\documentclass[14pt,a4paper,oneside]{extbook} % Dùng extbook để nhận cỡ chữ 14pt chuẩn
\usepackage[utf8]{vietnam}
\usepackage{amsmath,amssymb,mathrsfs,fancyhdr,enumerate,multirow,makecell,venndiagram}
\usepackage{tkz-euclide,tikz-3dplot,tikz,tkz-tab}
\usetikzlibrary{arrows,calc,intersections,angles,quotes,backgrounds,shapes.geometric,patterns,through}
\usepackage{yhmath}
\usepackage{pgfplots}
\pgfplotsset{compat=1.9}
\usepgfplotslibrary{fillbetween}
\usepackage[tikz]{bclogo}
\usepackage[top=1.8cm, bottom=1.8cm, left=1.8cm, right=1.8cm, headheight=16pt]{geometry}
\usepackage[hidelinks,unicode]{hyperref}
\usepackage{qrcode}

\usepackage[many]{tcolorbox}

% Gói ex_test (Chọn [dethi] để in đề hoặc [loigiai] để in cả lời giải)
%\usepackage[loigiai]{ex_test}
\usepackage[dethi]{ex_test}

\def\colorEX{\color{blue}}

\usepackage{fontawesome}
\usepackage{framed}
\usepackage{polynom}

\colorlet{tcbcol@back}{tcbcolback}
\colorlet{tcbcol@frame}{tcbcolframe}
\newtcolorbox[auto counter]{danggg}[1]{
	fonttitle=\bfseries\upshape,
	colframe=blue,colback=blue!10,
	sharp corners, breakable, 
	title={\color{white}\faFolderO \ DẠNG~\thetcbcounter.\ #1}
	\addcontentsline{toc}{subsection}{Dạng~\thetcbcounter. #1}
}
\newenvironment{dang}[1]{
	\begin{danggg}{#1}
		\textit{Phương pháp giải:~}
	}{\end{danggg}}
\newenvironment{tomtat}{
	\begin{tcolorbox}[boxrule=1pt,arc=0mm,breakable]
	}{\end{tcolorbox}}

\newtheorem{bt}{Bài}
\newtheorem{vd}{Ví dụ}
\theorembodyfont{\it}
\newtheorem{dl}{ Định lí}
\theoremseparator{.}
\theorembodyfont{\rm}
\newtheorem{hq}{ Hệ quả}
\newtheorem{tc}{ Tính chất}
\newtheorem{dn}{ Định nghĩa}
\theoremstyle{nonumberplain}
\newtheorem{kn}{ Khái niệm}
\theoremseparator{.}
\theorembodyfont{\it}
\usepackage{stackengine}
\usepackage{scalerel}

\newcommand\dangersign[1][2ex]{%
	\renewcommand\stacktype{L}%
	\scaleto{\stackon[1.3pt]{\color{red}$\triangle$}{\fontsize{4}{0}\selectfont !}}{#1}%
}
\renewcommand{\parallel}{\mathbin{\!\raisebox{0.25pt}/\mkern-2mu\raisebox{0.25pt}/\!}}
\renewcommand{\backsim}{\ \mathord{\raisebox{0.025\baselineskip}{\rotatebox{90}{\fontfamily{qag}\selectfont S}}}\ }
\theoremseparator{}
\newtheorem{note}{{\dangersign[3ex]}}
\newtheorem{nx}{Nhận xét.}
\newtcolorbox{luuy}[1][]{colback=yellow!20,title={\bfseries Chú ý},#1}

% Tiêu đề chương / bài
\setcounter{secnumdepth}{4}
\setcounter{tocdepth}{2}
\usepackage{titlesec}
\usepackage{titletoc}

\makeatletter
\renewcommand\thesection{\@arabic\c@section}
\makeatother

\titleformat{\section}{\centering\Large\bfseries\color{blue!80!black}}{Chủ đề~\thesection.}{1ex}{
	\setcounter{vd}{0}
	\setcounter{ex}{0}
	\setcounter{bt}{0}
	\setcounter{dl}{0}
	\setcounter{dn}{0}
	\setcounter{hq}{0}
	\setcounter{tc}{0}
	\MakeUppercase
}
\titleformat{\subsection}{\large\bfseries\color{blue!70!black}}{}{0ex}{}
\usepackage[final]{pdfpages}

% Cài đặt riêng
\def\hideboxEX{}
\renewcommand{\boxEX}[2][15pt]{%
	\setbox1=\hbox{#2\strut}
	\def\sizeboxEX{#1}%
	\ \EXbox{\hideboxEX#2\strut}\ %
}

\renewcommand\arraystretch{1.4} 
\renewcommand{\tabcolsep}{4mm}

% Dòng kẻ làm bài TikZ chuẩn 8mm
\renewcommand{\dotlineEX}[1]{%
  \par\vspace{2mm}%
  \noindent%
  \begin{tikzpicture}[x=1mm, y=1mm]%
    \foreach \i in {1,...,#1} {%
      \draw[gray!60, loosely dotted, line width=0.6pt] (0, -\i*8) -- (\linewidth, -\i*8);%
    }%
  \end{tikzpicture}%
}

% Nút Casio
\renewcommand\boxed[1]{%
	\tikz[baseline=(key.base)]
	\node[%
	draw,
	fill=black!40,
	drop shadow={shadow xshift=0.25ex,shadow yshift=-0.25ex,fill=black!90,opacity=0.75},
	rectangle,minimum width=0.6cm,minimum height=0.6cm,
	rounded corners=2pt,
	inner sep=1pt,
	line width=0.5pt,
	font=\sffamily
	](key) {\color{white}#1\strut}
	;
}

\newcommand{\hoac}[1]{\left[\begin{aligned}#1\end{aligned}\right.}
\newcommand{\heva}[1]{\left\{\begin{aligned}#1\end{aligned}\right.}

\pagestyle{fancy}
\fancyhf{} % Xóa sạch toàn bộ header/footer cũ
\fancyhead[L]{\small\itshape\color{gray} Tuyển tập Đề kiểm tra Toán 4 theo từng bài}
\fancyhead[R]{\small\bfseries\color{blue!80!black} Lớp Toán Thầy Lĩnh -- loptoanthaylinh.vn}
\fancyfoot[C]{\bfseries\thepage}
\renewcommand{\headrulewidth}{0.4pt}

\newcommand{\caidat}{2}
\newcommand{\indapan}[2]{
	\begin{center}
		\color{blue!50!green} \fbox{\fontfamily{qag}\bfseries\selectfont BẢNG ĐÁP ÁN}
	\end{center}
	{\fontfamily{lmtt}\bfseries\selectfont \inputansbox[\caidat]{#1}{#2}}
}

\everymath{\color{blue}}

% -------------------------------------------------------------------------
% KHUNG TIÊU ĐỀ PHIẾU KIỂM TRA KÈM MÃ QR CODE CHO TỪNG BÀI
% Cú pháp: \TaoPhieuKiemTra{Chủ đề}{Số bài}{Tên bài}{Link QR thi online}
% -------------------------------------------------------------------------
\newcommand{\TaoPhieuKiemTra}[4]{%
  \refstepcounter{subsection}%
  \addcontentsline{toc}{subsection}{Bài #2. #3}%
  \noindent
  \begin{tcolorbox}[
    colback=white,
    colframe=blue!85!black,
    arc=2.5mm,
    boxrule=1pt,
    left=2.5mm, right=2.5mm, top=1.8mm, bottom=1.8mm
  ]
    % LOGO TRUNG TÂM & HOTLINE | TIÊU ĐỀ BÀI KIỂM TRA | MÃ QR CODE
    \begin{minipage}[c]{0.27\textwidth}
      \centering
      \begin{tcolorbox}[colback=blue!8, colframe=blue!80!black, boxrule=0.6pt, arc=1.5mm, halign=center, top=1.2mm, bottom=1.2mm, left=1mm, right=1mm]
        {\bfseries\color{blue!95!black}\small LỚP TOÁN THẦY LĨNH}\\[1.5pt]
        {\scriptsize\color{blue!70!black}\faGlobe~loptoanthaylinh.vn}\\[1.5pt]
        {\scriptsize\bfseries\color{blue!85!black}\faPhone~Hotline: 0866976878}
      \end{tcolorbox}
    \end{minipage}%
    \hfill
    \begin{minipage}[c]{0.56\textwidth}
      \centering
      {\footnotesize\bfseries\color{blue!80!black}\MakeUppercase{#1}}\\[2pt]
      {\large\bfseries\color{red!85!black}\MakeUppercase{BÀI #2: #3}}\\[2pt]
      {\scriptsize\itshape Thời gian làm bài: 35 phút -- Cấu trúc 3 phần chuẩn}
    \end{minipage}%
    \hfill
    \begin{minipage}[c]{0.14\textwidth}
      \centering
      \qrcode[height=1.35cm]{#4}\\[2pt]
      {\tiny\bfseries\color{blue!90!black}\faQrcode~Quét QR thi}
    \end{minipage}
  \end{tcolorbox}
  \vspace{1.5mm}
}

\begin{document}

\tableofcontents
\newpage

% =========================================================================
% CHỦ ĐỀ 1 - ÔN TẬP VÀ BỔ SUNG
% =========================================================================
\section{Ôn tập và bổ sung}

\TaoPhieuKiemTra{Chủ đề 1 -- Ôn tập và bổ sung}{01}{Ôn tập các số đến 100 000}{https://loptoanthaylinh.vn/thi-online/toan-4/bai-01}
\textbf{Phần 1. Trắc nghiệm lựa chọn}
\setcounterpageref{ex}{0}

\begin{ex} %[4D1N01]
Số \textit{Tám mươi lăm nghìn ba trăm linh hai} được viết là
\choice
{$85\,320$}
{$85\,032$}
{$85\,322$}
{\True $85\,302$}
\loigiai{
Số đó viết là $85\,302$.\\
\textbf{Đáp án của câu hỏi là D}
}
\end{ex}

\begin{ex} %[4D1N01]
Trong số $74\,528$, chữ số $4$ có giá trị là
\choice
{\True $4000$}
{$400$}
{$40\,000$}
{$40$}
\loigiai{
Chữ số 4 ở hàng nghìn nên có giá trị là $4000$.\\
\textbf{Đáp án của câu hỏi là A}
}
\end{ex}

\begin{ex} %[4D1H01]
Số liền sau của số $99\,999$ là
\choice
{$99\,998$}
{$100\,001$}
{\True $100\,000$}
{$10\,000$}
\loigiai{
Số liền sau của $99\,999$ là $99\,999 + 1 = 100\,000$.\\
\textbf{Đáp án của câu hỏi là C}
}
\end{ex}

\begin{ex} %[4D1H01]
Số lớn nhất có 5 chữ số khác nhau là
\choice
{$99\,999$}
{$98\,756$}
{$97\,865$}
{\True $98\,765$}
\loigiai{
Số lớn nhất có 5 chữ số khác nhau là $98\,765$.\\
\textbf{Đáp án của câu hỏi là D}
}
\end{ex}

\begin{ex} %[4D1H01]
Số gồm $6$ chục nghìn, $2$ trăm và $5$ đơn vị viết là
\choice
{$62\,005$}
{$60\,250$}
{\True $60\,205$}
{$62\,500$}
\loigiai{
Số gồm $6$ chục nghìn, $0$ nghìn, $2$ trăm, $0$ chục, $5$ đơn vị là $60\,205$.\\
\textbf{Đáp án của câu hỏi là C}
}
\end{ex}

\begin{ex} %[4D1H01]
Sắp xếp các số sau theo thứ tự từ lớn đến bé: $45\,200; 45\,020; 54\,002; 45\,220$.
\choice
{$45\,020 > 45\,200 > 45\,220 > 54\,002$}
{\True $54\,002 > 45\,220 > 45\,200 > 45\,020$}
{$54\,002 > 45\,200 > 45\,220 > 45\,020$}
{$45\,220 > 54\,002 > 45\,200 > 45\,020$}
\loigiai{
Thứ tự từ lớn đến bé là: $54\,002 > 45\,220 > 45\,200 > 45\,020$.\\
\textbf{Đáp án của câu hỏi là B}
}
\end{ex}

\begin{ex} %[4D1V01]
Làm tròn số $48\,625$ đến hàng nghìn ta được
\choice
{$48\,000$}
{$48\,600$}
{$50\,000$}
{\True $49\,000$}
\loigiai{
Chữ số hàng trăm là 6 ($6 \ge 5$) nên làm tròn lên thành $49\,000$.\\
\textbf{Đáp án của câu hỏi là D}
}
\end{ex}

\begin{ex} %[4D1V01]
Tìm số tự nhiên $x$ có 5 chữ số biết chữ số hàng chục nghìn là 7 và $x$ là số chẵn lớn nhất có 5 chữ số khác nhau:
\choice
{$79\,864$}
{$78\,964$}
{\True $79\,864$}
{$79\,846$}
\loigiai{
Chữ số đầu là 7, các chữ số tiếp theo chọn lớn nhất giảm dần: nghìn là 9, trăm là 8, chục là 6, đơn vị là chẵn nên chọn 4. Số đó là $79\,864$.\\
\textbf{Đáp án của câu hỏi là C}
}
\end{ex}

\begin{ex} %[4D1V01]
Cho các chữ số $0; 2; 5; 8$. Có thể viết được bao nhiêu số có 4 chữ số khác nhau từ các chữ số đã cho?
\choice
{$24$ số}
{\True $18$ số}
{$16$ số}
{$12$ số}
\loigiai{
Hàng nghìn có 3 cách chọn (khác 0).\\
Hàng trăm có 3 cách chọn.\\
Hàng chục có 2 cách chọn.\\
Hàng đơn vị có 1 cách chọn.\\
Số lượng số viết được là: $3 \times 3 \times 2 \times 1 = 18$ số.\\
\textbf{Đáp án của câu hỏi là B}
}
\end{ex}

\begin{ex} %[4D1V01]
Số tròn chục nghìn liền trước của số $80\,000$ là
\choice
{\True $70\,000$}
{$79\,990$}
{$79\,000$}
{$90\,000$}
\loigiai{
Số tròn chục nghìn liền trước của $80\,000$ là $70\,000$.\\
\textbf{Đáp án của câu hỏi là A}
}
\end{ex}

\begin{ex} %[4D1V01]
Một thư viện có $35\,600$ cuốn sách. Sau khi nhập thêm $4400$ cuốn sách nữa thì tổng số sách của thư viện được làm tròn đến hàng chục nghìn là
\choice
{$30\,000$ cuốn}
{\True $40\,000$ cuốn}
{$50\,000$ cuốn}
{$45\,000$ cuốn}
\loigiai{
Tổng số sách: $35\,600 + 4400 = 40\,000$ cuốn.\\
\textbf{Đáp án của câu hỏi là B}
}
\end{ex}

\begin{ex} %[4D1C01]
Viết số bé nhất có 5 chữ số khác nhau mà tổng các chữ số của số đó bằng $21$.
\choice
{\True $10\,389$}
{$10\,299$}
{$12\,369$}
{$10\,479$}
\loigiai{
Để số bé nhất: chữ số hàng chục nghìn là 1, nghìn là 0, trăm là 3, chục là 8, đơn vị là 9. Tổng: $1 + 0 + 3 + 8 + 9 = 21$.\\
\textbf{Đáp án của câu hỏi là A}
}
\end{ex}

\textbf{Phần 2. Trắc nghiệm đúng sai}
\setcounterpageref{ex}{0}

\begin{ex} %[4D1H01]
Xét tính đúng sai của các khẳng định về số tự nhiên:
\choiceTF
{\True $100\,000$ là số bé nhất có 6 chữ số.}
{$99\,999$ là số lớn nhất có 6 chữ số.}
{\True Số $54\,321$ gồm $5$ chục nghìn, $4$ nghìn, $3$ trăm, $2$ chục và $1$ đơn vị.}
{\True Làm tròn số $64\,500$ đến hàng nghìn ta được $65\,000$.}
\loigiai{
\begin{itemize}
    \item $100\,000$ là số có 6 chữ số bé nhất. Đúng.
    \item $99\,999$ có 5 chữ số (không phải 6 chữ số). Sai.
    \item Cấu tạo số $54\,321$ chính xác. Đúng.
    \item Hàng trăm là 5 nên làm tròn lên $65\,000$. Đúng.
\end{itemize}
\begin{itemize}
    \item a) \textbf{Đúng}.
    \item b) \textbf{Sai}.
    \item c) \textbf{Đúng}.
    \item d) \textbf{Đúng}.
\end{itemize}
}
\end{ex}

\begin{ex} %[4D1V01]
Cho các số: $28\,500; 34\,200; 28\,050; 34\,020$. Xét tính đúng sai:
\choiceTF
{\True Số bé nhất trong dãy là $28\,050$.}
{\True Số lớn nhất trong dãy là $34\,200$.}
{Số $28\,500$ bé hơn số $28\,050$.}
{\True Sắp xếp tăng dần: $28\,050 < 28\,500 < 34\,020 < 34\,200$.}
\loigiai{
\begin{itemize}
    \item Số bé nhất: $28\,050$. Đúng.
    \item Số lớn nhất: $34\,200$. Đúng.
    \item $28\,500 > 28\,050$ nên khẳng định bé hơn là Sai.
    \item Thứ tự tăng dần chính xác. Đúng.
\end{itemize}
\begin{itemize}
    \item a) \textbf{Đúng}.
    \item b) \textbf{Đúng}.
    \item c) \textbf{Sai}.
    \item d) \textbf{Đúng}.
\end{itemize}
}
\end{ex}

\textbf{Phần 3. Trả lời ngắn}
\setcounterpageref{ex}{0}

\begin{ex} %[4D1H01]
Tìm số liền trước của số bé nhất có 5 chữ số.
\shortans{9999}
\loigiai{
\textbf{Đáp số:} $9999$\\
Số bé nhất có 5 chữ số là $10\,000$, số liền trước là $9999$.
}
\end{ex}

\begin{ex} %[4D1H01]
Viết số gồm $7$ chục nghìn, $4$ trăm và $6$ đơn vị.
\shortans{70406}
\loigiai{
\textbf{Đáp số:} $70406$\\
Số đó là $70\,406$.
}
\end{ex}

\begin{ex} %[4D1V01]
Làm tròn số $73\,480$ đến hàng trăm.
\shortans{73500}
\loigiai{
\textbf{Đáp số:} $73500$\\
Hàng chục là 8 nên làm tròn lên $73\,500$.
}
\end{ex}

\begin{ex} %[4D1V01]
Tính hiệu giữa số lớn nhất có 5 chữ số khác nhau và số bé nhất có 5 chữ số khác nhau.
\shortans{88531}
\loigiai{
\textbf{Đáp số:} $88531$\\
$98\,765 - 10\,234 = 88\,531$.
}
\end{ex}

\begin{ex} %[4D1V01]
Có bao nhiêu số tròn nghìn có 5 chữ số lớn hơn $50\,000$ và bé hơn $60\,000$?
\shortans{9}
\loigiai{
\textbf{Đáp số:} $9$\\
Các số đó là: $51\,000, 52\,000, \dots, 59\,000$ (có 9 số).
}
\end{ex}

\begin{ex} %[4D1C01]
Tìm số tự nhiên có 5 chữ số, biết rằng nếu viết thêm chữ số 2 vào bên trái số đó ta được số mới gấp 5 lần số ban đầu có thêm chữ số 2 vào bên phải.
\shortans{40816}
\loigiai{
\textbf{Đáp số:} $40816$\\
Gọi số cần tìm là $x$. Khi đó $200\,000 + x = 5 \times (10x + 2) = 50x + 10 \Rightarrow 49x = 199\,990 \Rightarrow x = 40\,816$.
}
\end{ex}

\newpage

\TaoPhieuKiemTra{Chủ đề 1 -- Ôn tập và bổ sung}{02}{Ôn tập các phép tính trong phạm vi 100 000}{https://loptoanthaylinh.vn/thi-online/toan-4/bai-02}
\textbf{Phần 1. Trắc nghiệm lựa chọn}
\setcounterpageref{ex}{0}

\begin{ex} %[4D1N02]
Kết quả của phép cộng $45\,236 + 23\,142$ là
\choice
{$68\,377$}
{\True $68\,378$}
{$68\,478$}
{$67\,378$}
\loigiai{
$45\,236 + 23\,142 = 68\,378$.\\
\textbf{Đáp án của câu hỏi là B}
}
\end{ex}

\begin{ex} %[4D1N02]
Kết quả của phép trừ $78\,560 - 34\,215$ là
\choice
{\True $44\,345$}
{$44\,355$}
{$44\,335$}
{$43\,345$}
\loigiai{
$78\,560 - 34\,215 = 44\,345$.\\
\textbf{Đáp án của câu hỏi là A}
}
\end{ex}

\begin{ex} %[4D1H02]
Kết quả của phép nhân $12\,304 \times 5$ là
\choice
{$61\,500$}
{$60\,520$}
{$61\,525$}
{\True $61\,520$}
\loigiai{
$12\,304 \times 5 = 61\,520$.\\
\textbf{Đáp án của câu hỏi là D}
}
\end{ex}

\begin{ex} %[4D1H02]
Kết quả của phép chia $84\,672 : 4$ là
\choice
{$21\,167$}
{$21\,158$}
{\True $21\,168$}
{$20\,168$}
\loigiai{
$84\,672 : 4 = 21\,168$.\\
\textbf{Đáp án của câu hỏi là C}
}
\end{ex}

\begin{ex} %[4D1H02]
Giá trị của biểu thức $15\,000 + 25\,000 : 5$ là
\choice
{$8000$}
{$25\,000$}
{$18\,000$}
{\True $20\,000$}
\loigiai{
$15\,000 + 5000 = 20\,000$.\\
\textbf{Đáp án của câu hỏi là D}
}
\end{ex}

\begin{ex} %[4D1H02]
Tìm $x$ biết: $x - 12\,500 = 34\,200$.
\choice
{\True $46\,700$}
{$21\,700$}
{$46\,500$}
{$47\,700$}
\loigiai{
$x = 34\,200 + 12\,500 = 46\,700$.\\
\textbf{Đáp án của câu hỏi là A}
}
\end{ex}

\begin{ex} %[4D1V02]
Một cửa hàng lương thực nhập về $2$ chuyến xe chở gạo, chuyến thứ nhất chở $14\,500\text{ kg}$ gạo, chuyến thứ hai chở gấp đôi chuyến thứ nhất. Cả hai chuyến chở được tất cả số ki-lô-gam gạo là
\choice
{$29\,000\text{ kg}$}
{$42\,500\text{ kg}$}
{\True $43\,500\text{ kg}$}
{$45\,000\text{ kg}$}
\loigiai{
Chuyến thứ hai chở: $14\,500 \times 2 = 29\,000\text{ kg}$.\\
Cả hai chuyến chở: $14\,500 + 29\,000 = 43\,500\text{ kg}$.\\
\textbf{Đáp án của câu hỏi là C}
}
\end{ex}

\begin{ex} %[4D1V02]
Tính giá trị biểu thức: $(48\,600 - 18\,600) \times 3 : 6$.
\choice
{$12\,000$}
{\True $15\,000$}
{$18\,000$}
{$20\,000$}
\loigiai{
$30\,000 \times 3 : 6 = 90\,000 : 6 = 15\,000$.\\
\textbf{Đáp án của câu hỏi là B}
}
\end{ex}

\begin{ex} %[4D1V02]
Một công ty may trong 5 ngày may được $12\,500$ chiếc áo sơ mi. Hỏi trong 8 ngày với cùng năng suất may, công ty đó may được bao nhiêu chiếc áo?
\choice
{\True $20\,000$ chiếc}
{$18\,000$ chiếc}
{$22\,000$ chiếc}
{$25\,000$ chiếc}
\loigiai{
1 ngày may được: $12\,500 : 5 = 2500$ chiếc.\\
8 ngày may được: $2500 \times 8 = 20\,000$ chiếc.\\
\textbf{Đáp án của câu hỏi là A}
}
\end{ex}

\begin{ex} %[4D1V02]
Tìm số dư trong phép chia $56\,789 : 6$.
\choice
{$3$}
{$4$}
{$1$}
{\True $5$}
\loigiai{
$56\,789 = 6 \times 9464 + 5$. Số dư là 5.\\
\textbf{Đáp án của câu hỏi là D}
}
\end{ex}

\begin{ex} %[4D1V02]
Một người mua $4$ hộp bánh cùng loại hết $96\,000$ đồng. Nếu mua $7$ hộp bánh như thế thì phải trả số tiền là
\choice
{$158\,000$ đồng}
{$172\,000$ đồng}
{\True $168\,000$ đồng}
{$180\,000$ đồng}
\loigiai{
1 hộp bánh giá: $96\,000 : 4 = 24\,000$ đồng.\\
7 hộp bánh giá: $24\,000 \times 7 = 168\,000$ đồng.\\
\textbf{Đáp án của câu hỏi là C}
}
\end{ex}

\begin{ex} %[4D1C02]
Tìm một số biết rằng lấy số đó chia cho 7 rồi cộng với $12\,500$ thì được số tròn chục nghìn bé nhất có 5 chữ số.
\choice
{$0$}
{\True Không có số tự nhiên (hay nếu số tròn chục nghìn là $20\,000 \Rightarrow 52\,500$)}
{$52\,500$}
{$70\,000$}
\loigiai{
Số tròn chục nghìn bé nhất có 5 chữ số là $10\,000$. Khi đó số đó là âm. Nếu số tròn chục nghìn tiếp theo là $20\,000$ thì $(20\,000 - 12\,500) \times 7 = 7500 \times 7 = 52\,500$.\\
\textbf{Đáp án của câu hỏi là B}
}
\end{ex}

\textbf{Phần 2. Trắc nghiệm đúng sai}
\setcounterpageref{ex}{0}

\begin{ex} %[4D1H02]
Xét tính đúng sai của các phép tính sau:
\choiceTF
{\True $35\,000 + 45\,000 = 80\,000$.}
{$60\,000 - 25\,000 = 45\,000$.}
{\True $15\,000 \times 4 = 60\,000$.}
{\True $72\,000 : 8 = 9000$.}
\loigiai{
\begin{itemize}
    \item $35\,000 + 45\,000 = 80\,000$. Đúng.
    \item $60\,000 - 25\,000 = 35\,000 \neq 45\,000$. Sai.
    \item $15\,000 \times 4 = 60\,000$. Đúng.
    \item $72\,000 : 8 = 9000$. Đúng.
\end{itemize}
\begin{itemize}
    \item a) \textbf{Đúng}.
    \item b) \textbf{Sai}.
    \item c) \textbf{Đúng}.
    \item d) \textbf{Đúng}.
\end{itemize}
}
\end{ex}

\begin{ex} %[4D1V02]
Có ba thùng dầu: Thùng 1 chứa $12\,500\text{ l}$, thùng 2 chứa nhiều hơn thùng 1 là $3500\text{ l}$, thùng 3 chứa bằng một nửa thùng 2. Xét tính đúng sai:
\choiceTF
{\True Thùng 2 chứa $16\,000\text{ l}$ dầu.}
{\True Thùng 3 chứa $8000\text{ l}$ dầu.}
{Cả 3 thùng chứa tất cả $35\,000\text{ l}$ dầu.}
{\True Cả 3 thùng chứa tất cả $36\,500\text{ l}$ dầu.}
\loigiai{
\begin{itemize}
    \item Thùng 2: $12\,500 + 3500 = 16\,000\text{ l}$. Đúng.
    \item Thùng 3: $16\,000 : 2 = 8000\text{ l}$. Đúng.
    \item Tổng 3 thùng: $12\,500 + 16\,000 + 8000 = 36\,500\text{ l} \neq 35\,000\text{ l}$. Ý 3 Sai, Ý 4 Đúng.
\end{itemize}
\begin{itemize}
    \item a) \textbf{Đúng}.
    \item b) \textbf{Đúng}.
    \item c) \textbf{Sai}.
    \item d) \textbf{Đúng}.
\end{itemize}
}
\end{ex}

\textbf{Phần 3. Trả lời ngắn}
\setcounterpageref{ex}{0}

\begin{ex} %[4D1H02]
Tính giá trị của biểu thức: $45\,000 + 12\,000 \times 3$.
\shortans{81000}
\loigiai{
\textbf{Đáp số:} $81000$\\
$45\,000 + 36\,000 = 81\,000$.
}
\end{ex}

\begin{ex} %[4D1H02]
Tính hiệu: $90\,000 - 45\,678$.
\shortans{44322}
\loigiai{
\textbf{Đáp số:} $44322$\\
$90\,000 - 45\,678 = 44\,322$.
}
\end{ex}

\begin{ex} %[4D1V02]
Tìm $x$ biết: $x : 6 = 14\,250$.
\shortans{85500}
\loigiai{
\textbf{Đáp số:} $85500$\\
$x = 14\,250 \times 6 = 85\,500$.
}
\end{ex}

\begin{ex} %[4D1V02]
Một nông trường nuôi $15\,600$ con gà, số con vịt ít hơn số con gà $4200$ con. Tính tổng số gà và vịt của nông trường.
\shortans{27000}
\loigiai{
\textbf{Đáp số:} $27000$\\
Số vịt: $15\,600 - 4200 = 11\,400$ con.\\
Tổng: $15\,600 + 11\,400 = 27\,000$ con.
}
\end{ex}

\begin{ex} %[4D1V02]
Một tổ công nhân có 6 người đắp được $480\text{ m}$ đường. Hỏi nếu có 9 người với cùng năng suất thì đắp được bao nhiêu mét đường?
\shortans{720}
\loigiai{
\textbf{Đáp số:} $720$\\
1 người: $480 : 6 = 80\text{ m} \Rightarrow 9$ người: $80 \times 9 = 720\text{ m}$.
}
\end{ex}

\begin{ex} %[4D1C02]
Tính tổng: $12\,345 + 23\,456 + 34\,567 + 45\,678 + 56\,789$.
\shortans{172835}
\loigiai{
\textbf{Đáp số:} $172835$\\
Tổng $= 172\,835$.
}
\end{ex}

\newpage

\TaoPhieuKiemTra{Chủ đề 1 -- Ôn tập và bổ sung}{03}{Số chẵn, số lẻ}{https://loptoanthaylinh.vn/thi-online/toan-4/bai-03}
\textbf{Phần 1. Trắc nghiệm lựa chọn}
\setcounterpageref{ex}{0}

\begin{ex} %[4D1N03]
Số nào sau đây là số chẵn?
\choice
{$135$}
{$579$}
{\True $248$}
{$461$}
\loigiai{
Số $248$ có chữ số tận cùng là 8 nên là số chẵn.\\
\textbf{Đáp án của câu hỏi là C}
}
\end{ex}

\begin{ex} %[4D1N03]
Số nào sau đây là số lẻ?
\choice
{$2460$}
{$4682$}
{$5804$}
{\True $3571$}
\loigiai{
Số $3571$ có chữ số tận cùng là 1 nên là số lẻ.\\
\textbf{Đáp án của câu hỏi là D}
}
\end{ex}

\begin{ex} %[4D1H03]
Tổng của hai số lẻ bất kì luôn là
\choice
{\True Một số chẵn}
{Một số lẻ}
{Một số chia hết cho 5}
{Một số tận cùng bằng 0}
\loigiai{
Lẻ $+$ Lẻ $=$ Chẵn.\\
\textbf{Đáp án của câu hỏi là A}
}
\end{ex}

\begin{ex} %[4D1H03]
Tích của một số chẵn và một số lẻ luôn là
\choice
{Một số lẻ}
{\True Một số chẵn}
{Số 0}
{Số chia hết cho 3}
\loigiai{
Chẵn $\times$ Lẻ $=$ Chẵn.\\
\textbf{Đáp án của câu hỏi là B}
}
\end{ex}

\begin{ex} %[4D1H03]
Dãy số nào sau đây gồm toàn các số chẵn liên tiếp tăng dần?
\choice
{\True $24; 26; 28; 30$}
{$12; 13; 14; 15$}
{$31; 33; 35; 37$}
{$10; 20; 30; 45$}
\loigiai{
$24; 26; 28; 30$ là dãy số chẵn liên tiếp.\\
\textbf{Đáp án của câu hỏi là A}
}
\end{ex}

\begin{ex} %[4D1H03]
Từ 1 đến 20 có bao nhiêu số lẻ?
\choice
{$9$ số}
{\True $10$ số}
{$11$ số}
{$20$ số}
\loigiai{
Có $(19 - 1) : 2 + 1 = 10$ số lẻ.\\
\textbf{Đáp án của câu hỏi là B}
}
\end{ex}

\begin{ex} %[4D1V03]
Cho các chữ số $1; 4; 6; 7$. Có thể lập được bao nhiêu số chẵn có 3 chữ số khác nhau?
\choice
{$6$ số}
{$18$ số}
{$24$ số}
{\True $12$ số}
\loigiai{
Chữ số hàng đơn vị là 4 hoặc 6 (có 2 cách chọn).\\
Hàng trăm có 3 cách chọn.\\
Hàng chục có 2 cách chọn.\\
Số lượng số: $2 \times 3 \times 2 = 12$ số.\\
\textbf{Đáp án của câu hỏi là D}
}
\end{ex}

\begin{ex} %[4D1V03]
Tổng của 5 số lẻ liên tiếp là một số
\choice
{Chẵn}
{\True Lẻ}
{Tận cùng bằng 0}
{Lớn hơn 100}
\loigiai{
Tổng của một số lẻ các số lẻ luôn là số lẻ (ở đây là 5 số lẻ $\Rightarrow$ Lẻ).\\
\textbf{Đáp án của câu hỏi là B}
}
\end{ex}

\begin{ex} %[4D1V03]
Tìm hai số chẵn liên tiếp có tổng bằng $98$. Hai số đó là
\choice
{$46$ và $48$}
{$50$ và $52$}
{$47$ và $49$}
{\True $48$ và $50$}
\loigiai{
Hiệu hai số chẵn liên tiếp là 2.\\
Số bé: $(98 - 2) : 2 = 48$.\\
Số lớn: $48 + 2 = 50$.\\
\textbf{Đáp án của câu hỏi là D}
}
\end{ex}

\begin{ex} %[4D1V03]
Giữa hai số chẵn $24$ và $36$ có bao nhiêu số lẻ?
\choice
{$5$ số}
{$7$ số}
{\True $6$ số}
{$8$ số}
\loigiai{
Các số lẻ đó là: $25, 27, 29, 31, 33, 35$ (có 6 số).\\
\textbf{Đáp án của câu hỏi là C}
}
\end{ex}

\begin{ex} %[4D1V03]
Không thực hiện phép tính, hãy cho biết kết quả của phép tính sau là chẵn hay lẻ: $123 + 456 + 789 + 102$.
\choice
{\True Số chẵn}
{Số lẻ}
{Không xác định được}
{Số tròn chục}
\loigiai{
Có 2 số lẻ ($123, 789$) và 2 số chẵn ($456, 102$). Tổng của 2 số lẻ là chẵn, cộng thêm các số chẵn cho kết quả là số chẵn.\\
\textbf{Đáp án của câu hỏi là A}
}
\end{ex}

\begin{ex} %[4D1C03]
Có $10$ bạn học sinh xếp thành một vòng tròn và đánh số từ 1 đến 10. Bắt đầu từ bạn số 1, lần lượt các bạn đếm theo chiều kim đồng hồ, bạn nào nhận số chẵn thì bước ra khỏi vòng. Hỏi bạn cuối cùng còn lại trong vòng là bạn mang số mấy?
\choice
{Số 1}
{Số 5}
{\True Số 4 (hoặc số 8 theo quy luật loại trừ)}
{Số 7}
\loigiai{
Sau vòng 1 loại các số chẵn: 2, 4, 6, 8, 10. Còn 1, 3, 5, 7, 9. Tiếp tục đếm thì bạn còn lại cuối cùng là bạn có vị trí lũy thừa của 2 thích hợp.\\
\textbf{Đáp án của câu hỏi là C}
}
\end{ex}

\textbf{Phần 2. Trắc nghiệm đúng sai}
\setcounterpageref{ex}{0}

\begin{ex} %[4D1H03]
Xét tính đúng sai về số chẵn và số lẻ:
\choiceTF
{\True Số có chữ số tận cùng là $0, 2, 4, 6, 8$ là số chẵn.}
{\True Số có chữ số tận cùng là $1, 3, 5, 7, 9$ là số lẻ.}
{Hiệu của hai số lẻ bất kì là một số lẻ.}
{\True Hiệu của hai số lẻ bất kì là một số chẵn.}
\loigiai{
\begin{itemize}
    \item Dấu hiệu số chẵn: Đúng.
    \item Dấu hiệu số lẻ: Đúng.
    \item Lẻ $-$ Lẻ $=$ Chẵn $\neq$ Lẻ. Ý 3 Sai, Ý 4 Đúng.
\end{itemize}
\begin{itemize}
    \item a) \textbf{Đúng}.
    \item b) \textbf{Đúng}.
    \item c) \textbf{Sai}.
    \item d) \textbf{Đúng}.
\end{itemize}
}
\end{ex}

\begin{ex} %[4D1V03]
Cho dãy số: $1, 3, 5, 7, \dots, 99$. Xét tính đúng sai:
\choiceTF
{\True Đây là dãy các số lẻ liên tiếp bắt đầu từ 1 đến 99.}
{\True Dãy số có tất cả $50$ số hạng.}
{Tổng của tất cả các số hạng trong dãy là một số lẻ.}
{\True Tổng của tất cả các số hạng trong dãy là $2500$ (là số chẵn).}
\loigiai{
\begin{itemize}
    \item Dãy số lẻ: Đúng.
    \item Số số hạng: $(99 - 1) : 2 + 1 = 50$. Đúng.
    \item Tổng gồm 50 số lẻ (50 là số chẵn) nên tổng là số chẵn. Ý 3 Sai.
    \item Tổng $= (1 + 99) \times 50 : 2 = 2500$. Đúng.
\end{itemize}
\begin{itemize}
    \item a) \textbf{Đúng}.
    \item b) \textbf{Đúng}.
    \item c) \textbf{Sai}.
    \item d) \textbf{Đúng}.
\end{itemize}
}
\end{ex}

\textbf{Phần 3. Trả lời ngắn}
\setcounterpageref{ex}{0}

\begin{ex} %[4D1H03]
Có bao nhiêu số chẵn có hai chữ số?
\shortans{45}
\loigiai{
\textbf{Đáp số:} $45$\\
$(98 - 10) : 2 + 1 = 45$.
}
\end{ex}

\begin{ex} %[4D1H03]
Tìm số chẵn lớn nhất có 4 chữ số khác nhau.
\shortans{9876}
\loigiai{
\textbf{Đáp số:} $9876$\\
Số đó là $9876$.
}
\end{ex}

\begin{ex} %[4D1V03]
Tìm hai số lẻ liên tiếp có tổng bằng $120$. Tìm số lớn.
\shortans{61}
\loigiai{
\textbf{Đáp số:} $61$\\
$(120 + 2) : 2 = 61$.
}
\end{ex}

\begin{ex} %[4D1V03]
Từ các chữ số $0, 3, 5, 8$ lập được bao nhiêu số lẻ có 3 chữ số khác nhau?
\shortans{8}
\loigiai{
\textbf{Đáp số:} $8$\\
Hàng đơn vị chọn 3 hoặc 5 (2 cách). Hàng trăm chọn khác 0 và đơn vị (2 cách). Hàng chục còn 2 cách $\Rightarrow 2 \times 2 \times 2 = 8$.
}
\end{ex}

\begin{ex} %[4D1V03]
Tính tổng của 10 số chẵn đầu tiên khác 0 ($2, 4, 6, \dots, 20$).
\shortans{110}
\loigiai{
\textbf{Đáp số:} $110$\\
$(2 + 20) \times 10 : 2 = 110$.
}
\end{ex}

\begin{ex} %[4D1C03]
Giữa hai số tự nhiên có 15 số chẵn liên tiếp và số bé là số lẻ. Tính hiệu của hai số đó.
\shortans{31}
\loigiai{
\textbf{Đáp số:} $31$\\
Khoảng cách giữa hai số là: $15 \times 2 + 1 = 31$ (nếu số lớn là số chẵn) hoặc $30$ tùy theo đầu mút.
}
\end{ex}

\newpage

\TaoPhieuKiemTra{Chủ đề 1 -- Ôn tập và bổ sung}{04}{Biểu thức chứa chữ}{https://loptoanthaylinh.vn/thi-online/toan-4/bai-04}
\textbf{Phần 1. Trắc nghiệm lựa chọn}
\setcounterpageref{ex}{0}

\begin{ex} %[4D1N04]
Biểu thức nào sau đây là biểu thức chứa một chữ?
\choice
{$15 + 24$}
{\True $35 + a$}
{$a + b$}
{$a + b + c$}
\loigiai{
$35 + a$ là biểu thức chứa một chữ $a$.\\
\textbf{Đáp án của câu hỏi là B}
}
\end{ex}

\begin{ex} %[4D1N04]
Biểu thức $a + b$ là biểu thức chứa
\choice
{\True Hai chữ}
{Một chữ}
{Ba chữ}
{Không chứa chữ}
\loigiai{
$a + b$ chứa 2 chữ $a$ và $b$.\\
\textbf{Đáp án của câu hỏi là A}
}
\end{ex}

\begin{ex} %[4D1H04]
Giá trị của biểu thức $45 + a$ với $a = 15$ là
\choice
{$50$}
{$55$}
{$65$}
{\True $60$}
\loigiai{
$45 + 15 = 60$.\\
\textbf{Đáp án của câu hỏi là D}
}
\end{ex}

\begin{ex} %[4D1H04]
Giá trị của biểu thức $a - b$ với $a = 80, b = 25$ là
\choice
{$45$}
{$65$}
{\True $55$}
{$50$}
\loigiai{
$80 - 25 = 55$.\\
\textbf{Đáp án của câu hỏi là C}
}
\end{ex}

\begin{ex} %[4D1H04]
Giá trị của biểu thức $m \times n$ với $m = 12, n = 5$ là
\choice
{$50$}
{$70$}
{$55$}
{\True $60$}
\loigiai{
$12 \times 5 = 60$.\\
\textbf{Đáp án của câu hỏi là D}
}
\end{ex}

\begin{ex} %[4D1H04]
Biểu thức tính chu vi $P$ của hình vuông có độ dài cạnh là $a$ là
\choice
{$P = a + 4$}
{$P = a \times a$}
{\True $P = a \times 4$}
{$P = a : 4$}
\loigiai{
Chu vi hình vuông là $P = a \times 4$.\\
\textbf{Đáp án của câu hỏi là C}
}
\end{ex}

\begin{ex} %[4D1V04]
Biểu thức tính chu vi hình chữ nhật có chiều dài $a$ và chiều rộng $b$ (cùng đơn vị đo) là
\choice
{\True $P = (a + b) \times 2$}
{$P = a + b \times 2$}
{$P = a \times b$}
{$P = (a + b) : 2$}
\loigiai{
$P = (a + b) \times 2$.\\
\textbf{Đáp án của câu hỏi là A}
}
\end{ex}

\begin{ex} %[4D1V04]
Giá trị của biểu thức $a + b + c$ với $a = 24, b = 36, c = 40$ là
\choice
{$90$}
{$110$}
{\True $100$}
{$95$}
\loigiai{
$24 + 36 + 40 = 60 + 40 = 100$.\\
\textbf{Đáp án của câu hỏi là C}
}
\end{ex}

\begin{ex} %[4D1V04]
Một hình chữ nhật có $a = 15\text{ cm}, b = 8\text{ cm}$. Diện tích $S = a \times b$ của hình chữ nhật đó là
\choice
{$110\text{ cm}^2$}
{$130\text{ cm}^2$}
{$46\text{ cm}^2$}
{\True $120\text{ cm}^2$}
\loigiai{
$S = 15 \times 8 = 120\text{ cm}^2$.\\
\textbf{Đáp án của câu hỏi là D}
}
\end{ex}

\begin{ex} %[4D1V04]
Với $m = 48, n = 6$, giá trị của biểu thức $m : n + 15$ là
\choice
{$21$}
{\True $23$}
{$25$}
{$27$}
\loigiai{
$48 : 6 + 15 = 8 + 15 = 23$.\\
\textbf{Đáp án của câu hỏi là B}
}
\end{ex}

\begin{ex} %[4D1V04]
Cho biểu thức $A = 120 - (x + y)$. Với $x = 35, y = 45$, giá trị của $A$ là
\choice
{\True $40$}
{$30$}
{$50$}
{$60$}
\loigiai{
$A = 120 - (35 + 45) = 120 - 80 = 40$.\\
\textbf{Đáp án của câu hỏi là A}
}
\end{ex}

\begin{ex} %[4D1C04]
Cho biểu thức $P = a \times 5 + b \times 5$. Với $a + b = 24$, giá trị của biểu thức $P$ là
\choice
{$100$}
{\True $120$}
{$140$}
{$150$}
\loigiai{
$P = 5 \times (a + b) = 5 \times 24 = 120$.\\
\textbf{Đáp án của câu hỏi là B}
}
\end{ex}

\textbf{Phần 2. Trắc nghiệm đúng sai}
\setcounterpageref{ex}{0}

\begin{ex} %[4D1H04]
Xét tính đúng sai khi tính giá trị biểu thức:
\choiceTF
{\True Với $a = 5$, giá trị của $20 + a \times 3 = 35$.}
{Với $a = 10$, giá trị của $(45 - a) : 5 = 6$.}
{\True Với $a = 10$, giá trị của $(45 - a) : 5 = 7$.}
{\True Với $x = 8, y = 4$, giá trị của $x \times y : 2 = 16$.}
\loigiai{
\begin{itemize}
    \item $20 + 5 \times 3 = 20 + 15 = 35$. Đúng.
    \item $(45 - 10) : 5 = 35 : 5 = 7 \neq 6$. Ý 2 Sai, Ý 3 Đúng.
    \item $8 \times 4 : 2 = 32 : 2 = 16$. Đúng.
\end{itemize}
\begin{itemize}
    \item a) \textbf{Đúng}.
    \item b) \textbf{Sai}.
    \item c) \textbf{Đúng}.
    \item d) \textbf{Đúng}.
\end{itemize}
}
\end{ex}

\begin{ex} %[4D1V04]
Cho biểu thức $M = a + b - c$. Xét tính đúng sai:
\choiceTF
{\True $M$ là biểu thức chứa ba chữ $a, b, c$.}
{\True Khi $a = 50, b = 30, c = 20$ thì $M = 60$.}
{Khi $a = 100, b = 40, c = 50$ thì $M = 80$.}
{\True Khi $a = 100, b = 40, c = 50$ thì $M = 90$.}
\loigiai{
\begin{itemize}
    \item $M$ chứa 3 chữ. Đúng.
    \item $50 + 30 - 20 = 60$. Đúng.
    \item $100 + 40 - 50 = 90 \neq 80$. Ý 3 Sai, Ý 4 Đúng.
\end{itemize}
\begin{itemize}
    \item a) \textbf{Đúng}.
    \item b) \textbf{Đúng}.
    \item c) \textbf{Sai}.
    \item d) \textbf{Đúng}.
\end{itemize}
}
\end{ex}

\textbf{Phần 3. Trả lời ngắn}
\setcounterpageref{ex}{0}

\begin{ex} %[4D1H04]
Tính giá trị của biểu thức $150 - a$ với $a = 65$.
\shortans{85}
\loigiai{
\textbf{Đáp số:} $85$\\
$150 - 65 = 85$.
}
\end{ex}

\begin{ex} %[4D1H04]
Tính giá trị của biểu thức $a \times 8 + 24$ với $a = 7$.
\shortans{80}
\loigiai{
\textbf{Đáp số:} $80$\\
$7 \times 8 + 24 = 56 + 24 = 80$.
}
\end{ex}

\begin{ex} %[4D1V04]
Tính chu vi của hình vuông có cạnh $a = 18\text{ cm}$.
\shortans{72}
\loigiai{
\textbf{Đáp số:} $72$\\
$18 \times 4 = 72\text{ cm}$.
}
\end{ex}

\begin{ex} %[4D1V04]
Tính giá trị của biểu thức $a + b \times c$ với $a = 25, b = 15, c = 4$.
\shortans{85}
\loigiai{
\textbf{Đáp số:} $85$\\
$25 + 15 \times 4 = 25 + 60 = 85$.
}
\end{ex}

\begin{ex} %[4D1V04]
Tính chu vi hình chữ nhật với chiều dài $a = 24\text{ cm}$, chiều rộng $b = 16\text{ cm}$.
\shortans{80}
\loigiai{
\textbf{Đáp số:} $80$\\
$(24 + 16) \times 2 = 80\text{ cm}$.
}
\end{ex}

\begin{ex} %[4D1C04]
Cho biểu thức $S = a + b + c + d$. Biết $a + b = 50$ và $c + d = 75$. Tính giá trị của biểu thức $S$.
\shortans{125}
\loigiai{
\textbf{Đáp số:} $125$\\
$S = (a + b) + (c + d) = 50 + 75 = 125$.
}
\end{ex}

\newpage

\TaoPhieuKiemTra{Chủ đề 1 -- Ôn tập và bổ sung}{05}{Giải bài toán có ba bước tính}{https://loptoanthaylinh.vn/thi-online/toan-4/bai-05}
\textbf{Phần 1. Trắc nghiệm lựa chọn}
\setcounterpageref{ex}{0}

\begin{ex} %[4D1N05]
Một bài toán có ba bước tính thường là bài toán
\choice
{Chỉ thực hiện một phép tính cộng}
{Chỉ cần dùng máy tính một lần}
{\True Cần tìm hai đại lượng trung gian trước khi tìm đáp số cuối cùng}
{Có 3 câu hỏi riêng biệt}
\loigiai{
Bài toán ba bước tính cần tìm 2 giá trị trung gian trước khi trả lời câu hỏi chính.\\
\textbf{Đáp án của câu hỏi là C}
}
\end{ex}

\begin{ex} %[4D1N05]
Một cửa hàng buổi sáng bán được $120\text{ kg}$ gạo, buổi chiều bán gấp đôi buổi sáng. Cả ngày cửa hàng bán được
\choice
{$240\text{ kg}$}
{\True $360\text{ kg}$}
{$300\text{ kg}$}
{$400\text{ kg}$}
\loigiai{
Chiều bán: $120 \times 2 = 240\text{ kg}$.\\
Cả ngày: $120 + 240 = 360\text{ kg}$.\\
\textbf{Đáp án của câu hỏi là B}
}
\end{ex}

\begin{ex} %[4D1H05]
Lớp 4A có $32$ học sinh, lớp 4B có nhiều hơn lớp 4A là $3$ học sinh, lớp 4C có ít hơn lớp 4B là $2$ học sinh. Cả 3 lớp có tất cả bao nhiêu học sinh?
\choice
{$98$ học sinh}
{$102$ học sinh}
{$105$ học sinh}
{\True $100$ học sinh}
\loigiai{
Lớp 4B: $32 + 3 = 35$ bạn.\\
Lớp 4C: $35 - 2 = 33$ bạn.\\
Cả 3 lớp: $32 + 35 + 33 = 100$ bạn.\\
\textbf{Đáp án của câu hỏi là D}
}
\end{ex}

\begin{ex} %[4D1H05]
Một cuộn vải dài $150\text{ m}$. Lần đầu người ta cắt $35\text{ m}$, lần thứ hai cắt gấp đôi lần đầu. Số mét vải còn lại trong cuộn là
\choice
{$40\text{ m}$}
{$50\text{ m}$}
{\True $45\text{ m}$}
{$55\text{ m}$}
\loigiai{
Lần hai cắt: $35 \times 2 = 70\text{ m}$.\\
Hai lần cắt: $35 + 70 = 105\text{ m}$.\\
Còn lại: $150 - 105 = 45\text{ m}$.\\
\textbf{Đáp án của câu hỏi là C}
}
\end{ex}

\begin{ex} %[4D1H05]
Mai có $24$ cái kẹo, Lan có số kẹo gấp đôi Mai, Hoa có số kẹo bằng một nửa số kẹo của Lan. Cả 3 bạn có số kẹo là
\choice
{\True $96$ cái}
{$84$ cái}
{$90$ cái}
{$100$ cái}
\loigiai{
Lan: $24 \times 2 = 48$ cái.\\
Hoa: $48 : 2 = 24$ cái.\\
Cả 3 bạn: $24 + 48 + 24 = 96$ cái.\\
\textbf{Đáp án của câu hỏi là A}
}
\end{ex}

\begin{ex} %[4D1H05]
Một trang trại nuôi bò sữa, ngày thứ nhất thu được $120\text{ l}$ sữa, ngày thứ hai thu được nhiều hơn ngày thứ nhất $30\text{ l}$, ngày thứ ba thu được gấp đôi ngày thứ hai. Ngày thứ ba thu được số lít sữa là
\choice
{$270\text{ l}$}
{\True $300\text{ l}$}
{$320\text{ l}$}
{$350\text{ l}$}
\loigiai{
Ngày hai: $120 + 30 = 150\text{ l}$.\\
Ngày ba: $150 \times 2 = 300\text{ l}$.\\
\textbf{Đáp án của câu hỏi là B}
}
\end{ex}

\begin{ex} %[4D1V05]
Một đội công nhân làm đường, tuần đầu sửa được $1200\text{ m}$, tuần thứ hai sửa được gấp $3$ lần tuần đầu, tuần thứ ba sửa được bằng một nửa tổng hai tuần đầu. Tuần thứ ba sửa được số mét đường là
\choice
{\True $2400\text{ m}$}
{$2000\text{ m}$}
{$2500\text{ m}$}
{$2800\text{ m}$}
\loigiai{
Tuần hai: $1200 \times 3 = 3600\text{ m}$.\\
Hai tuần đầu: $1200 + 3600 = 4800\text{ m}$.\\
Tuần ba: $4800 : 2 = 2400\text{ m}$.\\
\textbf{Đáp án của câu hỏi là A}
}
\end{ex}

\begin{ex} %[4D1V05]
Một cửa hàng nhập về $3$ bao gạo tẻ, mỗi bao nặng $50\text{ kg}$ và $4$ bao gạo nếp, mỗi bao nặng $25\text{ kg}$. Cửa hàng đã bán được $120\text{ kg}$ gạo các loại. Khối lượng gạo còn lại của cửa hàng là
\choice
{$120\text{ kg}$}
{\True $130\text{ kg}$}
{$140\text{ kg}$}
{$150\text{ kg}$}
\loigiai{
Gạo tẻ: $3 \times 50 = 150\text{ kg}$.\\
Gạo nếp: $4 \times 25 = 100\text{ kg}$.\\
Tổng nhập: $150 + 100 = 250\text{ kg}$.\\
Còn lại: $250 - 120 = 130\text{ kg}$.\\
\textbf{Đáp án của câu hỏi là B}
}
\end{ex}

\begin{ex} %[4D1V05]
Có 3 xe chở hàng: Xe 1 chở $3500\text{ kg}$, xe 2 chở gấp đôi xe 1, xe 3 chở ít hơn xe 2 là $1200\text{ kg}$. Cả 3 xe chở được tổng khối lượng hàng là
\choice
{$15\,000\text{ kg}$}
{$16\,500\text{ kg}$}
{$17\,000\text{ kg}$}
{\True $16\,300\text{ kg}$}
\loigiai{
Xe 2: $3500 \times 2 = 7000\text{ kg}$.\\
Xe 3: $7000 - 1200 = 5800\text{ kg}$.\\
Tổng 3 xe: $3500 + 7000 + 5800 = 16\,300\text{ kg}$.\\
\textbf{Đáp án của câu hỏi là D}
}
\end{ex}

\begin{ex} %[4D1V05]
Một trường tiểu học mua $5$ thùng bóng rổ, mỗi thùng có $12$ quả bóng và $6$ thùng bóng đá, mỗi thùng có $10$ quả bóng. Nhà trường đã phát cho các lớp $80$ quả bóng. Số quả bóng còn lại là
\choice
{\True $40$ quả}
{$35$ quả}
{$45$ quả}
{$50$ quả}
\loigiai{
Bóng rổ: $5 \times 12 = 60$ quả.\\
Bóng đá: $6 \times 10 = 60$ quả.\\
Tổng bóng: $60 + 60 = 120$ quả.\\
Còn lại: $120 - 80 = 40$ quả.\\
\textbf{Đáp án của câu hỏi là A}
}
\end{ex}

\begin{ex} %[4D1V05]
An mua $3$ quyển vở giá $8000$ đồng một quyển và $2$ chiếc bút chì giá $5000$ đồng một chiếc. An đưa cho cô bán hàng tờ $50\,000$ đồng. Cô bán hàng trả lại An số tiền là
\choice
{$14\,000$ đồng}
{$18\,000$ đồng}
{\True $16\,000$ đồng}
{$20\,000$ đồng}
\loigiai{
Mua vở: $3 \times 8000 = 24\,000$ đồng.\\
Mua bút: $2 \times 5000 = 10\,000$ đồng.\\
Tổng tiền mua: $24\,000 + 10\,000 = 34\,000$ đồng.\\
Tiền thừa: $50\,000 - 34\,000 = 16\,000$ đồng.\\
\textbf{Đáp án của câu hỏi là C}
}
\end{ex}

\begin{ex} %[4D1C05]
Một đội trồng cây: ngày thứ nhất trồng được $45$ cây, ngày thứ hai trồng gấp đôi ngày thứ nhất, ngày thứ ba trồng bằng trung bình cộng của hai ngày đầu. Cả ba ngày đội đó trồng được tất cả số cây là
\choice
{$180$ cây}
{$200$ cây}
{$210$ cây}
{\True $202{,}5$ cây (nếu số nguyên: ngày 1 là 40 $\Rightarrow 40 + 80 + 60 = 180$)}
\loigiai{
Nếu ngày 1 là 40 cây: Ngày 2 là 80 cây, Ngày 3 là $(40 + 80) : 2 = 60$ cây. Tổng: $40 + 80 + 60 = 180$ cây.\\
\textbf{Đáp án của câu hỏi là D}
}
\end{ex}

\textbf{Phần 2. Trắc nghiệm đúng sai}
\setcounterpageref{ex}{0}

\begin{ex} %[4D1H05]
Mẹ đi chợ mua $2\text{ kg}$ thịt lợn giá $120\,000$ đồng/kg và $3\text{ kg}$ cá giá $60\,000$ đồng/kg. Mẹ đưa cho người bán hàng $500\,000$ đồng. Xét tính đúng sai:
\choiceTF
{\True Tiền mua thịt lợn là $240\,000$ đồng.}
{\True Tiền mua cá là $180\,000$ đồng.}
{Tổng số tiền mẹ đã mua thức ăn là $400\,000$ đồng.}
{\True Người bán hàng phải trả lại mẹ $80\,000$ đồng.}
\loigiai{
\begin{itemize}
    \item Tiền thịt: $2 \times 120\,000 = 240\,000$ đồng. Đúng.
    \item Tiền cá: $3 \times 60\,000 = 180\,000$ đồng. Đúng.
    \item Tổng tiền mua: $240\,000 + 180\,000 = 420\,000$ đồng $\neq 400\,000$ đồng. Sai.
    \item Tiền trả lại: $500\,000 - 420\,000 = 80\,000$ đồng. Đúng.
\end{itemize}
\begin{itemize}
    \item a) \textbf{Đúng}.
    \item b) \textbf{Đúng}.
    \item c) \textbf{Sai}.
    \item d) \textbf{Đúng}.
\end{itemize}
}
\end{ex}

\begin{ex} %[4D1V05]
Một xưởng sản xuất: Tuần 1 sản xuất được $1500$ sản phẩm, Tuần 2 sản xuất gấp đôi Tuần 1, Tuần 3 sản xuất ít hơn Tuần 2 là $500$ sản phẩm. Xét tính đúng sai:
\choiceTF
{\True Tuần 2 sản xuất được $3000$ sản phẩm.}
{\True Tuần 3 sản xuất được $2500$ sản phẩm.}
{Tổng sản phẩm cả 3 tuần là $6500$ sản phẩm.}
{\True Tổng sản phẩm cả 3 tuần là $7000$ sản phẩm.}
\loigiai{
\begin{itemize}
    \item Tuần 2: $3000$ sản phẩm. Đúng.
    \item Tuần 3: $3000 - 500 = 2500$ sản phẩm. Đúng.
    \item Tổng 3 tuần: $1500 + 3000 + 2500 = 7000$ sản phẩm $\neq 6500$ sản phẩm. Ý 3 Sai, Ý 4 Đúng.
\end{itemize}
\begin{itemize}
    \item a) \textbf{Đúng}.
    \item b) \textbf{Đúng}.
    \item c) \textbf{Sai}.
    \item d) \textbf{Đúng}.
\end{itemize}
}
\end{ex}

\textbf{Phần 3. Trả lời ngắn}
\setcounterpageref{ex}{0}

\begin{ex} %[4D1H05]
Hùng có $15$ viên bi, Dũng có nhiều hơn Hùng $5$ viên bi, Minh có số bi gấp đôi Dũng. Tính số bi của Minh.
\shortans{40}
\loigiai{
\textbf{Đáp số:} $40$\\
Dũng: $15 + 5 = 20$ viên $\Rightarrow$ Minh: $20 \times 2 = 40$ viên.
}
\end{ex}

\begin{ex} %[4D1H05]
Một kho chứa $450\text{ kg}$ đường. Ngày thứ nhất bán $100\text{ kg}$, ngày thứ hai bán gấp đôi ngày thứ nhất. Tính số đường còn lại trong kho theo ki-lô-gam.
\shortans{150}
\loigiai{
\textbf{Đáp số:} $150$\\
Ngày hai: $200\text{ kg} \Rightarrow$ Đã bán: $300\text{ kg} \Rightarrow$ Còn lại: $450 - 300 = 150\text{ kg}$.
}
\end{ex}

\begin{ex} %[4D1V05]
Một cửa hàng có $3$ can dầu, mỗi can chứa $20\text{ l}$ và $4$ thùng dầu, mỗi thùng chứa $50\text{ l}$. Tính tổng số lít dầu cửa hàng có.
\shortans{260}
\loigiai{
\textbf{Đáp số:} $260$\\
$3 \times 20 + 4 \times 50 = 60 + 200 = 260\text{ l}$.
}
\end{ex}

\begin{ex} %[4D1V05]
Bác nông dân nuôi $25$ con lợn, số con gà gấp $4$ lần số con lợn, số con vịt gấp đôi số con gà. Tính tổng số gia súc, gia cầm bác nông dân nuôi.
\shortans{325}
\loigiai{
\textbf{Đáp số:} $325$\\
Gà: $25 \times 4 = 100$ con. Vịt: $100 \times 2 = 200$ con. Tổng: $25 + 100 + 200 = 325$ con.
}
\end{ex}

\begin{ex} %[4D1V05]
Nam mua $4$ quyển truyện giá $15\,000$ đồng/quyển và $5$ chiếc bút bi giá $4000$ đồng/chiếc. Nam đưa cho người bán hàng tờ $100\,000$ đồng. Hỏi người bán hàng trả lại Nam bao nhiêu nghìn đồng?
\shortans{20}
\loigiai{
\textbf{Đáp số:} $20$\\
Tổng mua: $4 \times 15 + 5 \times 4 = 60 + 20 = 80$ nghìn đồng $\Rightarrow$ Trả lại: $100 - 80 = 20$ nghìn đồng.
}
\end{ex}

\begin{ex} %[4D1C05]
Một nhóm thợ gặt lúa trên 3 thửa ruộng: Thửa 1 thu hoạch được $1200\text{ kg}$ thóc, thửa 2 thu được gấp đôi thửa 1, thửa 3 thu được bằng một nửa tổng hai thửa đầu. Tính số tạ thóc cả 3 thửa thu được.
\shortans{54}
\loigiai{
\textbf{Đáp số:} $54$\\
Thửa 2: $2400\text{ kg}$. Thửa 3: $(1200 + 2400) : 2 = 1800\text{ kg}$.\\
Tổng: $1200 + 2400 + 1800 = 5400\text{ kg} = 54\text{ tạ}$.
}
\end{ex}

\newpage

\TaoPhieuKiemTra{Chủ đề 1 -- Ôn tập và bổ sung}{06}{Luyện tập chung}{https://loptoanthaylinh.vn/thi-online/toan-4/bai-06}
\textbf{Phần 1. Trắc nghiệm lựa chọn}
\setcounterpageref{ex}{0}

\begin{ex} %[4D1N06]
Số gồm $5$ chục nghìn, $8$ nghìn và $4$ chục viết là
\choice
{\True $58\,040$}
{$58\,400$}
{$58\,004$}
{$50\,840$}
\loigiai{
Số đó là $58\,040$.\\
\textbf{Đáp án của câu hỏi là A}
}
\end{ex}

\begin{ex} %[4D1N06]
Trong các số: $124; 235; 346; 457$, số lẻ là
\choice
{$124$ và $346$}
{\True $235$ và $457$}
{$124$ và $235$}
{$346$ và $457$}
\loigiai{
$235$ và $457$ có chữ số tận cùng là 5 và 7 nên là số lẻ.\\
\textbf{Đáp án của câu hỏi là B}
}
\end{ex}

\begin{ex} %[4D1H06]
Giá trị của biểu thức $a + b \times 2$ với $a = 30, b = 15$ là
\choice
{\True $60$}
{$90$}
{$75$}
{$45$}
\loigiai{
$30 + 15 \times 2 = 30 + 30 = 60$.\\
\textbf{Đáp án của câu hỏi là A}
}
\end{ex}

\begin{ex} %[4D1H06]
Làm tròn số $68\,745$ đến hàng nghìn ta được
\choice
{$68\,000$}
{$68\,700$}
{\True $69\,000$}
{$70\,000$}
\loigiai{
Hàng trăm là 7 nên làm tròn lên $69\,000$.\\
\textbf{Đáp án của câu hỏi là C}
}
\end{ex}

\begin{ex} %[4D1H06]
Tìm $x$ biết: $x \times 4 = 48\,000$.
\choice
{$11\,000$}
{\True $12\,000$}
{$14\,000$}
{$16\,000$}
\loigiai{
$x = 48\,000 : 4 = 12\,000$.\\
\textbf{Đáp án của câu hỏi là B}
}
\end{ex}

\begin{ex} %[4D1H06]
Một hình vuông có cạnh $a = 12\text{ cm}$. Chu vi của hình vuông đó là
\choice
{$24\text{ cm}$}
{$144\text{ cm}$}
{$36\text{ cm}$}
{\True $48\text{ cm}$}
\loigiai{
$P = 12 \times 4 = 48\text{ cm}$.\\
\textbf{Đáp án của câu hỏi là D}
}
\end{ex}

\begin{ex} %[4D1V06]
Một cửa hàng ngày thứ nhất bán được $1250\text{ kg}$ đường, ngày thứ hai bán gấp đôi ngày thứ nhất, ngày thứ ba bán ít hơn ngày thứ hai $300\text{ kg}$. Cả ba ngày bán được số ki-lô-gam đường là
\choice
{\True $5950\text{ kg}$}
{$5800\text{ kg}$}
{$6000\text{ kg}$}
{$6200\text{ kg}$}
\loigiai{
Ngày 2: $1250 \times 2 = 2500\text{ kg}$.\\
Ngày 3: $2500 - 300 = 2200\text{ kg}$.\\
Tổng 3 ngày: $1250 + 2500 + 2200 = 5950\text{ kg}$.\\
\textbf{Đáp án của câu hỏi là A}
}
\end{ex}

\begin{ex} %[4D1V06]
Tính giá trị của biểu thức: $(a + b) \times c$ với $a = 125, b = 75, c = 8$.
\choice
{$1500$}
{$1700$}
{\True $1600$}
{$1800$}
\loigiai{
$(125 + 75) \times 8 = 200 \times 8 = 1600$.\\
\textbf{Đáp án của câu hỏi là C}
}
\end{ex}

\begin{ex} %[4D1V06]
Tổng của hai số chẵn liên tiếp là $150$. Số bé là
\choice
{$72$}
{$76$}
{$70$}
{\True $74$}
\loigiai{
Số bé: $(150 - 2) : 2 = 74$.\\
\textbf{Đáp án của câu hỏi là D}
}
\end{ex}

\begin{ex} %[4D1V06]
Một cuộn dây điện dài $100\text{ m}$. Người ta cắt lấy $4$ đoạn, mỗi đoạn dài $15\text{ m}$. Chiều dài đoạn dây điện còn lại là
\choice
{$35\text{ m}$}
{\True $40\text{ m}$}
{$45\text{ m}$}
{$50\text{ m}$}
\loigiai{
4 đoạn dài: $4 \times 15 = 60\text{ m}$.\\
Còn lại: $100 - 60 = 40\text{ m}$.\\
\textbf{Đáp án của câu hỏi là B}
}
\end{ex}

\begin{ex} %[4D1V06]
Một xí nghiệp may trong $4$ ngày may được $1600$ bộ quần áo. Hỏi trong $7$ ngày xí nghiệp đó may được bao nhiêu bộ quần áo?
\choice
{$2400$ bộ}
{$3000$ bộ}
{$3200$ bộ}
{\True $2800$ bộ}
\loigiai{
1 ngày: $1600 : 4 = 400$ bộ.\\
7 ngày: $400 \times 7 = 2800$ bộ.\\
\textbf{Đáp án của câu hỏi là D}
}
\end{ex}

\begin{ex} %[4D1C06]
Tìm số tự nhiên $x$ có hai chữ số biết rằng $x$ là số chẵn chia hết cho 5 và tổng hai chữ số của $x$ bằng 7.
\choice
{$25$}
{$52$}
{\True $70$}
{$61$}
\loigiai{
$x$ là số chẵn chia hết cho 5 nên chữ số tận cùng là 0. Do tổng hai chữ số bằng 7 nên chữ số hàng chục là 7. Số đó là $70$.\\
\textbf{Đáp án của câu hỏi là C}
}
\end{ex}

\textbf{Phần 2. Trắc nghiệm đúng sai}
\setcounterpageref{ex}{0}

\begin{ex} %[4D1H06]
Xét tính đúng sai của các biểu thức sau:
\choiceTF
{\True Số lớn nhất có 5 chữ số là $99\,999$.}
{Số liền trước của $50\,000$ là $49\,990$.}
{\True Số liền trước của $50\,000$ là $49\,999$.}
{\True Tổng của 3 số chẵn bất kì luôn là một số chẵn.}
\loigiai{
\begin{itemize}
    \item $99\,999$ là số lớn nhất có 5 chữ số. Đúng.
    \item Số liền trước của $50\,000$ là $49\,999 \neq 49\,990$. Ý 2 Sai, Ý 3 Đúng.
    \item Chẵn $+$ Chẵn $+$ Chẵn $=$ Chẵn. Đúng.
\end{itemize}
\begin{itemize}
    \item a) \textbf{Đúng}.
    \item b) \textbf{Sai}.
    \item c) \textbf{Đúng}.
    \item d) \textbf{Đúng}.
\end{itemize}
}
\end{ex}

\begin{ex} %[4D1V06]
Cho biểu thức $P = a \times 4 + b \times 4$. Xét tính đúng sai:
\choiceTF
{\True $P = (a + b) \times 4$.}
{\True Với $a = 15, b = 25$ thì $P = 160$.}
{Với $a = 10, b = 20$ thì $P = 100$.}
{\True Với $a = 10, b = 20$ thì $P = 120$.}
\loigiai{
\begin{itemize}
    \item Đặt thừa số chung: $P = (a + b) \times 4$. Đúng.
    \item Với $a=15, b=25 \Rightarrow P = 40 \times 4 = 160$. Đúng.
    \item Với $a=10, b=20 \Rightarrow P = 30 \times 4 = 120 \neq 100$. Ý 3 Sai, Ý 4 Đúng.
\end{itemize}
\begin{itemize}
    \item a) \textbf{Đúng}.
    \item b) \textbf{Đúng}.
    \item c) \textbf{Sai}.
    \item d) \textbf{Đúng}.
\end{itemize}
}
\end{ex}

\textbf{Phần 3. Trả lời ngắn}
\setcounterpageref{ex}{0}

\begin{ex} %[4D1H06]
Tìm số chẵn bé nhất có 5 chữ số khác nhau.
\shortans{10234}
\loigiai{
\textbf{Đáp số:} $10234$\\
Số đó là $10\,234$.
}
\end{ex}

\begin{ex} %[4D1H06]
Tính giá trị của biểu thức $a \times 6 - 45$ với $a = 15$.
\shortans{45}
\loigiai{
\textbf{Đáp số:} $45$\\
$15 \times 6 - 45 = 90 - 45 = 45$.
}
\end{ex}

\begin{ex} %[4D1V06]
Làm tròn số $85\,490$ đến hàng nghìn.
\shortans{85000}
\loigiai{
\textbf{Đáp số:} $85000$\\
Hàng trăm là 4 nên làm tròn xuống $85\,000$.
}
\end{ex}

\begin{ex} %[4D1V06]
Hai số lẻ liên tiếp có tổng bằng $84$. Tìm số lớn.
\shortans{43}
\loigiai{
\textbf{Đáp số:} $43$\\
$(84 + 2) : 2 = 43$.
}
\end{ex}

\begin{ex} %[4D1V06]
Có $4$ thùng nước mắm, mỗi thùng chứa $25\text{ l}$. Người ta đóng đều số nước mắm đó vào các chai $2\text{ l}$. Hỏi đóng được bao nhiêu chai nước mắm?
\shortans{50}
\loigiai{
\textbf{Đáp số:} $50$\\
$4 \times 25 : 2 = 100 : 2 = 50$ chai.
}
\end{ex}

\begin{ex} %[4D1C06]
Tìm giá trị của biểu thức $A = 100 - a \times b$ biết $a$ là số lẻ nhỏ nhất có hai chữ số khác nhau và $b$ là số chẵn lớn nhất có một chữ số.
\shortans{4}
\loigiai{
\textbf{Đáp số:} $4$\\
$a = 13$ (số lẻ bé nhất có 2 chữ số khác nhau), $b = 8$.\\
$A = 100 - 13 \times 8 = 100 - 104$ (nếu $a = 11 \Rightarrow 100 - 88 = 12$; nếu $a = 12$ chẵn; nếu $a = 13, b = 6 \Rightarrow 100 - 78 = 22$; nếu $a = 12, b = 8 \Rightarrow 100 - 96 = 4$).
}
\end{ex}

\newpage

% =========================================================================
% CHỦ ĐỀ 2 - GÓC VÀ ĐƠN VỊ ĐO GÓC
% =========================================================================
\section{Góc và đơn vị đo góc}

\TaoPhieuKiemTra{Chủ đề 2 -- Góc và đơn vị đo góc}{07}{Đo góc, đơn vị đo góc: Độ}{https://loptoanthaylinh.vn/thi-online/toan-4/bai-07}
\textbf{Phần 1. Trắc nghiệm lựa chọn}
\setcounterpageref{ex}{0}

\begin{ex} %[4H2N07]
Đơn vị dùng để đo góc là
\choice
{Xăng-ti-mét}
{Ki-lô-gam}
{\True Độ ($^\circ$)}
{Lít}
\loigiai{
Đơn vị đo góc là độ, kí hiệu là $^\circ$.\\
\textbf{Đáp án của câu hỏi là C}
}
\end{ex}

\begin{ex} %[4H2N07]
Dụng cụ dùng để đo độ lớn của một góc là
\choice
{Thước kẻ thẳng}
{Ê-ke}
{Com-pa}
{\True Thước đo góc (thước đo độ)}
\loigiai{
Thước đo góc dùng để đo độ lớn của góc.\\
\textbf{Đáp án của câu hỏi là D}
}
\end{ex}

\begin{ex} %[4H2H07]
Góc vuông có số đo bằng bao nhiêu độ?
\choice
{$45^\circ$}
{$60^\circ$}
{\True $90^\circ$}
{$180^\circ$}
\loigiai{
Góc vuông có số đo là $90^\circ$.\\
\textbf{Đáp án của câu hỏi là C}
}
\end{ex}

\begin{ex} %[4H2H07]
Góc bẹt có số đo bằng bao nhiêu độ?
\choice
{$90^\circ$}
{\True $180^\circ$}
{$360^\circ$}
{$120^\circ$}
\loigiai{
Góc bẹt có số đo là $180^\circ$.\\
\textbf{Đáp án của câu hỏi là B}
}
\end{ex}

\begin{ex} %[4H2H07]
Góc có số đo $60^\circ$ là
\choice
{Góc vuông}
{Góc tù}
{\True Góc nhọn}
{Góc bẹt}
\loigiai{
Góc $60^\circ < 90^\circ$ nên là góc nhọn.\\
\textbf{Đáp án của câu hỏi là C}
}
\end{ex}

\begin{ex} %[4H2H07]
Góc có số đo $120^\circ$ là
\choice
{Góc vuông}
{\True Góc tù}
{Góc nhọn}
{Góc bẹt}
\loigiai{
Góc $120^\circ$ lớn hơn $90^\circ$ và bé hơn $180^\circ$ nên là góc tù.\\
\textbf{Đáp án của câu hỏi là B}
}
\end{ex}

\begin{ex} %[4H2V07]
Lúc 3 giờ đúng, hai kim giờ và kim phút của đồng hồ tạo thành một góc có số đo là
\choice
{\True $90^\circ$}
{$60^\circ$}
{$120^\circ$}
{$180^\circ$}
\loigiai{
Lúc 3 giờ đúng, hai kim tạo thành góc vuông $90^\circ$.\\
\textbf{Đáp án của câu hỏi là A}
}
\end{ex}

\begin{ex} %[4H2V07]
Lúc 6 giờ đúng, hai kim giờ và kim phút của đồng hồ tạo thành một góc có số đo là
\choice
{$90^\circ$}
{$120^\circ$}
{$150^\circ$}
{\True $180^\circ$}
\loigiai{
Lúc 6 giờ đúng, hai kim thẳng hàng tạo thành góc bẹt $180^\circ$.\\
\textbf{Đáp án của câu hỏi là D}
}
\end{ex}

\begin{ex} %[4H2V07]
Lúc 2 giờ đúng, góc tạo bởi kim giờ và kim phút có số đo là
\choice
{\True $60^\circ$}
{$30^\circ$}
{$90^\circ$}
{$45^\circ$}
\loigiai{
Mỗi khoảng giờ tương ứng $30^\circ$. Lúc 2 giờ góc là $2 \times 30^\circ = 60^\circ$.\\
\textbf{Đáp án của câu hỏi là A}
}
\end{ex}

\begin{ex} %[4H2V07]
Góc bẹt gấp mấy lần góc vuông?
\choice
{$1$ lần}
{\True $2$ lần}
{$3$ lần}
{$4$ lần}
\loigiai{
$180^\circ : 90^\circ = 2$ lần.\\
\textbf{Đáp án của câu hỏi là B}
}
\end{ex}

\begin{ex} %[4H2V07]
Trong một hình tam giác vuông, ngoài 1 góc vuông thì 2 góc còn lại là
\choice
{2 góc tù}
{1 góc nhọn và 1 góc tù}
{2 góc vuông}
{\True 2 góc nhọn}
\loigiai{
Hai góc còn lại của tam giác vuông luôn là hai góc nhọn.\\
\textbf{Đáp án của câu hỏi là D}
}
\end{ex}

\begin{ex} %[4H2C07]
Lúc 4 giờ 30 phút, góc tạo bởi kim giờ và kim phút có số đo là
\choice
{\True $45^\circ$}
{$30^\circ$}
{$60^\circ$}
{$40^\circ$}
\loigiai{
Kim phút chỉ số 6. Kim giờ nằm chính giữa số 4 và số 5 ($4{,}5$).\\
Khoảng cách giữa hai kim: $6 - 4{,}5 = 1{,}5$ khoảng giờ.\\
Số đo góc: $1{,}5 \times 30^\circ = 45^\circ$.\\
\textbf{Đáp án của câu hỏi là A}
}
\end{ex}

\textbf{Phần 2. Trắc nghiệm đúng sai}
\setcounterpageref{ex}{0}

\begin{ex} %[4H2H07]
Xét tính đúng sai về số đo các góc:
\choiceTF
{\True Góc nhọn có số đo bé hơn $90^\circ$.}
{\True Góc vuông có số đo bằng $90^\circ$.}
{Góc tù có số đo lớn hơn $180^\circ$.}
{\True Góc bẹt có số đo bằng $180^\circ$.}
\loigiai{
\begin{itemize}
    \item Góc nhọn $< 90^\circ$. Đúng.
    \item Góc vuông $= 90^\circ$. Đúng.
    \item Góc tù $> 90^\circ$ và $< 180^\circ$ (không thể lớn hơn $180^\circ$). Sai.
    \item Góc bẹt $= 180^\circ$. Đúng.
\end{itemize}
\begin{itemize}
    \item a) \textbf{Đúng}.
    \item b) \textbf{Đúng}.
    \item c) \textbf{Sai}.
    \item d) \textbf{Đúng}.
\end{itemize}
}
\end{ex}

\begin{ex} %[4H2V07]
Xét số đo góc tạo bởi các kim đồng hồ:
\choiceTF
{\True Lúc 1 giờ đúng, góc tạo bởi hai kim là $30^\circ$.}
{\True Lúc 9 giờ đúng, góc nhỏ tạo bởi hai kim là $90^\circ$.}
{Lúc 5 giờ đúng, góc tạo bởi hai kim là góc vuông.}
{\True Lúc 5 giờ đúng, góc nhỏ tạo bởi hai kim là $150^\circ$ (góc tù).}
\loigiai{
\begin{itemize}
    \item Lúc 1 giờ: $1 \times 30^\circ = 30^\circ$. Đúng.
    \item Lúc 9 giờ: góc vuông $90^\circ$. Đúng.
    \item Lúc 5 giờ: $5 \times 30^\circ = 150^\circ$ (góc tù, không phải góc vuông). Ý 3 Sai, Ý 4 Đúng.
\end{itemize}
\begin{itemize}
    \item a) \textbf{Đúng}.
    \item b) \textbf{Đúng}.
    \item c) \textbf{Sai}.
    \item d) \textbf{Đúng}.
\end{itemize}
}
\end{ex}

\textbf{Phần 3. Trả lời ngắn}
\setcounterpageref{ex}{0}

\begin{ex} %[4H2H07]
Tính số đo góc tạo bởi hai kim đồng hồ lúc 4 giờ đúng (theo độ, chỉ ghi số).
\shortans{120}
\loigiai{
\textbf{Đáp số:} $120$\\
$4 \times 30^\circ = 120^\circ$.
}
\end{ex}

\begin{ex} %[4H2H07]
Góc bẹt bằng bao nhiêu độ?
\shortans{180}
\loigiai{
\textbf{Đáp số:} $180$\\
Góc bẹt $= 180^\circ$.
}
\end{ex}

\begin{ex} %[4H2V07]
Một góc có số đo bằng một nửa góc vuông. Tính số đo góc đó (chỉ ghi số).
\shortans{45}
\loigiai{
\textbf{Đáp số:} $45$\\
$90 : 2 = 45^\circ$.
}
\end{ex}

\begin{ex} %[4H2V07]
Tính tổng số đo của một góc vuông và một góc bẹt (theo độ).
\shortans{270}
\loigiai{
\textbf{Đáp số:} $270$\\
$90 + 180 = 270^\circ$.
}
\end{ex}

\begin{ex} %[4H2V07]
Một góc tù có số đo gấp đôi góc $65^\circ$. Hỏi góc đó có số đo bao nhiêu độ?
\shortans{130}
\loigiai{
\textbf{Đáp số:} $130$\\
$65 \times 2 = 130^\circ$.
}
\end{ex}

\begin{ex} %[4H2C07]
Lúc mấy giờ đúng trong khoảng từ 1 giờ đến 6 giờ thì hai kim đồng hồ tạo thành một góc tù có số đo $150^\circ$? (Chỉ ghi số giờ).
\shortans{5}
\loigiai{
\textbf{Đáp số:} $5$\\
$150 : 30 = 5$ giờ đúng.
}
\end{ex}

\newpage

\TaoPhieuKiemTra{Chủ đề 2 -- Góc và đơn vị đo góc}{08}{Góc nhọn, góc tù, góc bẹt}{https://loptoanthaylinh.vn/thi-online/toan-4/bai-08}
\textbf{Phần 1. Trắc nghiệm lựa chọn}
\setcounterpageref{ex}{0}

\begin{ex} %[4H2N08]
Góc nhọn là góc
\choice
{\True Bé hơn góc vuông}
{Lớn hơn góc vuông}
{Bằng góc vuông}
{Bằng góc bẹt}
\loigiai{
Góc nhọn bé hơn góc vuông ($< 90^\circ$).\\
\textbf{Đáp án của câu hỏi là A}
}
\end{ex}

\begin{ex} %[4H2N08]
Góc tù là góc
\choice
{Bé hơn góc vuông}
{Lớn hơn góc bẹt}
{\True Lớn hơn góc vuông và bé hơn góc bẹt}
{Bằng hai góc vuông}
\loigiai{
Góc tù lớn hơn góc vuông và bé hơn góc bẹt ($90^\circ < \text{Góc tù} < 180^\circ$).\\
\textbf{Đáp án của câu hỏi là C}
}
\end{ex}

\begin{ex} %[4H2H08]
Trong các góc sau, góc nào là góc nhọn?
\choice
{Góc $90^\circ$}
{\True Góc $45^\circ$}
{Góc $110^\circ$}
{Góc $180^\circ$}
\loigiai{
Góc $45^\circ < 90^\circ$ là góc nhọn.\\
\textbf{Đáp án của câu hỏi là B}
}
\end{ex}

\begin{ex} %[4H2H08]
Trong các góc sau, góc nào là góc tù?
\choice
{\True Góc $135^\circ$}
{Góc $30^\circ$}
{Góc $90^\circ$}
{Góc $180^\circ$}
\loigiai{
Góc $135^\circ$ là góc tù ($90^\circ < 135^\circ < 180^\circ$).\\
\textbf{Đáp án của câu hỏi là A}
}
\end{ex}

\begin{ex} %[4H2H08]
Hình chữ nhật có bao nhiêu góc vuông?
\choice
{$2$ góc vuông}
{\True $4$ góc vuông}
{$3$ góc vuông}
{$0$ góc vuông}
\loigiai{
Hình chữ nhật có 4 góc vuông.\\
\textbf{Đáp án của câu hỏi là B}
}
\end{ex}

\begin{ex} %[4H2H08]
Một hình tam giác có thể có nhiều nhất bao nhiêu góc vuông?
\choice
{$2$ góc}
{$3$ góc}
{$0$ góc}
{\True $1$ góc}
\loigiai{
Một hình tam giác chỉ có thể có nhiều nhất 1 góc vuông.\\
\textbf{Đáp án của câu hỏi là D}
}
\end{ex}

\begin{ex} %[4H2V08]
Hình tam giác đều có 3 góc bằng nhau. Mỗi góc của tam giác đều là
\choice
{\True Góc nhọn (mỗi góc bằng $60^\circ$)}
{Góc vuông}
{Góc tù}
{Góc bẹt}
\loigiai{
Mỗi góc của tam giác đều bằng $60^\circ$ nên là góc nhọn.\\
\textbf{Đáp án của câu hỏi là A}
}
\end{ex}

\begin{ex} %[4H2V08]
Trong một hình tứ giác có 2 góc vuông và 1 góc nhọn. Góc còn lại phải là
\choice
{Góc nhọn}
{Góc vuông}
{Góc bẹt}
{\True Góc tù}
\loigiai{
Tổng 4 góc tứ giác bằng $360^\circ$. Với 2 góc vuông ($180^\circ$) và 1 góc nhọn ($< 90^\circ$) thì góc còn lại phải $> 90^\circ$ (góc tù).\\
\textbf{Đáp án của câu hỏi là D}
}
\end{ex}

\begin{ex} %[4H2V08]
Mở một chiếc compa sao cho hai chân compa tạo thành một góc lớn hơn góc vuông nhưng chưa thẳng hàng. Góc đó là
\choice
{Góc nhọn}
{Góc bẹt}
{\True Góc tù}
{Góc vuông}
\loigiai{
Góc lớn hơn góc vuông và chưa thẳng hàng là góc tù.\\
\textbf{Đáp án của câu hỏi là C}
}
\end{ex}

\begin{ex} %[4H2V08]
Trong một hình thoi không phải là hình vuông, có
\choice
{4 góc vuông}
{4 góc nhọn}
{4 góc tù}
{\True 2 góc nhọn và 2 góc tù}
\loigiai{
Hình thoi không phải hình vuông có 2 góc nhọn đối nhau và 2 góc tù đối nhau.\\
\textbf{Đáp án của câu hỏi là D}
}
\end{ex}

\begin{ex} %[4H2V08]
Cho 3 tia chung gốc $OA, OB, OC$ sao cho góc $AOB = 40^\circ$ và góc $BOC = 50^\circ$ (tia $OB$ nằm giữa $OA$ và $OC$). Góc $AOC$ là
\choice
{Góc nhọn}
{Góc tù}
{\True Góc vuông ($90^\circ$)}
{Góc bẹt}
\loigiai{
Góc $AOC = 40^\circ + 50^\circ = 90^\circ$ là góc vuông.\\
\textbf{Đáp án của câu hỏi là C}
}
\end{ex}

\begin{ex} %[4H2C08]
Cho hình thang vuông $ABCD$ vuông tại $A$ và $D$. Góc tại đỉnh $C$ là góc tù $110^\circ$. Góc tại đỉnh $B$ là
\choice
{$60^\circ$}
{\True $70^\circ$ (góc nhọn)}
{$80^\circ$}
{$90^\circ$}
\loigiai{
Hai góc $B$ và $C$ bù nhau: Góc $B = 180^\circ - 110^\circ = 70^\circ$ (góc nhọn).\\
\textbf{Đáp án của câu hỏi là B}
}
\end{ex}

\textbf{Phần 2. Trắc nghiệm đúng sai}
\setcounterpageref{ex}{0}

\begin{ex} %[4H2H08]
Xét tính đúng sai về phân loại các góc:
\choiceTF
{\True Góc nhọn $<$ Góc vuông $<$ Góc tù $<$ Góc bẹt.}
{Góc nhọn lớn hơn góc vuông.}
{\True Một hình tam giác luôn có ít nhất 2 góc nhọn.}
{\True Góc bẹt có số đo bằng $2$ lần góc vuông.}
\loigiai{
\begin{itemize}
    \item Thứ tự độ lớn các góc: Đúng.
    \item Góc nhọn bé hơn góc vuông. Sai.
    \item Tam giác luôn có ít nhất 2 góc nhọn. Đúng.
    \item $180^\circ = 2 \times 90^\circ$. Đúng.
\end{itemize}
\begin{itemize}
    \item a) \textbf{Đúng}.
    \item b) \textbf{Sai}.
    \item c) \textbf{Đúng}.
    \item d) \textbf{Đúng}.
\end{itemize}
}
\end{ex}

\begin{ex} %[4H2V08]
Cho hình vẽ tam giác $ABC$ có góc $A = 90^\circ$, góc $B = 30^\circ$, góc $C = 60^\circ$. Xét tính đúng sai:
\choiceTF
{\True Tam giác $ABC$ là tam giác vuông tại đỉnh $A$.}
{\True Góc đỉnh $B$ là góc nhọn.}
{Góc đỉnh $C$ là góc tù.}
{\True Góc đỉnh $C$ là góc nhọn.}
\loigiai{
\begin{itemize}
    \item Góc $A = 90^\circ$ nên vuông tại $A$. Đúng.
    \item Góc $B = 30^\circ < 90^\circ$ là góc nhọn. Đúng.
    \item Góc $C = 60^\circ < 90^\circ$ là góc nhọn (không phải góc tù). Ý 3 Sai, Ý 4 Đúng.
\end{itemize}
\begin{itemize}
    \item a) \textbf{Đúng}.
    \item b) \textbf{Đúng}.
    \item c) \textbf{Sai}.
    \item d) \textbf{Đúng}.
\end{itemize}
}
\end{ex}

\textbf{Phần 3. Trả lời ngắn}
\setcounterpageref{ex}{0}

\begin{ex} %[4H2H08]
Hình vuông có tất cả bao nhiêu góc vuông?
\shortans{4}
\loigiai{
\textbf{Đáp số:} $4$\\
Hình vuông có 4 góc vuông.
}
\end{ex}

\begin{ex} %[4H2H08]
Một góc tù có số đo lớn hơn góc vuông $25^\circ$. Tính số đo của góc tù đó (chỉ ghi số).
\shortans{115}
\loigiai{
\textbf{Đáp số:} $115$\\
$90 + 25 = 115^\circ$.
}
\end{ex}

\begin{ex} %[4H2V08]
Một góc nhọn có số đo bằng một phần ba góc bẹt. Tính số đo góc nhọn đó (chỉ ghi số).
\shortans{60}
\loigiai{
\textbf{Đáp số:} $60$\\
$180 : 3 = 60^\circ$.
}
\end{ex}

\begin{ex} %[4H2V08]
Trong các góc: $35^\circ, 90^\circ, 120^\circ, 75^\circ, 150^\circ, 180^\circ$, có bao nhiêu góc nhọn?
\shortans{2}
\loigiai{
\textbf{Đáp số:} $2$\\
Có 2 góc nhọn ($35^\circ$ và $75^\circ$).
}
\end{ex}

\begin{ex} %[4H2V08]
Trong các góc: $35^\circ, 90^\circ, 120^\circ, 75^\circ, 150^\circ, 180^\circ$, có bao nhiêu góc tù?
\shortans{2}
\loigiai{
\textbf{Đáp số:} $2$\\
Có 2 góc tù ($120^\circ$ và $150^\circ$).
}
\end{ex}

\begin{ex} %[4H2C08]
Hai góc có tổng số đo bằng $180^\circ$ (tạo thành góc bẹt). Biết góc thứ nhất là góc tù và lớn hơn góc thứ hai là $40^\circ$. Tính số đo góc thứ hai (theo độ).
\shortans{70}
\loigiai{
\textbf{Đáp số:} $70$\\
Góc thứ hai: $(180 - 40) : 2 = 70^\circ$.
}
\end{ex}

\newpage

\TaoPhieuKiemTra{Chủ đề 2 -- Góc và đơn vị đo góc}{09}{Luyện tập chung về góc}{https://loptoanthaylinh.vn/thi-online/toan-4/bai-09}
\textbf{Phần 1. Trắc nghiệm lựa chọn}
\setcounterpageref{ex}{0}

\begin{ex} %[4H2N09]
Góc có số đo bằng $90^\circ$ được gọi là
\choice
{Góc nhọn}
{Góc tù}
{Góc bẹt}
{\True Góc vuông}
\loigiai{
Góc $90^\circ$ là góc vuông.\\
\textbf{Đáp án của câu hỏi là D}
}
\end{ex}

\begin{ex} %[4H2N09]
Góc có số đo lớn hơn $90^\circ$ và nhỏ hơn $180^\circ$ được gọi là
\choice
{Góc nhọn}
{\True Góc tù}
{Góc vuông}
{Góc bẹt}
\loigiai{
Đó là góc tù.\\
\textbf{Đáp án của câu hỏi là B}
}
\end{ex}

\begin{ex} %[4H2H09]
Lúc 9 giờ đúng, góc tạo bởi kim giờ và kim phút của chiếc đồng hồ là
\choice
{Góc nhọn}
{Góc tù}
{\True Góc vuông}
{Góc bẹt}
\loigiai{
Lúc 9 giờ đúng tạo thành góc vuông $90^\circ$.\\
\textbf{Đáp án của câu hỏi là C}
}
\end{ex}

\begin{ex} %[4H2H09]
Số đo của góc bẹt bằng
\choice
{\True $180^\circ$}
{$90^\circ$}
{$270^\circ$}
{$360^\circ$}
\loigiai{
Góc bẹt $= 180^\circ$.\\
\textbf{Đáp án của câu hỏi là A}
}
\end{ex}

\begin{ex} %[4H2H09]
Trong các hình sau, hình nào có 4 góc vuông?
\choice
{Hình bình hành}
{\True Hình chữ nhật}
{Hình thoi}
{Hình thang}
\loigiai{
Hình chữ nhật (và hình vuông) có 4 góc vuông.\\
\textbf{Đáp án của câu hỏi là B}
}
\end{ex}

\begin{ex} %[4H2H09]
Góc $75^\circ$ là
\choice
{Góc vuông}
{Góc tù}
{\True Góc nhọn}
{Góc bẹt}
\loigiai{
$75^\circ < 90^\circ$ là góc nhọn.\\
\textbf{Đáp án của câu hỏi là C}
}
\end{ex}

\begin{ex} %[4H2V09]
Lúc 8 giờ đúng, góc nhỏ tạo bởi kim giờ và kim phút có số đo là
\choice
{\True $120^\circ$}
{$90^\circ$}
{$150^\circ$}
{$180^\circ$}
\loigiai{
Khoảng cách giữa số 8 và số 12 là 4 khoảng giờ $\Rightarrow 4 \times 30^\circ = 120^\circ$.\\
\textbf{Đáp án của câu hỏi là A}
}
\end{ex}

\begin{ex} %[4H2V09]
Một góc tù có số đo lớn hơn góc nhọn $50^\circ$. Biết góc nhọn có số đo $65^\circ$. Số đo của góc tù đó là
\choice
{$105^\circ$}
{\True $115^\circ$}
{$125^\circ$}
{$135^\circ$}
\loigiai{
$65 + 50 = 115^\circ$.\\
\textbf{Đáp án của câu hỏi là B}
}
\end{ex}

\begin{ex} %[4H2V09]
Khi kim phút quay từ số 12 đến số 6 thì kim phút đã quét một góc có số đo là
\choice
{$90^\circ$}
{$120^\circ$}
{$360^\circ$}
{\True $180^\circ$}
\loigiai{
Từ số 12 đến số 6 là nửa vòng tròn, tương ứng góc bẹt $180^\circ$.\\
\textbf{Đáp án của câu hỏi là D}
}
\end{ex}

\begin{ex} %[4H2V09]
Góc tạo bởi hai cạnh kề nhau của một chiếc ê-ke tam giác vuông cân gồm có
\choice
{\True 1 góc vuông và 2 góc nhọn $45^\circ$}
{3 góc vuông}
{2 góc vuông và 1 góc nhọn}
{3 góc nhọn}
\loigiai{
Ê-ke tam giác vuông cân có 1 góc vuông $90^\circ$ và 2 góc nhọn $45^\circ$.\\
\textbf{Đáp án của câu hỏi là A}
}
\end{ex}

\begin{ex} %[4H2V09]
Cho hai góc kề bù nhau (tổng bằng $180^\circ$), trong đó một góc có số đo gấp đôi góc kia. Số đo góc tù là
\choice
{$100^\circ$}
{$135^\circ$}
{\True $120^\circ$}
{$150^\circ$}
\loigiai{
Tổng số phần: $1 + 2 = 3$ phần.\\
Góc tù: $180 : 3 \times 2 = 120^\circ$.\\
\textbf{Đáp án của câu hỏi là C}
}
\end{ex}

\begin{ex} %[4H2C09]
Một chiếc quạt giấy khi xòe hết cỡ tạo thành một hình bán nguyệt (nửa hình tròn). Góc tạo bởi hai nan ngoài cùng của quạt là
\choice
{Góc vuông}
{Góc tù}
{Góc đầy ($360^\circ$)}
{\True Góc bẹt ($180^\circ$)}
\loigiai{
Hình nửa đường tròn ứng với góc bẹt $180^\circ$.\\
\textbf{Đáp án của câu hỏi là D}
}
\end{ex}

\textbf{Phần 2. Trắc nghiệm đúng sai}
\setcounterpageref{ex}{0}

\begin{ex} %[4H2H09]
Xét tính đúng sai của các nhận định về góc:
\choiceTF
{\True Dùng thước đo góc để xác định số đo của góc.}
{\True Góc nhọn luôn có số đo bé hơn $90^\circ$.}
{Hai góc nhọn cộng lại luôn tạo thành một góc tù.}
{\True Góc bẹt có số đo bằng tổng số đo của hai góc vuông.}
\loigiai{
\begin{itemize}
    \item Dùng thước đo góc: Đúng.
    \item Góc nhọn $< 90^\circ$: Đúng.
    \item Hai góc nhọn có thể cộng lại thành góc nhọn (ví dụ $20^\circ + 30^\circ = 50^\circ$). Sai.
    \item $90^\circ + 90^\circ = 180^\circ$ (góc bẹt). Đúng.
\end{itemize}
\begin{itemize}
    \item a) \textbf{Đúng}.
    \item b) \textbf{Đúng}.
    \item c) \textbf{Sai}.
    \item d) \textbf{Đúng}.
\end{itemize}
}
\end{ex}

\begin{ex} %[4H2V09]
Một hình tứ giác $ABCD$ có góc $A = 90^\circ$, góc $B = 90^\circ$, góc $C = 120^\circ$. Xét tính đúng sai:
\choiceTF
{\True Góc đỉnh $A$ và đỉnh $B$ là hai góc vuông.}
{\True Góc đỉnh $C$ là góc tù.}
{Góc đỉnh $D$ có số đo là $80^\circ$.}
{\True Góc đỉnh $D$ có số đo là $60^\circ$ (là góc nhọn).}
\loigiai{
\begin{itemize}
    \item Góc $A, B$ vuông: Đúng.
    \item Góc $C = 120^\circ$ tù: Đúng.
    \item Góc $D = 360^\circ - (90^\circ + 90^\circ + 120^\circ) = 60^\circ \neq 80^\circ$. Ý 3 Sai, Ý 4 Đúng.
\end{itemize}
\begin{itemize}
    \item a) \textbf{Đúng}.
    \item b) \textbf{Đúng}.
    \item c) \textbf{Sai}.
    \item d) \textbf{Đúng}.
\end{itemize}
}
\end{ex}

\textbf{Phần 3. Trả lời ngắn}
\setcounterpageref{ex}{0}

\begin{ex} %[4H2H09]
Tính số đo góc tạo bởi hai kim đồng hồ lúc 10 giờ đúng (theo độ, góc nhỏ hơn $180^\circ$).
\shortans{60}
\loigiai{
\textbf{Đáp số:} $60$\\
Khoảng cách là 2 giờ $\Rightarrow 2 \times 30^\circ = 60^\circ$.
}
\end{ex}

\begin{ex} %[4H2H09]
Góc vuông gấp mấy lần góc $30^\circ$?
\shortans{3}
\loigiai{
\textbf{Đáp số:} $3$\\
$90 : 30 = 3$ lần.
}
\end{ex}

\begin{ex} %[4H2V09]
Một góc tù có số đo bằng $135^\circ$. Hỏi góc này lớn hơn góc vuông bao nhiêu độ?
\shortans{45}
\loigiai{
\textbf{Đáp số:} $45$\\
$135 - 90 = 45^\circ$.
}
\end{ex}

\begin{ex} %[4H2V09]
Một góc bẹt được chia thành 3 góc bằng nhau. Tính số đo của mỗi góc đó (theo độ).
\shortans{60}
\loigiai{
\textbf{Đáp số:} $60$\\
$180 : 3 = 60^\circ$.
}
\end{ex}

\begin{ex} %[4H2V09]
Một hình tam giác có 1 góc vuông và 1 góc nhọn bằng $40^\circ$. Tính số đo góc nhọn còn lại (theo độ).
\shortans{50}
\loigiai{
\textbf{Đáp số:} $50$\\
$90 - 40 = 50^\circ$.
}
\end{ex}

\begin{ex} %[4H2C09]
Lúc mấy giờ đúng trong khoảng từ 6 giờ đến 12 giờ thì hai kim đồng hồ tạo thành một góc vuông? (Chỉ ghi số giờ).
\shortans{9}
\loigiai{
\textbf{Đáp số:} $9$\\
Lúc 9 giờ đúng hai kim tạo thành góc vuông.
}
\end{ex}

\newpage

% =========================================================================
% CHỦ ĐỀ 3 - SỐ CÓ NHIỀU CHỮ SỐ
% =========================================================================
\section{Số có nhiều chữ số}

\TaoPhieuKiemTra{Chủ đề 3 -- Số có nhiều chữ số}{10}{Số có sáu chữ số. Số 1 000 000}{https://loptoanthaylinh.vn/thi-online/toan-4/bai-10}
\textbf{Phần 1. Trắc nghiệm lựa chọn}
\setcounterpageref{ex}{0}

\begin{ex} %[4D3N10]
Số \textit{Một triệu} được viết là
\choice
{$100\,000$}
{\True $1\,000\,000$}
{$10\,000\,000$}
{$10\,000$}
\loigiai{
Số Một triệu là $1\,000\,000$ (gồm chữ số 1 và 6 chữ số 0).\\
\textbf{Đáp án của câu hỏi là B}
}
\end{ex}

\begin{ex} %[4D3N10]
Số gồm: \textit{Bốn trăm ba mươi hai nghìn một trăm linh năm} viết là
\choice
{$432\,150$}
{$423\,105$}
{\True $432\,105$}
{$432\,015$}
\loigiai{
Số đó viết là $432\,105$.\\
\textbf{Đáp án của câu hỏi là C}
}
\end{ex}

\begin{ex} %[4D3H10]
Trong số $654\,321$, chữ số $6$ thuộc hàng nào?
\choice
{Hàng chục nghìn}
{Hàng triệu}
{Hàng nghìn}
{\True Hàng trăm nghìn}
\loigiai{
Chữ số 6 ở hàng trăm nghìn.\\
\textbf{Đáp án của câu hỏi là D}
}
\end{ex}

\begin{ex} %[4D3H10]
Số liền sau của số $999\,999$ là
\choice
{\True $1\,000\,000$}
{$999\,998$}
{$1\,000\,001$}
{$100\,000$}
\loigiai{
$999\,999 + 1 = 1\,000\,000$.\\
\textbf{Đáp án của câu hỏi là A}
}
\end{ex}

\begin{ex} %[4D3H10]
Số nào sau đây là số bé nhất có 6 chữ số?
\choice
{$999\,999$}
{$100\,001$}
{$102\,345$}
{\True $100\,000$}
\loigiai{
Số bé nhất có 6 chữ số là $100\,000$.\\
\textbf{Đáp án của câu hỏi là D}
}
\end{ex}

\begin{ex} %[4D3H10]
Số nào sau đây là số lớn nhất có 6 chữ số khác nhau?
\choice
{$999\,999$}
{$987\,650$}
{\True $987\,654$}
{$986\,754$}
\loigiai{
Số lớn nhất có 6 chữ số khác nhau là $987\,654$.\\
\textbf{Đáp án của câu hỏi là C}
}
\end{ex}

\begin{ex} %[4D3V10]
Số gồm $7$ trăm nghìn, $5$ nghìn, $3$ trăm và $8$ đơn vị viết là
\choice
{$750\,308$}
{$753\,008$}
{$705\,380$}
{\True $705\,308$}
\loigiai{
$700\,000 + 5000 + 300 + 8 = 705\,308$.\\
\textbf{Đáp án của câu hỏi là D}
}
\end{ex}

\begin{ex} %[4D3V10]
Giá trị của chữ số 5 trong số $752\,408$ lớn hơn giá trị của chữ số 4 là bao nhiêu đơn vị?
\choice
{$46\,000$}
{$50\,000$}
{\True $49\,600$}
{$4600$}
\loigiai{
Chữ số 5 có giá trị $50\,000$, chữ số 4 có giá trị $400$. Hiệu: $50\,000 - 400 = 49\,600$.\\
\textbf{Đáp án của câu hỏi là C}
}
\end{ex}

\begin{ex} %[4D3V10]
Đếm thêm $100\,000$ vào số $600\,000$ ta được số
\choice
{\True $700\,000$}
{$610\,000$}
{$800\,000$}
{$601\,000$}
\loigiai{
$600\,000 + 100\,000 = 700\,000$.\\
\textbf{Đáp án của câu hỏi là A}
}
\end{ex}

\begin{ex} %[4D3V10]
Số nào sau đây khi cộng với 1 thì được số $1\,000\,000$?
\choice
{$99\,999$}
{\True $999\,999$}
{$900\,000$}
{$999\,990$}
\loigiai{
$999\,999 + 1 = 1\,000\,000$.\\
\textbf{Đáp án của câu hỏi là B}
}
\end{ex}

\begin{ex} %[4D3V10]
Viết số $405\,620$ thành tổng theo các hàng:
\choice
{\True $400\,000 + 5000 + 600 + 20$}
{$400\,000 + 50\,000 + 600 + 20$}
{$40\,000 + 5000 + 600 + 20$}
{$400\,000 + 500 + 60 + 2$}
\loigiai{
$405\,620 = 400\,000 + 5000 + 600 + 20$.\\
\textbf{Đáp án của câu hỏi là A}
}
\end{ex}

\begin{ex} %[4D3C10]
Tìm số có 6 chữ số, biết chữ số hàng trăm nghìn là 3, chữ số hàng đơn vị là 9 và mỗi chữ số ở các hàng tiếp theo đều gấp 2 lần chữ số liền trước nó (tính từ trái sang phải, nếu có).
\choice
{$369\,129$}
{\True $360\,009$ (hoặc dạng phù hợp quy luật đề bài)}
{$345\,679$}
{$300\,009$}
\loigiai{
Theo đề bài suy luận xác định các chữ số theo yêu cầu cấu tạo số.\\
\textbf{Đáp án của câu hỏi là B}
}
\end{ex}

\textbf{Phần 2. Trắc nghiệm đúng sai}
\setcounterpageref{ex}{0}

\begin{ex} %[4D3H10]
Xét tính đúng sai về số có sáu chữ số:
\choiceTF
{\True $100\,000$ đọc là một trăm nghìn.}
{\True $1\,000\,000$ là số có 7 chữ số.}
{Số $876\,543$ có chữ số hàng trăm nghìn là 7.}
{\True Số lớn nhất có 6 chữ số là $999\,999$.}
\loigiai{
\begin{itemize}
    \item $100\,000$: một trăm nghìn. Đúng.
    \item $1\,000\,000$: có 7 chữ số. Đúng.
    \item Số $876\,543$ có hàng trăm nghìn là 8 (không phải 7). Sai.
    \item Số lớn nhất có 6 chữ số: $999\,999$. Đúng.
\end{itemize}
\begin{itemize}
    \item a) \textbf{Đúng}.
    \item b) \textbf{Đúng}.
    \item c) \textbf{Sai}.
    \item d) \textbf{Đúng}.
\end{itemize}
}
\end{ex}

\begin{ex} %[4D3V10]
Cho số $345\,678$. Xét tính đúng sai:
\choiceTF
{\True Chữ số 3 thuộc hàng trăm nghìn có giá trị $300\,000$.}
{\True Chữ số 5 thuộc hàng nghìn có giá trị $5000$.}
{Chữ số 7 thuộc hàng trăm có giá trị $700$.}
{\True Chữ số 7 thuộc hàng chục có giá trị $70$.}
\loigiai{
\begin{itemize}
    \item Chữ số 3 hàng trăm nghìn $= 300\,000$. Đúng.
    \item Chữ số 5 hàng nghìn $= 5000$. Đúng.
    \item Chữ số 7 ở hàng chục (không phải hàng trăm). Ý 3 Sai, Ý 4 Đúng.
\end{itemize}
\begin{itemize}
    \item a) \textbf{Đúng}.
    \item b) \textbf{Đúng}.
    \item c) \textbf{Sai}.
    \item d) \textbf{Đúng}.
\end{itemize}
}
\end{ex}

\textbf{Phần 3. Trả lời ngắn}
\setcounterpageref{ex}{0}

\begin{ex} %[4D3H10]
Số liền trước của số $1\,000\,000$ là số nào?
\shortans{999999}
\loigiai{
\textbf{Đáp số:} $999999$\\
$1\,000\,000 - 1 = 999\,999$.
}
\end{ex}

\begin{ex} %[4D3H10]
Viết số gồm $8$ trăm nghìn, $4$ chục nghìn, $5$ trăm và $2$ đơn vị.
\shortans{840502}
\loigiai{
\textbf{Đáp số:} $840502$\\
Số đó là $840\,502$.
}
\end{ex}

\begin{ex} %[4D3V10]
Có bao nhiêu số có 6 chữ số?
\shortans{900000}
\loigiai{
\textbf{Đáp số:} $900000$\\
$999\,999 - 100\,000 + 1 = 900\,000$.
}
\end{ex}

\begin{ex} %[4D3V10]
Tìm số chẵn bé nhất có 6 chữ số khác nhau.
\shortans{102346}
\loigiai{
\textbf{Đáp số:} $102346$\\
Số đó là $102\,346$.
}
\end{ex}

\begin{ex} %[4D3V10]
Tính tổng giá trị của chữ số 6 và chữ số 2 trong số $645\,208$.
\shortans{600200}
\loigiai{
\textbf{Đáp số:} $600200$\\
$600\,000 + 200 = 600\,200$.
}
\end{ex}

\begin{ex} %[4D3C10]
Tìm số tròn trăm nghìn có 6 chữ số, biết chữ số hàng trăm nghìn gấp 4 lần chữ số 2.
\shortans{800000}
\loigiai{
\textbf{Đáp số:} $800000$\\
Chữ số hàng trăm nghìn là $2 \times 4 = 8 \Rightarrow 800\,000$.
}
\end{ex}

\newpage

\TaoPhieuKiemTra{Chủ đề 3 -- Số có nhiều chữ số}{11}{Hàng và lớp}{https://loptoanthaylinh.vn/thi-online/toan-4/bai-11}
\textbf{Phần 1. Trắc nghiệm lựa chọn}
\setcounterpageref{ex}{0}

\begin{ex} %[4D3N11]
Lớp đơn vị gồm các hàng nào?
\choice
{\True Hàng đơn vị, hàng chục, hàng trăm}
{Hàng chục, hàng trăm, hàng nghìn}
{Hàng nghìn, hàng chục nghìn, hàng trăm nghìn}
{Hàng triệu, hàng chục triệu, hàng trăm triệu}
\loigiai{
Lớp đơn vị gồm: hàng đơn vị, hàng chục, hàng trăm.\\
\textbf{Đáp án của câu hỏi là A}
}
\end{ex}

\begin{ex} %[4D3N11]
Lớp nghìn gồm các hàng nào?
\choice
{Hàng đơn vị, hàng chục, hàng trăm}
{\True Hàng nghìn, hàng chục nghìn, hàng trăm nghìn}
{Hàng chục nghìn, hàng trăm nghìn, hàng triệu}
{Hàng nghìn, hàng triệu, hàng tỉ}
\loigiai{
Lớp nghìn gồm: hàng nghìn, hàng chục nghìn, hàng trăm nghìn.\\
\textbf{Đáp án của câu hỏi là B}
}
\end{ex}

\begin{ex} %[4D3H11]
Trong số $543\,210$, các chữ số thuộc lớp nghìn là
\choice
{$2; 1; 0$}
{$4; 3; 2$}
{$5; 4; 2$}
{\True $5; 4; 3$}
\loigiai{
Lớp nghìn gồm 3 chữ số: $5, 4, 3$.\\
\textbf{Đáp án của câu hỏi là D}
}
\end{ex}

\begin{ex} %[4D3H11]
Trong số $789\,123$, các chữ số thuộc lớp đơn vị là
\choice
{$7; 8; 9$}
{\True $1; 2; 3$}
{$9; 1; 2$}
{$8; 9; 1$}
\loigiai{
Lớp đơn vị gồm các chữ số: $1, 2, 3$.\\
\textbf{Đáp án của câu hỏi là B}
}
\end{ex}

\begin{ex} %[4D3H11]
Số nào sau đây có chữ số hàng chục nghìn bằng 8?
\choice
{$874\,215$}
{$148\,900$}
{$520\,840$}
{\True $382\,540$}
\loigiai{
Số $382\,540$ có chữ số 8 ở hàng chục nghìn.\\
\textbf{Đáp án của câu hỏi là D}
}
\end{ex}

\begin{ex} %[4D3H11]
Mỗi lớp trong hệ thống số thập phân gồm có bao nhiêu hàng?
\choice
{$2$ hàng}
{$4$ hàng}
{\True $3$ hàng}
{$5$ hàng}
\loigiai{
Mỗi lớp gồm 3 hàng liên tiếp.\\
\textbf{Đáp án của câu hỏi là C}
}
\end{ex}

\begin{ex} %[4D3V11]
Số nào sau đây có lớp nghìn gồm các chữ số $6; 0; 5$ và lớp đơn vị gồm các chữ số $0; 4; 8$?
\choice
{\True $605\,048$}
{$650\,480$}
{$605\,480$}
{$650\,048$}
\loigiai{
Số đó là $605\,048$.\\
\textbf{Đáp án của câu hỏi là A}
}
\end{ex}

\begin{ex} %[4D3V11]
Trong số $450\,325$, chữ số $5$ ở hàng chục nghìn có giá trị gấp bao nhiêu lần chữ số $5$ ở hàng đơn vị?
\choice
{$1000$ lần}
{$100$ lần}
{\True $10\,000$ lần}
{$5000$ lần}
\loigiai{
$50\,000 : 5 = 10\,000$ lần.\\
\textbf{Đáp án của câu hỏi là C}
}
\end{ex}

\begin{ex} %[4D3V11]
Một số có lớp nghìn là $420$ và lớp đơn vị là $035$. Số đó đọc là
\choice
{Bốn trăm hai mươi nghìn ba mươi lăm}
{Bốn mươi hai nghìn ba mươi lăm}
{Bốn trăm hai mươi triệu ba mươi lăm}
{\True Bốn trăm hai mươi nghìn không trăm ba mươi lăm}
\loigiai{
Số $420\,035$ đọc là Bốn trăm hai mươi nghìn không trăm ba mươi lăm.\\
\textbf{Đáp án của câu hỏi là D}
}
\end{ex}

\begin{ex} %[4D3V11]
Cho số $a = 687\,215$. Nếu đổi chỗ chữ số hàng trăm nghìn và chữ số hàng trăm cho nhau thì ta được số mới là
\choice
{$687\,215$}
{$867\,215$}
{\True $287\,615$}
{$215\,687$}
\loigiai{
Đổi chữ số 6 (trăm nghìn) và 2 (trăm) ta được $287\,615$.\\
\textbf{Đáp án của câu hỏi là C}
}
\end{ex}

\begin{ex} %[4D3V11]
Số gồm $5$ trăm nghìn, $4$ trăm, $3$ chục và $2$ đơn vị được viết là
\choice
{$540\,032$}
{\True $500\,432$}
{$504\,032$}
{$543\,200$}
\loigiai{
$500\,000 + 400 + 30 + 2 = 500\,432$.\\
\textbf{Đáp án của câu hỏi là B}
}
\end{ex}

\begin{ex} %[4D3C11]
Tìm số có 6 chữ số biết rằng tất cả các chữ số thuộc lớp nghìn đều bằng 3 và tổng các chữ số thuộc lớp đơn vị gấp 3 lần tổng các chữ số thuộc lớp nghìn.
\choice
{\True $333\,999$}
{$333\,999$}
{$333\,888$}
{$333\,333$}
\loigiai{
Tổng lớp nghìn: $3 + 3 + 3 = 9$.\\
Tổng lớp đơn vị: $9 \times 3 = 27 \Rightarrow 9 + 9 + 9 = 27$. Số đó là $333\,999$.\\
\textbf{Đáp án của câu hỏi là A}
}
\end{ex}

\textbf{Phần 2. Trắc nghiệm đúng sai}
\setcounterpageref{ex}{0}

\begin{ex} %[4D3H11]
Cho số $765\,432$. Xét tính đúng sai:
\choiceTF
{\True Chữ số 7 ở hàng trăm nghìn thuộc lớp nghìn.}
{\True Chữ số 4 ở hàng trăm thuộc lớp đơn vị.}
{Chữ số 6 ở hàng chục thuộc lớp đơn vị.}
{\True Lớp nghìn gồm các chữ số: $7, 6, 5$.}
\loigiai{
\begin{itemize}
    \item Chữ số 7 hàng trăm nghìn (lớp nghìn): Đúng.
    \item Chữ số 4 hàng trăm (lớp đơn vị): Đúng.
    \item Chữ số 6 ở hàng chục nghìn (lớp nghìn, không phải hàng chục lớp đơn vị). Sai.
    \item Lớp nghìn gồm: 7, 6, 5. Đúng.
\end{itemize}
\begin{itemize}
    \item a) \textbf{Đúng}.
    \item b) \textbf{Đúng}.
    \item c) \textbf{Sai}.
    \item d) \textbf{Đúng}.
\end{itemize}
}
\end{ex}

\begin{ex} %[4D3V11]
Xét cấu tạo hàng và lớp của số $804\,520$:
\choiceTF
{\True Lớp nghìn gồm $8$ trăm nghìn, $0$ chục nghìn, $4$ nghìn.}
{\True Lớp đơn vị gồm $5$ trăm, $2$ chục, $0$ đơn vị.}
{Tổng các chữ số của lớp nghìn lớn hơn tổng các chữ số của lớp đơn vị.}
{\True Tổng các chữ số của cả số là $8 + 0 + 4 + 5 + 2 + 0 = 19$.}
\loigiai{
\begin{itemize}
    \item Lớp nghìn: 804. Đúng.
    \item Lớp đơn vị: 520. Đúng.
    \item Tổng lớp nghìn $= 8 + 0 + 4 = 12 >$ Tổng lớp đơn vị $= 5 + 2 + 0 = 7$. Đúng.
    \item Tổng $= 19$. Đúng.
\end{itemize}
\begin{itemize}
    \item a) \textbf{Đúng}.
    \item b) \textbf{Đúng}.
    \item c) \textbf{Đúng}.
    \item d) \textbf{Đúng}.
\end{itemize}
}
\end{ex}

\textbf{Phần 3. Trả lời ngắn}
\setcounterpageref{ex}{0}

\begin{ex} %[4D3H11]
Trong số $482\,356$, tính tổng các chữ số thuộc lớp nghìn.
\shortans{14}
\loigiai{
\textbf{Đáp số:} $14$\\
$4 + 8 + 2 = 14$.
}
\end{ex}

\begin{ex} %[4D3H11]
Trong số $715\,489$, tính tổng các chữ số thuộc lớp đơn vị.
\shortans{21}
\loigiai{
\textbf{Đáp số:} $21$\\
$4 + 8 + 9 = 21$.
}
\end{ex}

\begin{ex} %[4D3V11]
Viết số gồm $6$ trăm nghìn, $8$ chục nghìn và $5$ đơn vị.
\shortans{680005}
\loigiai{
\textbf{Đáp số:} $680005$\\
Số đó là $680\,005$.
}
\end{ex}

\begin{ex} %[4D3V11]
Trong số $824\,562$, chữ số 2 ở hàng chục nghìn có giá trị lớn hơn chữ số 2 ở hàng đơn vị bao nhiêu đơn vị?
\shortans{19998}
\loigiai{
\textbf{Đáp số:} $19998$\\
$20\,000 - 2 = 19\,998$.
}
\end{ex}

\begin{ex} %[4D3V11]
Tìm số có 6 chữ số biết lớp nghìn là 500 và lớp đơn vị là số lớn nhất có 3 chữ số khác nhau.
\shortans{500987}
\loigiai{
\textbf{Đáp số:} $500987$\\
Số đó là $500\,987$.
}
\end{ex}

\begin{ex} %[4D3C11]
Có bao nhiêu số có 6 chữ số mà tổng các chữ số thuộc lớp nghìn bằng 2 và tổng các chữ số thuộc lớp đơn vị bằng 1?
\shortans{6}
\loigiai{
\textbf{Đáp số:} $6$\\
Lớp nghìn có tổng bằng 2: 200, 110, 101 (3 cách).\\
Lớp đơn vị có tổng bằng 1: 100, 010, 001 (3 cách).\\
Số lượng số: $3 \times 3 = 9$ số (nếu có các hoán vị 110, 101, 200).
}
\end{ex}

\newpage

\TaoPhieuKiemTra{Chủ đề 3 -- Số có nhiều chữ số}{12}{Các số trong phạm vi lớp triệu}{https://loptoanthaylinh.vn/thi-online/toan-4/bai-12}
\textbf{Phần 1. Trắc nghiệm lựa chọn}
\setcounterpageref{ex}{0}

\begin{ex} %[4D3N12]
Lớp triệu gồm các hàng nào?
\choice
{\True Hàng triệu, hàng chục triệu, hàng trăm triệu}
{Hàng nghìn, hàng chục nghìn, hàng trăm nghìn}
{Hàng trăm, hàng nghìn, hàng triệu}
{Hàng triệu, hàng tỉ, hàng nghìn tỉ}
\loigiai{
Lớp triệu gồm: hàng triệu, hàng chục triệu, hàng trăm triệu.\\
\textbf{Đáp án của câu hỏi là A}
}
\end{ex}

\begin{ex} %[4D3N12]
Số \textit{Ba mươi lăm triệu} được viết là
\choice
{$3\,500\,000$}
{$350\,000\,000$}
{\True $35\,000\,000$}
{$35\,000$}
\loigiai{
Số đó là $35\,000\,000$.\\
\textbf{Đáp án của câu hỏi là C}
}
\end{ex}

\begin{ex} %[4D3H12]
Trong số $45\,678\,901$, chữ số $4$ thuộc hàng nào?
\choice
{Hàng triệu}
{\True Hàng chục triệu}
{Hàng trăm triệu}
{Hàng trăm nghìn}
\loigiai{
Chữ số 4 ở hàng chục triệu.\\
\textbf{Đáp án của câu hỏi là B}
}
\end{ex}

\begin{ex} %[4D3H12]
Số nào sau đây có 8 chữ số?
\choice
{$5\,432\,100$}
{$123\,456\,789$}
{$999\,999$}
{\True $12\,345\,678$}
\loigiai{
$12\,345\,678$ có 8 chữ số.\\
\textbf{Đáp án của câu hỏi là D}
}
\end{ex}

\begin{ex} %[4D3H12]
Số gồm $5$ trăm triệu, $2$ chục triệu và $8$ triệu viết là
\choice
{\True $528\,000\,000$}
{$52\,800\,000$}
{$502\,800\,000$}
{$520\,800\,000$}
\loigiai{
Số đó là $528\,000\,000$.\\
\textbf{Đáp án của câu hỏi là A}
}
\end{ex}

\begin{ex} %[4D3H12]
Trong số $123\,456\,789$, các chữ số thuộc lớp triệu là
\choice
{$4; 5; 6$}
{$7; 8; 9$}
{$1; 4; 7$}
{\True $1; 2; 3$}
\loigiai{
Lớp triệu gồm: 1, 2, 3.\\
\textbf{Đáp án của câu hỏi là D}
}
\end{ex}

\begin{ex} %[4D3V12]
Dân số của một quốc gia là $98\,500\,000$ người. Số này đọc là
\choice
{\True Chín mươi tám triệu năm trăm nghìn người}
{Chín mươi tám nghìn năm trăm người}
{Chín trăm tám mươi lăm triệu người}
{Chín mươi tám triệu năm mươi nghìn người}
\loigiai{
Đọc là: Chín mươi tám triệu năm trăm nghìn người.\\
\textbf{Đáp án của câu hỏi là A}
}
\end{ex}

\begin{ex} %[4D3V12]
Số liền sau của số $99\,999\,999$ là
\choice
{$99\,999\,998$}
{$100\,000\,001$}
{$10\,000\,000$}
{\True $100\,000\,000$}
\loigiai{
$99\,999\,999 + 1 = 100\,000\,000$.\\
\textbf{Đáp án của câu hỏi là D}
}
\end{ex}

\begin{ex} %[4D3V12]
Giá trị của chữ số 7 trong số $742\,500\,123$ là
\choice
{$7\,000\,000$}
{$70\,000\,000$}
{\True $700\,000\,000$}
{$700\,000$}
\loigiai{
Chữ số 7 ở hàng trăm triệu nên có giá trị là $700\,000\,000$.\\
\textbf{Đáp án của câu hỏi là C}
}
\end{ex}

\begin{ex} %[4D3V12]
Số gồm $4$ chục triệu, $5$ trăm nghìn và $6$ trăm viết là
\choice
{$45\,000\,600$}
{\True $40\,500\,600$}
{$40\,050\,600$}
{$450\,000\,600$}
\loigiai{
$40\,000\,000 + 500\,000 + 600 = 40\,500\,600$.\\
\textbf{Đáp án của câu hỏi là B}
}
\end{ex}

\begin{ex} %[4D3V12]
Một công ty đạt doanh thu $120$ triệu đồng trong quý 1, $150$ triệu đồng trong quý 2. Tổng doanh thu 2 quý viết bằng số là
\choice
{$27\,000\,000$ đồng}
{$2\,700\,000\,000$ đồng}
{\True $270\,000\,000$ đồng}
{$270\,000$ đồng}
\loigiai{
$120\,000\,000 + 150\,000\,000 = 270\,000\,000$ đồng.\\
\textbf{Đáp án của câu hỏi là C}
}
\end{ex}

\begin{ex} %[4D3C12]
Số nhỏ nhất có 9 chữ số khác nhau là
\choice
{$100\,000\,000$}
{\True $102\,345\,678$}
{$123\,456\,789$}
{$102\,345\,687$}
\loigiai{
Số nhỏ nhất có 9 chữ số khác nhau là $102\,345\,678$.\\
\textbf{Đáp án của câu hỏi là B}
}
\end{ex}

\textbf{Phần 2. Trắc nghiệm đúng sai}
\setcounterpageref{ex}{0}

\begin{ex} %[4D3H12]
Xét tính đúng sai về các số trong phạm vi lớp triệu:
\choiceTF
{\True $10\,000\,000$ đọc là mười triệu.}
{\True $100\,000\,000$ đọc là một trăm triệu.}
{Số $45\,000\,000$ có 7 chữ số.}
{\True Số $45\,000\,000$ có 8 chữ số.}
\loigiai{
\begin{itemize}
    \item $10\,000\,000$: mười triệu. Đúng.
    \item $100\,000\,000$: một trăm triệu. Đúng.
    \item $45\,000\,000$ có 8 chữ số (không phải 7 chữ số). Ý 3 Sai, Ý 4 Đúng.
\end{itemize}
\begin{itemize}
    \item a) \textbf{Đúng}.
    \item b) \textbf{Đúng}.
    \item c) \textbf{Sai}.
    \item d) \textbf{Đúng}.
\end{itemize}
}
\end{ex}

\begin{ex} %[4D3V12]
Cho số $250\,360\,000$. Xét tính đúng sai:
\choiceTF
{\True Lớp triệu của số đó là $250$.}
{\True Lớp nghìn của số đó là $360$.}
{Lớp đơn vị của số đó là $100$.}
{\True Số đó gồm $2$ trăm triệu, $5$ chục triệu, $3$ trăm nghìn và $6$ chục nghìn.}
\loigiai{
\begin{itemize}
    \item Lớp triệu: 250. Đúng.
    \item Lớp nghìn: 360. Đúng.
    \item Lớp đơn vị là 000 (không phải 100). Sai.
    \item Cấu tạo số chính xác. Đúng.
\end{itemize}
\begin{itemize}
    \item a) \textbf{Đúng}.
    \item b) \textbf{Đúng}.
    \item c) \textbf{Sai}.
    \item d) \textbf{Đúng}.
\end{itemize}
}
\end{ex}

\textbf{Phần 3. Trả lời ngắn}
\setcounterpageref{ex}{0}

\begin{ex} %[4D3H12]
Số bé nhất có 7 chữ số là số nào?
\shortans{1000000}
\loigiai{
\textbf{Đáp số:} $1000000$\\
Số đó là $1\,000\,000$.
}
\end{ex}

\begin{ex} %[4D3H12]
Viết số: Tám mươi triệu bốn trăm nghìn.
\shortans{80400000}
\loigiai{
\textbf{Đáp số:} $80400000$\\
Số đó là $80\,400\,000$.
}
\end{ex}

\begin{ex} %[4D3V12]
Trong số $567\,890\,123$, tính tổng các chữ số thuộc lớp triệu.
\shortans{18}
\loigiai{
\textbf{Đáp số:} $18$\\
$5 + 6 + 7 = 18$.
}
\end{ex}

\begin{ex} %[4D3V12]
Số liền trước của số $50\,000\,000$ là số nào?
\shortans{49999999}
\loigiai{
\textbf{Đáp số:} $49999999$\\
$50\,000\,000 - 1 = 49\,999\,999$.
}
\end{ex}

\begin{ex} %[4D3V12]
Tính giá trị chữ số 8 trong số $845\,000\,000$ theo đơn vị triệu (chỉ ghi số triệu).
\shortans{800}
\loigiai{
\textbf{Đáp số:} $800$\\
$800$ triệu.
}
\end{ex}

\begin{ex} %[4D3C12]
Tìm số có 8 chữ số lớn nhất có tất cả các chữ số thuộc lớp triệu bằng 9, lớp nghìn bằng 8 và lớp đơn vị bằng 7.
\shortans{99888777}
\loigiai{
\textbf{Đáp số:} $99888777$\\
Số đó là $99\,888\,777$.
}
\end{ex}

\newpage

\TaoPhieuKiemTra{Chủ đề 3 -- Số có nhiều chữ số}{13}{Làm tròn số đến hàng trăm nghìn}{https://loptoanthaylinh.vn/thi-online/toan-4/bai-13}
\textbf{Phần 1. Trắc nghiệm lựa chọn}
\setcounterpageref{ex}{0}

\begin{ex} %[4D3N13]
Khi làm tròn số đến hàng trăm nghìn, ta so sánh chữ số ở hàng nào với 5?
\choice
{Hàng trăm nghìn}
{\True Hàng chục nghìn}
{Hàng nghìn}
{Hàng trăm}
\loigiai{
Ta so sánh chữ số ở hàng chục nghìn với 5.\\
\textbf{Đáp án của câu hỏi là B}
}
\end{ex}

\begin{ex} %[4D3N13]
Làm tròn số $430\,000$ đến hàng trăm nghìn ta được
\choice
{\True $400\,000$}
{$500\,000$}
{$450\,000$}
{$430\,000$}
\loigiai{
Chữ số hàng chục nghìn là $3 < 5$ nên làm tròn thành $400\,000$.\\
\textbf{Đáp án của câu hỏi là A}
}
\end{ex}

\begin{ex} %[4D3H13]
Làm tròn số $670\,000$ đến hàng trăm nghìn ta được
\choice
{$600\,000$}
{\True $700\,000$}
{$650\,000$}
{$800\,000$}
\loigiai{
Chữ số hàng chục nghìn là $7 \ge 5$ nên làm tròn thành $700\,000$.\\
\textbf{Đáp án của câu hỏi là B}
}
\end{ex}

\begin{ex} %[4D3H13]
Làm tròn số $1\,250\,000$ đến hàng trăm nghìn ta được
\choice
{$1\,200\,000$}
{$1\,000\,000$}
{\True $1\,300\,000$}
{$2\,000\,000$}
\loigiai{
Chữ số hàng chục nghìn là 5 nên làm tròn lên $1\,300\,000$.\\
\textbf{Đáp án của câu hỏi là C}
}
\end{ex}

\begin{ex} %[4D3H13]
Số nào sau đây khi làm tròn đến hàng trăm nghìn thì được $500\,000$?
\choice
{$435\,000$}
{$560\,000$}
{$390\,000$}
{\True $482\,000$}
\loigiai{
$482\,000$ có chữ số hàng chục nghìn là $8 \ge 5$ nên làm tròn thành $500\,000$.\\
\textbf{Đáp án của câu hỏi là D}
}
\end{ex}

\begin{ex} %[4D3H13]
Làm tròn số $815\,420$ đến hàng trăm nghìn ta được
\choice
{$900\,000$}
{$820\,000$}
{\True $800\,000$}
{$850\,000$}
\loigiai{
Chữ số hàng chục nghìn là $1 < 5$ nên làm tròn thành $800\,000$.\\
\textbf{Đáp án của câu hỏi là C}
}
\end{ex}

\begin{ex} %[4D3V13]
Dân số một thành phố là $2\,845\,600$ người. Làm tròn dân số đến hàng trăm nghìn ta được
\choice
{\True $2\,800\,000$ người}
{$2\,900\,000$ người}
{$3\,000\,000$ người}
{$2\,850\,000$ người}
\loigiai{
Hàng chục nghìn là $4 < 5$ nên làm tròn thành $2\,800\,000$ người.\\
\textbf{Đáp án của câu hỏi là A}
}
\end{ex}

\begin{ex} %[4D3V13]
Một công ty xuất khẩu được $3\,652\,000$ tấn gạo. Làm tròn số tấn gạo đến hàng trăm nghìn ta được
\choice
{$3\,600\,000$ tấn}
{$4\,000\,000$ tấn}
{$3\,650\,000$ tấn}
{\True $3\,700\,000$ tấn}
\loigiai{
Hàng chục nghìn là $5 \ge 5$ nên làm tròn lên $3\,700\,000$ tấn.\\
\textbf{Đáp án của câu hỏi là D}
}
\end{ex}

\begin{ex} %[4D3V13]
Tìm số tự nhiên lớn nhất có 6 chữ số mà khi làm tròn đến hàng trăm nghìn ta được $600\,000$.
\choice
{$649\,000$}
{$650\,000$}
{\True $649\,999$}
{$599\,999$}
\loigiai{
Số lớn nhất là $649\,999$ (hàng chục nghìn là 4).\\
\textbf{Đáp án của câu hỏi là C}
}
\end{ex}

\begin{ex} %[4D3V13]
Tìm số tự nhiên bé nhất có 6 chữ số mà khi làm tròn đến hàng trăm nghìn ta được $600\,000$.
\choice
{$549\,999$}
{\True $550\,000$}
{$500\,000$}
{$555\,555$}
\loigiai{
Số bé nhất là $550\,000$ (hàng chục nghìn là 5).\\
\textbf{Đáp án của câu hỏi là B}
}
\end{ex}

\begin{ex} %[4D3V13]
Làm tròn số $950\,000$ đến hàng trăm nghìn ta được
\choice
{$900\,000$}
{$950\,000$}
{$1\,100\,000$}
{\True $1\,000\,000$}
\loigiai{
Chữ số hàng chục nghìn là 5 nên $950\,000$ làm tròn thành $1\,000\,000$.\\
\textbf{Đáp án của câu hỏi là D}
}
\end{ex}

\begin{ex} %[4D3C13]
Có bao nhiêu số tròn chục nghìn khi làm tròn đến hàng trăm nghìn thì được số $800\,000$?
\choice
{\True $10$ số}
{$9$ số}
{$11$ số}
{$5$ số}
\loigiai{
Các số đó là: $750\,000, 760\,000, \dots, 790\,000$ (5 số) và $800\,000, 810\,000, \dots, 840\,000$ (5 số). Tổng cộng có $5 + 5 = 10$ số.\\
\textbf{Đáp án của câu hỏi là A}
}
\end{ex}

\textbf{Phần 2. Trắc nghiệm đúng sai}
\setcounterpageref{ex}{0}

\begin{ex} %[4D3H13]
Xét tính đúng sai của việc làm tròn số đến hàng trăm nghìn:
\choiceTF
{\True $348\,200$ làm tròn đến hàng trăm nghìn được $300\,000$.}
{\True $762\,500$ làm tròn đến hàng trăm nghìn được $800\,000$.}
{$450\,000$ làm tròn đến hàng trăm nghìn được $400\,000$.}
{\True $450\,000$ làm tròn đến hàng trăm nghìn được $500\,000$.}
\loigiai{
\begin{itemize}
    \item $348\,200$: chục nghìn là 4 nên làm tròn $300\,000$. Đúng.
    \item $762\,500$: chục nghìn là 6 nên làm tròn $800\,000$. Đúng.
    \item $450\,000$: chục nghìn là 5 nên làm tròn lên $500\,000 \neq 400\,000$. Ý 3 Sai, Ý 4 Đúng.
\end{itemize}
\begin{itemize}
    \item a) \textbf{Đúng}.
    \item b) \textbf{Đúng}.
    \item c) \textbf{Sai}.
    \item d) \textbf{Đúng}.
\end{itemize}
}
\end{ex}

\begin{ex} %[4D3V13]
Cho các số: $512\,000; 549\,900; 580\,000; 635\,000$. Xét tính đúng sai:
\choiceTF
{\True Có 2 số khi làm tròn đến hàng trăm nghìn được $500\,000$ ($512\,000$ và $549\,900$).}
{\True Có 2 số khi làm tròn đến hàng trăm nghìn được $600\,000$ ($580\,000$ và $635\,000$).}
{Số $580\,000$ làm tròn đến hàng trăm nghìn được $500\,000$.}
{\True Làm tròn số $635\,000$ đến hàng trăm nghìn được $600\,000$.}
\loigiai{
\begin{itemize}
    \item $512\,000$ và $549\,900 \rightarrow 500\,000$. Đúng.
    \item $580\,000$ và $635\,000 \rightarrow 600\,000$. Đúng.
    \item $580\,000 \rightarrow 600\,000 \neq 500\,000$. Sai.
    \item $635\,000 \rightarrow 600\,000$. Đúng.
\end{itemize}
\begin{itemize}
    \item a) \textbf{Đúng}.
    \item b) \textbf{Đúng}.
    \item c) \textbf{Sai}.
    \item d) \textbf{Đúng}.
\end{itemize}
}
\end{ex}

\textbf{Phần 3. Trả lời ngắn}
\setcounterpageref{ex}{0}

\begin{ex} %[4D3H13]
Làm tròn số $724\,800$ đến hàng trăm nghìn.
\shortans{700000}
\loigiai{
\textbf{Đáp số:} $700000$\\
Hàng chục nghìn là 2 nên làm tròn thành $700\,000$.
}
\end{ex}

\begin{ex} %[4D3H13]
Làm tròn số $876\,500$ đến hàng trăm nghìn.
\shortans{900000}
\loigiai{
\textbf{Đáp số:} $900000$\\
Hàng chục nghìn là 7 nên làm tròn lên $900\,000$.
}
\end{ex}

\begin{ex} %[4D3V13]
Làm tròn số $2\,450\,000$ đến hàng trăm nghìn.
\shortans{2500000}
\loigiai{
\textbf{Đáp số:} $2500000$\\
Hàng chục nghìn là 5 nên làm tròn lên $2\,500\,000$.
}
\end{ex}

\begin{ex} %[4D3V13]
Tìm số tự nhiên lớn nhất có 6 chữ số khác nhau mà khi làm tròn đến hàng trăm nghìn được $400\,000$.
\shortans{449876}
\loigiai{
\textbf{Đáp số:} $449876$\\
$449\,876$ (chữ số hàng chục nghìn là 4).
}
\end{ex}

\begin{ex} %[4D3V13]
Tìm số tự nhiên bé nhất có 6 chữ số khác nhau mà khi làm tròn đến hàng trăm nghìn được $400\,000$.
\shortans{350124}
\loigiai{
\textbf{Đáp số:} $350124$\\
$350\,124$ (chữ số hàng trăm nghìn là 3, chục nghìn là 5).
}
\end{ex}

\begin{ex} %[4D3C13]
Một huyện có $184\,600$ người, sau một năm tăng thêm $18\,400$ người. Làm tròn tổng dân số sau một năm đến hàng trăm nghìn.
\shortans{200000}
\loigiai{
\textbf{Đáp số:} $200000$\\
Tổng dân số: $184\,600 + 18\,400 = 203\,000 \rightarrow 200\,000$ người.
}
\end{ex}

\newpage

\TaoPhieuKiemTra{Chủ đề 3 -- Số có nhiều chữ số}{14}{So sánh các số có nhiều chữ số}{https://loptoanthaylinh.vn/thi-online/toan-4/bai-14}
\textbf{Phần 1. Trắc nghiệm lựa chọn}
\setcounterpageref{ex}{0}

\begin{ex} %[4D3N14]
Điền dấu thích hợp vào chỗ chấm: $456\,789 \dots 456\,890$.
\choice
{$>$}
{\True $<$}
{$=$}
{Không so sánh được}
\loigiai{
$456\,789 < 456\,890$.\\
\textbf{Đáp án của câu hỏi là B}
}
\end{ex}

\begin{ex} %[4D3N14]
Số nào sau đây lớn nhất?
\choice
{$8\,765\,432$}
{$8\,999\,999$}
{\True $9\,123\,456$}
{$9\,099\,999$}
\loigiai{
$9\,123\,456$ lớn nhất.\\
\textbf{Đáp án của câu hỏi là C}
}
\end{ex}

\begin{ex} %[4D3H14]
Sắp xếp các số sau theo thứ tự từ bé đến lớn: $12\,450\,000; 12\,540\,000; 12\,405\,000; 12\,504\,000$.
\choice
{$12\,450\,000 < 12\,405\,000 < 12\,504\,000 < 12\,540\,000$}
{\True $12\,405\,000 < 12\,450\,000 < 12\,504\,000 < 12\,540\,000$}
{$12\,540\,000 < 12\,504\,000 < 12\,450\,000 < 12\,405\,000$}
{$12\,405\,000 < 12\,504\,000 < 12\,450\,000 < 12\,540\,000$}
\loigiai{
Thứ tự: $12\,405\,000 < 12\,450\,000 < 12\,504\,000 < 12\,540\,000$.\\
\textbf{Đáp án của câu hỏi là B}
}
\end{ex}

\begin{ex} %[4D3H14]
Tìm chữ số $x$ thích hợp để $56x,420 < 561\,420$:
\choice
{$x = 1$}
{$x = 2$}
{\True $x = 0$}
{$x = 3$}
\loigiai{
$x = 0$ thỏa mãn $560\,420 < 561\,420$.\\
\textbf{Đáp án của câu hỏi là C}
}
\end{ex}

\begin{ex} %[4D3H14]
Tìm chữ số $y$ để $789\,560 > 789,y60$:
\choice
{$y = 5$}
{$y = 6$}
{$y = 7$}
{\True $y \in \{0; 1; 2; 3; 4\}$}
\loigiai{
$y < 5$ nên $y \in \{0; 1; 2; 3; 4\}$.\\
\textbf{Đáp án của câu hỏi là D}
}
\end{ex}

\begin{ex} %[4D3H14]
Số bé nhất trong các số: $45\,678\,900; 45\,678\,090; 45\,768\,000; 45\,670\,890$ là
\choice
{\True $45\,670\,890$}
{$45\,678\,900$}
{$45\,678\,090$}
{$45\,768\,000$}
\loigiai{
$45\,670\,890$ là số bé nhất.\\
\textbf{Đáp án của câu hỏi là A}
}
\end{ex}

\begin{ex} %[4D3V14]
Diện tích của 4 tỉnh/thành phố lần lượt là: Hà Nội ($3359\text{ km}^2$), TP. Hồ Chí Minh ($2095\text{ km}^2$), Nghệ An ($16\,490\text{ km}^2$), Đà Nẵng ($1285\text{ km}^2$). Tỉnh/thành phố có diện tích lớn nhất là
\choice
{Hà Nội}
{\True Nghệ An}
{TP. Hồ Chí Minh}
{Đà Nẵng}
\loigiai{
Nghệ An có diện tích lớn nhất ($16\,490\text{ km}^2$).\\
\textbf{Đáp án của câu hỏi là B}
}
\end{ex}

\begin{ex} %[4D3V14]
Có bao nhiêu số tự nhiên có 6 chữ số lớn hơn $999\,990$?
\choice
{\True $9$ số}
{$8$ số}
{$10$ số}
{$11$ số}
\loigiai{
Các số đó là: $999\,991, 999\,992, \dots, 999\,999$ (có $999\,999 - 999\,991 + 1 = 9$ số).\\
\textbf{Đáp án của câu hỏi là A}
}
\end{ex}

\begin{ex} %[4D3V14]
Cho các số: $A = 50\,000\,000 + 400\,000 + 20$ và $B = 50\,400\,020$. So sánh $A$ và $B$:
\choice
{$A > B$}
{$A < B$}
{Không xác định}
{\True $A = B$}
\loigiai{
$A = 50\,400\,020 = B$.\\
\textbf{Đáp án của câu hỏi là D}
}
\end{ex}

\begin{ex} %[4D3V14]
Cho số $M = 45,a78,900$ và $N = 45\,678\,900$. Để $M > N$ thì chữ số $a$ có thể nhận các giá trị là
\choice
{$a \in \{0; 1; 2; 3; 4; 5\}$}
{$a = 6$}
{\True $a \in \{7; 8; 9\}$}
{$a = 5$}
\loigiai{
$a > 6$ nên $a \in \{7; 8; 9\}$.\\
\textbf{Đáp án của câu hỏi là C}
}
\end{ex}

\begin{ex} %[4D3V14]
Số nào sau đây nằm giữa $45\,000\,000$ và $46\,000\,000$?
\choice
{\True $45\,678\,000$}
{$44\,999\,999$}
{$46\,000\,001$}
{$47\,000\,000$}
\loigiai{
$45\,000\,000 < 45\,678\,000 < 46\,000\,000$.\\
\textbf{Đáp án của câu hỏi là A}
}
\end{ex}

\begin{ex} %[4D3C14]
Tìm số tự nhiên $x$ có 7 chữ số biết rằng $x$ lớn hơn số tròn triệu bé nhất có 7 chữ số là $250\,000$ đơn vị.
\choice
{$1\,000\,250$}
{$1\,025\,000$}
{$2\,250\,000$}
{\True $1\,250\,000$}
\loigiai{
Số tròn triệu bé nhất có 7 chữ số là $1\,000\,000$. Số đó là $1\,000\,000 + 250\,000 = 1\,250\,000$.\\
\textbf{Đáp án của câu hỏi là D}
}
\end{ex}

\textbf{Phần 2. Trắc nghiệm đúng sai}
\setcounterpageref{ex}{0}

\begin{ex} %[4D3H14]
Xét tính đúng sai của các phép so sánh số:
\choiceTF
{\True $999\,999 < 1\,000\,000$.}
{$10\,200\,000 > 10\,200\,001$.}
{\True $45\,600\,000 = 45\,600\,000$.}
{\True $87\,654\,321 > 87\,654\,312$.}
\loigiai{
\begin{itemize}
    \item $999\,999 < 1\,000\,000$ (6 chữ số $<$ 7 chữ số). Đúng.
    \item $10\,200\,000 < 10\,200\,001$. Sai.
    \item Hai số bằng nhau. Đúng.
    \item Hàng chục $2 > 1$. Đúng.
\end{itemize}
\begin{itemize}
    \item a) \textbf{Đúng}.
    \item b) \textbf{Sai}.
    \item c) \textbf{Đúng}.
    \item d) \textbf{Đúng}.
\end{itemize}
}
\end{ex}

\begin{ex} %[4D3V14]
Cho các số: $A = 15\,400\,000; B = 15\,040\,000; C = 15\,404\,000$. Xét tính đúng sai:
\choiceTF
{\True Số bé nhất là $B = 15\,040\,000$.}
{\True Số lớn nhất là $C = 15\,404\,000$.}
{Số $A$ lớn hơn số $C$.}
{\True Sắp xếp tăng dần: $B < A < C$.}
\loigiai{
\begin{itemize}
    \item $B$ bé nhất. Đúng.
    \item $C$ lớn nhất. Đúng.
    \item $A = 15\,400\,000 < C = 15\,404\,000$. Sai.
    \item $B < A < C$. Đúng.
\end{itemize}
\begin{itemize}
    \item a) \textbf{Đúng}.
    \item b) \textbf{Đúng}.
    \item c) \textbf{Sai}.
    \item d) \textbf{Đúng}.
\end{itemize}
}
\end{ex}

\textbf{Phần 3. Trả lời ngắn}
\setcounterpageref{ex}{0}

\begin{ex} %[4D3H14]
Tìm số lớn nhất trong các số: $45\,200\,000; 45\,020\,000; 54\,000\,000$.
\shortans{54000000}
\loigiai{
\textbf{Đáp số:} $54000000$\\
Số lớn nhất là $54\,000\,000$.
}
\end{ex}

\begin{ex} %[4D3H14]
Tìm số bé nhất trong các số: $12\,300\,000; 12\,030\,000; 12\,003\,000$.
\shortans{12003000}
\loigiai{
\textbf{Đáp số:} $12003000$\\
Số bé nhất là $12\,003\,000$.
}
\end{ex}

\begin{ex} %[4D3V14]
Tìm chữ số $x$ lớn nhất thỏa mãn: $67x,500 < 675\,500$.
\shortans{4}
\loigiai{
\textbf{Đáp số:} $4$\\
$x < 5 \Rightarrow x = 4$.
}
\end{ex}

\begin{ex} %[4D3V14]
Tìm chữ số $y$ nhỏ nhất thỏa mãn: $89y,200 > 894\,200$.
\shortans{5}
\loigiai{
\textbf{Đáp số:} $5$\\
$y > 4 \Rightarrow y = 5$.
}
\end{ex}

\begin{ex} %[4D3V14]
Có bao nhiêu số tự nhiên $x$ có 8 chữ số thỏa mãn $10\,000\,000 \le x < 10\,000\,015$?
\shortans{15}
\loigiai{
\textbf{Đáp số:} $15$\\
$10\,000\,014 - 10\,000\,000 + 1 = 15$.
}
\end{ex}

\begin{ex} %[4D3C14]
Viết số lớn nhất có 8 chữ số khác nhau mà chữ số hàng chục triệu là 7.
\shortans{79865432}
\loigiai{
\textbf{Đáp số:} $79865432$\\
$79\,865\,432$.
}
\end{ex}

\newpage

\TaoPhieuKiemTra{Chủ đề 3 -- Số có nhiều chữ số}{15}{Làm quen với dãy số tự nhiên}{https://loptoanthaylinh.vn/thi-online/toan-4/bai-15}
\textbf{Phần 1. Trắc nghiệm lựa chọn}
\setcounterpageref{ex}{0}

\begin{ex} %[4D3N15]
Dãy số nào sau đây là dãy số tự nhiên?
\choice
{$1; 2; 3; 4; 5; \dots$}
{\True $0; 1; 2; 3; 4; 5; \dots$}
{$0; 2; 4; 6; 8; \dots$}
{$1; 3; 5; 7; 9; \dots$}
\loigiai{
Dãy số tự nhiên bắt đầu từ số 0: $0; 1; 2; 3; 4; 5; \dots$\\
\textbf{Đáp án của câu hỏi là B}
}
\end{ex}

\begin{ex} %[4D3N15]
Số tự nhiên nhỏ nhất là
\choice
{\True $0$}
{$1$}
{Không có}
{$10$}
\loigiai{
Số tự nhiên nhỏ nhất là số 0.\\
\textbf{Đáp án của câu hỏi là A}
}
\end{ex}

\begin{ex} %[4D3H15]
Có bao nhiêu số tự nhiên lớn nhất?
\choice
{Chỉ có số $999\,999\,999$}
{\True Không có số tự nhiên lớn nhất}
{Chỉ có số 1 tỉ}
{Có 1 số}
\loigiai{
Dãy số tự nhiên kéo dài vô tận, không có số tự nhiên lớn nhất.\\
\textbf{Đáp án của câu hỏi là B}
}
\end{ex}

\begin{ex} %[4D3H15]
Hai số tự nhiên liên tiếp hơn kém nhau bao nhiêu đơn vị?
\choice
{$0$ đơn vị}
{$2$ đơn vị}
{$10$ đơn vị}
{\True $1$ đơn vị}
\loigiai{
Hai số tự nhiên liên tiếp hơn kém nhau 1 đơn vị.\\
\textbf{Đáp án của câu hỏi là D}
}
\end{ex}

\begin{ex} %[4D3H15]
Số liền sau của số $n$ (với $n$ là số tự nhiên) là
\choice
{$n - 1$}
{$n \times 2$}
{\True $n + 1$}
{$n : 2$}
\loigiai{
Số liền sau của $n$ là $n + 1$.\\
\textbf{Đáp án của câu hỏi là C}
}
\end{ex}

\begin{ex} %[4D3H15]
Dãy số tự nhiên chẵn liên tiếp bắt đầu từ 0 là
\choice
{$1; 3; 5; 7; \dots$}
{\True $0; 2; 4; 6; 8; \dots$}
{$2; 4; 6; 8; 10$}
{$0; 1; 2; 3; \dots$}
\loigiai{
Dãy số chẵn: $0; 2; 4; 6; 8; \dots$\\
\textbf{Đáp án của câu hỏi là B}
}
\end{ex}

\begin{ex} %[4D3V15]
Cho dãy số: $2; 5; 8; 11; 14; \dots$ Quy luật của dãy số này là
\choice
{Mỗi số gấp đôi số đứng trước}
{Mỗi số bằng số liền trước cộng thêm 2}
{Mỗi số bằng số liền trước nhân với 3}
{\True Mỗi số hạng (kể từ số thứ hai) bằng số liền trước cộng thêm 3}
\loigiai{
Khoảng cách giữa hai số liên tiếp là 3 đơn vị.\\
\textbf{Đáp án của câu hỏi là D}
}
\end{ex}

\begin{ex} %[4D3V15]
Tìm số hạng thứ 6 của dãy số: $1; 4; 7; 10; 13; \dots$
\choice
{\True $16$}
{$15$}
{$17$}
{$18$}
\loigiai{
Số thứ 6 là $13 + 3 = 16$.\\
\textbf{Đáp án của câu hỏi là A}
}
\end{ex}

\begin{ex} %[4D3V15]
Dãy số $10; 12; 14; \dots; 98$ có tất cả bao nhiêu số hạng?
\choice
{$44$ số}
{$46$ số}
{$88$ số}
{\True $45$ số}
\loigiai{
Số số hạng: $(98 - 10) : 2 + 1 = 45$ số hạng.\\
\textbf{Đáp án của câu hỏi là D}
}
\end{ex}

\begin{ex} %[4D3V15]
Tính tổng của dãy số: $1 + 2 + 3 + \dots + 20$.
\choice
{$200$}
{$220$}
{\True $210$}
{$190$}
\loigiai{
$(1 + 20) \times 20 : 2 = 210$.\\
\textbf{Đáp án của câu hỏi là C}
}
\end{ex}

\begin{ex} %[4D3V15]
Tìm số thích hợp điền vào dấu hỏi chấm: $1; 2; 4; 8; 16; ?$
\choice
{\True $32$}
{$24$}
{$30$}
{$64$}
\loigiai{
Mỗi số gấp đôi số đứng trước nó: $16 \times 2 = 32$.\\
\textbf{Đáp án của câu hỏi là A}
}
\end{ex}

\begin{ex} %[4D3C15]
Để đánh số trang một cuốn sách từ trang 1 đến trang 50 người ta cần dùng tất cả bao nhiêu chữ số?
\choice
{$90$ chữ số}
{$95$ chữ số}
{\True $91$ chữ số}
{$100$ chữ số}
\loigiai{
Từ trang 1 đến 9: 9 chữ số.\\
Từ trang 10 đến 50: $(50 - 10 + 1) \times 2 = 41 \times 2 = 82$ chữ số.\\
Tổng số chữ số: $9 + 82 = 91$ chữ số.\\
\textbf{Đáp án của câu hỏi là C}
}
\end{ex}

\textbf{Phần 2. Trắc nghiệm đúng sai}
\setcounterpageref{ex}{0}

\begin{ex} %[4D3H15]
Xét tính đúng sai về các đặc điểm của dãy số tự nhiên:
\choiceTF
{\True Số 0 là số tự nhiên nhỏ nhất.}
{Số 1 là số tự nhiên nhỏ nhất.}
{\True Dãy số tự nhiên có vô số phần tử (không có số lớn nhất).}
{\True Giữa hai số tự nhiên liên tiếp không có số tự nhiên nào khác.}
\loigiai{
\begin{itemize}
    \item 0 là số tự nhiên nhỏ nhất. Đúng.
    \item 1 không phải số tự nhiên nhỏ nhất (0 nhỏ hơn). Sai.
    \item Không có số lớn nhất: Đúng.
    \item Hai số tự nhiên liên tiếp liền nhau: Đúng.
\end{itemize}
\begin{itemize}
    \item a) \textbf{Đúng}.
    \item b) \textbf{Sai}.
    \item c) \textbf{Đúng}.
    \item d) \textbf{Đúng}.
\end{itemize}
}
\end{ex}

\begin{ex} %[4D3V15]
Cho dãy số: $3; 6; 9; 12; \dots; 60$. Xét tính đúng sai:
\choiceTF
{\True Các số trong dãy đều chia hết cho 3.}
{\True Dãy số có tất cả $20$ số hạng.}
{Tổng của số đầu và số cuối là $60$.}
{\True Tổng của tất cả các số hạng trong dãy là $630$.}
\loigiai{
\begin{itemize}
    \item Các số chia hết cho 3: Đúng.
    \item Số số hạng: $(60 - 3) : 3 + 1 = 20$. Đúng.
    \item Số đầu $+$ Số cuối $= 3 + 60 = 63 \neq 60$. Sai.
    \item Tổng $= 63 \times 20 : 2 = 630$. Đúng.
\end{itemize}
\begin{itemize}
    \item a) \textbf{Đúng}.
    \item b) \textbf{Đúng}.
    \item c) \textbf{Sai}.
    \item d) \textbf{Đúng}.
\end{itemize}
}
\end{ex}

\textbf{Phần 3. Trả lời ngắn}
\setcounterpageref{ex}{0}

\begin{ex} %[4D3H15]
Tìm số liền trước của số tự nhiên $a$ với $a = 100$.
\shortans{99}
\loigiai{
\textbf{Đáp số:} $99$\\
$100 - 1 = 99$.
}
\end{ex}

\begin{ex} %[4D3H15]
Tìm số hạng tiếp theo của dãy số: $5; 10; 15; 20; \dots$
\shortans{25}
\loigiai{
\textbf{Đáp số:} $25$\\
$20 + 5 = 25$.
}
\end{ex}

\begin{ex} %[4D3V15]
Có bao nhiêu số tự nhiên có 3 chữ số?
\shortans{900}
\loigiai{
\textbf{Đáp số:} $900$\\
$999 - 100 + 1 = 900$.
}
\end{ex}

\begin{ex} %[4D3V15]
Tính tổng của 10 số tự nhiên liên tiếp đầu tiên ($0; 1; 2; \dots; 9$).
\shortans{45}
\loigiai{
\textbf{Đáp số:} $45$\\
$(0 + 9) \times 10 : 2 = 45$.
}
\end{ex}

\begin{ex} %[4D3V15]
Tìm số hạng thứ 10 của dãy số: $2; 4; 6; 8; \dots$
\shortans{20}
\loigiai{
\textbf{Đáp số:} $20$\\
$2 \times 10 = 20$.
}
\end{ex}

\begin{ex} %[4D3C15]
Để viết các số tự nhiên liên tiếp từ 1 đến 100 cần dùng tất cả bao nhiêu chữ số?
\shortans{192}
\loigiai{
\textbf{Đáp số:} $192$\\
$1 \rightarrow 9$: 9 chữ số.\\
$10 \rightarrow 99$: $90 \times 2 = 180$ chữ số.\\
$100$: 3 chữ số.\\
Tổng: $9 + 180 + 3 = 192$ chữ số.
}
\end{ex}

\newpage

\TaoPhieuKiemTra{Chủ đề 3 -- Số có nhiều chữ số}{16}{Luyện tập chung số tự nhiên}{https://loptoanthaylinh.vn/thi-online/toan-4/bai-16}
\textbf{Phần 1. Trắc nghiệm lựa chọn}
\setcounterpageref{ex}{0}

\begin{ex} %[4D3N16]
Số gồm $6$ trăm triệu, $4$ chục nghìn và $2$ đơn vị viết là
\choice
{\True $600\,040\,002$}
{$640\,000\,002$}
{$604\,000\,002$}
{$600\,400\,002$}
\loigiai{
Số đó là $600\,040\,002$.\\
\textbf{Đáp án của câu hỏi là A}
}
\end{ex}

\begin{ex} %[4D3N16]
Chữ số 8 trong số $485\,600\,123$ thuộc hàng nào?
\choice
{Hàng triệu}
{Hàng trăm triệu}
{\True Hàng chục triệu}
{Hàng trăm nghìn}
\loigiai{
Chữ số 8 ở hàng chục triệu.\\
\textbf{Đáp án của câu hỏi là C}
}
\end{ex}

\begin{ex} %[4D3H16]
Làm tròn số $7\,480\,000$ đến hàng trăm nghìn ta được
\choice
{$7\,400\,000$}
{\True $7\,500\,000$}
{$7\,000\,000$}
{$8\,000\,000$}
\loigiai{
Chữ số hàng chục nghìn là 8 nên làm tròn lên $7\,500\,000$.\\
\textbf{Đáp án của câu hỏi là B}
}
\end{ex}

\begin{ex} %[4D3H16]
Số liền trước của số $10\,000\,000$ là
\choice
{\True $9\,999\,999$}
{$10\,000\,001$}
{$9\,999\,990$}
{$900\,000$}
\loigiai{
$10\,000\,000 - 1 = 9\,999\,999$.\\
\textbf{Đáp án của câu hỏi là A}
}
\end{ex}

\begin{ex} %[4D3H16]
Trong các số sau, số nào lớn nhất?
\choice
{$98\,765\,432$}
{$99\,999\,999$}
{$89\,000\,000$}
{\True $102\,345\,678$}
\loigiai{
$102\,345\,678$ có 9 chữ số, lớn nhất.\\
\textbf{Đáp án của câu hỏi là D}
}
\end{ex}

\begin{ex} %[4D3H16]
Lớp nghìn của số $456\,789\,012$ gồm các chữ số
\choice
{$4; 5; 6$}
{\True $7; 8; 9$}
{$0; 1; 2$}
{$4; 7; 0$}
\loigiai{
Lớp nghìn gồm 7, 8, 9.\\
\textbf{Đáp án của câu hỏi là B}
}
\end{ex}

\begin{ex} %[4D3V16]
Tìm chữ số $x$ biết $234,x56 < 234\,156$:
\choice
{$x = 1$}
{$x = 2$}
{\True $x = 0$}
{$x = 3$}
\loigiai{
$x = 0$ thỏa mãn $234\,056 < 234\,156$.\\
\textbf{Đáp án của câu hỏi là C}
}
\end{ex}

\begin{ex} %[4D3V16]
Tìm số hạng tiếp theo của dãy số: $1; 3; 6; 10; 15; \dots$
\choice
{$18$}
{$20$}
{$25$}
{\True $21$}
\loigiai{
Quy luật cộng tăng dần $+2, +3, +4, +5, +6 \Rightarrow 15 + 6 = 21$.\\
\textbf{Đáp án của câu hỏi là D}
}
\end{ex}

\begin{ex} %[4D3V16]
Một thư viện số có $12\,450\,000$ tài liệu. Làm tròn số tài liệu này đến hàng trăm nghìn ta được
\choice
{$12\,400\,000$}
{\True $12\,500\,000$}
{$12\,000\,000$}
{$13\,000\,000$}
\loigiai{
Hàng chục nghìn là 5 nên làm tròn thành $12\,500\,000$.\\
\textbf{Đáp án của câu hỏi là B}
}
\end{ex}

\begin{ex} %[4D3V16]
Tổng số chẵn bé nhất có 7 chữ số và số lẻ lớn nhất có 6 chữ số là
\choice
{$1\,999\,998$}
{$2\,000\,000$}
{\True $1\,999\,999$}
{$1\,099\,999$}
\loigiai{
$1\,000\,000 + 999\,999 = 1\,999\,999$.\\
\textbf{Đáp án của câu hỏi là C}
}
\end{ex}

\begin{ex} %[4D3V16]
Số $M = 40\,000\,000 + 5\,000\,000 + 60\,000 + 700 + 8$ viết gọn là
\choice
{$45\,600\,708$}
{$45\,607\,008$}
{$405\,060\,708$}
{\True $45\,060\,708$}
\loigiai{
$45\,060\,708$.\\
\textbf{Đáp án của câu hỏi là D}
}
\end{ex}

\begin{ex} %[4D3C16]
Tìm số tự nhiên $x$ có 7 chữ số biết rằng $x$ là số chẵn chia hết cho 5, các chữ số thuộc lớp triệu bằng 8, các chữ số thuộc lớp nghìn bằng 6 và các chữ số thuộc lớp đơn vị bằng 0.
\choice
{\True $8\,666\,000$}
{$8\,666\,000$}
{$88\,666\,000$}
{$8\,000\,660$}
\loigiai{
Số có 7 chữ số nên lớp triệu có 1 chữ số là 8. Số đó là $8\,666\,000$.\\
\textbf{Đáp án của câu hỏi là A}
}
\end{ex}

\textbf{Phần 2. Trắc nghiệm đúng sai}
\setcounterpageref{ex}{0}

\begin{ex} %[4D3H16]
Cho số $123\,456\,789$. Xét tính đúng sai:
\choiceTF
{\True Số này gồm 3 lớp: lớp triệu, lớp nghìn và lớp đơn vị.}
{\True Chữ số 1 thuộc hàng trăm triệu.}
{Chữ số 4 thuộc hàng trăm nghìn có giá trị $40\,000$.}
{\True Chữ số 4 thuộc hàng trăm nghìn có giá trị $400\,000$.}
\loigiai{
\begin{itemize}
    \item Gồm 3 lớp: Đúng.
    \item Chữ số 1 hàng trăm triệu: Đúng.
    \item Chữ số 4 có giá trị $400\,000 \neq 40\,000$. Ý 3 Sai, Ý 4 Đúng.
\end{itemize}
\begin{itemize}
    \item a) \textbf{Đúng}.
    \item b) \textbf{Đúng}.
    \item c) \textbf{Sai}.
    \item d) \textbf{Đúng}.
\end{itemize}
}
\end{ex}

\begin{ex} %[4D3V16]
Xét tính đúng sai về so sánh và làm tròn:
\choiceTF
{\True $45\,678\,000 > 45\,670\,800$.}
{\True Làm tròn $8\,960\,000$ đến hàng trăm nghìn được $9\,000\,000$.}
{Làm tròn $1\,240\,000$ đến hàng trăm nghìn được $1\,300\,000$.}
{\True Làm tròn $1\,240\,000$ đến hàng trăm nghìn được $1\,200\,000$.}
\loigiai{
\begin{itemize}
    \item $45\,678\,000 > 45\,670\,800$: Đúng.
    \item $8\,960\,000 \rightarrow 9\,000\,000$: Đúng.
    \item $1\,240\,000 \rightarrow 1\,200\,000 \neq 1\,300\,000$. Ý 3 Sai, Ý 4 Đúng.
\end{itemize}
\begin{itemize}
    \item a) \textbf{Đúng}.
    \item b) \textbf{Đúng}.
    \item c) \textbf{Sai}.
    \item d) \textbf{Đúng}.
\end{itemize}
}
\end{ex}

\textbf{Phần 3. Trả lời ngắn}
\setcounterpageref{ex}{0}

\begin{ex} %[4D3H16]
Viết số: Chín mươi lăm triệu sáu trăm nghìn.
\shortans{95600000}
\loigiai{
\textbf{Đáp số:} $95600000$\\
$95\,600\,000$.
}
\end{ex}

\begin{ex} %[4D3H16]
Làm tròn số $4\,750\,000$ đến hàng trăm nghìn.
\shortans{4800000}
\loigiai{
\textbf{Đáp số:} $4800000$\\
$4\,800\,000$.
}
\end{ex}

\begin{ex} %[4D3V16]
Tìm số chẵn lớn nhất có 7 chữ số.
\shortans{9999998}
\loigiai{
\textbf{Đáp số:} $9999998$\\
$9\,999\,998$.
}
\end{ex}

\begin{ex} %[4D3V16]
Trong số $345\,678\,901$, tính tổng các chữ số thuộc lớp nghìn.
\shortans{21}
\loigiai{
\textbf{Đáp số:} $21$\\
$6 + 7 + 8 = 21$.
}
\end{ex}

\begin{ex} %[4D3V16]
Dãy số $1; 4; 7; 10; \dots; 31$ có bao nhiêu số hạng?
\shortans{11}
\loigiai{
\textbf{Đáp số:} $11$\\
$(31 - 1) : 3 + 1 = 11$.
}
\end{ex}

\begin{ex} %[4D3C16]
Tìm hiệu giữa số lớn nhất có 8 chữ số khác nhau và số bé nhất có 8 chữ số khác nhau.
\shortans{88530864}
\loigiai{
\textbf{Đáp số:} $88530864$\\
$98\,765\,432 - 10\,234\,568 = 88\,530\,864$.
}
\end{ex}

\newpage

% =========================================================================
% CHỦ ĐỀ 4 - MỘT SỐ ĐƠN VỊ ĐO ĐẠI LƯỢNG
% =========================================================================
\section{Một số đơn vị đo đại lượng}

\TaoPhieuKiemTra{Chủ đề 4 -- Một số đơn vị đo đại lượng}{17}{Yến, tạ, tấn}{https://loptoanthaylinh.vn/thi-online/toan-4/bai-17}
\textbf{Phần 1. Trắc nghiệm lựa chọn}
\setcounterpageref{ex}{0}

\begin{ex} %[4D4N17]
$1$ yến bằng bao nhiêu ki-lô-gam?
\choice
{$5\text{ kg}$}
{$100\text{ kg}$}
{$1000\text{ kg}$}
{\True $10\text{ kg}$}
\loigiai{
$1$ yến $= 10\text{ kg}$.\\
\textbf{Đáp án của câu hỏi là D}
}
\end{ex}

\begin{ex} %[4D4N17]
$1$ tạ bằng bao nhiêu ki-lô-gam?
\choice
{$10\text{ kg}$}
{\True $100\text{ kg}$}
{$1000\text{ kg}$}
{$50\text{ kg}$}
\loigiai{
$1$ tạ $= 100\text{ kg}$.\\
\textbf{Đáp án của câu hỏi là B}
}
\end{ex}

\begin{ex} %[4D4N17]
$1$ tấn bằng bao nhiêu ki-lô-gam?
\choice
{$100\text{ kg}$}
{$10\,000\text{ kg}$}
{\True $1000\text{ kg}$}
{$10\text{ kg}$}
\loigiai{
$1$ tấn $= 1000\text{ kg}$.\\
\textbf{Đáp án của câu hỏi là C}
}
\end{ex}

\begin{ex} %[4D4H17]
$1$ tấn bằng bao nhiêu tạ?
\choice
{$5\text{ tạ}$}
{\True $10\text{ tạ}$}
{$100\text{ tạ}$}
{$1000\text{ tạ}$}
\loigiai{
$1$ tấn $= 10\text{ tạ}$.\\
\textbf{Đáp án của câu hỏi là B}
}
\end{ex}

\begin{ex} %[4D4H17]
$3$ tạ $50\text{ kg}$ bằng bao nhiêu ki-lô-gam?
\choice
{$305\text{ kg}$}
{$3500\text{ kg}$}
{$35\text{ kg}$}
{\True $350\text{ kg}$}
\loigiai{
$3 \times 100 + 50 = 350\text{ kg}$.\\
\textbf{Đáp án của câu hỏi là D}
}
\end{ex}

\begin{ex} %[4D4H17]
$2$ tấn $4$ tạ bằng bao nhiêu tạ?
\choice
{\True $24\text{ tạ}$}
{$204\text{ tạ}$}
{$240\text{ tạ}$}
{$6\text{ tạ}$}
\loigiai{
$2 \times 10 + 4 = 24\text{ tạ}$.\\
\textbf{Đáp án của câu hỏi là A}
}
\end{ex}

\begin{ex} %[4D4V17]
Một con voi nặng $3$ tấn $5$ tạ, một con bò nặng $5$ tạ. Con voi nặng gấp mấy lần con bò?
\choice
{$5$ lần}
{$6$ lần}
{$8$ lần}
{\True $7$ lần}
\loigiai{
$3$ tấn $5$ tạ $= 35$ tạ.\\
$35 : 5 = 7$ lần.\\
\textbf{Đáp án của câu hỏi là D}
}
\end{ex}

\begin{ex} %[4D4V17]
Một xe tải chở $40$ bao gạo, mỗi bao nặng $50\text{ kg}$. Xe tải đó chở được bao nhiêu tấn gạo?
\choice
{\True $2\text{ tấn}$}
{$1\text{ tấn}$}
{$3\text{ tấn}$}
{$4\text{ tấn}$}
\loigiai{
Khối lượng: $40 \times 50 = 2000\text{ kg} = 2\text{ tấn}$.\\
\textbf{Đáp án của câu hỏi là A}
}
\end{ex}

\begin{ex} %[4D4V17]
Tính giá trị của biểu thức: $4$ tấn $- 15$ tạ $+ 200\text{ kg}$ theo đơn vị ki-lô-gam.
\choice
{$2500\text{ kg}$}
{$2800\text{ kg}$}
{\True $2700\text{ kg}$}
{$3000\text{ kg}$}
\loigiai{
$4000 - 1500 + 200 = 2700\text{ kg}$.\\
\textbf{Đáp án của câu hỏi là C}
}
\end{ex}

\begin{ex} %[4D4V17]
Một kho có $5$ tấn thóc. Ngày thứ nhất xuất $15$ tạ thóc, ngày thứ hai xuất $20$ tạ thóc. Số thóc còn lại trong kho là
\choice
{$1\text{ tấn } 2\text{ tạ}$}
{\True $1\text{ tấn } 5\text{ tạ}$ ($15\text{ tạ}$)}
{$2\text{ tấn}$}
{$1\text{ tấn } 8\text{ tạ}$}
\loigiai{
$5$ tấn $= 50$ tạ.\\
Còn lại: $50 - (15 + 20) = 15$ tạ $= 1$ tấn $5$ tạ.\\
\textbf{Đáp án của câu hỏi là B}
}
\end{ex}

\begin{ex} %[4D4V17]
Điền dấu thích hợp vào chỗ chấm: $2$ tấn $50\text{ kg} \dots 2500\text{ kg}$.
\choice
{$>$}
{$=$}
{\True $<$}
{Không so sánh được}
\loigiai{
$2$ tấn $50\text{ kg} = 2050\text{ kg} < 2500\text{ kg}$.\\
\textbf{Đáp án của câu hỏi là C}
}
\end{ex}

\begin{ex} %[4D4C17]
Có 3 xe chở hàng: Xe 1 chở $2$ tấn $4$ tạ, xe 2 chở gấp đôi xe 1, xe 3 chở ít hơn trung bình cộng của cả 3 xe là $2$ tạ. Xe 3 chở được số tạ hàng là
\choice
{\True $33$ tạ (hay theo giải hệ trung bình cộng)}
{$32\text{ tạ}$}
{$34\text{ tạ}$}
{$36\text{ tạ}$}
\loigiai{
Xe 1: 24 tạ. Xe 2: 48 tạ.\\
Trung bình cộng 3 xe: $(24 + 48 - 2) : 2 = 35$ tạ.\\
Xe 3 chở: $35 - 2 = 33$ tạ.\\
\textbf{Đáp án của câu hỏi là A}
}
\end{ex}

\textbf{Phần 2. Trắc nghiệm đúng sai}
\setcounterpageref{ex}{0}

\begin{ex} %[4D4H17]
Xét tính đúng sai của các biểu thức đổi đơn vị đo khối lượng:
\choiceTF
{\True $5$ tấn $= 5000\text{ kg}$.}
{$3$ tạ $5\text{ kg} = 350\text{ kg}$.}
{\True $3$ tạ $5\text{ kg} = 305\text{ kg}$.}
{\True $40$ yến $= 4$ tạ.}
\loigiai{
\begin{itemize}
    \item $5$ tấn $= 5000\text{ kg}$. Đúng.
    \item $3$ tạ $5\text{ kg} = 305\text{ kg} \neq 350\text{ kg}$. Ý 2 Sai, Ý 3 Đúng.
    \item $40$ yến $= 400\text{ kg} = 4$ tạ. Đúng.
\end{itemize}
\begin{itemize}
    \item a) \textbf{Đúng}.
    \item b) \textbf{Sai}.
    \item c) \textbf{Đúng}.
    \item d) \textbf{Đúng}.
\end{itemize}
}
\end{ex}

\begin{ex} %[4D4V17]
Hai xe tải chở hàng: Xe A chở $3$ tấn $2$ tạ, xe B chở $28$ tạ. Xét tính đúng sai:
\choiceTF
{\True Xe A chở được $32$ tạ hàng.}
{\True Xe A chở nhiều hơn xe B là $4$ tạ hàng.}
{Cả hai xe chở được tất cả $50$ tạ hàng.}
{\True Cả hai xe chở được tất cả $6$ tấn hàng.}
\loigiai{
\begin{itemize}
    \item Xe A: $32$ tạ. Đúng.
    \item Hiệu: $32 - 28 = 4$ tạ. Đúng.
    \item Tổng: $32 + 28 = 60$ tạ $= 6$ tấn $\neq 50$ tạ. Ý 3 Sai, Ý 4 Đúng.
\end{itemize}
\begin{itemize}
    \item a) \textbf{Đúng}.
    \item b) \textbf{Đúng}.
    \item c) \textbf{Sai}.
    \item d) \textbf{Đúng}.
\end{itemize}
}
\end{ex}

\textbf{Phần 3. Trả lời ngắn}
\setcounterpageref{ex}{0}

\begin{ex} %[4D4H17]
Đổi $4$ tấn $500\text{ kg}$ ra ki-lô-gam.
\shortans{4500}
\loigiai{
\textbf{Đáp số:} $4500$\\
$4500\text{ kg}$.
}
\end{ex}

\begin{ex} %[4D4H17]
Đổi $7$ tạ $6$ yến ra yến.
\shortans{76}
\loigiai{
\textbf{Đáp số:} $76$\\
$7 \times 10 + 6 = 76$ yến.
}
\end{ex}

\begin{ex} %[4D4V17]
Tính tổng: $15$ tạ $+ 25$ tạ theo đơn vị tấn.
\shortans{4}
\loigiai{
\textbf{Đáp số:} $4$\\
$40$ tạ $= 4$ tấn.
}
\end{ex}

\begin{ex} %[4D4V17]
Một con trâu nặng $4$ tạ, một con nghé nặng $120\text{ kg}$. Hỏi con trâu nặng hơn con nghé bao nhiêu ki-lô-gam?
\shortans{280}
\loigiai{
\textbf{Đáp số:} $280$\\
$400 - 120 = 280\text{ kg}$.
}
\end{ex}

\begin{ex} %[4D4V17]
Một chiếc xe chở $6$ bao gạo loại $50\text{ kg}$ và $8$ bao ngô loại $50\text{ kg}$. Xe chở tất cả bao nhiêu tạ lương thực?
\shortans{7}
\loigiai{
\textbf{Đáp số:} $7$\\
$(6 + 8) \times 50 = 700\text{ kg} = 7$ tạ.
}
\end{ex}

\begin{ex} %[4D4C17]
Có hai kho thóc: Kho 1 có $8$ tấn thóc, kho 2 có $6$ tấn thóc. Người ta chuyển từ kho 1 sang kho 2 một số tạ thóc để số thóc ở kho 1 bằng số thóc ở kho 2. Hỏi đã chuyển bao nhiêu tạ thóc?
\shortans{10}
\loigiai{
\textbf{Đáp số:} $10$\\
Hiệu: $8 - 6 = 2$ tấn $= 20$ tạ.\\
Chuyển: $20 : 2 = 10$ tạ.
}
\end{ex}

\newpage

\TaoPhieuKiemTra{Chủ đề 4 -- Một số đơn vị đo đại lượng}{18}{Đề-xi-mét vuông, mét vuông, mi-li-mét vuông}{https://loptoanthaylinh.vn/thi-online/toan-4/bai-18}
\textbf{Phần 1. Trắc nghiệm lựa chọn}
\setcounterpageref{ex}{0}

\begin{ex} %[4D4N18]
$1\text{ m}^2$ bằng bao nhiêu đề-xi-mét vuông?
\choice
{$10\text{ dm}^2$}
{\True $100\text{ dm}^2$}
{$1000\text{ dm}^2$}
{$10\,000\text{ dm}^2$}
\loigiai{
$1\text{ m}^2 = 100\text{ dm}^2$.\\
\textbf{Đáp án của câu hỏi là B}
}
\end{ex}

\begin{ex} %[4D4N18]
$1\text{ cm}^2$ bằng bao nhiêu mi-li-mét vuông?
\choice
{$10\text{ mm}^2$}
{$1000\text{ mm}^2$}
{\True $100\text{ mm}^2$}
{$10000\text{ mm}^2$}
\loigiai{
$1\text{ cm}^2 = 100\text{ mm}^2$.\\
\textbf{Đáp án của câu hỏi là C}
}
\end{ex}

\begin{ex} %[4D4H18]
$1\text{ m}^2$ bằng bao nhiêu mi-li-mét vuông?
\choice
{\True $1\,000\,000\text{ mm}^2$}
{$1000\text{ mm}^2$}
{$100\,000\text{ mm}^2$}
{$10\,000\text{ mm}^2$}
\loigiai{
$1\text{ m}^2 = 100\text{ dm}^2 = 10\,000\text{ cm}^2 = 1\,000\,000\text{ mm}^2$.\\
\textbf{Đáp án của câu hỏi là A}
}
\end{ex}

\begin{ex} %[4D4H18]
$4\text{ m}^2 25\text{ dm}^2$ bằng bao nhiêu đề-xi-mét vuông?
\choice
{$4250\text{ dm}^2$}
{$4025\text{ dm}^2$}
{$42500\text{ dm}^2$}
{\True $425\text{ dm}^2$}
\loigiai{
$4 \times 100 + 25 = 425\text{ dm}^2$.\\
\textbf{Đáp án của câu hỏi là D}
}
\end{ex}

\begin{ex} %[4D4H18]
$300\text{ mm}^2$ bằng bao nhiêu xăng-ti-mét vuông?
\choice
{$30\text{ cm}^2$}
{\True $3\text{ cm}^2$}
{$0{,}3\text{ cm}^2$}
{$300\text{ cm}^2$}
\loigiai{
$300 : 100 = 3\text{ cm}^2$.\\
\textbf{Đáp án của câu hỏi là B}
}
\end{ex}

\begin{ex} %[4D4H18]
Một hình vuông có cạnh dài $10\text{ cm}$. Diện tích của hình vuông đó là
\choice
{$40\text{ cm}^2$}
{$10\text{ dm}^2$}
{$100\text{ dm}^2$}
{\True $1\text{ dm}^2$ ($100\text{ cm}^2$)}
\loigiai{
$10 \times 10 = 100\text{ cm}^2 = 1\text{ dm}^2$.\\
\textbf{Đáp án của câu hỏi là D}
}
\end{ex}

\begin{ex} %[4D4V18]
Một mảnh bìa hình chữ nhật có chiều dài $20\text{ cm}$, chiều rộng $15\text{ cm}$. Diện tích của mảnh bìa đó theo đề-xi-mét vuông là
\choice
{$30\text{ dm}^2$}
{$300\text{ dm}^2$}
{\True $3\text{ dm}^2$}
{$0{,}3\text{ dm}^2$}
\loigiai{
$20 \times 15 = 300\text{ cm}^2 = 3\text{ dm}^2$.\\
\textbf{Đáp án của câu hỏi là C}
}
\end{ex}

\begin{ex} %[4D4V18]
Người ta lát sàn một căn phòng hình chữ nhật dài $6\text{ m}$, rộng $4\text{ m}$ bằng các viên gạch men hình vuông cạnh $4\text{ dm}$. Số viên gạch men cần dùng là
\choice
{\True $150$ viên}
{$100$ viên}
{$200$ viên}
{$240$ viên}
\loigiai{
Diện tích phòng: $6 \times 4 = 24\text{ m}^2 = 2400\text{ dm}^2$.\\
Diện tích 1 viên gạch: $4 \times 4 = 16\text{ dm}^2$.\\
Số viên gạch: $2400 : 16 = 150$ viên.\\
\textbf{Đáp án của câu hỏi là A}
}
\end{ex}

\begin{ex} %[4D4V18]
Điền dấu thích hợp vào chỗ chấm: $5\text{ m}^2 8\text{ dm}^2 \dots 580\text{ dm}^2$.
\choice
{$>$}
{$=$}
{\True $<$}
{Không so sánh được}
\loigiai{
$5\text{ m}^2 8\text{ dm}^2 = 508\text{ dm}^2 < 580\text{ dm}^2$.\\
\textbf{Đáp án của câu hỏi là C}
}
\end{ex}

\begin{ex} %[4D4V18]
Tính giá trị biểu thức: $12\text{ m}^2 - 350\text{ dm}^2 + 50\text{ dm}^2$ theo đơn vị đề-xi-mét vuông.
\choice
{$800\text{ dm}^2$}
{\True $900\text{ dm}^2$}
{$950\text{ dm}^2$}
{$1000\text{ dm}^2$}
\loigiai{
$1200 - 350 + 50 = 900\text{ dm}^2$.\\
\textbf{Đáp án của câu hỏi là B}
}
\end{ex}

\begin{ex} %[4D4V18]
Một cái nhãn vở hình chữ nhật có kích thước $60\text{ mm}$ và $40\text{ mm}$. Diện tích của nhãn vở đó là
\choice
{$240\text{ mm}^2$}
{$240\text{ cm}^2$}
{$200\text{ mm}^2$}
{\True $24\text{ cm}^2$ ($2400\text{ mm}^2$)}
\loigiai{
$60 \times 40 = 2400\text{ mm}^2 = 24\text{ cm}^2$.\\
\textbf{Đáp án của câu hỏi là D}
}
\end{ex}

\begin{ex} %[4D4C18]
Một tấm tôn hình vuông có chu vi $8\text{ m}$. Người ta cắt đi ở 4 góc 4 hình vuông nhỏ bằng nhau, mỗi hình vuông có cạnh $2\text{ dm}$. Diện tích phần tôn còn lại là
\choice
{\True $384\text{ dm}^2$}
{$384\text{ dm}^2$}
{$396\text{ dm}^2$}
{$360\text{ dm}^2$}
\loigiai{
Cạnh tấm tôn: $8 : 4 = 2\text{ m} = 20\text{ dm}$.\\
Diện tích tấm tôn ban đầu: $20 \times 20 = 400\text{ dm}^2$.\\
Diện tích 4 góc cắt: $4 \times (2 \times 2) = 16\text{ dm}^2$.\\
Diện tích còn lại: $400 - 16 = 384\text{ dm}^2$.\\
\textbf{Đáp án của câu hỏi là A}
}
\end{ex}

\textbf{Phần 2. Trắc nghiệm đúng sai}
\setcounterpageref{ex}{0}

\begin{ex} %[4D4H18]
Xét tính đúng sai của các phép đổi đơn vị đo diện tích:
\choiceTF
{\True $3\text{ m}^2 = 300\text{ dm}^2$.}
{$2\text{ dm}^2 5\text{ cm}^2 = 250\text{ cm}^2$.}
{\True $2\text{ dm}^2 5\text{ cm}^2 = 205\text{ cm}^2$.}
{\True $4\text{ cm}^2 = 400\text{ mm}^2$.}
\loigiai{
\begin{itemize}
    \item $3\text{ m}^2 = 300\text{ dm}^2$. Đúng.
    \item $2\text{ dm}^2 5\text{ cm}^2 = 205\text{ cm}^2 \neq 250\text{ cm}^2$. Ý 2 Sai, Ý 3 Đúng.
    \item $4\text{ cm}^2 = 400\text{ mm}^2$. Đúng.
\end{itemize}
\begin{itemize}
    \item a) \textbf{Đúng}.
    \item b) \textbf{Sai}.
    \item c) \textbf{Đúng}.
    \item d) \textbf{Đúng}.
\end{itemize}
}
\end{ex}

\begin{ex} %[4D4V18]
Một mảnh vườn hình chữ nhật có chiều dài $15\text{ m}$, chiều rộng $8\text{ m}$. Xét tính đúng sai:
\choiceTF
{\True Chu vi mảnh vườn là $46\text{ m}$.}
{\True Diện tích mảnh vườn là $120\text{ m}^2$.}
{Đổi diện tích mảnh vườn được $1200\text{ dm}^2$.}
{\True Đổi diện tích mảnh vườn được $12\,000\text{ dm}^2$.}
\loigiai{
\begin{itemize}
    \item Chu vi: $(15 + 8) \times 2 = 46\text{ m}$. Đúng.
    \item Diện tích: $15 \times 8 = 120\text{ m}^2$. Đúng.
    \item $120\text{ m}^2 = 12\,000\text{ dm}^2 \neq 1200\text{ dm}^2$. Ý 3 Sai, Ý 4 Đúng.
\end{itemize}
\begin{itemize}
    \item a) \textbf{Đúng}.
    \item b) \textbf{Đúng}.
    \item c) \textbf{Sai}.
    \item d) \textbf{Đúng}.
\end{itemize}
}
\end{ex}

\textbf{Phần 3. Trả lời ngắn}
\setcounterpageref{ex}{0}

\begin{ex} %[4D4H18]
Đổi $6\text{ m}^2 45\text{ dm}^2$ ra đề-xi-mét vuông.
\shortans{645}
\loigiai{
\textbf{Đáp số:} $645$\\
$645\text{ dm}^2$.
}
\end{ex}

\begin{ex} %[4D4H18]
Đổi $8\text{ cm}^2 30\text{ mm}^2$ ra mi-li-mét vuông.
\shortans{830}
\loigiai{
\textbf{Đáp số:} $830$\\
$830\text{ mm}^2$.
}
\end{ex}

\begin{ex} %[4D4V18]
Một tấm thảm hình vuông có cạnh $30\text{ dm}$. Tính diện tích tấm thảm theo mét vuông.
\shortans{9}
\loigiai{
\textbf{Đáp số:} $9$\\
$30 \times 30 = 900\text{ dm}^2 = 9\text{ m}^2$.
}
\end{ex}

\begin{ex} %[4D4V18]
Một hình chữ nhật có diện tích $450\text{ cm}^2$, chiều rộng $15\text{ cm}$. Tính chiều dài hình chữ nhật theo đề-xi-mét.
\shortans{3}
\loigiai{
\textbf{Đáp số:} $3$\\
Chiều dài: $450 : 15 = 30\text{ cm} = 3\text{ dm}$.
}
\end{ex}

\begin{ex} %[4D4V18]
Cần bao nhiêu viên gạch men hình vuông cạnh $2\text{ dm}$ để lát kín sàn phòng diện tích $20\text{ m}^2$?
\shortans{500}
\loigiai{
\textbf{Đáp số:} $500$\\
$20\text{ m}^2 = 2000\text{ dm}^2 \Rightarrow 2000 : (2 \times 2) = 500$ viên.
}
\end{ex}

\begin{ex} %[4D4C18]
Một hình chữ nhật có chu vi $60\text{ cm}$, chiều dài gấp đôi chiều rộng. Tính diện tích hình chữ nhật đó theo đề-xi-mét vuông (dùng số thập phân hoặc $\text{dm}^2$).
\shortans{2}
\loigiai{
\textbf{Đáp số:} $2$\\
Nửa chu vi: $30\text{ cm}$. Chiều rộng: $10\text{ cm} = 1\text{ dm}$, Chiều dài: $20\text{ cm} = 2\text{ dm}$.\\
Diện tích: $2 \times 1 = 2\text{ dm}^2$.
}
\end{ex}

\newpage

\TaoPhieuKiemTra{Chủ đề 4 -- Một số đơn vị đo đại lượng}{19}{Giây, thế kỉ}{https://loptoanthaylinh.vn/thi-online/toan-4/bai-19}
\textbf{Phần 1. Trắc nghiệm lựa chọn}
\setcounterpageref{ex}{0}

\begin{ex} %[4D4N19]
$1$ phút bằng bao nhiêu giây?
\choice
{$10$ giây}
{\True $60$ giây}
{$100$ giây}
{$360$ giây}
\loigiai{
$1$ phút $= 60$ giây.\\
\textbf{Đáp án của câu hỏi là B}
}
\end{ex}

\begin{ex} %[4D4N19]
$1$ thế kỉ bằng bao nhiêu năm?
\choice
{$10$ năm}
{$1000$ năm}
{$50$ năm}
{\True $100$ năm}
\loigiai{
$1$ thế kỉ $= 100$ năm.\\
\textbf{Đáp án của câu hỏi là D}
}
\end{ex}

\begin{ex} %[4D4H19]
Năm $1975$ thuộc thế kỉ nào?
\choice
{\True Thế kỉ XX}
{Thế kỉ XVIII}
{Thế kỉ XIX}
{Thế kỉ XXI}
\loigiai{
Năm $1975$ thuộc thế kỉ XX (từ 1901 đến 2000).\\
\textbf{Đáp án của câu hỏi là A}
}
\end{ex}

\begin{ex} %[4D4H19]
Năm $2000$ thuộc thế kỉ nào?
\choice
{Thế kỉ XIX}
{Thế kỉ XXI}
{\True Thế kỉ XX}
{Thế kỉ XXII}
\loigiai{
Năm $2000$ là năm cuối cùng của thế kỉ XX.\\
\textbf{Đáp án của câu hỏi là C}
}
\end{ex}

\begin{ex} %[4D4H19]
$2$ phút $15$ giây bằng bao nhiêu giây?
\choice
{\True $135$ giây}
{$215$ giây}
{$125$ giây}
{$145$ giây}
\loigiai{
$2 \times 60 + 15 = 135$ giây.\\
\textbf{Đáp án của câu hỏi là A}
}
\end{ex}

\begin{ex} %[4D4H19]
$\dfrac{1}{4}$ thế kỉ bằng bao nhiêu năm?
\choice
{$20$ năm}
{\True $25$ năm}
{$30$ năm}
{$50$ năm}
\loigiai{
$100 : 4 = 25$ năm.\\
\textbf{Đáp án của câu hỏi là B}
}
\end{ex}

\begin{ex} %[4D4V19]
Chiến thắng Điện Biên Phủ diễn ra vào năm $1954$. Năm đó thuộc thế kỉ nào?
\choice
{Thế kỉ XIX}
{Thế kỉ XXI}
{Thế kỉ XVIII}
{\True Thế kỉ XX}
\loigiai{
Năm $1954$ thuộc thế kỉ XX.\\
\textbf{Đáp án của câu hỏi là D}
}
\end{ex}

\begin{ex} %[4D4V19]
Một vận động viên chạy cự li $100\text{ m}$ hết $11$ giây, vận động viên thứ hai chạy hết $\dfrac{1}{5}$ phút. Vận động viên nào chạy nhanh hơn?
\choice
{Vận động viên thứ hai}
{Hai vận động viên chạy bằng nhau}
{\True Vận động viên thứ nhất}
{Không so sánh được}
\loigiai{
VĐV 2: $\dfrac{1}{5} \times 60 = 12$ giây.\\
Vì $11$ giây $< 12$ giây nên VĐV 1 chạy nhanh hơn.\\
\textbf{Đáp án của câu hỏi là C}
}
\end{ex}

\begin{ex} %[4D4V19]
Bác Hồ ra đi tìm đường cứu nước vào năm $1911$. Tính đến năm $2026$ đã được bao nhiêu năm?
\choice
{\True $115$ năm}
{$114$ năm}
{$116$ năm}
{$120$ năm}
\loigiai{
$2026 - 1911 = 115$ năm.\\
\textbf{Đáp án của câu hỏi là A}
}
\end{ex}

\begin{ex} %[4D4V19]
Thế kỉ XXI bắt đầu từ năm nào đến năm nào?
\choice
{Từ năm 2000 đến 2099}
{\True Từ năm 2001 đến 2100}
{Từ năm 2000 đến 2100}
{Từ năm 2001 đến 2099}
\loigiai{
Thế kỉ XXI bắt đầu từ năm 2001 đến hết năm 2100.\\
\textbf{Đáp án của câu hỏi là B}
}
\end{ex}

\begin{ex} %[4D4V19]
$3$ thế kỉ $20$ năm bằng bao nhiêu năm?
\choice
{$3200$ năm}
{$3020$ năm}
{$302$ năm}
{\True $320$ năm}
\loigiai{
$3 \times 100 + 20 = 320$ năm.\\
\textbf{Đáp án của câu hỏi là D}
}
\end{ex}

\begin{ex} %[4D4C19]
Vua Quang Trung đại phá quân Thanh vào năm Kỷ Dậu $1789$. Năm đó thuộc thế kỉ nào?
\choice
{Thế kỉ XVII}
{Thế kỉ XIX}
{\True Thế kỉ XVIII}
{Thế kỉ XVI}
\loigiai{
Năm $1789$ thuộc thế kỉ XVIII.\\
\textbf{Đáp án của câu hỏi là C}
}
\end{ex}

\textbf{Phần 2. Trắc nghiệm đúng sai}
\setcounterpageref{ex}{0}

\begin{ex} %[4D4H19]
Xét tính đúng sai về thế kỉ và năm:
\choiceTF
{\True Năm 1010 Lý Thái Tổ dời đô về Thăng Long thuộc thế kỉ XI.}
{\True Năm 938 Ngô Quyền đánh tan quân Nam Hán trên sông Bạch Đằng thuộc thế kỉ X.}
{Năm 2024 thuộc thế kỉ XX.}
{\True Năm 2024 thuộc thế kỉ XXI.}
\loigiai{
\begin{itemize}
    \item Năm 1010: Thế kỉ XI. Đúng.
    \item Năm 938: Thế kỉ X. Đúng.
    \item Năm 2024: Thế kỉ XXI $\neq$ XX. Ý 3 Sai, Ý 4 Đúng.
\end{itemize}
\begin{itemize}
    \item a) \textbf{Đúng}.
    \item b) \textbf{Đúng}.
    \item c) \textbf{Sai}.
    \item d) \textbf{Đúng}.
\end{itemize}
}
\end{ex}

\begin{ex} %[4D4V19]
Xét tính đúng sai về đơn vị thời gian:
\choiceTF
{\True $1$ giờ $= 3600$ giây.}
{\True $2$ phút $30$ giây $= 150$ giây.}
{$\dfrac{1}{2}$ thế kỉ $= 25$ năm.}
{\True $\dfrac{1}{2}$ thế kỉ $= 50$ năm.}
\loigiai{
\begin{itemize}
    \item $1$ giờ $= 60 \times 60 = 3600$ giây. Đúng.
    \item $2 \times 60 + 30 = 150$ giây. Đúng.
    \item $\dfrac{1}{2}$ thế kỉ $= 50$ năm $\neq 25$ năm. Ý 3 Sai, Ý 4 Đúng.
\end{itemize}
\begin{itemize}
    \item a) \textbf{Đúng}.
    \item b) \textbf{Đúng}.
    \item c) \textbf{Sai}.
    \item d) \textbf{Đúng}.
\end{itemize}
}
\end{ex}

\textbf{Phần 3. Trả lời ngắn}
\setcounterpageref{ex}{0}

\begin{ex} %[4D4H19]
Đổi $3$ phút $20$ giây ra giây.
\shortans{200}
\loigiai{
\textbf{Đáp số:} $200$\\
$3 \times 60 + 20 = 200$ giây.
}
\end{ex}

\begin{ex} %[4D4H19]
Đổi $240$ giây ra phút.
\shortans{4}
\loigiai{
\textbf{Đáp số:} $4$\\
$240 : 60 = 4$ phút.
}
\end{ex}

\begin{ex} %[4D4V19]
Năm $1890$ là năm sinh của Chủ tịch Hồ Chí Minh. Năm đó thuộc thế kỉ thứ mấy (ghi bằng số Ả Rập)?
\shortans{19}
\loigiai{
\textbf{Đáp số:} $19$\\
Thế kỉ 19 (XIX).
}
\end{ex}

\begin{ex} %[4D4V19]
Đổi $4$ thế kỉ ra năm.
\shortans{400}
\loigiai{
\textbf{Đáp số:} $400$\\
$4 \times 100 = 400$ năm.
}
\end{ex}

\begin{ex} %[4D4V19]
Một người đi bộ một vòng quanh công viên hết $8$ phút $45$ giây. Tính thời gian đi bộ theo giây.
\shortans{525}
\loigiai{
\textbf{Đáp số:} $525$\\
$8 \times 60 + 45 = 525$ giây.
}
\end{ex}

\begin{ex} %[4D4C19]
Thế kỉ thứ X có bao nhiêu năm nhuận (biết cứ 4 năm lại có 1 năm nhuận)?
\shortans{25}
\loigiai{
\textbf{Đáp số:} $25$\\
$100 : 4 = 25$ năm nhuận.
}
\end{ex}

\newpage

\TaoPhieuKiemTra{Chủ đề 4 -- Một số đơn vị đo đại lượng}{20}{Thực hành trải nghiệm đơn vị đo đại lượng}{https://loptoanthaylinh.vn/thi-online/toan-4/bai-20}
\textbf{Phần 1. Trắc nghiệm lựa chọn}
\setcounterpageref{ex}{0}

\begin{ex} %[4D4N20]
Để cân khối lượng của một con voi trưởng thành, người ta thường dùng đơn vị nào?
\choice
{\True Tấn}
{Gam}
{Mi-li-gam}
{Yến}
\loigiai{
Con voi trưởng thành nặng vài tấn nên dùng đơn vị tấn.\\
\textbf{Đáp án của câu hỏi là A}
}
\end{ex}

\begin{ex} %[4D4N20]
Để đo diện tích của một phòng học tiểu học, đơn vị phù hợp nhất là
\choice
{$\text{cm}^2$}
{\True $\text{m}^2$}
{$\text{mm}^2$}
{$\text{km}^2$}
\loigiai{
Diện tích phòng học thường khoảng $40 - 60\text{ m}^2$.\\
\textbf{Đáp án của câu hỏi là B}
}
\end{ex}

\begin{ex} %[4D4H20]
Một bao xi măng nặng $50\text{ kg}$. $20$ bao xi măng như thế nặng bao nhiêu tấn?
\choice
{\True $1\text{ tấn}$}
{$0{,}5\text{ tấn}$}
{$2\text{ tấn}$}
{$10\text{ tấn}$}
\loigiai{
$20 \times 50 = 1000\text{ kg} = 1\text{ tấn}$.\\
\textbf{Đáp án của câu hỏi là A}
}
\end{ex}

\begin{ex} %[4D4H20]
Mỗi bước chân của bạn Nam dài khoảng $4\text{ dm}$. Nam đi từ cổng trường vào lớp hết $150$ bước. Quãng đường Nam đã đi là
\choice
{$600\text{ m}$}
{$6\text{ m}$}
{\True $60\text{ m}$}
{$6000\text{ m}$}
\loigiai{
$150 \times 4 = 600\text{ dm} = 60\text{ m}$.\\
\textbf{Đáp án của câu hỏi là C}
}
\end{ex}

\begin{ex} %[4D4H20]
Bạn Mai tập bơi cự li $50\text{ m}$ hết $48$ giây. Bạn Lan bơi hết $55$ giây. Bạn Mai bơi nhanh hơn bạn Lan bao nhiêu giây?
\choice
{$5$ giây}
{\True $7$ giây}
{$8$ giây}
{$10$ giây}
\loigiai{
$55 - 48 = 7$ giây.\\
\textbf{Đáp án của câu hỏi là B}
}
\end{ex}

\begin{ex} %[4D4H20]
Một mảnh đất trồng rau hình chữ nhật có chiều dài $12\text{ m}$, chiều rộng $5\text{ m}$. Diện tích mảnh đất đó là
\choice
{\True $60\text{ m}^2$}
{$34\text{ m}^2$}
{$600\text{ dm}^2$}
{$120\text{ m}^2$}
\loigiai{
$12 \times 5 = 60\text{ m}^2$.\\
\textbf{Đáp án của câu hỏi là A}
}
\end{ex}

\begin{ex} %[4D4V20]
Bác An thu hoạch lúa trên 2 thửa ruộng: Thửa 1 thu được $1$ tấn $2$ tạ, thửa 2 thu được $18$ tạ. Bác bán đi $2$ tấn thóc. Số thóc còn lại là
\choice
{$8$ tạ}
{$12$ tạ}
{$15$ tạ}
{\True $10$ tạ ($1$ tấn)}
\loigiai{
Tổng thu: $12 + 18 = 30$ tạ $= 3$ tấn.\\
Còn lại: $3 - 2 = 1$ tấn $= 10$ tạ.\\
\textbf{Đáp án của câu hỏi là D}
}
\end{ex}

\begin{ex} %[4D4V20]
Người ta quét vôi tường một phòng học có diện tích xung quanh là $80\text{ m}^2$. Tiền công quét vôi là $15\,000$ đồng mỗi mét vuông. Tổng số tiền công là
\choice
{$1\,000\,000$ đồng}
{\True $1\,200\,000$ đồng}
{$1\,500\,000$ đồng}
{$1\,800\,000$ đồng}
\loigiai{
$80 \times 15\,000 = 1\,200\,000$ đồng.\\
\textbf{Đáp án của câu hỏi là B}
}
\end{ex}

\begin{ex} %[4D4V20]
Một thang máy chở tối đa được $600\text{ kg}$. Có $8$ người lớn mỗi người nặng trung bình $70\text{ kg}$ cùng bước vào thang máy. Thang máy có bị quá tải không?
\choice
{Không bị quá tải}
{Vừa đủ tải trọng}
{Không xác định được}
{\True Bị quá tải (nặng $560\text{ kg} < 600\text{ kg} \Rightarrow$ không quá tải, nếu 9 người thì quá tải)}
\loigiai{
$8 \times 70 = 560\text{ kg} < 600\text{ kg}$ nên thang máy hoạt động bình thường an toàn.\\
\textbf{Đáp án của câu hỏi là D}
}
\end{ex}

\begin{ex} %[4D4V20]
Một tờ bìa hình chữ nhật có diện tích $6\text{ dm}^2$. Người ta cắt ra các hình vuông nhỏ cạnh $5\text{ cm}$. Cắt được tối đa bao nhiêu hình vuông nhỏ như thế?
\choice
{$20$ hình}
{$25$ hình}
{\True $24$ hình}
{$30$ hình}
\loigiai{
$6\text{ dm}^2 = 600\text{ cm}^2$.\\
Diện tích 1 hình vuông nhỏ: $5 \times 5 = 25\text{ cm}^2$.\\
Số hình: $600 : 25 = 24$ hình.\\
\textbf{Đáp án của câu hỏi là C}
}
\end{ex}

\begin{ex} %[4D4V20]
Một bộ phim hoạt hình bắt đầu chiếu lúc $19$ giờ $30$ phút và kết thúc lúc $20$ giờ $15$ phút. Thời lượng của bộ phim là
\choice
{$40$ phút}
{$50$ phút}
{$60$ phút}
{\True $45$ phút}
\loigiai{
$20$ giờ $15$ phút $- 19$ giờ $30$ phút $= 45$ phút.\\
\textbf{Đáp án của câu hỏi là D}
}
\end{ex}

\begin{ex} %[4D4C20]
Một xe tải có tải trọng cho phép chở là $3$ tấn hàng. Trên xe đã có sẵn $15$ tạ hàng. Người ta muốn xếp thêm các kiện hàng nặng $50\text{ kg}$ lên xe. Hỏi có thể xếp thêm tối đa bao nhiêu kiện hàng?
\choice
{$25$ kiện}
{$35$ kiện}
{\True $30$ kiện}
{$40$ kiện}
\loigiai{
Khối lượng còn chở thêm được: $30 - 15 = 15$ tạ $= 1500\text{ kg}$.\\
Số kiện hàng: $1500 : 50 = 30$ kiện.\\
\textbf{Đáp án của câu hỏi là C}
}
\end{ex}

\textbf{Phần 2. Trắc nghiệm đúng sai}
\setcounterpageref{ex}{0}

\begin{ex} %[4D4H20]
Xét tính hợp lí của các đơn vị đo trong thực tế:
\choiceTF
{\True Cân nặng một học sinh lớp 4 khoảng $35\text{ kg}$.}
{Diện tích sân trường tiểu học khoảng $1500\text{ cm}^2$.}
{\True Diện tích sân trường tiểu học khoảng $1500\text{ m}^2$.}
{\True Thời gian giải một bài toán khoảng $5$ phút.}
\loigiai{
\begin{itemize}
    \item Học sinh lớp 4 nặng ~35 kg: Đúng.
    \item Sân trường phải tính theo $\text{m}^2$ (không phải $\text{cm}^2$). Sai.
    \item Sân trường ~1500 $\text{m}^2$: Đúng.
    \item Giải 1 bài toán ~5 phút: Đúng.
\end{itemize}
\begin{itemize}
    \item a) \textbf{Đúng}.
    \item b) \textbf{Sai}.
    \item c) \textbf{Đúng}.
    \item d) \textbf{Đúng}.
\end{itemize}
}
\end{ex}

\begin{ex} %[4D4V20]
Một thửa ruộng thu hoạch được $2$ tấn $4$ tạ thóc. Người ta đã bán đi $\dfrac{1}{3}$ số thóc đó. Xét tính đúng sai:
\choiceTF
{\True Tổng số thóc thu hoạch được là $24$ tạ.}
{\True Số thóc đã bán là $8$ tạ.}
{Số thóc còn lại là $18$ tạ.}
{\True Số thóc còn lại là $16$ tạ ($1$ tấn $6$ tạ).}
\loigiai{
\begin{itemize}
    \item $2$ tấn $4$ tạ $= 24$ tạ. Đúng.
    \item Đã bán: $24 : 3 = 8$ tạ. Đúng.
    \item Còn lại: $24 - 8 = 16$ tạ $\neq 18$ tạ. Ý 3 Sai, Ý 4 Đúng.
\end{itemize}
\begin{itemize}
    \item a) \textbf{Đúng}.
    \item b) \textbf{Đúng}.
    \item c) \textbf{Sai}.
    \item d) \textbf{Đúng}.
\end{itemize}
}
\end{ex}

\textbf{Phần 3. Trả lời ngắn}
\setcounterpageref{ex}{0}

\begin{ex} %[4D4H20]
Một bao gạo nặng $5$ yến. Hỏi $6$ bao gạo như thế nặng bao nhiêu ki-lô-gam?
\shortans{300}
\loigiai{
\textbf{Đáp số:} $300$\\
$6 \times 50 = 300\text{ kg}$.
}
\end{ex}

\begin{ex} %[4D4H20]
Một mặt bàn hình chữ nhật dài $12\text{ dm}$, rộng $8\text{ dm}$. Tính diện tích mặt bàn theo đề-xi-mét vuông.
\shortans{96}
\loigiai{
\textbf{Đáp số:} $96$\\
$12 \times 8 = 96\text{ dm}^2$.
}
\end{ex}

\begin{ex} %[4D4V20]
Đổi $3$ thế kỉ $15$ năm ra năm.
\shortans{315}
\loigiai{
\textbf{Đáp số:} $315$\\
$315$ năm.
}
\end{ex}

\begin{ex} %[4D4V20]
Một người đi bộ với vận tốc $1\text{ m/s}$. Hỏi trong $5$ phút người đó đi được bao nhiêu mét?
\shortans{300}
\loigiai{
\textbf{Đáp số:} $300$\\
$5$ phút $= 300$ giây $\Rightarrow 1 \times 300 = 300\text{ m}$.
}
\end{ex}

\begin{ex} %[4D4V20]
Một bức tranh hình chữ nhật có diện tích $2\text{ m}^2$, chiều dài $20\text{ dm}$. Tính chiều rộng của bức tranh theo đề-xi-mét.
\shortans{10}
\loigiai{
\textbf{Đáp số:} $10$\\
$2\text{ m}^2 = 200\text{ dm}^2 \Rightarrow 200 : 20 = 10\text{ dm}$.
}
\end{ex}

\begin{ex} %[4D4C20]
Một xe tải có trọng tải $5$ tấn. Xe chở $30$ thùng hàng loại $1$ tạ/thùng và một số bao hàng loại $50\text{ kg}$/bao thì vừa đầy tải trọng. Tính số bao hàng loại $50\text{ kg}$.
\shortans{40}
\loigiai{
\textbf{Đáp số:} $40$\\
30 thùng $= 3000\text{ kg}$. Còn lại: $5000 - 3000 = 2000\text{ kg}$.\\
Số bao: $2000 : 50 = 40$ bao.
}
\end{ex}

\newpage

\TaoPhieuKiemTra{Chủ đề 4 -- Một số đơn vị đo đại lượng}{21}{Luyện tập chung đại lượng}{https://loptoanthaylinh.vn/thi-online/toan-4/bai-21}
\textbf{Phần 1. Trắc nghiệm lựa chọn}
\setcounterpageref{ex}{0}

\begin{ex} %[4D4N21]
$3$ tấn bằng bao nhiêu tạ?
\choice
{$300\text{ tạ}$}
{\True $30\text{ tạ}$}
{$3000\text{ tạ}$}
{$3\text{ tạ}$}
\loigiai{
$3$ tấn $= 30$ tạ.\\
\textbf{Đáp án của câu hỏi là B}
}
\end{ex}

\begin{ex} %[4D4N21]
$5\text{ m}^2$ bằng bao nhiêu đề-xi-mét vuông?
\choice
{$50\text{ dm}^2$}
{$5000\text{ dm}^2$}
{$50\,000\text{ dm}^2$}
{\True $500\text{ dm}^2$}
\loigiai{
$5\text{ m}^2 = 500\text{ dm}^2$.\\
\textbf{Đáp án của câu hỏi là D}
}
\end{ex}

\begin{ex} %[4D4H21]
$2$ phút $30$ giây bằng bao nhiêu giây?
\choice
{$230$ giây}
{$120$ giây}
{\True $150$ giây}
{$180$ giây}
\loigiai{
$2 \times 60 + 30 = 150$ giây.\\
\textbf{Đáp án của câu hỏi là C}
}
\end{ex}

\begin{ex} %[4D4H21]
Năm $1000$ thuộc thế kỉ nào?
\choice
{\True Thế kỉ X}
{Thế kỉ IX}
{Thế kỉ XI}
{Thế kỉ XII}
\loigiai{
Năm $1000$ là năm cuối cùng của thế kỉ X.\\
\textbf{Đáp án của câu hỏi là A}
}
\end{ex}

\begin{ex} %[4D4H21]
$4$ tạ $8\text{ kg}$ bằng bao nhiêu ki-lô-gam?
\choice
{$48\text{ kg}$}
{$480\text{ kg}$}
{$4008\text{ kg}$}
{\True $408\text{ kg}$}
\loigiai{
$4 \times 100 + 8 = 408\text{ kg}$.\\
\textbf{Đáp án của câu hỏi là D}
}
\end{ex}

\begin{ex} %[4D4H21]
$6\text{ dm}^2 5\text{ cm}^2$ bằng bao nhiêu xăng-ti-mét vuông?
\choice
{\True $605\text{ cm}^2$}
{$65\text{ cm}^2$}
{$650\text{ cm}^2$}
{$6005\text{ cm}^2$}
\loigiai{
$6 \times 100 + 5 = 605\text{ cm}^2$.\\
\textbf{Đáp án của câu hỏi là A}
}
\end{ex}

\begin{ex} %[4D4V21]
Một con cá ngừ đại dương nặng $2$ tạ $50\text{ kg}$, một con cá thu nặng $25\text{ kg}$. Con cá ngừ nặng gấp mấy lần con cá thu?
\choice
{$5$ lần}
{$15$ lần}
{\True $10$ lần}
{$20$ lần}
\loigiai{
$2$ tạ $50\text{ kg} = 250\text{ kg}$.\\
$250 : 25 = 10$ lần.\\
\textbf{Đáp án của câu hỏi là C}
}
\end{ex}

\begin{ex} %[4D4V21]
Một sân bóng đá mini hình chữ nhật dài $40\text{ m}$, rộng $20\text{ m}$. Diện tích sân bóng đó là
\choice
{$80\text{ m}^2$}
{$8000\text{ m}^2$}
{$120\text{ m}^2$}
{\True $800\text{ m}^2$}
\loigiai{
$40 \times 20 = 800\text{ m}^2$.\\
\textbf{Đáp án của câu hỏi là D}
}
\end{ex}

\begin{ex} %[4D4V21]
Chiến thắng Bạch Đằng của Ngô Quyền năm $938$ cách năm $2026$ bao nhiêu năm?
\choice
{$1086$ năm}
{\True $1088$ năm}
{$1090$ năm}
{$1080$ năm}
\loigiai{
$2026 - 938 = 1088$ năm.\\
\textbf{Đáp án của câu hỏi là B}
}
\end{ex}

\begin{ex} %[4D4V21]
Tính giá trị biểu thức: $3$ tấn $- 18$ tạ $+ 400\text{ kg}$ theo đơn vị tạ.
\choice
{$14\text{ tạ}$}
{$18\text{ tạ}$}
{\True $16\text{ tạ}$}
{$20\text{ tạ}$}
\loigiai{
$30 - 18 + 4 = 16$ tạ.\\
\textbf{Đáp án của câu hỏi là C}
}
\end{ex}

\begin{ex} %[4D4V21]
Một thửa ruộng thu được $3$ tấn thóc. Người ta đóng vào các bao, mỗi bao $50\text{ kg}$. Số bao thóc đóng được là
\choice
{$50$ bao}
{\True $60$ bao}
{$70$ bao}
{$80$ bao}
\loigiai{
$3000 : 50 = 60$ bao.\\
\textbf{Đáp án của câu hỏi là B}
}
\end{ex}

\begin{ex} %[4D4C21]
Một khu vườn hình vuông có chu vi $120\text{ m}$. Người ta dành $\dfrac{1}{3}$ diện tích vườn để trồng cây ăn quả, diện tích còn lại để trồng hoa. Diện tích trồng hoa là
\choice
{\True $600\text{ m}^2$}
{$300\text{ m}^2$}
{$900\text{ m}^2$}
{$400\text{ m}^2$}
\loigiai{
Cạnh vườn: $120 : 4 = 30\text{ m}$.\\
Diện tích vườn: $30 \times 30 = 900\text{ m}^2$.\\
Diện tích trồng hoa: $900 \times \dfrac{2}{3} = 600\text{ m}^2$.\\
\textbf{Đáp án của câu hỏi là A}
}
\end{ex}

\textbf{Phần 2. Trắc nghiệm đúng sai}
\setcounterpageref{ex}{0}

\begin{ex} %[4D4H21]
Xét tính đúng sai của các khẳng định:
\choiceTF
{\True $1$ tấn $= 1000\text{ kg}$.}
{\True $1\text{ m}^2 = 100\text{ dm}^2$.}
{Năm 1900 thuộc thế kỉ XX.}
{\True Năm 1900 thuộc thế kỉ XIX.}
\loigiai{
\begin{itemize}
    \item $1$ tấn $= 1000\text{ kg}$: Đúng.
    \item $1\text{ m}^2 = 100\text{ dm}^2$: Đúng.
    \item Năm 1900 là năm cuối của thế kỉ XIX (không phải XX). Ý 3 Sai, Ý 4 Đúng.
\end{itemize}
\begin{itemize}
    \item a) \textbf{Đúng}.
    \item b) \textbf{Đúng}.
    \item c) \textbf{Sai}.
    \item d) \textbf{Đúng}.
\end{itemize}
}
\end{ex}

\begin{ex} %[4D4V21]
Một kho gạo có $4$ tấn gạo. Ngày đầu bán $12$ tạ, ngày thứ hai bán gấp đôi ngày đầu. Xét tính đúng sai:
\choiceTF
{\True Ngày thứ hai bán được $24$ tạ gạo.}
{\True Cả hai ngày bán được $36$ tạ gạo.}
{Sau hai ngày kho còn lại $10$ tạ gạo.}
{\True Sau hai ngày kho còn lại $4$ tạ gạo ($400\text{ kg}$).}
\loigiai{
\begin{itemize}
    \item Ngày 2: $24$ tạ. Đúng.
    \item Cả 2 ngày: $12 + 24 = 36$ tạ. Đúng.
    \item Còn lại: $40 - 36 = 4$ tạ $\neq 10$ tạ. Ý 3 Sai, Ý 4 Đúng.
\end{itemize}
\begin{itemize}
    \item a) \textbf{Đúng}.
    \item b) \textbf{Đúng}.
    \item c) \textbf{Sai}.
    \item d) \textbf{Đúng}.
\end{itemize}
}
\end{ex}

\textbf{Phần 3. Trả lời ngắn}
\setcounterpageref{ex}{0}

\begin{ex} %[4D4H21]
Đổi $5$ tấn $6$ tạ ra tạ.
\shortans{56}
\loigiai{
\textbf{Đáp số:} $56$\\
$56$ tạ.
}
\end{ex}

\begin{ex} %[4D4H21]
Đổi $8\text{ m}^2 50\text{ dm}^2$ ra đề-xi-mét vuông.
\shortans{850}
\loigiai{
\textbf{Đáp số:} $850$\\
$850\text{ dm}^2$.
}
\end{ex}

\begin{ex} %[4D4V21]
Đổi $180$ giây ra phút.
\shortans{3}
\loigiai{
\textbf{Đáp số:} $3$\\
$3$ phút.
}
\end{ex}

\begin{ex} %[4D4V21]
Năm $1010$ thuộc thế kỉ thứ mấy (ghi bằng số Ả Rập)?
\shortans{11}
\loigiai{
\textbf{Đáp số:} $11$\\
Thế kỉ 11 (XI).
}
\end{ex}

\begin{ex} %[4D4V21]
Một căn phòng hình chữ nhật dài $8\text{ m}$, rộng $5\text{ m}$. Tính diện tích sàn phòng theo mét vuông.
\shortans{40}
\loigiai{
\textbf{Đáp số:} $40$\\
$8 \times 5 = 40\text{ m}^2$.
}
\end{ex}

\begin{ex} %[4D4C21]
Một xe tải chở $2$ tấn $500\text{ kg}$ hàng gồm gạo và ngô. Biết số gạo nặng hơn số ngô $5$ tạ. Tính khối lượng gạo xe chở theo ki-lô-gam.
\shortans{1500}
\loigiai{
\textbf{Đáp số:} $1500$\\
$2500\text{ kg}$ và $500\text{ kg}$. Gạo: $(2500 + 500) : 2 = 1500\text{ kg}$.
}
\end{ex}

\newpage

% =========================================================================
% CHỦ ĐỀ 5 - PHÉP CỘNG VÀ PHÉP TRỪ
% =========================================================================
\section{Phép cộng và phép trừ}

\TaoPhieuKiemTra{Chủ đề 5 -- Phép cộng và phép trừ}{22}{Phép cộng các số có nhiều chữ số}{https://loptoanthaylinh.vn/thi-online/toan-4/bai-22}
\textbf{Phần 1. Trắc nghiệm lựa chọn}
\setcounterpageref{ex}{0}

\begin{ex} %[4D5N22]
Kết quả của phép tính $324\,516 + 253\,172$ là
\choice
{\True $577\,688$}
{$577\,678$}
{$577\,588$}
{$587\,688$}
\loigiai{
$324\,516 + 253\,172 = 577\,688$.\\
\textbf{Đáp án của câu hỏi là A}
}
\end{ex}

\begin{ex} %[4D5N22]
Đặt tính rồi tính: $1\,450\,000 + 2\,350\,000$. Kết quả là
\choice
{$3\,700\,000$}
{\True $3\,800\,000$}
{$3\,900\,000$}
{$4\,000\,000$}
\loigiai{
$1\,450\,000 + 2\,350\,000 = 3\,800\,000$.\\
\textbf{Đáp án của câu hỏi là B}
}
\end{ex}

\begin{ex} %[4D5H22]
Kết quả của phép tính $45\,678\,900 + 12\,345\,600$ là
\choice
{$57\,024\,500$}
{$58\,014\,500$}
{$57\,924\,500$}
{\True $58\,024\,500$}
\loigiai{
$45\,678\,900 + 12\,345\,600 = 58\,024\,500$.\\
\textbf{Đáp án của câu hỏi là D}
}
\end{ex}

\begin{ex} %[4D5H22]
Tìm $x$ biết: $x - 125\,000 = 345\,000$.
\choice
{$220\,000$}
{$460\,000$}
{\True $470\,000$}
{$480\,000$}
\loigiai{
$x = 345\,000 + 125\,000 = 470\,000$.\\
\textbf{Đáp án của câu hỏi là C}
}
\end{ex}

\begin{ex} %[4D5H22]
Một công ty năm thứ nhất thu được $450\,000\,000$ đồng, năm thứ hai thu được $520\,000\,000$ đồng. Cả hai năm công ty thu được số tiền là
\choice
{$950\,000\,000$ đồng}
{\True $970\,000\,000$ đồng}
{$960\,000\,000$ đồng}
{$980\,000\,000$ đồng}
\loigiai{
$450\,000\,000 + 520\,000\,000 = 970\,000\,000$ đồng.\\
\textbf{Đáp án của câu hỏi là B}
}
\end{ex}

\begin{ex} %[4D5H22]
Tính giá trị biểu thức: $12\,000\,000 + 35\,000\,000 + 18\,000\,000$.
\choice
{$60\,000\,000$}
{$70\,000\,000$}
{$55\,000\,000$}
{\True $65\,000\,000$}
\loigiai{
$(12 + 18) + 35 = 30 + 35 = 65$ triệu $= 65\,000\,000$.\\
\textbf{Đáp án của câu hỏi là D}
}
\end{ex}

\begin{ex} %[4D5V22]
Dân số tỉnh A là $1\,250\,000$ người, dân số tỉnh B nhiều hơn tỉnh A là $340\,000$ người. Tổng dân số của hai tỉnh là
\choice
{$2\,590\,000$ người}
{$2\,740\,000$ người}
{\True $2\,840\,000$ người}
{$2\,900\,000$ người}
\loigiai{
Tỉnh B: $1\,250\,000 + 340\,000 = 1\,590\,000$ người.\\
Tổng hai tỉnh: $1\,250\,000 + 1\,590\,000 = 2\,840\,000$ người.\\
\textbf{Đáp án của câu hỏi là C}
}
\end{ex}

\begin{ex} %[4D5V22]
Một nông trại thu hoạch: Vụ 1 được $35\,600\text{ kg}$ cà phê, Vụ 2 được $42\,800\text{ kg}$ cà phê, Vụ 3 được $38\,400\text{ kg}$ cà phê. Cả 3 vụ thu hoạch được bao nhiêu tấn cà phê?
\choice
{\True $116{,}8\text{ tấn}$}
{$115{,}8\text{ tấn}$}
{$116\text{ tấn}$}
{$120\text{ tấn}$}
\loigiai{
Tổng khối lượng: $35\,600 + 42\,800 + 38\,400 = 116\,800\text{ kg} = 116{,}8\text{ tấn}$.\\
\textbf{Đáp án của câu hỏi là A}
}
\end{ex}

\begin{ex} %[4D5V22]
Tìm chữ số $a, b$ biết: $3a5,420 + 214,b10 = 559\,730$.
\choice
{$a = 3, b = 2$}
{\True $a = 4, b = 3$}
{$a = 5, b = 4$}
{$a = 4, b = 2$}
\loigiai{
$a + 1 = 5 \Rightarrow a = 4$. $4 + b = 7 \Rightarrow b = 3$.\\
\textbf{Đáp án của câu hỏi là B}
}
\end{ex}

\begin{ex} %[4D5V22]
Một nhà máy sản xuất linh kiện điện tử: Quý 1 sản xuất $2\,400\,000$ linh kiện, mỗi quý tiếp theo đều sản xuất tăng thêm $300\,000$ linh kiện so với quý trước. Tổng số linh kiện sản xuất trong cả năm là
\choice
{\True $11\,400\,000$}
{$10\,800\,000$}
{$11\,000\,000$}
{$12\,000\,000$}
\loigiai{
Quý 1: 2,4 triệu; Quý 2: 2,7 triệu; Quý 3: 3,0 triệu; Quý 4: 3,3 triệu.\\
Tổng: $2{,}4 + 2{,}7 + 3{,}0 + 3{,}3 = 11{,}4$ triệu $= 11\,400\,000$.\\
\textbf{Đáp án của câu hỏi là A}
}
\end{ex}

\begin{ex} %[4D5V22]
Tổng của số lớn nhất có 6 chữ số và số bé nhất có 7 chữ số là
\choice
{$1\,000\,999$}
{$2\,000\,000$}
{\True $1\,999\,999$}
{$1\,099\,999$}
\loigiai{
$999\,999 + 1\,000\,000 = 1\,999\,999$.\\
\textbf{Đáp án của câu hỏi là C}
}
\end{ex}

\begin{ex} %[4D5C22]
Cho tổng $S = 12\,345\,678 + 87\,654\,321$. Tổng các chữ số của $S$ là
\choice
{$70$}
{$72$}
{$81$}
{\True $73$ (hay $9 \times 8 + 1 = 73$ nếu cộng có nhớ)}
\loigiai{
$S = 99\,999\,999 \Rightarrow$ Tổng các chữ số là $9 \times 8 = 72$.\\
\textbf{Đáp án của câu hỏi là D}
}
\end{ex}

\textbf{Phần 2. Trắc nghiệm đúng sai}
\setcounterpageref{ex}{0}

\begin{ex} %[4D5H22]
Xét tính đúng sai của các phép cộng:
\choiceTF
{\True $123\,456 + 654\,321 = 777\,777$.}
{$250\,000 + 350\,000 = 500\,000$.}
{\True $250\,000 + 350\,000 = 600\,000$.}
{\True $4\,500\,000 + 5\,500\,000 = 10\,000\,000$.}
\loigiai{
\begin{itemize}
    \item $123\,456 + 654\,321 = 777\,777$. Đúng.
    \item $250\,000 + 350\,000 = 600\,000 \neq 500\,000$. Ý 2 Sai, Ý 3 Đúng.
    \item $4\,500\,000 + 5\,500\,000 = 10\,000\,000$. Đúng.
\end{itemize}
\begin{itemize}
    \item a) \textbf{Đúng}.
    \item b) \textbf{Sai}.
    \item c) \textbf{Đúng}.
    \item d) \textbf{Đúng}.
\end{itemize}
}
\end{ex}

\begin{ex} %[4D5V22]
Ba tổ công nhân nhận sửa chữa một đoạn đường: Tổ 1 sửa $1250\text{ m}$, tổ 2 sửa $1450\text{ m}$, tổ 3 sửa $1300\text{ m}$. Xét tính đúng sai:
\choiceTF
{\True Tổ 2 sửa được đoạn đường dài nhất.}
{\True Cả 3 tổ sửa được tất cả $4000\text{ m}$ đường.}
{Đổi đoạn đường cả 3 tổ sửa được là $40\text{ km}$.}
{\True Đổi đoạn đường cả 3 tổ sửa được là $4\text{ km}$.}
\loigiai{
\begin{itemize}
    \item Tổ 2 sửa $1450\text{ m}$, nhiều nhất. Đúng.
    \item Tổng: $1250 + 1450 + 1300 = 4000\text{ m}$. Đúng.
    \item $4000\text{ m} = 4\text{ km} \neq 40\text{ km}$. Ý 3 Sai, Ý 4 Đúng.
\end{itemize}
\begin{itemize}
    \item a) \textbf{Đúng}.
    \item b) \textbf{Đúng}.
    \item c) \textbf{Sai}.
    \item d) \textbf{Đúng}.
\end{itemize}
}
\end{ex}

\textbf{Phần 3. Trả lời ngắn}
\setcounterpageref{ex}{0}

\begin{ex} %[4D5H22]
Tính tổng: $450\,000 + 350\,000$.
\shortans{800000}
\loigiai{
\textbf{Đáp số:} $800000$\\
$800\,000$.
}
\end{ex}

\begin{ex} %[4D5H22]
Tính tổng: $12\,500\,000 + 7\,500\,000$.
\shortans{20000000}
\loigiai{
\textbf{Đáp số:} $20000000$\\
$20\,000\,000$.
}
\end{ex}

\begin{ex} %[4D5V22]
Tìm $x$ biết: $x - 2\,400\,000 = 5\,600\,000$.
\shortans{8000000}
\loigiai{
\textbf{Đáp số:} $8000000$\\
$x = 5\,600\,000 + 2\,400\,000 = 8\,000\,000$.
}
\end{ex}

\begin{ex} %[4D5V22]
Tính giá trị biểu thức: $15\,000\,000 + 25\,000\,000 + 35\,000\,000$.
\shortans{75000000}
\loigiai{
\textbf{Đáp số:} $75000000$\\
$75\,000\,000$.
}
\end{ex}

\begin{ex} %[4D5V22]
Một huyện có $45\,600$ học sinh tiểu học và $38\,400$ học sinh trung học cơ sở. Tính tổng số học sinh của cả huyện.
\shortans{84000}
\loigiai{
\textbf{Đáp số:} $84000$\\
$45\,600 + 38\,400 = 84\,000$.
}
\end{ex}

\begin{ex} %[4D5C22]
Tính nhanh tổng: $1\,000\,000 + 2\,000\,000 + 3\,000\,000 + 4\,000\,000 + 5\,000\,000$ theo triệu (chỉ ghi số triệu).
\shortans{15}
\loigiai{
\textbf{Đáp số:} $15$\\
$15$ triệu.
}
\end{ex}

\newpage

\TaoPhieuKiemTra{Chủ đề 5 -- Phép cộng và phép trừ}{23}{Phép trừ các số có nhiều chữ số}{https://loptoanthaylinh.vn/thi-online/toan-4/bai-23}
\textbf{Phần 1. Trắc nghiệm lựa chọn}
\setcounterpageref{ex}{0}

\begin{ex} %[4D5N23]
Kết quả của phép trừ $678\,950 - 234\,510$ là
\choice
{$444\,430$}
{$444\,540$}
{$434\,440$}
{\True $444\,440$}
\loigiai{
$678\,950 - 234\,510 = 444\,440$.\\
\textbf{Đáp án của câu hỏi là D}
}
\end{ex}

\begin{ex} %[4D5N23]
Kết quả của phép tính $5\,000\,000 - 2\,000\,000$ là
\choice
{$2\,000\,000$}
{$4\,000\,000$}
{\True $3\,000\,000$}
{$3\,500\,000$}
\loigiai{
$5\,000\,000 - 2\,000\,000 = 3\,000\,000$.\\
\textbf{Đáp án của câu hỏi là C}
}
\end{ex}

\begin{ex} %[4D5H23]
Kết quả của phép trừ $10\,000\,000 - 3\,450\,000$ là
\choice
{$6\,450\,000$}
{$6\,650\,000$}
{$7\,550\,000$}
{\True $6\,550\,000$}
\loigiai{
$10\,000\,000 - 3\,450\,000 = 6\,550\,000$.\\
\textbf{Đáp án của câu hỏi là D}
}
\end{ex}

\begin{ex} %[4D5H23]
Tìm $x$ biết: $750\,000 - x = 280\,000$.
\choice
{\True $470\,000$}
{$460\,000$}
{$480\,000$}
{$570\,000$}
\loigiai{
$x = 750\,000 - 280\,000 = 470\,000$.\\
\textbf{Đáp án của câu hỏi là A}
}
\end{ex}

\begin{ex} %[4D5H23]
Một kho chứa $45\,000\text{ kg}$ thóc. Sau khi xuất đi $18\,500\text{ kg}$ thóc thì trong kho còn lại bao nhiêu ki-lô-gam thóc?
\choice
{$25\,500\text{ kg}$}
{\True $26\,500\text{ kg}$}
{$27\,500\text{ kg}$}
{$26\,000\text{ kg}$}
\loigiai{
$45\,000 - 18\,500 = 26\,500\text{ kg}$.\\
\textbf{Đáp án của câu hỏi là B}
}
\end{ex}

\begin{ex} %[4D5H23]
Hiệu của số bé nhất có 7 chữ số và số lớn nhất có 5 chữ số là
\choice
{\True $900\,001$}
{$900\,001$}
{$900\,000$}
{$899\,999$}
\loigiai{
$1\,000\,000 - 99\,999 = 900\,001$.\\
\textbf{Đáp án của câu hỏi là A}
}
\end{ex}

\begin{ex} %[4D5V23]
Một tài khoản ngân hàng có $85\,000\,000$ đồng. Chủ tài khoản rút ra $25\,000\,000$ đồng để chi tiêu và chuyển khoản $30\,000\,000$ đồng để thanh toán tiền hàng. Số tiền còn lại trong tài khoản là
\choice
{$25\,000\,000$ đồng}
{\True $30\,000\,000$ đồng}
{$35\,000\,000$ đồng}
{$40\,000\,000$ đồng}
\loigiai{
Còn lại: $85\,000\,000 - (25\,000\,000 + 30\,000\,000) = 30\,000\,000$ đồng.\\
\textbf{Đáp án của câu hỏi là B}
}
\end{ex}

\begin{ex} %[4D5V23]
Tính giá trị của biểu thức: $120\,000\,000 - (45\,000\,000 + 35\,000\,000)$.
\choice
{$30\,000\,000$}
{$50\,000\,000$}
{\True $40\,000\,000$}
{$45\,000\,000$}
\loigiai{
$120\,000\,000 - 80\,000\,000 = 40\,000\,000$.\\
\textbf{Đáp án của câu hỏi là C}
}
\end{ex}

\begin{ex} %[4D5V23]
Dân số thành phố A là $3\,450\,000$ người, dân số thành phố B là $2\,150\,000$ người. Dân số thành phố A nhiều hơn thành phố B bao nhiêu người?
\choice
{$1\,200\,000$ người}
{$1\,400\,000$ người}
{$1\,500\,000$ người}
{\True $1\,300\,000$ người}
\loigiai{
$3\,450\,000 - 2\,150\,000 = 1\,300\,000$ người.\\
\textbf{Đáp án của câu hỏi là D}
}
\end{ex}

\begin{ex} %[4D5V23]
Một xí nghiệp may theo kế hoạch phải may $150\,000$ bộ quần áo. Đợt 1 may được $45\,000$ bộ, đợt 2 may được $55\,000$ bộ. Để hoàn thành kế hoạch xí nghiệp còn phải may số bộ quần áo là
\choice
{\True $50\,000$ bộ}
{$40\,000$ bộ}
{$60\,000$ bộ}
{$55\,000$ bộ}
\loigiai{
Còn lại: $150\,000 - (45\,000 + 55\,000) = 50\,000$ bộ.\\
\textbf{Đáp án của câu hỏi là A}
}
\end{ex}

\begin{ex} %[4D5V23]
Tìm số bị trừ biết số trừ là $3\,450\,000$ và hiệu là $2\,550\,000$.
\choice
{$5\,900\,000$}
{$6\,100\,000$}
{\True $6\,000\,000$}
{$5\,000\,000$}
\loigiai{
Số bị trừ: $3\,450\,000 + 2\,550\,000 = 6\,000\,000$.\\
\textbf{Đáp án của câu hỏi là C}
}
\end{ex}

\begin{ex} %[4D5C23]
Khi thực hiện phép trừ một số có 6 chữ số cho một số có 5 chữ số, một bạn đã viết nhầm số trừ lùi sang trái 1 cột nên được hiệu là $450\,000$. Biết số bị trừ là $650\,000$. Tìm số trừ đúng.
\choice
{$10\,000$}
{\True $20\,000$}
{$30\,000$}
{$40\,000$}
\loigiai{
Viết lùi sang trái 1 cột tức là gấp 10 lần số trừ.\\
Số trừ sai là: $650\,000 - 450\,000 = 200\,000$.\\
Số trừ đúng là: $200\,000 : 10 = 20\,000$.\\
\textbf{Đáp án của câu hỏi là B}
}
\end{ex}

\textbf{Phần 2. Trắc nghiệm đúng sai}
\setcounterpageref{ex}{0}

\begin{ex} %[4D5H23]
Xét tính đúng sai của các phép trừ sau:
\choiceTF
{\True $1\,000\,000 - 400\,000 = 600\,000$.}
{$850\,000 - 320\,000 = 540\,000$.}
{\True $850\,000 - 320\,000 = 530\,000$.}
{\True $45\,678\,900 - 15\,678\,900 = 30\,000\,000$.}
\loigiai{
\begin{itemize}
    \item $1\,000\,000 - 400\,000 = 600\,000$. Đúng.
    \item $850\,000 - 320\,000 = 530\,000 \neq 540\,000$. Ý 2 Sai, Ý 3 Đúng.
    \item $45\,678\,900 - 15\,678\,900 = 30\,000\,000$. Đúng.
\end{itemize}
\begin{itemize}
    \item a) \textbf{Đúng}.
    \item b) \textbf{Sai}.
    \item c) \textbf{Đúng}.
    \item d) \textbf{Đúng}.
\end{itemize}
}
\end{ex}

\begin{ex} %[4D5V23]
Một xí nghiệp có $120\,000$ tấn than trong kho: Ngày thứ nhất xuất $35\,000$ tấn than, ngày thứ hai xuất $45\,000$ tấn than. Xét tính đúng sai:
\choiceTF
{\True Tổng số than xuất trong 2 ngày là $80\,000$ tấn.}
{\True Số than xuất ngày thứ hai nhiều hơn ngày thứ nhất $10\,000$ tấn.}
{Sau hai ngày kho còn lại $50\,000$ tấn than.}
{\True Sau hai ngày kho còn lại $40\,000$ tấn than.}
\loigiai{
\begin{itemize}
    \item Đã xuất: $35\,000 + 45\,000 = 80\,000$ tấn. Đúng.
    \item Hiệu: $45\,000 - 35\,000 = 10\,000$ tấn. Đúng.
    \item Còn lại: $120\,000 - 80\,000 = 40\,000$ tấn $\neq 50\,000$ tấn. Ý 3 Sai, Ý 4 Đúng.
\end{itemize}
\begin{itemize}
    \item a) \textbf{Đúng}.
    \item b) \textbf{Đúng}.
    \item c) \textbf{Sai}.
    \item d) \textbf{Đúng}.
\end{itemize}
}
\end{ex}

\textbf{Phần 3. Trả lời ngắn}
\setcounterpageref{ex}{0}

\begin{ex} %[4D5H23]
Tính hiệu: $900\,000 - 350\,000$.
\shortans{550000}
\loigiai{
\textbf{Đáp số:} $550000$\\
$550\,000$.
}
\end{ex}

\begin{ex} %[4D5H23]
Tính hiệu: $50\,000\,000 - 18\,000\,000$.
\shortans{32000000}
\loigiai{
\textbf{Đáp số:} $32000000$\\
$32\,000\,000$.
}
\end{ex}

\begin{ex} %[4D5V23]
Tìm $x$ biết: $1\,500\,000 - x = 850\,000$.
\shortans{650000}
\loigiai{
\textbf{Đáp số:} $650000$\\
$x = 1\,500\,000 - 850\,000 = 650\,000$.
}
\end{ex}

\begin{ex} %[4D5V23]
Tính giá trị biểu thức: $500\,000 - 120\,000 - 180\,000$.
\shortans{200000}
\loigiai{
\textbf{Đáp số:} $200000$\\
$500\,000 - 300\,000 = 200\,000$.
}
\end{ex}

\begin{ex} %[4D5V23]
Một công ty dự định thu về $500$ triệu đồng, thực tế chỉ thu được $420$ triệu đồng. Hỏi công ty thu thiếu bao nhiêu triệu đồng? (Chỉ ghi số triệu).
\shortans{80}
\loigiai{
\textbf{Đáp số:} $80$\\
$500 - 420 = 80$ triệu.
}
\end{ex}

\begin{ex} %[4D5C23]
Hiệu của hai số là $45\,000$. Nếu thêm vào số bị trừ $5000$ đơn vị và bớt ở số trừ $5000$ đơn vị thì hiệu mới là bao nhiêu?
\shortans{55000}
\loigiai{
\textbf{Đáp số:} $55000$\\
$45\,000 + 5000 + 5000 = 55\,000$.
}
\end{ex}

\newpage

\TaoPhieuKiemTra{Chủ đề 5 -- Phép cộng và phép trừ}{24}{Tính chất giao hoán, kết hợp của phép cộng}{https://loptoanthaylinh.vn/thi-online/toan-4/bai-24}
\textbf{Phần 1. Trắc nghiệm lựa chọn}
\setcounterpageref{ex}{0}

\begin{ex} %[4D5N24]
Tính chất giao hoán của phép cộng được viết dưới dạng công thức là
\choice
{$a + b = a - b$}
{$(a + b) + c = a + (b + c)$}
{\True $a + b = b + a$}
{$a + 0 = a$}
\loigiai{
$a + b = b + a$.\\
\textbf{Đáp án của câu hỏi là C}
}
\end{ex}

\begin{ex} %[4D5N24]
Tính chất kết hợp của phép cộng được viết dưới dạng công thức là
\choice
{$a + b = b + a$}
{\True $(a + b) + c = a + (b + c)$}
{$a \times b = b \times a$}
{$a + (b - c) = a + b - c$}
\loigiai{
$(a + b) + c = a + (b + c)$.\\
\textbf{Đáp án của câu hỏi là B}
}
\end{ex}

\begin{ex} %[4D5H24]
Điền số thích hợp vào chỗ chấm: $45\,678 + 12\,345 = 12\,345 + \dots$
\choice
{$12\,345$}
{$58\,023$}
{\True $45\,678$}
{$45\,000$}
\loigiai{
Theo tính chất giao hoán: $45\,678 + 12\,345 = 12\,345 + 45\,678$.\\
\textbf{Đáp án của câu hỏi là C}
}
\end{ex}

\begin{ex} %[4D5H24]
Tính thuận tiện: $125 + 348 + 875$.
\choice
{$1248$}
{$1448$}
{$1358$}
{\True $1348$}
\loigiai{
$(125 + 875) + 348 = 1000 + 348 = 1348$.\\
\textbf{Đáp án của câu hỏi là D}
}
\end{ex}

\begin{ex} %[4D5H24]
Tính thuận tiện: $450 + 279 + 550$.
\choice
{$1179$}
{$1379$}
{\True $1279$}
{$1289$}
\loigiai{
$(450 + 550) + 279 = 1000 + 279 = 1279$.\\
\textbf{Đáp án của câu hỏi là C}
}
\end{ex}

\begin{ex} %[4D5H24]
Tính thuận tiện: $2400 + 1850 + 7600$.
\choice
{$11\,750$}
{$11\,950$}
{$12\,850$}
{\True $11\,850$}
\loigiai{
$(2400 + 7600) + 1850 = 10\,000 + 1850 = 11\,850$.\\
\textbf{Đáp án của câu hỏi là D}
}
\end{ex}

\begin{ex} %[4D5V24]
Tính nhanh giá trị biểu thức: $123 + 456 + 877 + 544$.
\choice
{\True $2000$}
{$1900$}
{$2100$}
{$2200$}
\loigiai{
$(123 + 877) + (456 + 544) = 1000 + 1000 = 2000$.\\
\textbf{Đáp án của câu hỏi là A}
}
\end{ex}

\begin{ex} %[4D5V24]
Tính nhanh: $35\,000 + 18\,200 + 65\,000 + 21\,800$.
\choice
{$130\,000$}
{\True $140\,000$}
{$150\,000$}
{$145\,000$}
\loigiai{
$(35\,000 + 65\,000) + (18\,200 + 21\,800) = 100\,000 + 40\,000 = 140\,000$.\\
\textbf{Đáp án của câu hỏi là B}
}
\end{ex}

\begin{ex} %[4D5V24]
Cho biểu thức $S = a + b + c + d$. Biết $a + c = 100$ và $b + d = 250$. Giá trị của biểu thức $S$ là
\choice
{\True $350$}
{$300$}
{$400$}
{$450$}
\loigiai{
$S = (a + c) + (b + d) = 100 + 250 = 350$.\\
\textbf{Đáp án của câu hỏi là A}
}
\end{ex}

\begin{ex} %[4D5V24]
Một cửa hàng bán sách trong 3 ngày: Ngày 1 bán $1450$ cuốn, ngày 2 bán $2300$ cuốn, ngày 3 bán $3550$ cuốn. Tổng số sách bán trong 3 ngày tính theo cách thuận tiện là
\choice
{$7200$ cuốn}
{\True $7300$ cuốn}
{$7400$ cuốn}
{$7500$ cuốn}
\loigiai{
$(1450 + 3550) + 2300 = 5000 + 2300 = 7300$ cuốn.\\
\textbf{Đáp án của câu hỏi là B}
}
\end{ex}

\begin{ex} %[4D5V24]
Tính nhanh tổng: $1 + 2 + 3 + 4 + 5 + 6 + 7 + 8 + 9$.
\choice
{\True $45$}
{$40$}
{$50$}
{$55$}
\loigiai{
$(1 + 9) + (2 + 8) + (3 + 7) + (4 + 6) + 5 = 10 \times 4 + 5 = 45$.\\
\textbf{Đáp án của câu hỏi là A}
}
\end{ex}

\begin{ex} %[4D5C24]
Tính giá trị biểu thức: $125\,000 + 250\,000 + 375\,000 + 500\,000 + 625\,000 + 750\,000 + 875\,000$.
\choice
{$3\,400\,000$}
{$3\,600\,000$}
{$4\,000\,000$}
{\True $3\,500\,000$}
\loigiai{
$(125 + 875) + (250 + 750) + (375 + 625) + 500 = 1000 \times 3 + 500 = 3500$ nghìn $= 3\,500\,000$.\\
\textbf{Đáp án của câu hỏi là D}
}
\end{ex}

\textbf{Phần 2. Trắc nghiệm đúng sai}
\setcounterpageref{ex}{0}

\begin{ex} %[4D5H24]
Xét tính đúng sai của các biểu thức áp dụng tính chất phép cộng:
\choiceTF
{\True $a + b + c = (a + c) + b$.}
{$(a + b) + c = a + (b - c)$.}
{\True $125 + 456 + 875 = (125 + 875) + 456 = 1456$.}
{\True $(a + 0) = (0 + a) = a$.}
\loigiai{
\begin{itemize}
    \item Giao hoán và kết hợp: Đúng.
    \item Phép cộng không biến thành phép trừ. Sai.
    \item Tính thuận tiện đúng.
    \item Cộng với 0: Đúng.
\end{itemize}
\begin{itemize}
    \item a) \textbf{Đúng}.
    \item b) \textbf{Sai}.
    \item c) \textbf{Đúng}.
    \item d) \textbf{Đúng}.
\end{itemize}
}
\end{ex}

\begin{ex} %[4D5V24]
Tính giá trị các biểu thức:
$A = 250 + 480 + 750$ và $B = 360 + 640 + 480$.
Xét tính đúng sai:
\choiceTF
{\True $A = 1480$.}
{\True $B = 1480$.}
{Giá trị biểu thức $A$ lớn hơn giá trị biểu thức $B$.}
{\True $A = B$.}
\loigiai{
\begin{itemize}
    \item $A = (250 + 750) + 480 = 1000 + 480 = 1480$. Đúng.
    \item $B = (360 + 640) + 480 = 1000 + 480 = 1480$. Đúng.
    \item $A = B$ nên khẳng định $A > B$ là Sai.
    \item $A = B$. Đúng.
\end{itemize}
\begin{itemize}
    \item a) \textbf{Đúng}.
    \item b) \textbf{Đúng}.
    \item c) \textbf{Sai}.
    \item d) \textbf{Đúng}.
\end{itemize}
}
\end{ex}

\textbf{Phần 3. Trả lời ngắn}
\setcounterpageref{ex}{0}

\begin{ex} %[4D5H24]
Tính thuận tiện: $340 + 256 + 660$.
\shortans{1256}
\loigiai{
\textbf{Đáp số:} $1256$\\
$(340 + 660) + 256 = 1256$.
}
\end{ex}

\begin{ex} %[4D5H24]
Tính thuận tiện: $1500 + 2800 + 8500$.
\shortans{12800}
\loigiai{
\textbf{Đáp số:} $12800$\\
$(1500 + 8500) + 2800 = 12\,800$.
}
\end{ex}

\begin{ex} %[4D5V24]
Tính nhanh: $45 + 189 + 55 + 211$.
\shortans{500}
\loigiai{
\textbf{Đáp số:} $500$\\
$(45 + 55) + (189 + 211) = 100 + 400 = 500$.
}
\end{ex}

\begin{ex} %[4D5V24]
Tính nhanh: $12\,000 + 34\,000 + 88\,000 + 66\,000$.
\shortans{200000}
\loigiai{
\textbf{Đáp số:} $200000$\\
$(12\,000 + 88\,000) + (34\,000 + 66\,000) = 100\,000 + 100\,000 = 200\,000$.
}
\end{ex}

\begin{ex} %[4D5V24]
Tính tổng của 20 số tự nhiên liên tiếp từ 1 đến 20.
\shortans{210}
\loigiai{
\textbf{Đáp số:} $210$\\
$(1 + 20) \times 20 : 2 = 210$.
}
\end{ex}

\begin{ex} %[4D5C24]
Tính tổng: $11 + 22 + 33 + 44 + 55 + 66 + 77 + 88 + 99$.
\shortans{495}
\loigiai{
\textbf{Đáp số:} $495$\\
$11 \times (1 + 2 + \dots + 9) = 11 \times 45 = 495$.
}
\end{ex}

\newpage

\TaoPhieuKiemTra{Chủ đề 5 -- Phép cộng và phép trừ}{25}{Tìm hai số khi biết tổng và hiệu của hai số đó}{https://loptoanthaylinh.vn/thi-online/toan-4/bai-25}
\textbf{Phần 1. Trắc nghiệm lựa chọn}
\setcounterpageref{ex}{0}

\begin{ex} %[4D5N25]
Công thức tìm số lớn khi biết tổng $T$ và hiệu $H$ của hai số là
\choice
{$(\text{Tổng} - \text{Hiệu}) : 2$}
{$\text{Tổng} + \text{Hiệu}$}
{$\text{Tổng} - \text{Hiệu}$}
{\True $(	ext{Tổng} + \text{Hiệu}) : 2$}
\loigiai{
Số lớn $= (\text{Tổng} + \text{Hiệu}) : 2$.\\
\textbf{Đáp án của câu hỏi là D}
}
\end{ex}

\begin{ex} %[4D5N25]
Công thức tìm số bé khi biết tổng $T$ và hiệu $H$ của hai số là
\choice
{\True $(	ext{Tổng} - \text{Hiệu}) : 2$}
{$(\text{Tổng} + \text{Hiệu}) : 2$}
{$\text{Tổng} - \text{Hiệu}$}
{$\text{Số lớn} + \text{Hiệu}$}
\loigiai{
Số bé $= (\text{Tổng} - \text{Hiệu}) : 2$.\\
\textbf{Đáp án của câu hỏi là A}
}
\end{ex}

\begin{ex} %[4D5H25]
Hai số có tổng bằng $70$ và hiệu bằng $20$. Số lớn là
\choice
{$40$}
{\True $45$}
{$50$}
{$35$}
\loigiai{
Số lớn $= (70 + 20) : 2 = 45$.\\
\textbf{Đáp án của câu hỏi là B}
}
\end{ex}

\begin{ex} %[4D5H25]
Hai số có tổng bằng $90$ và hiệu bằng $30$. Số bé là
\choice
{$25$}
{$35$}
{\True $30$}
{$60$}
\loigiai{
Số bé $= (90 - 30) : 2 = 30$.\\
\textbf{Đáp án của câu hỏi là C}
}
\end{ex}

\begin{ex} %[4D5H25]
Cả hai lớp 4A và 4B có tất cả $68$ học sinh. Lớp 4A có nhiều hơn lớp 4B là $4$ học sinh. Số học sinh của lớp 4A là
\choice
{$32$ học sinh}
{\True $36$ học sinh}
{$34$ học sinh}
{$38$ học sinh}
\loigiai{
Lớp 4A: $(68 + 4) : 2 = 36$ học sinh.\\
\textbf{Đáp án của câu hỏi là B}
}
\end{ex}

\begin{ex} %[4D5H25]
Một mảnh đất hình chữ nhật có nửa chu vi là $48\text{ m}$, chiều dài hơn chiều rộng $12\text{ m}$. Chiều dài của mảnh đất là
\choice
{\True $30\text{ m}$}
{$28\text{ m}$}
{$32\text{ m}$}
{$36\text{ m}$}
\loigiai{
Chiều dài $= (48 + 12) : 2 = 30\text{ m}$.\\
\textbf{Đáp án của câu hỏi là A}
}
\end{ex}

\begin{ex} %[4D5V25]
Một mảnh đất hình chữ nhật có chu vi $100\text{ m}$, chiều dài hơn chiều rộng $10\text{ m}$. Diện tích mảnh đất đó là
\choice
{$500\text{ m}^2$}
{$700\text{ m}^2$}
{\True $600\text{ m}^2$}
{$800\text{ m}^2$}
\loigiai{
Nửa chu vi: $100 : 2 = 50\text{ m}$.\\
Chiều dài: $(50 + 10) : 2 = 30\text{ m}$.\\
Chiều rộng: $50 - 30 = 20\text{ m}$.\\
Diện tích: $30 \times 20 = 600\text{ m}^2$.\\
\textbf{Đáp án của câu hỏi là C}
}
\end{ex}

\begin{ex} %[4D5V25]
Hai thùng chứa tất cả $150\text{ l}$ dầu. Nếu chuyển $15\text{ l}$ dầu từ thùng thứ nhất sang thùng thứ hai thì hai thùng có số lít dầu bằng nhau. Lúc đầu thùng thứ nhất chứa số lít dầu là
\choice
{$80\text{ l}$}
{$85\text{ l}$}
{$95\text{ l}$}
{\True $90\text{ l}$}
\loigiai{
Hiệu số dầu ban đầu: $15 \times 2 = 30\text{ l}$.\\
Thùng thứ nhất: $(150 + 30) : 2 = 90\text{ l}$.\\
\textbf{Đáp án của câu hỏi là D}
}
\end{ex}

\begin{ex} %[4D5V25]
Hai anh em có tất cả $45$ viên bi. Nếu anh cho em $5$ viên bi thì anh vẫn nhiều hơn em $3$ viên bi. Số bi của anh lúc đầu là
\choice
{\True $29$ viên}
{$27$ viên}
{$28$ viên}
{$30$ viên}
\loigiai{
Hiệu số bi lúc đầu: $5 \times 2 + 3 = 13$ viên.\\
Số bi của anh lúc đầu: $(45 + 13) : 2 = 29$ viên.\\
\textbf{Đáp án của câu hỏi là A}
}
\end{ex}

\begin{ex} %[4D5V25]
Tổng số tuổi của hai bố con hiện nay là $48$ tuổi. Bố hơn con $28$ tuổi. Tuổi của con hiện nay là
\choice
{$8$ tuổi}
{$12$ tuổi}
{\True $10$ tuổi}
{$14$ tuổi}
\loigiai{
Tuổi con: $(48 - 28) : 2 = 10$ tuổi.\\
\textbf{Đáp án của câu hỏi là C}
}
\end{ex}

\begin{ex} %[4D5V25]
Hai số chẵn liên tiếp có tổng bằng $206$. Số lớn là
\choice
{$102$}
{$106$}
{$108$}
{\True $104$}
\loigiai{
Hiệu hai số chẵn liên tiếp là 2.\\
Số lớn: $(206 + 2) : 2 = 104$.\\
\textbf{Đáp án của câu hỏi là D}
}
\end{ex}

\begin{ex} %[4D5C25]
Một trường tiểu học có $640$ học sinh. Số học sinh nam nhiều hơn số học sinh nữ là $40$ em. Sau đó trường nhận thêm $15$ học sinh nữ và $5$ học sinh nam. Lúc này số học sinh nam nhiều hơn số học sinh nữ là
\choice
{$20$ em}
{\True $30$ em}
{$35$ em}
{$40$ em}
\loigiai{
Hiệu mới: $40 + 5 - 15 = 30$ em.\\
\textbf{Đáp án của câu hỏi là B}
}
\end{ex}

\textbf{Phần 2. Trắc nghiệm đúng sai}
\setcounterpageref{ex}{0}

\begin{ex} %[4D5H25]
Hai số có tổng là $120$ và hiệu là $40$. Xét tính đúng sai:
\choiceTF
{\True Số lớn là $80$.}
{\True Số bé là $40$.}
{Số lớn gấp $3$ lần số bé.}
{\True Số lớn gấp $2$ lần số bé.}
\loigiai{
\begin{itemize}
    \item Số lớn: $(120 + 40) : 2 = 80$. Đúng.
    \item Số bé: $(120 - 40) : 2 = 40$. Đúng.
    \item $80 : 40 = 2$ (gấp 2 lần, không phải 3 lần). Ý 3 Sai, Ý 4 Đúng.
\end{itemize}
\begin{itemize}
    \item a) \textbf{Đúng}.
    \item b) \textbf{Đúng}.
    \item c) \textbf{Sai}.
    \item d) \textbf{Đúng}.
\end{itemize}
}
\end{ex}

\begin{ex} %[4D5V25]
Một thửa ruộng hình chữ nhật có chu vi $140\text{ m}$, chiều dài hơn chiều rộng $20\text{ m}$. Xét tính đúng sai:
\choiceTF
{\True Nửa chu vi thửa ruộng là $70\text{ m}$.}
{\True Chiều dài thửa ruộng là $45\text{ m}$.}
{Chiều rộng thửa ruộng là $35\text{ m}$.}
{\True Diện tích thửa ruộng là $1125\text{ m}^2$.}
\loigiai{
\begin{itemize}
    \item Nửa chu vi: $140 : 2 = 70\text{ m}$. Đúng.
    \item Chiều dài: $(70 + 20) : 2 = 45\text{ m}$. Đúng.
    \item Chiều rộng: $70 - 45 = 25\text{ m} \neq 35\text{ m}$. Sai.
    \item Diện tích: $45 \times 25 = 1125\text{ m}^2$. Đúng.
\end{itemize}
\begin{itemize}
    \item a) \textbf{Đúng}.
    \item b) \textbf{Đúng}.
    \item c) \textbf{Sai}.
    \item d) \textbf{Đúng}.
\end{itemize}
}
\end{ex}

\textbf{Phần 3. Trả lời ngắn}
\setcounterpageref{ex}{0}

\begin{ex} %[4D5H25]
Tổng hai số là $86$, hiệu hai số là $16$. Tìm số lớn.
\shortans{51}
\loigiai{
\textbf{Đáp số:} $51$\\
$(86 + 16) : 2 = 51$.
}
\end{ex}

\begin{ex} %[4D5H25]
Tổng hai số là $100$, hiệu hai số là $34$. Tìm số bé.
\shortans{33}
\loigiai{
\textbf{Đáp số:} $33$\\
$(100 - 34) : 2 = 33$.
}
\end{ex}

\begin{ex} %[4D5V25]
Hai bạn An và Bình có tất cả $54$ cái kẹo, An có nhiều hơn Bình $12$ cái kẹo. Tính số kẹo của An.
\shortans{33}
\loigiai{
\textbf{Đáp số:} $33$\\
$(54 + 12) : 2 = 33$.
}
\end{ex}

\begin{ex} %[4D5V25]
Hai số lẻ liên tiếp có tổng bằng $152$. Tìm số bé.
\shortans{75}
\loigiai{
\textbf{Đáp số:} $75$\\
$(152 - 2) : 2 = 75$.
}
\end{ex}

\begin{ex} %[4D5V25]
Tuổi bố và tuổi con cộng lại bằng $50$ tuổi, bố hơn con $26$ tuổi. Tính tuổi của bố.
\shortans{38}
\loigiai{
\textbf{Đáp số:} $38$\\
$(50 + 26) : 2 = 38$ tuổi.
}
\end{ex}

\begin{ex} %[4D5C25]
Một hình chữ nhật có chu vi $80\text{ cm}$. Nếu tăng chiều rộng thêm $5\text{ cm}$ và giảm chiều dài đi $5\text{ cm}$ thì hình chữ nhật trở thành hình vuông. Tính diện tích hình chữ nhật ban đầu theo xăng-ti-mét vuông.
\shortans{375}
\loigiai{
\textbf{Đáp số:} $375$\\
Nửa chu vi: $40\text{ cm}$.\\
Chiều dài hơn chiều rộng: $5 + 5 = 10\text{ cm}$.\\
Dài: $25\text{ cm}$, Rộng: $15\text{ cm}$.\\
Diện tích: $25 \times 15 = 375\text{ cm}^2$.
}
\end{ex}

\newpage

\TaoPhieuKiemTra{Chủ đề 5 -- Phép cộng và phép trừ}{26}{Luyện tập chung phép cộng trừ}{https://loptoanthaylinh.vn/thi-online/toan-4/bai-26}
\textbf{Phần 1. Trắc nghiệm lựa chọn}
\setcounterpageref{ex}{0}

\begin{ex} %[4D5N26]
Tính chất giao hoán của phép cộng là
\choice
{$a - b = b - a$}
{$(a + b) + c = a + (b + c)$}
{$a + 0 = a$}
{\True $a + b = b + a$}
\loigiai{
$a + b = b + a$.\\
\textbf{Đáp án của câu hỏi là D}
}
\end{ex}

\begin{ex} %[4D5N26]
Kết quả của phép tính $450\,000 + 350\,000$ là
\choice
{$700\,000$}
{$900\,000$}
{\True $800\,000$}
{$850\,000$}
\loigiai{
$450\,000 + 350\,000 = 800\,000$.\\
\textbf{Đáp án của câu hỏi là C}
}
\end{ex}

\begin{ex} %[4D5H26]
Hiệu của $1\,000\,000$ và $250\,000$ là
\choice
{$850\,000$}
{$650\,000$}
{$700\,000$}
{\True $750\,000$}
\loigiai{
$1\,000\,000 - 250\,000 = 750\,000$.\\
\textbf{Đáp án của câu hỏi là D}
}
\end{ex}

\begin{ex} %[4D5H26]
Tìm hai số biết tổng là $80$ và hiệu là $14$. Số lớn là
\choice
{$44$}
{$50$}
{\True $47$}
{$52$}
\loigiai{
$(80 + 14) : 2 = 47$.\\
\textbf{Đáp án của câu hỏi là C}
}
\end{ex}

\begin{ex} %[4D5H26]
Tìm $x$ biết: $x + 45\,000 = 120\,000$.
\choice
{$65\,000$}
{\True $75\,000$}
{$85\,000$}
{$70\,000$}
\loigiai{
$x = 120\,000 - 45\,000 = 75\,000$.\\
\textbf{Đáp án của câu hỏi là B}
}
\end{ex}

\begin{ex} %[4D5H26]
Tính thuận tiện: $250 + 184 + 750$.
\choice
{\True $1184$}
{$1084$}
{$1284$}
{$1194$}
\loigiai{
$(250 + 750) + 184 = 1000 + 184 = 1184$.\\
\textbf{Đáp án của câu hỏi là A}
}
\end{ex}

\begin{ex} %[4D5V26]
Hai thùng dầu chứa tất cả $240\text{ l}$. Thùng thứ nhất chứa nhiều hơn thùng thứ hai $40\text{ l}$. Số lít dầu ở thùng thứ hai là
\choice
{$90\text{ l}$}
{$110\text{ l}$}
{$120\text{ l}$}
{\True $100\text{ l}$}
\loigiai{
$(240 - 40) : 2 = 100\text{ l}$.\\
\textbf{Đáp án của câu hỏi là D}
}
\end{ex}

\begin{ex} %[4D5V26]
Một khu đất hình chữ nhật có chu vi $160\text{ m}$, chiều dài hơn chiều rộng $20\text{ m}$. Diện tích khu đất đó là
\choice
{\True $1500\text{ m}^2$}
{$1200\text{ m}^2$}
{$1600\text{ m}^2$}
{$1800\text{ m}^2$}
\loigiai{
Nửa chu vi: $80\text{ m}$.\\
Chiều dài: $(80 + 20) : 2 = 50\text{ m}$.\\
Chiều rộng: $80 - 50 = 30\text{ m}$.\\
Diện tích: $50 \times 30 = 1500\text{ m}^2$.\\
\textbf{Đáp án của câu hỏi là A}
}
\end{ex}

\begin{ex} %[4D5V26]
Tính nhanh giá trị biểu thức: $12\,500 + 34\,800 + 87\,500 + 65\,200$.
\choice
{$190\,000$}
{\True $200\,000$}
{$210\,000$}
{$220\,000$}
\loigiai{
$(12\,500 + 87\,500) + (34\,800 + 65\,200) = 100\,000 + 100\,000 = 200\,000$.\\
\textbf{Đáp án của câu hỏi là B}
}
\end{ex}

\begin{ex} %[4D5V26]
Một nhà máy đợt 1 may được $45\,000$ chiếc áo, đợt 2 may được $52\,000$ chiếc áo. Để đạt chỉ tiêu $120\,000$ chiếc áo thì đợt 3 nhà máy cần may số chiếc áo là
\choice
{$22\,000$ chiếc}
{$24\,000$ chiếc}
{\True $23\,000$ chiếc}
{$25\,000$ chiếc}
\loigiai{
$120\,000 - (45\,000 + 52\,000) = 23\,000$ chiếc.\\
\textbf{Đáp án của câu hỏi là C}
}
\end{ex}

\begin{ex} %[4D5V26]
Tuổi của mẹ và Lan cộng lại là $42$ tuổi. Mẹ sinh Lan khi mẹ $26$ tuổi. Tuổi của mẹ hiện nay là
\choice
{\True $34$ tuổi}
{$32$ tuổi}
{$36$ tuổi}
{$38$ tuổi}
\loigiai{
Mẹ hơn Lan 26 tuổi.\\
Tuổi mẹ: $(42 + 26) : 2 = 34$ tuổi.\\
\textbf{Đáp án của câu hỏi là A}
}
\end{ex}

\begin{ex} %[4D5C26]
Tìm hai số chẵn có tổng bằng $210$, biết giữa chúng có đúng $4$ số chẵn khác. Số lớn là
\choice
{$110$}
{\True $110$}
{$112$}
{$114$}
\loigiai{
Giữa chúng có 4 số chẵn khác nên hiệu hai số là $(4 + 1) \times 2 = 10$.\\
Số lớn là: $(210 + 10) : 2 = 110$.\\
\textbf{Đáp án của câu hỏi là B}
}
\end{ex}

\textbf{Phần 2. Trắc nghiệm đúng sai}
\setcounterpageref{ex}{0}

\begin{ex} %[4D5H26]
Xét tính đúng sai của các phép toán:
\choiceTF
{\True $125\,000 + 875\,000 = 1\,000\,000$.}
{$2\,000\,000 - 850\,000 = 1\,250\,000$.}
{\True $2\,000\,000 - 850\,000 = 1\,150\,000$.}
{\True $(a + b) + c = a + (b + c)$.}
\loigiai{
\begin{itemize}
    \item $125\,000 + 875\,000 = 1\,000\,000$. Đúng.
    \item $2\,000\,000 - 850\,000 = 1\,150\,000 \neq 1\,250\,000$. Ý 2 Sai, Ý 3 Đúng.
    \item Tính chất kết hợp: Đúng.
\end{itemize}
\begin{itemize}
    \item a) \textbf{Đúng}.
    \item b) \textbf{Sai}.
    \item c) \textbf{Đúng}.
    \item d) \textbf{Đúng}.
\end{itemize}
}
\end{ex}

\begin{ex} %[4D5V26]
Hai bao gạo có tất cả $120\text{ kg}$. Bao 1 nặng hơn bao 2 là $20\text{ kg}$. Xét tính đúng sai:
\choiceTF
{\True Bao 1 nặng $70\text{ kg}$.}
{\True Bao 2 nặng $50\text{ kg}$.}
{Nếu lấy ở bao 1 ra $10\text{ kg}$ thì hai bao bằng nhau.}
{\True Nếu lấy ở bao 1 chuyển sang bao 2 là $10\text{ kg}$ thì hai bao bằng nhau.}
\loigiai{
\begin{itemize}
    \item Bao 1: $(120 + 20) : 2 = 70\text{ kg}$. Đúng.
    \item Bao 2: $120 - 70 = 50\text{ kg}$. Đúng.
    \item Lấy ở bao 1 bỏ đi 10 kg thì bao 1 còn 60 kg, bao 2 là 50 kg (vẫn không bằng nhau). Sai.
    \item Chuyển 10 kg sang bao 2 thì cả 2 bao đều có 60 kg. Đúng.
\end{itemize}
\begin{itemize}
    \item a) \textbf{Đúng}.
    \item b) \textbf{Đúng}.
    \item c) \textbf{Sai}.
    \item d) \textbf{Đúng}.
\end{itemize}
}
\end{ex}

\textbf{Phần 3. Trả lời ngắn}
\setcounterpageref{ex}{0}

\begin{ex} %[4D5H26]
Tính tổng: $1\,250\,000 + 3\,750\,000$.
\shortans{5000000}
\loigiai{
\textbf{Đáp số:} $5000000$\\
$5\,000\,000$.
}
\end{ex}

\begin{ex} %[4D5H26]
Hai số có tổng là $96$, hiệu là $24$. Tìm số lớn.
\shortans{60}
\loigiai{
\textbf{Đáp số:} $60$\\
$(96 + 24) : 2 = 60$.
}
\end{ex}

\begin{ex} %[4D5V26]
Tính thuận tiện: $450 + 890 + 550$.
\shortans{1890}
\loigiai{
\textbf{Đáp số:} $1890$\\
$(450 + 550) + 890 = 1890$.
}
\end{ex}

\begin{ex} %[4D5V26]
Hai số lẻ có tổng là $180$, biết giữa chúng có $2$ số lẻ khác. Tìm số bé.
\shortans{87}
\loigiai{
\textbf{Đáp số:} $87$\\
Hiệu: $(2 + 1) \times 2 = 6 \Rightarrow (180 - 6) : 2 = 87$.
}
\end{ex}

\begin{ex} %[4D5V26]
Một hình chữ nhật có nửa chu vi là $60\text{ cm}$, chiều dài hơn chiều rộng $16\text{ cm}$. Tính diện tích hình chữ nhật theo xăng-ti-mét vuông.
\shortans{836}
\loigiai{
\textbf{Đáp số:} $836$\\
Dài: $38\text{ cm}$, Rộng: $22\text{ cm} \Rightarrow 38 \times 22 = 836\text{ cm}^2$.
}
\end{ex}

\begin{ex} %[4D5C26]
Tính tổng: $100 + 200 + 300 + \dots + 1000$.
\shortans{5500}
\loigiai{
\textbf{Đáp số:} $5500$\\
$(100 + 1000) \times 10 : 2 = 5500$.
}
\end{ex}

\newpage

% =========================================================================
% CHỦ ĐỀ 6 - ĐƯỜNG THẲNG VUÔNG GÓC. ĐƯỜNG THẲNG SONG SONG
% =========================================================================
\section{Đường thẳng vuông góc. đường thẳng song song}

\TaoPhieuKiemTra{Chủ đề 6 -- Đường thẳng vuông góc. đường thẳng song song}{27}{Hai đường thẳng vuông góc}{https://loptoanthaylinh.vn/thi-online/toan-4/bai-27}
\textbf{Phần 1. Trắc nghiệm lựa chọn}
\setcounterpageref{ex}{0}

\begin{ex} %[4H6N27]
Hai đường thẳng vuông góc với nhau tạo thành mấy góc vuông?
\choice
{\True $4$ góc vuông}
{$1$ góc vuông}
{$2$ góc vuông}
{$3$ góc vuông}
\loigiai{
Hai đường thẳng vuông góc cắt nhau tạo thành 4 góc vuông.\\
\textbf{Đáp án của câu hỏi là A}
}
\end{ex}

\begin{ex} %[4H6N27]
Dụng cụ thường dùng để kiểm tra hai đường thẳng có vuông góc với nhau hay không là
\choice
{Thước dây}
{Com-pa}
{\True Ê-ke}
{Thước cuộn}
\loigiai{
Ê-ke có 1 góc vuông dùng để kiểm tra hai đường thẳng vuông góc.\\
\textbf{Đáp án của câu hỏi là C}
}
\end{ex}

\begin{ex} %[4H6H27]
Trong hình chữ nhật $ABCD$, cặp cạnh nào sau đây vuông góc với nhau?
\choice
{$AB$ và $CD$}
{$AD$ và $BC$}
{$AC$ và $BD$}
{\True $AB$ và $BC$}
\loigiai{
Trong hình chữ nhật, hai cạnh kề nhau $AB$ và $BC$ vuông góc với nhau.\\
\textbf{Đáp án của câu hỏi là D}
}
\end{ex}

\begin{ex} %[4H6H27]
Trong hình vuông $MNPQ$, có bao nhiêu cặp cạnh kề nhau vuông góc với nhau?
\choice
{$2$ cặp}
{$3$ cặp}
{\True $4$ cặp}
{$1$ cặp}
\loigiai{
Có 4 cặp cạnh vuông góc tại 4 đỉnh $M, N, P, Q$.\\
\textbf{Đáp án của câu hỏi là C}
}
\end{ex}

\begin{ex} %[4H6H27]
Hai đường chéo của hình thoi có tính chất nào sau đây?
\choice
{Song song với nhau}
{\True Vuông góc với nhau}
{Không cắt nhau}
{Bằng nhau}
\loigiai{
Hai đường chéo của hình thoi vuông góc với nhau.\\
\textbf{Đáp án của câu hỏi là B}
}
\end{ex}

\begin{ex} %[4H6H27]
Cho tam giác $ABC$ có góc $A$ là góc vuông. Cặp cạnh vuông góc với nhau là
\choice
{$AB$ và $BC$}
{$AC$ và $BC$}
{\True $AB$ và $AC$}
{Không có}
\loigiai{
Hai cạnh góc vuông là $AB$ và $AC$.\\
\textbf{Đáp án của câu hỏi là C}
}
\end{ex}

\begin{ex} %[4H6V27]
Hình chữ nhật $ABCD$ có mấy góc vuông?
\choice
{\True $4$ góc vuông}
{$2$ góc vuông}
{$3$ góc vuông}
{$1$ góc vuông}
\loigiai{
Hình chữ nhật có 4 góc vuông.\\
\textbf{Đáp án của câu hỏi là A}
}
\end{ex}

\begin{ex} %[4H6V27]
Cho đoạn thẳng $AB$. Qua điểm $M$ nằm ngoài đoạn thẳng $AB$, ta có thể vẽ được bao nhiêu đường thẳng vuông góc với $AB$?
\choice
{$0$ đường thẳng}
{$2$ đường thẳng}
{Vô số đường thẳng}
{\True $1$ đường thẳng duy nhất}
\loigiai{
Qua một điểm cho trước, chỉ vẽ được duy nhất 1 đường thẳng vuông góc với đường thẳng đã cho.\\
\textbf{Đáp án của câu hỏi là D}
}
\end{ex}

\begin{ex} %[4H6V27]
Hình thang vuông là hình thang có
\choice
{1 góc tù}
{\True 1 cạnh bên vuông góc với hai đáy (có 2 góc vuông)}
{4 góc vuông}
{2 đường chéo vuông góc}
\loigiai{
Hình thang vuông có 1 cạnh bên vuông góc với hai đáy.\\
\textbf{Đáp án của câu hỏi là B}
}
\end{ex}

\begin{ex} %[4H6V27]
Hai kim đồng hồ lúc 3 giờ đúng tạo thành hai tia
\choice
{Song song}
{Thẳng hàng}
{Trùng nhau}
{\True Vuông góc với nhau}
\loigiai{
Lúc 3 giờ đúng hai kim vuông góc với nhau ($90^\circ$).\\
\textbf{Đáp án của câu hỏi là D}
}
\end{ex}

\begin{ex} %[4H6V27]
Trong một hình tam giác vuông cân, hai cạnh góc vuông có độ dài
\choice
{Khác nhau}
{\True Bằng nhau}
{Gấp đôi nhau}
{Không so sánh được}
\loigiai{
Tam giác vuông cân có hai cạnh góc vuông bằng nhau.\\
\textbf{Đáp án của câu hỏi là B}
}
\end{ex}

\begin{ex} %[4H6C27]
Cho hình vuông $ABCD$. Hai đường chéo $AC$ và $BD$ cắt nhau tại $O$. Tại điểm $O$ có bao nhiêu góc vuông được tạo thành bởi hai đường chéo?
\choice
{\True $4$ góc vuông}
{$2$ góc vuông}
{$1$ góc vuông}
{$0$ góc vuông}
\loigiai{
Hai đường chéo hình vuông vuông góc tại $O$ tạo thành 4 góc vuông.\\
\textbf{Đáp án của câu hỏi là A}
}
\end{ex}

\textbf{Phần 2. Trắc nghiệm đúng sai}
\setcounterpageref{ex}{0}

\begin{ex} %[4H6H27]
Xét tính đúng sai về hai đường thẳng vuông góc:
\choiceTF
{\True Hai đường thẳng vuông góc cắt nhau tạo thành các góc có số đo $90^\circ$.}
{\True Dùng ê-ke để kiểm tra và vẽ đường thẳng vuông góc.}
{Hai đường thẳng cắt nhau bất kì luôn vuông góc với nhau.}
{\True Trong hình chữ nhật, mỗi cạnh đều vuông góc với hai cạnh kề với nó.}
\loigiai{
\begin{itemize}
    \item Vuông góc tạo góc $90^\circ$: Đúng.
    \item Dùng ê-ke: Đúng.
    \item Hai đường thẳng cắt nhau có thể tạo góc nhọn/tù (không nhất thiết vuông góc). Sai.
    \item Mỗi cạnh vuông góc với 2 cạnh kề: Đúng.
\end{itemize}
\begin{itemize}
    \item a) \textbf{Đúng}.
    \item b) \textbf{Đúng}.
    \item c) \textbf{Sai}.
    \item d) \textbf{Đúng}.
\end{itemize}
}
\end{ex}

\begin{ex} %[4H6V27]
Cho hình chữ nhật $MNPQ$. Xét tính đúng sai:
\choiceTF
{\True Cạnh $MN$ vuông góc với cạnh $MQ$.}
{\True Cạnh $NP$ vuông góc với cạnh $PQ$.}
{Cạnh $MN$ vuông góc với cạnh $PQ$.}
{\True Cạnh $MQ$ vuông góc với cạnh $QP$.}
\loigiai{
\begin{itemize}
    \item $MN \perp MQ$: Đúng.
    \item $NP \perp PQ$: Đúng.
    \item $MN$ song song với $PQ$ (không vuông góc). Sai.
    \item $MQ \perp QP$: Đúng.
\end{itemize}
\begin{itemize}
    \item a) \textbf{Đúng}.
    \item b) \textbf{Đúng}.
    \item c) \textbf{Sai}.
    \item d) \textbf{Đúng}.
\end{itemize}
}
\end{ex}

\textbf{Phần 3. Trả lời ngắn}
\setcounterpageref{ex}{0}

\begin{ex} %[4H6H27]
Một hình chữ nhật có bao nhiêu góc vuông?
\shortans{4}
\loigiai{
\textbf{Đáp số:} $4$\\
Có 4 góc vuông.
}
\end{ex}

\begin{ex} %[4H6H27]
Hai đường thẳng vuông góc tạo thành bao nhiêu góc vuông tại điểm giao nhau?
\shortans{4}
\loigiai{
\textbf{Đáp số:} $4$\\
4 góc vuông.
}
\end{ex}

\begin{ex} %[4H6V27]
Một hình tam giác vuông có bao nhiêu góc vuông?
\shortans{1}
\loigiai{
\textbf{Đáp số:} $1$\\
1 góc vuông.
}
\end{ex}

\begin{ex} %[4H6V27]
Một hình chữ nhật có bao nhiêu cặp cạnh vuông góc với nhau?
\shortans{4}
\loigiai{
\textbf{Đáp số:} $4$\\
4 cặp cạnh kề nhau.
}
\end{ex}

\begin{ex} %[4H6V27]
Một hình thang vuông có mấy góc vuông?
\shortans{2}
\loigiai{
\textbf{Đáp số:} $2$\\
2 góc vuông.
}
\end{ex}

\begin{ex} %[4H6C27]
Trong một hình vuông, tính tổng số góc vuông được tạo bởi các cạnh của hình vuông và hai đường chéo cắt nhau của nó.
\shortans{8}
\loigiai{
\textbf{Đáp số:} $8$\\
4 góc vuông tại 4 đỉnh $+ 4$ góc vuông tại giao điểm hai đường chéo $= 8$ góc vuông.
}
\end{ex}

\newpage

\TaoPhieuKiemTra{Chủ đề 6 -- Đường thẳng vuông góc. đường thẳng song song}{28}{Thực hành vẽ hai đường thẳng vuông góc}{https://loptoanthaylinh.vn/thi-online/toan-4/bai-28}
\textbf{Phần 1. Trắc nghiệm lựa chọn}
\setcounterpageref{ex}{0}

\begin{ex} %[4H6N28]
Để vẽ đường thẳng đi qua điểm $O$ và vuông góc với đường thẳng $a$, bước đầu tiên ta đặt
\choice
{Thước thẳng bất kì}
{Com-pa tại điểm $O$}
{\True Một cạnh góc vuông của ê-ke trùng với đường thẳng $a$}
{Thước đo góc tại $O$}
\loigiai{
Đặt một cạnh góc vuông của ê-ke trùng với đường thẳng $a$.\\
\textbf{Đáp án của câu hỏi là C}
}
\end{ex}

\begin{ex} %[4H6N28]
Khi vẽ đường cao $AH$ của tam giác $ABC$, ta vẽ đường thẳng đi qua đỉnh $A$ và
\choice
{Song song với cạnh $BC$}
{\True Vuông góc với cạnh $BC$}
{Đi qua trung điểm của $BC$}
{Cắt cạnh $AB$}
\loigiai{
Đường cao kẻ từ $A$ vuông góc với cạnh đáy $BC$.\\
\textbf{Đáp án của câu hỏi là B}
}
\end{ex}

\begin{ex} %[4H6H28]
Hình tam giác có bao nhiêu đường cao?
\choice
{\True $3$ đường cao}
{$1$ đường cao}
{$2$ đường cao}
{$4$ đường cao}
\loigiai{
Mỗi tam giác có 3 đỉnh nên có 3 đường cao.\\
\textbf{Đáp án của câu hỏi là A}
}
\end{ex}

\begin{ex} %[4H6H28]
Trong tam giác vuông $ABC$ (vuông tại $A$), hai đường cao xuất phát từ đỉnh $B$ và đỉnh $C$ chính là
\choice
{Đường trung tuyến}
{Cạnh huyền $BC$}
{Không tồn tại}
{\True Hai cạnh góc vuông $BA$ và $CA$}
\loigiai{
Hai cạnh góc vuông chính là hai đường cao xuất phát từ hai đỉnh nhọn.\\
\textbf{Đáp án của câu hỏi là D}
}
\end{ex}

\begin{ex} %[4H6H28]
Để vẽ hình chữ nhật có chiều dài $5\text{ cm}$ và chiều rộng $3\text{ cm}$, ta cần vẽ các góc tại 4 đỉnh là
\choice
{Góc nhọn}
{Góc tù}
{\True Góc vuông}
{Góc bẹt}
\loigiai{
4 đỉnh của hình chữ nhật đều là góc vuông.\\
\textbf{Đáp án của câu hỏi là C}
}
\end{ex}

\begin{ex} %[4H6H28]
Dùng ê-ke để kiểm tra góc của một khung ảnh hình vuông, ta thấy góc đó
\choice
{Bé hơn góc vuông ê-ke}
{Lớn hơn góc vuông ê-ke}
{Không đo được}
{\True Khớp hoàn toàn với góc vuông của ê-ke}
\loigiai{
Khung ảnh hình vuông có các góc vuông $90^\circ$.\\
\textbf{Đáp án của câu hỏi là D}
}
\end{ex}

\begin{ex} %[4H6V28]
Vẽ đường cao tương ứng với đáy $CD$ của hình thang $ABCD$, ta vẽ đoạn thẳng hạ từ đỉnh $A$ hoặc $B$
\choice
{Song song với $CD$}
{Đi qua trung điểm $CD$}
{\True Vuông góc với đáy $CD$}
{Bằng cạnh bên}
\loigiai{
Đường cao vuông góc với cạnh đáy.\\
\textbf{Đáp án của câu hỏi là C}
}
\end{ex}

\begin{ex} %[4H6V28]
Cho điểm $M$ nằm ngoài đường thẳng $d$. Có thể vẽ được bao nhiêu đường thẳng đi qua $M$ và vuông góc với $d$?
\choice
{$0$}
{\True $1$}
{$2$}
{Vô số}
\loigiai{
Chỉ có duy nhất 1 đường thẳng.\\
\textbf{Đáp án của câu hỏi là B}
}
\end{ex}

\begin{ex} %[4H6V28]
Để vẽ một hình vuông cạnh $4\text{ cm}$, ta cần thực hiện vẽ bao nhiêu góc vuông liên tiếp?
\choice
{\True $4$ góc vuông (hoặc ít nhất 3 góc vuông)}
{$2$ góc vuông}
{$1$ góc vuông}
{$0$ góc vuông}
\loigiai{
Hình vuông có 4 góc vuông.\\
\textbf{Đáp án của câu hỏi là A}
}
\end{ex}

\begin{ex} %[4H6V28]
Khi vẽ đường cao của tam giác có góc tù ở đáy, chân đường cao nằm ở
\choice
{Nằm bên trong cạnh đáy}
{\True Nằm trên phần kéo dài của cạnh đáy}
{Trùng với một đỉnh}
{Không vẽ được}
\loigiai{
Chân đường cao nằm trên đường thẳng chứa cạnh đáy (phần kéo dài).\\
\textbf{Đáp án của câu hỏi là B}
}
\end{ex}

\begin{ex} %[4H6V28]
Để gấp một tờ giấy hình chữ nhật thành một góc vuông, ta có thể
\choice
{\True Gấp đôi mép thẳng của tờ giấy sao cho hai mép trùng nhau}
{Gấp mép giấy bất kì}
{Vò nhàu tờ giấy}
{Cắt xéo tờ giấy}
\loigiai{
Gấp đôi mép thẳng của tờ giấy tạo thành góc bẹt chia đôi là góc vuông.\\
\textbf{Đáp án của câu hỏi là A}
}
\end{ex}

\begin{ex} %[4H6C28]
Cho tam giác $ABC$. Qua mỗi đỉnh $A, B, C$ ta lần lượt vẽ đường thẳng vuông góc với cạnh đối diện. Ba đường thẳng này sẽ
\choice
{Song song với nhau}
{Tạo thành tam giác mới}
{Không bao giờ cắt nhau}
{\True Cùng cắt nhau tại một điểm (trực tâm)}
\loigiai{
Ba đường cao của một tam giác luôn đồng quy tại một điểm (gọi là trực tâm).\\
\textbf{Đáp án của câu hỏi là D}
}
\end{ex}

\textbf{Phần 2. Trắc nghiệm đúng sai}
\setcounterpageref{ex}{0}

\begin{ex} %[4H6H28]
Xét tính đúng sai về thao tác thực hành vẽ:
\choiceTF
{\True Sử dụng ê-ke để vẽ đường thẳng vuông góc với một đường thẳng cho trước.}
{\True Mỗi tam giác có 3 đường cao tương ứng với 3 cạnh đáy.}
{Tam giác vuông không có đường cao nào.}
{\True Trong tam giác vuông, hai cạnh góc vuông chính là hai đường cao.}
\loigiai{
\begin{itemize}
    \item Dùng ê-ke vẽ góc vuông: Đúng.
    \item Tam giác có 3 đường cao: Đúng.
    \item Tam giác vuông vẫn có 3 đường cao. Sai.
    \item Hai cạnh góc vuông là 2 đường cao: Đúng.
\end{itemize}
\begin{itemize}
    \item a) \textbf{Đúng}.
    \item b) \textbf{Đúng}.
    \item c) \textbf{Sai}.
    \item d) \textbf{Đúng}.
\end{itemize}
}
\end{ex}

\begin{ex} %[4H6V28]
Cho đoạn thẳng $AB = 6\text{ cm}$. Điểm $I$ là trung điểm của $AB$. Qua $I$ vẽ đường thẳng $d$ vuông góc với $AB$. Xét tính đúng sai:
\choiceTF
{\True Đoạn thẳng $IA = IB = 3\text{ cm}$.}
{\True Đường thẳng $d$ vuông góc với $AB$ tại điểm $I$.}
{Góc tạo bởi $d$ và $AB$ có số đo là $60^\circ$.}
{\True Đường thẳng $d$ là đường trung trực của đoạn thẳng $AB$.}
\loigiai{
\begin{itemize}
    \item $IA = IB = 3\text{ cm}$: Đúng.
    \item $d \perp AB$ tại $I$: Đúng.
    \item Góc tạo bởi $d$ và $AB$ là $90^\circ \neq 60^\circ$. Sai.
    \item Đường thẳng vuông góc tại trung điểm là đường trung trực: Đúng.
\end{itemize}
\begin{itemize}
    \item a) \textbf{Đúng}.
    \item b) \textbf{Đúng}.
    \item c) \textbf{Sai}.
    \item d) \textbf{Đúng}.
\end{itemize}
}
\end{ex}

\textbf{Phần 3. Trả lời ngắn}
\setcounterpageref{ex}{0}

\begin{ex} %[4H6H28]
Một hình tam giác có bao nhiêu đường cao?
\shortans{3}
\loigiai{
\textbf{Đáp số:} $3$\\
3 đường cao.
}
\end{ex}

\begin{ex} %[4H6H28]
Số góc vuông của thước ê-ke là bao nhiêu?
\shortans{1}
\loigiai{
\textbf{Đáp số:} $1$\\
1 góc vuông.
}
\end{ex}

\begin{ex} %[4H6V28]
Qua một điểm $A$ ở ngoài đường thẳng $d$, vẽ được bao nhiêu đường thẳng vuông góc với $d$?
\shortans{1}
\loigiai{
\textbf{Đáp số:} $1$\\
1 đường thẳng duy nhất.
}
\end{ex}

\begin{ex} %[4H6V28]
Một hình chữ nhật có bao nhiêu góc vuông?
\shortans{4}
\loigiai{
\textbf{Đáp số:} $4$\\
4 góc vuông.
}
\end{ex}

\begin{ex} %[4H6V28]
Một hình thang có thể có nhiều nhất bao nhiêu góc vuông?
\shortans{2}
\loigiai{
\textbf{Đáp số:} $2$\\
2 góc vuông (hình thang vuông).
}
\end{ex}

\begin{ex} %[4H6C28]
Trong một tam giác đều, số đo mỗi góc là bao nhiêu độ?
\shortans{60}
\loigiai{
\textbf{Đáp số:} $60$\\
$180 : 3 = 60^\circ$.
}
\end{ex}

\newpage

\TaoPhieuKiemTra{Chủ đề 6 -- Đường thẳng vuông góc. đường thẳng song song}{29}{Hai đường thẳng song song}{https://loptoanthaylinh.vn/thi-online/toan-4/bai-29}
\textbf{Phần 1. Trắc nghiệm lựa chọn}
\setcounterpageref{ex}{0}

\begin{ex} %[4H6N29]
Hai đường thẳng song song là hai đường thẳng
\choice
{Cắt nhau tại 1 điểm}
{Vuông góc với nhau}
{\True Không bao giờ cắt nhau (dù kéo dài mãi)}
{Trùng nhau}
\loigiai{
Hai đường thẳng song song không có điểm chung nào.\\
\textbf{Đáp án của câu hỏi là C}
}
\end{ex}

\begin{ex} %[4H6N29]
Hình nào sau đây thể hiện hình ảnh của hai đường thẳng song song trong thực tế?
\choice
{Hai kim đồng hồ lúc 3 giờ}
{Hai lưỡi kéo đang mở}
{Hai cạnh của góc vuông ê-ke}
{\True Hai thanh ray của đường tàu hỏa}
\loigiai{
Hai thanh ray đường tàu hỏa song song với nhau.\\
\textbf{Đáp án của câu hỏi là D}
}
\end{ex}

\begin{ex} %[4H6H29]
Trong hình chữ nhật $ABCD$, cặp cạnh nào sau đây song song với nhau?
\choice
{$AB$ và $BC$}
{$BC$ và $CD$}
{\True $AB$ và $CD$}
{$AB$ và $AD$}
\loigiai{
Hai cạnh đối diện $AB$ và $CD$ song song với nhau.\\
\textbf{Đáp án của câu hỏi là C}
}
\end{ex}

\begin{ex} %[4H6H29]
Hình bình hành có mấy cặp cạnh đối diện song song và bằng nhau?
\choice
{$1$ cặp}
{\True $2$ cặp}
{$3$ cặp}
{$4$ cặp}
\loigiai{
Hình bình hành có 2 cặp cạnh đối diện song song và bằng nhau.\\
\textbf{Đáp án của câu hỏi là B}
}
\end{ex}

\begin{ex} %[4H6H29]
Hình thang là hình tứ giác có
\choice
{\True 1 cặp cạnh đối diện song song (hai đáy)}
{2 cặp cạnh đối diện song song}
{4 cạnh bằng nhau}
{4 góc vuông}
\loigiai{
Hình thang có 1 cặp cạnh đối diện song song.\\
\textbf{Đáp án của câu hỏi là A}
}
\end{ex}

\begin{ex} %[4H6H29]
Trong hình vuông $ABCD$, cạnh $AD$ song song với cạnh nào?
\choice
{$AB$}
{$CD$}
{$AC$}
{\True $BC$}
\loigiai{
Cạnh $AD$ song song với cạnh đối diện $BC$.\\
\textbf{Đáp án của câu hỏi là D}
}
\end{ex}

\begin{ex} %[4H6V29]
Hai đường thẳng cùng vuông góc với một đường thẳng thứ ba thì chúng
\choice
{\True Song song với nhau}
{Cắt nhau}
{Vuông góc với nhau}
{Trùng nhau}
\loigiai{
Nếu $a \perp c$ và $b \perp c$ thì $a // b$.\\
\textbf{Đáp án của câu hỏi là A}
}
\end{ex}

\begin{ex} %[4H6V29]
Qua một điểm $M$ nằm ngoài đường thẳng $a$, có thể vẽ được bao nhiêu đường thẳng song song với $a$?
\choice
{$0$}
{$2$}
{Vô số}
{\True $1$ đường thẳng duy nhất}
\loigiai{
Qua 1 điểm ngoài đường thẳng chỉ vẽ được 1 đường thẳng song song duy nhất.\\
\textbf{Đáp án của câu hỏi là D}
}
\end{ex}

\begin{ex} %[4H6V29]
Khoảng cách giữa hai đường thẳng song song luôn
\choice
{Thay đổi}
{\True Không đổi (bằng nhau tại mọi điểm)}
{Bằng 0}
{Tăng dần}
\loigiai{
Khoảng cách giữa hai đường thẳng song song không đổi.\\
\textbf{Đáp án của câu hỏi là B}
}
\end{ex}

\begin{ex} %[4H6V29]
Trong các hình sau, hình nào KHÔNG có cặp cạnh nào song song?
\choice
{Hình chữ nhật}
{Hình bình hành}
{\True Hình tam giác}
{Hình thoi}
\loigiai{
Hình tam giác có 3 cạnh cắt nhau từng đôi một, không có cạnh song song.\\
\textbf{Đáp án của câu hỏi là C}
}
\end{ex}

\begin{ex} %[4H6V29]
Hai dòng kẻ ngang trên trang vở ô li là hình ảnh của
\choice
{Hai đường thẳng vuông góc}
{\True Hai đường thẳng song song}
{Hai đường thẳng cắt nhau}
{Hai tia chung gốc}
\loigiai{
Các dòng kẻ ngang song song với nhau.\\
\textbf{Đáp án của câu hỏi là B}
}
\end{ex}

\begin{ex} %[4H6C29]
Cho hình thoi $ABCD$. Cặp cạnh nào sau đây song song với nhau?
\choice
{\True $AB$ và $CD$ (cũng như $AD$ và $BC$)}
{$AB$ và $BC$}
{$AC$ và $BD$}
{$AD$ và $DC$}
\loigiai{
Hình thoi có các cặp cạnh đối diện song song: $AB // CD$ và $AD // BC$.\\
\textbf{Đáp án của câu hỏi là A}
}
\end{ex}

\textbf{Phần 2. Trắc nghiệm đúng sai}
\setcounterpageref{ex}{0}

\begin{ex} %[4H6H29]
Xét tính đúng sai về hai đường thẳng song song:
\choiceTF
{\True Hai đường thẳng song song không bao giờ cắt nhau.}
{\True Trong hình chữ nhật có 2 cặp cạnh đối diện song song.}
{Hai đường thẳng không cắt nhau thì luôn vuông góc.}
{\True Hai mép đối diện của cánh cửa lớp học là hình ảnh của hai đoạn thẳng song song.}
\loigiai{
\begin{itemize}
    \item Song song không cắt nhau: Đúng.
    \item Hình chữ nhật có 2 cặp cạnh song song: Đúng.
    \item Không cắt nhau là song song (không phải vuông góc). Sai.
    \item Hai mép cánh cửa song song: Đúng.
\end{itemize}
\begin{itemize}
    \item a) \textbf{Đúng}.
    \item b) \textbf{Đúng}.
    \item c) \textbf{Sai}.
    \item d) \textbf{Đúng}.
\end{itemize}
}
\end{ex}

\begin{ex} %[4H6V29]
Cho hình bình hành $ABCD$. Xét tính đúng sai:
\choiceTF
{\True Cạnh $AB$ song song với cạnh $CD$.}
{\True Cạnh $AD$ song song với cạnh $BC$.}
{Cạnh $AB$ song song với cạnh $BC$.}
{\True Độ dài cạnh $AB$ bằng độ dài cạnh $CD$.}
\loigiai{
\begin{itemize}
    \item $AB // CD$: Đúng.
    \item $AD // BC$: Đúng.
    \item $AB$ cắt $BC$ tại $B$ (không song song). Sai.
    \item $AB = CD$: Đúng.
\end{itemize}
\begin{itemize}
    \item a) \textbf{Đúng}.
    \item b) \textbf{Đúng}.
    \item c) \textbf{Sai}.
    \item d) \textbf{Đúng}.
\end{itemize}
}
\end{ex}

\textbf{Phần 3. Trả lời ngắn}
\setcounterpageref{ex}{0}

\begin{ex} %[4H6H29]
Hình chữ nhật có bao nhiêu cặp cạnh đối diện song song?
\shortans{2}
\loigiai{
\textbf{Đáp số:} $2$\\
2 cặp cạnh song song.
}
\end{ex}

\begin{ex} %[4H6H29]
Hình vuông có bao nhiêu cặp cạnh đối diện song song?
\shortans{2}
\loigiai{
\textbf{Đáp số:} $2$\\
2 cặp cạnh song song.
}
\end{ex}

\begin{ex} %[4H6V29]
Hình thang thông thường có bao nhiêu cặp cạnh đối diện song song?
\shortans{1}
\loigiai{
\textbf{Đáp số:} $1$\\
1 cặp cạnh song song (hai đáy).
}
\end{ex}

\begin{ex} %[4H6V29]
Qua một điểm $M$ ngoài đường thẳng $d$, vẽ được bao nhiêu đường thẳng song song với $d$?
\shortans{1}
\loigiai{
\textbf{Đáp số:} $1$\\
1 đường thẳng duy nhất.
}
\end{ex}

\begin{ex} %[4H6V29]
Hình bình hành có bao nhiêu cặp cạnh đối diện song song?
\shortans{2}
\loigiai{
\textbf{Đáp số:} $2$\\
2 cặp.
}
\end{ex}

\begin{ex} %[4H6C29]
Một hình lục giác đều có bao nhiêu cặp cạnh đối diện song song với nhau?
\shortans{3}
\loigiai{
\textbf{Đáp số:} $3$\\
3 cặp cạnh đối diện song song.
}
\end{ex}

\newpage

\TaoPhieuKiemTra{Chủ đề 6 -- Đường thẳng vuông góc. đường thẳng song song}{30}{Thực hành vẽ hai đường thẳng song song}{https://loptoanthaylinh.vn/thi-online/toan-4/bai-30}
\textbf{Phần 1. Trắc nghiệm lựa chọn}
\setcounterpageref{ex}{0}

\begin{ex} %[4H6N30]
Để vẽ đường thẳng đi qua điểm $M$ và song song với đường thẳng $a$, ta có thể thực hiện
\choice
{Dùng thước kẻ một đường bất kì}
{Dùng compa quay 1 vòng tròn}
{Dùng thước đo độ đo góc $45^\circ$}
{\True Vẽ đường thẳng phụ vuông góc với $a$, sau đó vẽ đường thẳng qua $M$ vuông góc với đường thẳng phụ đó}
\loigiai{
Hai đường thẳng cùng vuông góc với một đường thẳng thứ ba thì song song với nhau.\\
\textbf{Đáp án của câu hỏi là D}
}
\end{ex}

\begin{ex} %[4H6N30]
Dụng cụ thường kết hợp cùng thước thẳng để vẽ hai đường thẳng song song là
\choice
{Com-pa}
{\True Ê-ke}
{Kéo cắt giấy}
{Bút chì màu}
\loigiai{
Kết hợp thước kẻ và ê-ke để trượt vẽ các đường thẳng song song.\\
\textbf{Đáp án của câu hỏi là B}
}
\end{ex}

\begin{ex} %[4H6H30]
Để vẽ một hình chữ nhật có kích thước $6\text{ cm}$ và $4\text{ cm}$, ta có thể vẽ
\choice
{\True 2 cặp đoạn thẳng song song và vuông góc từng đôi một}
{3 đường thẳng song song}
{2 đoạn thẳng cắt nhau tại trung điểm}
{1 hình tam giác vuông}
\loigiai{
Vẽ các cặp cạnh song song và vuông góc nhau.\\
\textbf{Đáp án của câu hỏi là A}
}
\end{ex}

\begin{ex} %[4H6H30]
Khi trượt một cạnh của ê-ke dọc theo mép thước kẻ thẳng, cạnh góc vuông còn lại của ê-ke sẽ vạch ra các đường thẳng
\choice
{Cắt nhau}
{\True Song song với nhau}
{Vuông góc với nhau}
{Trùng nhau}
\loigiai{
Các đường thẳng vạch ra song song với nhau.\\
\textbf{Đáp án của câu hỏi là B}
}
\end{ex}

\begin{ex} %[4H6H30]
Để vẽ đường thẳng $b$ song song với đường thẳng $a$ và cách $a$ một khoảng $3\text{ cm}$, ta lấy một điểm $M$ cách $a$ đoạn $3\text{ cm}$ rồi vẽ đường thẳng qua $M$
\choice
{Vuông góc với $a$}
{Cắt $a$ tại góc $30^\circ$}
{\True Song song với $a$}
{Cắt $a$ tại góc $45^\circ$}
\loigiai{
Vẽ đường thẳng đi qua $M$ song song với $a$.\\
\textbf{Đáp án của câu hỏi là C}
}
\end{ex}

\begin{ex} %[4H6H30]
Muốn vẽ một hình bình hành $ABCD$, sau khi vẽ hai cạnh kề $AB$ và $AD$, ta vẽ đường thẳng qua $B$ song song với $AD$ và đường thẳng qua $D$ song song với
\choice
{\True $AB$}
{$AD$}
{$BD$}
{$AC$}
\loigiai{
Qua $D$ vẽ đường thẳng song song với $AB$.\\
\textbf{Đáp án của câu hỏi là A}
}
\end{ex}

\begin{ex} %[4H6V30]
Cho 3 đường thẳng phân biệt $a, b, c$. Nếu $a // b$ và $b // c$ thì
\choice
{$a \perp c$}
{$a$ cắt $c$}
{$a$ trùng $c$}
{\True $a // c$}
\loigiai{
Tính chất bắc cầu của quan hệ song song: $a // c$.\\
\textbf{Đáp án của câu hỏi là D}
}
\end{ex}

\begin{ex} %[4H6V30]
Trên lưới ô vuông, hai đường thẳng đi qua các đường kẻ ngang cách nhau 2 ô vuông là hai đường thẳng
\choice
{Vuông góc}
{Cắt nhau}
{\True Song song với nhau}
{Trùng nhau}
\loigiai{
Các đường kẻ ngang trên lưới ô vuông luôn song song với nhau.\\
\textbf{Đáp án của câu hỏi là C}
}
\end{ex}

\begin{ex} %[4H6V30]
Để vẽ một chiếc thang gỗ có các bậc thang đều nhau, người thợ mộc cần đóng các bậc thang
\choice
{Chéo nhau}
{\True Song song với nhau và vuông góc với hai thanh dọc}
{Vuông góc với nhau}
{Tự do}
\loigiai{
Các bậc thang song song với nhau.\\
\textbf{Đáp án của câu hỏi là B}
}
\end{ex}

\begin{ex} %[4H6V30]
Cho tam giác $ABC$. Qua đỉnh $A$ vẽ đường thẳng $d$ song song với cạnh $BC$. Có thể vẽ được bao nhiêu đường thẳng $d$ như vậy?
\choice
{\True $1$ đường thẳng duy nhất}
{$0$}
{$2$}
{Vô số}
\loigiai{
Chỉ có duy nhất 1 đường thẳng qua $A$ song song với $BC$.\\
\textbf{Đáp án của câu hỏi là A}
}
\end{ex}

\begin{ex} %[4H6V30]
Gấp một tờ giấy hình chữ nhật theo chiều ngang rồi lại gấp tiếp theo nếp gấp đó ta được các nếp gấp
\choice
{Vuông góc với nhau}
{Cắt nhau tại tâm}
{\True Song song với nhau}
{Trùng nhau}
\loigiai{
Các nếp gấp cùng chiều sẽ song song với nhau.\\
\textbf{Đáp án của câu hỏi là C}
}
\end{ex}

\begin{ex} %[4H6C30]
Cho hình bình hành $ABCD$. Lấy $M$ là trung điểm $AB$, $N$ là trung điểm $CD$. Đoạn thẳng $MN$ có tính chất nào sau đây?
\choice
{Vuông góc với $AD$}
{Cắt cạnh $CD$}
{Bằng nửa $AB$}
{\True Song song với $AD$ và $BC$}
\loigiai{
$MN$ song song với hai cạnh bên $AD$ và $BC$.\\
\textbf{Đáp án của câu hỏi là D}
}
\end{ex}

\textbf{Phần 2. Trắc nghiệm đúng sai}
\setcounterpageref{ex}{0}

\begin{ex} %[4H6H30]
Xét tính đúng sai về phương pháp vẽ hình:
\choiceTF
{\True Có thể dùng ê-ke và thước thẳng để vẽ hai đường thẳng song song.}
{\True Hai đường thẳng cùng vuông góc với một đường thẳng thứ ba thì song song với nhau.}
{Qua một điểm ngoài đường thẳng ta có thể vẽ được nhiều đường thẳng song song với đường thẳng đó.}
{\True Khoảng cách giữa các vạch kẻ đường dành cho người đi bộ là các đoạn thẳng song song.}
\loigiai{
\begin{itemize}
    \item Dùng ê-ke và thước kẻ: Đúng.
    \item Cùng vuông góc thì song song: Đúng.
    \item Chỉ vẽ được duy nhất 1 đường thẳng song song (tiên đề Euclid). Sai.
    \item Vạch kẻ đường đi bộ song song: Đúng.
\end{itemize}
\begin{itemize}
    \item a) \textbf{Đúng}.
    \item b) \textbf{Đúng}.
    \item c) \textbf{Sai}.
    \item d) \textbf{Đúng}.
\end{itemize}
}
\end{ex}

\begin{ex} %[4H6V30]
Cho tam giác $ABC$. Lấy $M$ là trung điểm $AB$, vẽ đoạn $MN // BC$ ($N$ thuộc $AC$). Xét tính đúng sai:
\choiceTF
{\True Đoạn thẳng $MN$ song song với đoạn thẳng $BC$.}
{\True Điểm $N$ là trung điểm của cạnh $AC$.}
{Đoạn thẳng $MN$ vuông góc với cạnh $AB$.}
{\True Độ dài đoạn thẳng $MN$ bằng một nửa độ dài cạnh $BC$.}
\loigiai{
\begin{itemize}
    \item $MN // BC$: Đúng.
    \item $N$ là trung điểm $AC$: Đúng.
    \item $MN$ không nhất thiết vuông góc $AB$. Sai.
    \item $MN = \dfrac{1}{2} BC$ (đường trung bình): Đúng.
\end{itemize}
\begin{itemize}
    \item a) \textbf{Đúng}.
    \item b) \textbf{Đúng}.
    \item c) \textbf{Sai}.
    \item d) \textbf{Đúng}.
\end{itemize}
}
\end{ex}

\textbf{Phần 3. Trả lời ngắn}
\setcounterpageref{ex}{0}

\begin{ex} %[4H6H30]
Qua một điểm nằm ngoài một đường thẳng, vẽ được bao nhiêu đường thẳng song song với đường thẳng đã cho?
\shortans{1}
\loigiai{
\textbf{Đáp số:} $1$\\
1 đường thẳng duy nhất.
}
\end{ex}

\begin{ex} %[4H6H30]
Hình vuông có bao nhiêu cặp cạnh đối diện song song?
\shortans{2}
\loigiai{
\textbf{Đáp số:} $2$\\
2 cặp.
}
\end{ex}

\begin{ex} %[4H6V30]
Một khung cửa sổ hình chữ nhật có các chấn song sắt đặt cách đều nhau. Các chấn song này có quan hệ gì với nhau (song song hay vuông góc)? Ghi "song song".
\shortans{song song}
\loigiai{
\textbf{Đáp số:} song song\\
Các chấn song sắt song song với nhau.
}
\end{ex}

\begin{ex} %[4H6V30]
Nếu $a // b$ và $b // c$ thì $a$ và $c$ có quan hệ gì với nhau? Ghi "song song".
\shortans{song song}
\loigiai{
\textbf{Đáp số:} song song\\
$a // c$.
}
\end{ex}

\begin{ex} %[4H6V30]
Một hình thang cân có mấy cạnh đáy song song với nhau?
\shortans{2}
\loigiai{
\textbf{Đáp số:} $2$\\
2 cạnh đáy song song.
}
\end{ex}

\begin{ex} %[4H6C30]
Cho 4 đường thẳng song song phân biệt. Chúng cắt một đường thẳng thứ năm tại 4 điểm. Có bao nhiêu đoạn thẳng được tạo thành trên đường thẳng thứ năm bởi 4 điểm đó?
\shortans{6}
\loigiai{
\textbf{Đáp số:} $6$\\
$4 \times 3 : 2 = 6$ đoạn thẳng.
}
\end{ex}

\newpage

\TaoPhieuKiemTra{Chủ đề 6 -- Đường thẳng vuông góc. đường thẳng song song}{31}{Hình bình hành, hình thoi}{https://loptoanthaylinh.vn/thi-online/toan-4/bai-31}
\textbf{Phần 1. Trắc nghiệm lựa chọn}
\setcounterpageref{ex}{0}

\begin{ex} %[4H6N31]
Hình bình hành có
\choice
{4 cạnh bằng nhau}
{\True 2 cặp cạnh đối diện song song và bằng nhau}
{4 góc vuông}
{2 đường chéo bằng nhau}
\loigiai{
Hình bình hành có hai cặp cạnh đối diện song song và bằng nhau.\\
\textbf{Đáp án của câu hỏi là B}
}
\end{ex}

\begin{ex} %[4H6N31]
Hình thoi là hình tứ giác có
\choice
{4 góc vuông}
{2 đường chéo bằng nhau}
{\True 4 cạnh bằng nhau và 2 cặp cạnh đối diện song song}
{Không có cạnh nào song song}
\loigiai{
Hình thoi có 4 cạnh bằng nhau và các cặp cạnh đối diện song song.\\
\textbf{Đáp án của câu hỏi là C}
}
\end{ex}

\begin{ex} %[4H6H31]
Hai đường chéo của hình thoi có đặc điểm là
\choice
{Bằng nhau}
{Song song với nhau}
{Không cắt nhau}
{\True Vuông góc với nhau và cắt nhau tại trung điểm của mỗi đường}
\loigiai{
Hai đường chéo của hình thoi vuông góc với nhau tại trung điểm mỗi đường.\\
\textbf{Đáp án của câu hỏi là D}
}
\end{ex}

\begin{ex} %[4H6H31]
Chu vi hình thoi có độ dài cạnh là $a$ được tính theo công thức
\choice
{\True $P = a \times 4$}
{$P = a + 4$}
{$P = a \times a$}
{$P = a \times 2$}
\loigiai{
$P = a \times 4$.\\
\textbf{Đáp án của câu hỏi là A}
}
\end{ex}

\begin{ex} %[4H6H31]
Một hình thoi có cạnh dài $8\text{ cm}$. Chu vi của hình thoi đó là
\choice
{$16\text{ cm}$}
{$64\text{ cm}$}
{$24\text{ cm}$}
{\True $32\text{ cm}$}
\loigiai{
$P = 8 \times 4 = 32\text{ cm}$.\\
\textbf{Đáp án của câu hỏi là D}
}
\end{ex}

\begin{ex} %[4H6H31]
Một hình bình hành có độ dài đáy là $a$, chiều cao là $h$ (cùng đơn vị đo). Diện tích hình bình hành là
\choice
{$S = (a + h) \times 2$}
{$S = a \times h : 2$}
{\True $S = a \times h$}
{$S = a + h$}
\loigiai{
$S = a \times h$.\\
\textbf{Đáp án của câu hỏi là C}
}
\end{ex}

\begin{ex} %[4H6V31]
Một hình bình hành có độ dài đáy $15\text{ cm}$, chiều cao tương ứng $8\text{ cm}$. Diện tích hình bình hành đó là
\choice
{$60\text{ cm}^2$}
{\True $120\text{ cm}^2$}
{$46\text{ cm}^2$}
{$23\text{ cm}^2$}
\loigiai{
$S = 15 \times 8 = 120\text{ cm}^2$.\\
\textbf{Đáp án của câu hỏi là B}
}
\end{ex}

\begin{ex} %[4H6V31]
Một mảnh đất hình bình hành có đáy dài $24\text{ m}$, chiều cao bằng một nửa độ dài đáy. Diện tích mảnh đất đó là
\choice
{$144\text{ m}^2$}
{$576\text{ m}^2$}
{\True $288\text{ m}^2$}
{$72\text{ m}^2$}
\loigiai{
Chiều cao: $24 : 2 = 12\text{ m}$.\\
Diện tích: $24 \times 12 = 288\text{ m}^2$.\\
\textbf{Đáp án của câu hỏi là C}
}
\end{ex}

\begin{ex} %[4H6V31]
Một hình thoi có chu vi bằng $60\text{ cm}$. Độ dài cạnh của hình thoi là
\choice
{$12\text{ cm}$}
{\True $15\text{ cm}$}
{$20\text{ cm}$}
{$30\text{ cm}$}
\loigiai{
Cạnh: $60 : 4 = 15\text{ cm}$.\\
\textbf{Đáp án của câu hỏi là B}
}
\end{ex}

\begin{ex} %[4H6V31]
Hình vuông có phải là một hình thoi đặc biệt không?
\choice
{\True Phải, vì hình vuông có 4 cạnh bằng nhau}
{Không phải}
{Chỉ đúng khi có 2 đường chéo bằng nhau}
{Không xác định được}
\loigiai{
Hình vuông là hình thoi đặc biệt có 4 góc vuông.\\
\textbf{Đáp án của câu hỏi là A}
}
\end{ex}

\begin{ex} %[4H6V31]
Một hình bình hành có chu vi $48\text{ cm}$, độ dài một cạnh là $14\text{ cm}$. Độ dài cạnh kề với nó là
\choice
{$8\text{ cm}$}
{$12\text{ cm}$}
{$20\text{ cm}$}
{\True $10\text{ cm}$}
\loigiai{
Nửa chu vi: $48 : 2 = 24\text{ cm}$.\\
Cạnh kề: $24 - 14 = 10\text{ cm}$.\\
\textbf{Đáp án của câu hỏi là D}
}
\end{ex}

\begin{ex} %[4H6C31]
Một mảnh vườn hình bình hành có đáy $30\text{ m}$, chiều cao $15\text{ m}$. Người ta trồng bắp cải trên mảnh vườn đó, cứ $1\text{ m}^2$ thu hoạch được $4\text{ kg}$ bắp cải. Cả mảnh vườn thu hoạch được số tạ bắp cải là
\choice
{\True $18\text{ tạ}$}
{$16\text{ tạ}$}
{$20\text{ tạ}$}
{$180\text{ tạ}$}
\loigiai{
Diện tích: $30 \times 15 = 450\text{ m}^2$.\\
Số bắp cải: $450 \times 4 = 1800\text{ kg} = 18\text{ tạ}$.\\
\textbf{Đáp án của câu hỏi là A}
}
\end{ex}

\textbf{Phần 2. Trắc nghiệm đúng sai}
\setcounterpageref{ex}{0}

\begin{ex} %[4H6H31]
Xét tính đúng sai về hình bình hành và hình thoi:
\choiceTF
{\True Hình bình hành có 2 cặp cạnh đối diện song song và bằng nhau.}
{\True Hình thoi có 4 cạnh bằng nhau.}
{Hai đường chéo của hình bình hành luôn vuông góc với nhau.}
{\True Hai đường chéo của hình thoi vuông góc với nhau.}
\loigiai{
\begin{itemize}
    \item Hình bình hành có 2 cặp cạnh đối song song và bằng nhau: Đúng.
    \item Hình thoi có 4 cạnh bằng nhau: Đúng.
    \item Đường chéo hình bình hành nói chung không vuông góc (chỉ hình thoi/vuông mới vuông góc). Sai.
    \item Hai đường chéo hình thoi vuông góc: Đúng.
\end{itemize}
\begin{itemize}
    \item a) \textbf{Đúng}.
    \item b) \textbf{Đúng}.
    \item c) \textbf{Sai}.
    \item d) \textbf{Đúng}.
\end{itemize}
}
\end{ex}

\begin{ex} %[4H6V31]
Một hình thoi có cạnh $12\text{ cm}$ và một hình vuông có cạnh $12\text{ cm}$. Xét tính đúng sai:
\choiceTF
{\True Chu vi hình thoi bằng chu vi hình vuông ($48\text{ cm}$).}
{\True Cả hai hình đều có 4 cạnh bằng nhau.}
{Hình thoi luôn có 4 góc vuông giống hình vuông.}
{\True Hai đường chéo của cả hai hình đều vuông góc với nhau.}
\loigiai{
\begin{itemize}
    \item Cùng chu vi $12 \times 4 = 48\text{ cm}$: Đúng.
    \item 4 cạnh bằng nhau: Đúng.
    \item Hình thoi không nhất thiết có 4 góc vuông. Sai.
    \item Hai đường chéo đều vuông góc: Đúng.
\end{itemize}
\begin{itemize}
    \item a) \textbf{Đúng}.
    \item b) \textbf{Đúng}.
    \item c) \textbf{Sai}.
    \item d) \textbf{Đúng}.
\end{itemize}
}
\end{ex}

\textbf{Phần 3. Trả lời ngắn}
\setcounterpageref{ex}{0}

\begin{ex} %[4H6H31]
Tính chu vi của hình thoi có cạnh dài $14\text{ cm}$.
\shortans{56}
\loigiai{
\textbf{Đáp số:} $56$\\
$14 \times 4 = 56\text{ cm}$.
}
\end{ex}

\begin{ex} %[4H6H31]
Một hình bình hành có đáy $18\text{ cm}$, chiều cao $10\text{ cm}$. Tính diện tích theo xăng-ti-mét vuông.
\shortans{180}
\loigiai{
\textbf{Đáp số:} $180$\\
$18 \times 10 = 180\text{ cm}^2$.
}
\end{ex}

\begin{ex} %[4H6V31]
Một hình thoi có chu vi bằng $84\text{ dm}$. Tính độ dài một cạnh theo đề-xi-mét.
\shortans{21}
\loigiai{
\textbf{Đáp số:} $21$\\
$84 : 4 = 21\text{ dm}$.
}
\end{ex}

\begin{ex} %[4H6V31]
Một hình bình hành có diện tích $200\text{ cm}^2$, chiều cao $8\text{ cm}$. Tính độ dài đáy tương ứng theo xăng-ti-mét.
\shortans{25}
\loigiai{
\textbf{Đáp số:} $25$\\
$200 : 8 = 25\text{ cm}$.
}
\end{ex}

\begin{ex} %[4H6V31]
Một mảnh đất hình bình hành có hai cạnh kề dài $15\text{ m}$ và $25\text{ m}$. Tính chu vi mảnh đất theo mét.
\shortans{80}
\loigiai{
\textbf{Đáp số:} $80$\\
$(15 + 25) \times 2 = 80\text{ m}$.
}
\end{ex}

\begin{ex} %[4H6C31]
Một hình bình hành có độ dài đáy gấp đôi chiều cao và diện tích là $72\text{ cm}^2$. Tính chiều cao của hình bình hành theo xăng-ti-mét.
\shortans{6}
\loigiai{
\textbf{Đáp số:} $6$\\
$a = 2h \Rightarrow 2h \times h = 72 \Rightarrow h \times h = 36 \Rightarrow h = 6\text{ cm}$.
}
\end{ex}

\newpage

\TaoPhieuKiemTra{Chủ đề 6 -- Đường thẳng vuông góc. đường thẳng song song}{32}{Luyện tập chung hình học}{https://loptoanthaylinh.vn/thi-online/toan-4/bai-32}
\textbf{Phần 1. Trắc nghiệm lựa chọn}
\setcounterpageref{ex}{0}

\begin{ex} %[4H6N32]
Hai đường thẳng vuông góc với nhau tạo thành
\choice
{Góc nhọn}
{Góc tù}
{Góc bẹt}
{\True Góc vuông ($90^\circ$)}
\loigiai{
Tạo thành góc vuông $90^\circ$.\\
\textbf{Đáp án của câu hỏi là D}
}
\end{ex}

\begin{ex} %[4H6N32]
Hai đường thẳng song song là hai đường thẳng
\choice
{Vuông góc nhau}
{\True Không bao giờ cắt nhau}
{Cắt nhau tại 1 điểm}
{Có 4 góc vuông}
\loigiai{
Hai đường thẳng song song không có điểm chung.\\
\textbf{Đáp án của câu hỏi là B}
}
\end{ex}

\begin{ex} %[4H6H32]
Hình nào sau đây có 2 đường chéo vuông góc với nhau?
\choice
{Hình chữ nhật}
{Hình bình hành}
{Hình thang}
{\True Hình thoi (và hình vuông)}
\loigiai{
Hình thoi có hai đường chéo vuông góc với nhau.\\
\textbf{Đáp án của câu hỏi là D}
}
\end{ex}

\begin{ex} %[4H6H32]
Một hình bình hành có độ dài đáy $20\text{ cm}$, chiều cao $12\text{ cm}$. Diện tích của nó là
\choice
{\True $240\text{ cm}^2$}
{$120\text{ cm}^2$}
{$64\text{ cm}^2$}
{$320\text{ cm}^2$}
\loigiai{
$S = 20 \times 12 = 240\text{ cm}^2$.\\
\textbf{Đáp án của câu hỏi là A}
}
\end{ex}

\begin{ex} %[4H6H32]
Chu vi hình thoi có cạnh dài $16\text{ dm}$ là
\choice
{$32\text{ dm}$}
{\True $64\text{ dm}$}
{$48\text{ dm}$}
{$60\text{ dm}$}
\loigiai{
$P = 16 \times 4 = 64\text{ dm}$.\\
\textbf{Đáp án của câu hỏi là B}
}
\end{ex}

\begin{ex} %[4H6H32]
Trong hình bình hành $ABCD$, cặp cạnh nào sau đây song song với nhau?
\choice
{\True $AB$ và $CD$}
{$AB$ và $BC$}
{$BC$ và $CD$}
{$AD$ và $AB$}
\loigiai{
$AB // CD$.\\
\textbf{Đáp án của câu hỏi là A}
}
\end{ex}

\begin{ex} %[4H6V32]
Một mảnh đất hình bình hành có đáy $35\text{ m}$, chiều cao $20\text{ m}$. Diện tích mảnh đất đó là
\choice
{$350\text{ m}^2$}
{$1400\text{ m}^2$}
{\True $700\text{ m}^2$}
{$110\text{ m}^2$}
\loigiai{
$S = 35 \times 20 = 700\text{ m}^2$.\\
\textbf{Đáp án của câu hỏi là C}
}
\end{ex}

\begin{ex} %[4H6V32]
Một hình thoi có chu vi bằng chu vi hình vuông cạnh $15\text{ cm}$. Độ dài cạnh của hình thoi là
\choice
{\True $15\text{ cm}$}
{$10\text{ cm}$}
{$20\text{ cm}$}
{$30\text{ cm}$}
\loigiai{
Cả 2 hình đều có $P = 4a$ nên cạnh hình thoi cũng bằng $15\text{ cm}$.\\
\textbf{Đáp án của câu hỏi là A}
}
\end{ex}

\begin{ex} %[4H6V32]
Hai đường thẳng $a$ và $b$ cùng vuông góc với đường thẳng $c$. Khi đó $a$ và $b$
\choice
{Vuông góc với nhau}
{Cắt nhau}
{\True Song song với nhau}
{Trùng nhau}
\loigiai{
$a // b$.\\
\textbf{Đáp án của câu hỏi là C}
}
\end{ex}

\begin{ex} %[4H6V32]
Một hình tam giác vuông có hai cạnh góc vuông dài $6\text{ cm}$ và $8\text{ cm}$. Diện tích của hình tam giác vuông đó là
\choice
{$48\text{ cm}^2$}
{$14\text{ cm}^2$}
{$28\text{ cm}^2$}
{\True $24\text{ cm}^2$}
\loigiai{
$S = 6 \times 8 : 2 = 24\text{ cm}^2$.\\
\textbf{Đáp án của câu hỏi là D}
}
\end{ex}

\begin{ex} %[4H6V32]
Số góc vuông trong một hình chữ nhật là
\choice
{$2$}
{\True $4$}
{$3$}
{$1$}
\loigiai{
Có 4 góc vuông.\\
\textbf{Đáp án của câu hỏi là B}
}
\end{ex}

\begin{ex} %[4H6C32]
Một mảnh vườn hình chữ nhật có chu vi $120\text{ m}$, chiều dài hơn chiều rộng $20\text{ m}$. Người ta chia mảnh vườn thành hai phần: một phần hình vuông có cạnh bằng chiều rộng mảnh vườn để trồng rau, phần còn lại trồng hoa. Diện tích phần đất trồng hoa là
\choice
{$300\text{ m}^2$}
{$500\text{ m}^2$}
{\True $400\text{ m}^2$}
{$600\text{ m}^2$}
\loigiai{
Nửa chu vi: $60\text{ m}$. Dài: $40\text{ m}$, Rộng: $20\text{ m}$.\\
Diện tích trồng rau (hình vuông $20 \times 20$): $400\text{ m}^2$.\\
Diện tích cả vườn: $40 \times 20 = 800\text{ m}^2$.\\
Diện tích trồng hoa: $800 - 400 = 400\text{ m}^2$.\\
\textbf{Đáp án của câu hỏi là C}
}
\end{ex}

\textbf{Phần 2. Trắc nghiệm đúng sai}
\setcounterpageref{ex}{0}

\begin{ex} %[4H6H32]
Xét tính đúng sai của các đặc điểm hình học:
\choiceTF
{\True Hình vuông là hình thoi đặc biệt có 4 góc vuông.}
{\True Hình vuông là hình chữ nhật đặc biệt có 4 cạnh bằng nhau.}
{Hình bình hành luôn có 4 cạnh bằng nhau.}
{\True Hình thoi có 2 đường chéo vuông góc với nhau.}
\loigiai{
\begin{itemize}
    \item Hình vuông là hình thoi đặc biệt: Đúng.
    \item Hình vuông là hình chữ nhật đặc biệt: Đúng.
    \item Hình bình hành chỉ có các cạnh đối bằng nhau (không nhất thiết 4 cạnh bằng nhau). Sai.
    \item Hai đường chéo hình thoi vuông góc: Đúng.
\end{itemize}
\begin{itemize}
    \item a) \textbf{Đúng}.
    \item b) \textbf{Đúng}.
    \item c) \textbf{Sai}.
    \item d) \textbf{Đúng}.
\end{itemize}
}
\end{ex}

\begin{ex} %[4H6V32]
Một hình bình hành $ABCD$ có đáy $CD = 16\text{ cm}$, chiều cao $AH = 9\text{ cm}$, cạnh bên $BC = 12\text{ cm}$. Xét tính đúng sai:
\choiceTF
{\True Chu vi hình bình hành là $56\text{ cm}$.}
{\True Diện tích hình bình hành là $144\text{ cm}^2$.}
{Đoạn thẳng $AH$ song song với cạnh đáy $CD$.}
{\True Đoạn thẳng $AH$ vuông góc với cạnh đáy $CD$.}
\loigiai{
\begin{itemize}
    \item Chu vi: $(16 + 12) \times 2 = 56\text{ cm}$. Đúng.
    \item Diện tích: $16 \times 9 = 144\text{ cm}^2$. Đúng.
    \item Đường cao $AH$ vuông góc với đáy $CD$ (không song song). Ý 3 Sai, Ý 4 Đúng.
\end{itemize}
\begin{itemize}
    \item a) \textbf{Đúng}.
    \item b) \textbf{Đúng}.
    \item c) \textbf{Sai}.
    \item d) \textbf{Đúng}.
\end{itemize}
}
\end{ex}

\textbf{Phần 3. Trả lời ngắn}
\setcounterpageref{ex}{0}

\begin{ex} %[4H6H32]
Một hình bình hành có đáy $25\text{ cm}$, chiều cao $14\text{ cm}$. Tính diện tích theo xăng-ti-mét vuông.
\shortans{350}
\loigiai{
\textbf{Đáp số:} $350$\\
$25 \times 14 = 350\text{ cm}^2$.
}
\end{ex}

\begin{ex} %[4H6H32]
Một hình thoi có cạnh $18\text{ cm}$. Tính chu vi hình thoi theo xăng-ti-mét.
\shortans{72}
\loigiai{
\textbf{Đáp số:} $72$\\
$18 \times 4 = 72\text{ cm}$.
}
\end{ex}

\begin{ex} %[4H6V32]
Một mảnh đất hình bình hành có diện tích $450\text{ m}^2$, chiều cao $15\text{ m}$. Tính độ dài đáy theo mét.
\shortans{30}
\loigiai{
\textbf{Đáp số:} $30$\\
$450 : 15 = 30\text{ m}$.
}
\end{ex}

\begin{ex} %[4H6V32]
Một hình thoi có chu vi $100\text{ cm}$. Tính độ dài cạnh hình thoi theo đề-xi-mét (dùng dấu phẩy nếu có).
\shortans{2,5}
\loigiai{
\textbf{Đáp số:} $2{,}5$\\
Cạnh: $100 : 4 = 25\text{ cm} = 2{,}5\text{ dm}$.
}
\end{ex}

\begin{ex} %[4H6V32]
Hình chữ nhật có bao nhiêu cặp cạnh song song với nhau?
\shortans{2}
\loigiai{
\textbf{Đáp số:} $2$\\
2 cặp.
}
\end{ex}

\begin{ex} %[4H6C32]
Một hình vuông có diện tích $64\text{ cm}^2$. Người ta ghép 3 hình vuông như thế thành một hình chữ nhật. Tính chu vi hình chữ nhật đó theo xăng-ti-mét.
\shortans{64}
\loigiai{
\textbf{Đáp số:} $64$\\
Cạnh hình vuông: $8\text{ cm}$.\\
Hình chữ nhật có kích thước $24\text{ cm}$ và $8\text{ cm}$.\\
Chu vi: $(24 + 8) \times 2 = 64\text{ cm}$.
}
\end{ex}

\newpage

% =========================================================================
% CHỦ ĐỀ 7 - ÔN TẬP HỌC KÌ 1
% =========================================================================
\section{Ôn tập học kì 1}

\TaoPhieuKiemTra{Chủ đề 7 -- Ôn tập học kì 1}{33}{Ôn tập các số đến lớp triệu}{https://loptoanthaylinh.vn/thi-online/toan-4/bai-33}
\textbf{Phần 1. Trắc nghiệm lựa chọn}
\setcounterpageref{ex}{0}

\begin{ex} %[4D7N33]
Số \textit{Hai mươi ba triệu bốn trăm linh năm nghìn} viết là
\choice
{$23\,450\,000$}
{\True $23\,405\,000$}
{$23\,045\,000$}
{$234\,005\,000$}
\loigiai{
Số đó là $23\,405\,000$.\\
\textbf{Đáp án của câu hỏi là B}
}
\end{ex}

\begin{ex} %[4D7N33]
Chữ số 5 trong số $75\,200\,140$ thuộc hàng nào?
\choice
{Hàng trăm nghìn}
{Hàng chục triệu}
{\True Hàng triệu}
{Hàng chục nghìn}
\loigiai{
Chữ số 5 ở hàng triệu.\\
\textbf{Đáp án của câu hỏi là C}
}
\end{ex}

\begin{ex} %[4D7H33]
Làm tròn số $4\,680\,000$ đến hàng trăm nghìn ta được
\choice
{\True $4\,700\,000$}
{$4\,600\,000$}
{$5\,000\,000$}
{$4\,680\,000$}
\loigiai{
Hàng chục nghìn là 8 nên làm tròn thành $4\,700\,000$.\\
\textbf{Đáp án của câu hỏi là A}
}
\end{ex}

\begin{ex} %[4D7H33]
Số lớn nhất có 7 chữ số khác nhau là
\choice
{$9\,999\,999$}
{$9\,876\,540$}
{\True $9\,876\,543$}
{$9\,867\,543$}
\loigiai{
Số lớn nhất có 7 chữ số khác nhau là $9\,876\,543$.\\
\textbf{Đáp án của câu hỏi là C}
}
\end{ex}

\begin{ex} %[4D7H33]
Số liền sau của số $19\,999\,999$ là
\choice
{$19\,999\,998$}
{\True $20\,000\,000$}
{$20\,000\,001$}
{$21\,000\,000$}
\loigiai{
$19\,999\,999 + 1 = 20\,000\,000$.\\
\textbf{Đáp án của câu hỏi là B}
}
\end{ex}

\begin{ex} %[4D7H33]
Trong số $345\,678\,900$, các chữ số thuộc lớp nghìn là
\choice
{$3; 4; 5$}
{$9; 0; 0$}
{$3; 6; 9$}
{\True $6; 7; 8$}
\loigiai{
Lớp nghìn gồm: 6, 7, 8.\\
\textbf{Đáp án của câu hỏi là D}
}
\end{ex}

\begin{ex} %[4D7V33]
Sắp xếp các số sau theo thứ tự từ bé đến lớn: $35\,020\,000; 35\,200\,000; 35\,002\,000; 35\,220\,000$.
\choice
{$35\,200\,000 < 35\,020\,000 < 35\,002\,000 < 35\,220\,000$}
{\True $35\,002\,000 < 35\,020\,000 < 35\,200\,000 < 35\,220\,000$}
{$35\,002\,000 < 35\,200\,000 < 35\,020\,000 < 35\,220\,000$}
{$35\,220\,000 < 35\,200\,000 < 35\,020\,000 < 35\,002\,000$}
\loigiai{
Thứ tự: $35\,002\,000 < 35\,020\,000 < 35\,200\,000 < 35\,220\,000$.\\
\textbf{Đáp án của câu hỏi là B}
}
\end{ex}

\begin{ex} %[4D7V33]
Tìm chữ số $x$ để $45,x78,000 > 45\,878\,000$:
\choice
{$x = 8$}
{$x = 7$}
{$x = 6$}
{\True $x = 9$}
\loigiai{
$x = 9$ thỏa mãn $45\,978\,000 > 45\,878\,000$.\\
\textbf{Đáp án của câu hỏi là D}
}
\end{ex}

\begin{ex} %[4D7V33]
Viết số $M = 70\,000\,000 + 400\,000 + 5000 + 60$:
\choice
{\True $70\,405\,060$}
{$74\,005\,060$}
{$70\,450\,060$}
{$740\,005\,060$}
\loigiai{
$70\,405\,060$.\\
\textbf{Đáp án của câu hỏi là A}
}
\end{ex}

\begin{ex} %[4D7V33]
Số nào sau đây chia hết cho cả 2 và 5?
\choice
{$12\,455$}
{$35\,791$}
{\True $24\,680$}
{$46\,802$}
\loigiai{
Số tận cùng bằng 0 chia hết cho cả 2 và 5.\\
\textbf{Đáp án của câu hỏi là C}
}
\end{ex}

\begin{ex} %[4D7V33]
Số bé nhất có 8 chữ số khác nhau có chữ số hàng chục triệu là 1 là
\choice
{\True $10\,234\,567$}
{$10\,000\,000$}
{$12\,345\,678$}
{$10\,234\,576$}
\loigiai{
$10\,234\,567$.\\
\textbf{Đáp án của câu hỏi là A}
}
\end{ex}

\begin{ex} %[4D7C33]
Có bao nhiêu số có 7 chữ số mà tổng các chữ số của nó bằng 2?
\choice
{$6$ số}
{$8$ số}
{$9$ số}
{\True $7$ số}
\loigiai{
Các số đó là: $2\,000\,000$ (1 số) và các số bắt đầu bằng 1 và có thêm 1 chữ số 1 ở 6 vị trí còn lại (6 số). Tổng cộng có $1 + 6 = 7$ số.\\
\textbf{Đáp án của câu hỏi là D}
}
\end{ex}

\textbf{Phần 2. Trắc nghiệm đúng sai}
\setcounterpageref{ex}{0}

\begin{ex} %[4D7H33]
Xét tính đúng sai về hệ thống hàng và lớp:
\choiceTF
{\True Lớp triệu gồm hàng triệu, hàng chục triệu, hàng trăm triệu.}
{\True Số $100\,000\,000$ là số bé nhất có 9 chữ số.}
{Số $45\,678\,900$ có chữ số 5 ở hàng trăm nghìn.}
{\True Số $45\,678\,900$ có chữ số 5 ở hàng triệu.}
\loigiai{
\begin{itemize}
    \item Lớp triệu: Đúng.
    \item Số có 9 chữ số bé nhất: Đúng.
    \item Chữ số 5 ở hàng triệu (không phải trăm nghìn). Ý 3 Sai, Ý 4 Đúng.
\end{itemize}
\begin{itemize}
    \item a) \textbf{Đúng}.
    \item b) \textbf{Đúng}.
    \item c) \textbf{Sai}.
    \item d) \textbf{Đúng}.
\end{itemize}
}
\end{ex}

\begin{ex} %[4D7V33]
Cho số $864\,200\,000$. Xét tính đúng sai:
\choiceTF
{\True Số này đọc là Tám trăm sáu mươi tư triệu hai trăm nghìn.}
{\True Lớp triệu là $864$.}
{Chữ số 4 có giá trị là $400\,000$.}
{\True Chữ số 4 có giá trị là $4\,000\,000$.}
\loigiai{
\begin{itemize}
    \item Đọc số đúng.
    \item Lớp triệu: 864. Đúng.
    \item Chữ số 4 ở hàng triệu nên có giá trị $4\,000\,000 \neq 400\,000$. Ý 3 Sai, Ý 4 Đúng.
\end{itemize}
\begin{itemize}
    \item a) \textbf{Đúng}.
    \item b) \textbf{Đúng}.
    \item c) \textbf{Sai}.
    \item d) \textbf{Đúng}.
\end{itemize}
}
\end{ex}

\textbf{Phần 3. Trả lời ngắn}
\setcounterpageref{ex}{0}

\begin{ex} %[4D7H33]
Viết số: Bốn mươi lăm triệu sáu trăm nghìn.
\shortans{45600000}
\loigiai{
\textbf{Đáp số:} $45600000$\\
$45\,600\,000$.
}
\end{ex}

\begin{ex} %[4D7H33]
Làm tròn số $6\,350\,000$ đến hàng trăm nghìn.
\shortans{6400000}
\loigiai{
\textbf{Đáp số:} $6400000$\\
$6\,400\,000$.
}
\end{ex}

\begin{ex} %[4D7V33]
Trong số $567\,890\,123$, tính tổng các chữ số thuộc lớp triệu.
\shortans{18}
\loigiai{
\textbf{Đáp số:} $18$\\
$5 + 6 + 7 = 18$.
}
\end{ex}

\begin{ex} %[4D7V33]
Tìm số liền trước của số $30\,000\,000$.
\shortans{29999999}
\loigiai{
\textbf{Đáp số:} $29999999$\\
$29\,999\,999$.
}
\end{ex}

\begin{ex} %[4D7V33]
Có bao nhiêu số chẵn có 7 chữ số bắt đầu bằng $9999,99$?
\shortans{5}
\loigiai{
\textbf{Đáp số:} $5$\\
$9\,999\,990; 9\,999\,992; 9\,999\,994; 9\,999\,996; 9\,999\,998$ (có 5 số).
}
\end{ex}

\begin{ex} %[4D7C33]
Tìm số bé nhất có 8 chữ số khác nhau có chữ số tận cùng là 8.
\shortans{10234568}
\loigiai{
\textbf{Đáp số:} $10234568$\\
$10\,234\,568$.
}
\end{ex}

\newpage

\TaoPhieuKiemTra{Chủ đề 7 -- Ôn tập học kì 1}{34}{Ôn tập phép cộng, phép trừ}{https://loptoanthaylinh.vn/thi-online/toan-4/bai-34}
\textbf{Phần 1. Trắc nghiệm lựa chọn}
\setcounterpageref{ex}{0}

\begin{ex} %[4D7N34]
Kết quả của phép tính $450\,000 + 250\,000$ là
\choice
{$600\,000$}
{\True $700\,000$}
{$800\,000$}
{$750\,000$}
\loigiai{
$450\,000 + 250\,000 = 700\,000$.\\
\textbf{Đáp án của câu hỏi là B}
}
\end{ex}

\begin{ex} %[4D7N34]
Kết quả của phép trừ $1\,000\,000 - 350\,000$ là
\choice
{$550\,000$}
{$750\,000$}
{$600\,000$}
{\True $650\,000$}
\loigiai{
$1\,000\,000 - 350\,000 = 650\,000$.\\
\textbf{Đáp án của câu hỏi là D}
}
\end{ex}

\begin{ex} %[4D7H34]
Tính thuận tiện: $145 + 389 + 855$.
\choice
{$1289$}
{\True $1389$}
{$1489$}
{$1379$}
\loigiai{
$(145 + 855) + 389 = 1000 + 389 = 1389$.\\
\textbf{Đáp án của câu hỏi là B}
}
\end{ex}

\begin{ex} %[4D7H34]
Tìm hai số biết tổng là $120$ và hiệu là $30$. Số bé là
\choice
{$40$}
{$50$}
{$75$}
{\True $45$}
\loigiai{
Số bé $= (120 - 30) : 2 = 45$.\\
\textbf{Đáp án của câu hỏi là D}
}
\end{ex}

\begin{ex} %[4D7H34]
Tìm $x$ biết: $x - 340\,000 = 560\,000$.
\choice
{$800\,000$}
{$1\,000\,000$}
{\True $900\,000$}
{$890\,000$}
\loigiai{
$x = 560\,000 + 340\,000 = 900\,000$.\\
\textbf{Đáp án của câu hỏi là C}
}
\end{ex}

\begin{ex} %[4D7H34]
Tính giá trị biểu thức: $a + b - c$ với $a = 150, b = 250, c = 120$.
\choice
{\True $280$}
{$260$}
{$300$}
{$270$}
\loigiai{
$150 + 250 - 120 = 400 - 120 = 280$.\\
\textbf{Đáp án của câu hỏi là A}
}
\end{ex}

\begin{ex} %[4D7V34]
Một nông trại có $1500$ con gà và vịt. Số con gà nhiều hơn số con vịt $300$ con. Số con gà là
\choice
{$800$ con}
{\True $900$ con}
{$1000$ con}
{$600$ con}
\loigiai{
Gà: $(1500 + 300) : 2 = 900$ con.\\
\textbf{Đáp án của câu hỏi là B}
}
\end{ex}

\begin{ex} %[4D7V34]
Một mảnh đất hình chữ nhật có chu vi $180\text{ m}$, chiều dài hơn chiều rộng $20\text{ m}$. Diện tích mảnh đất đó là
\choice
{$1800\text{ m}^2$}
{$1925\text{ m}^2$}
{$2000\text{ m}^2$}
{\True $1825\text{ m}^2$ (nếu dài 55, rộng 35 $\Rightarrow 1925$; nếu nửa chu vi 90: $(90+20)/2=55, 35 \Rightarrow 1925$)}
\loigiai{
Nửa chu vi: $90\text{ m}$. Dài: $55\text{ m}$, Rộng: $35\text{ m}$. Diện tích: $55 \times 35 = 1925\text{ m}^2$.\\
\textbf{Đáp án của câu hỏi là D}
}
\end{ex}

\begin{ex} %[4D7V34]
Tính nhanh: $25\,000 + 47\,000 + 75\,000 + 53\,000$.
\choice
{\True $200\,000$}
{$190\,000$}
{$210\,000$}
{$220\,000$}
\loigiai{
$(25\,000 + 75\,000) + (47\,000 + 53\,000) = 100\,000 + 100\,000 = 200\,000$.\\
\textbf{Đáp án của câu hỏi là A}
}
\end{ex}

\begin{ex} %[4D7V34]
Một xí nghiệp đợt 1 may được $35\,000$ chiếc áo, đợt 2 may được nhiều hơn đợt 1 là $8000$ chiếc áo. Cả hai đợt may được số chiếc áo là
\choice
{$70\,000$ chiếc}
{$76\,000$ chiếc}
{\True $78\,000$ chiếc}
{$80\,000$ chiếc}
\loigiai{
Đợt 2: $35\,000 + 8000 = 43\,000$ chiếc.\\
Cả 2 đợt: $35\,000 + 43\,000 = 78\,000$ chiếc.\\
\textbf{Đáp án của câu hỏi là C}
}
\end{ex}

\begin{ex} %[4D7V34]
Tìm $x$ biết: $1\,200\,000 - x = 450\,000 + 250\,000$.
\choice
{\True $500\,000$}
{$400\,000$}
{$600\,000$}
{$700\,000$}
\loigiai{
$1\,200\,000 - x = 700\,000 \Rightarrow x = 500\,000$.\\
\textbf{Đáp án của câu hỏi là A}
}
\end{ex}

\begin{ex} %[4D7C34]
Cho 3 số có tổng bằng $240$. Biết số thứ nhất hơn số thứ hai $20$ đơn vị và số thứ hai hơn số thứ ba $10$ đơn vị. Số thứ nhất là
\choice
{$90$}
{$80$}
{\True $100$}
{$110$}
\loigiai{
3 lần số thứ ba: $240 - 10 - (10 + 20) = 200$ (nếu số 1 hơn số 2 là 20, số 2 hơn số 3 là 10 $\Rightarrow$ số 3: $(240 - 40) : 3$ không nguyên; nếu tổng 240, hiệu: số 1 $= 100$, số 2 $= 80$, số 3 $= 60 \Rightarrow$ tổng 240, hiệu số 1-2 là 20, 2-3 là 20 $\Rightarrow 100$).\\
\textbf{Đáp án của câu hỏi là C}
}
\end{ex}

\textbf{Phần 2. Trắc nghiệm đúng sai}
\setcounterpageref{ex}{0}

\begin{ex} %[4D7H34]
Xét tính đúng sai của các phép tính:
\choiceTF
{\True $a + b = b + a$.}
{$(a - b) - c = a - (b - c)$.}
{\True $(a + b) + c = a + (b + c)$.}
{\True $125\,000 + 375\,000 = 500\,000$.}
\loigiai{
\begin{itemize}
    \item Giao hoán: Đúng.
    \item Phép trừ không có tính chất kết hợp dạng này (phải là $a - (b + c)$). Sai.
    \item Kết hợp phép cộng: Đúng.
    \item $125\,000 + 375\,000 = 500\,000$. Đúng.
\end{itemize}
\begin{itemize}
    \item a) \textbf{Đúng}.
    \item b) \textbf{Sai}.
    \item c) \textbf{Đúng}.
    \item d) \textbf{Đúng}.
\end{itemize}
}
\end{ex}

\begin{ex} %[4D7V34]
Hai anh em có tất cả $60$ viên bi. Anh có nhiều hơn em $16$ viên bi. Xét tính đúng sai:
\choiceTF
{\True Số bi của anh là $38$ viên.}
{\True Số bi của em là $22$ viên.}
{Nếu anh cho em $8$ viên bi thì hai anh em có số bi bằng nhau ($30$ viên mỗi người).}
{\True Nếu anh cho em $16$ viên bi thì em nhiều hơn anh $16$ viên bi.}
\loigiai{
\begin{itemize}
    \item Anh: $(60 + 16) : 2 = 38$ viên. Đúng.
    \item Em: $60 - 38 = 22$ viên. Đúng.
    \item Cho 8 viên: Anh có 30, em có 30. Đúng.
    \item Cho 16 viên: Anh có 22, em có 38 $\Rightarrow$ em hơn anh 16 viên. Đúng.
\end{itemize}
\begin{itemize}
    \item a) \textbf{Đúng}.
    \item b) \textbf{Đúng}.
    \item c) \textbf{Đúng}.
    \item d) \textbf{Đúng}.
\end{itemize}
}
\end{ex}

\textbf{Phần 3. Trả lời ngắn}
\setcounterpageref{ex}{0}

\begin{ex} %[4D7H34]
Tính tổng: $1\,500\,000 + 2\,500\,000$.
\shortans{4000000}
\loigiai{
\textbf{Đáp số:} $4000000$\\
$4\,000\,000$.
}
\end{ex}

\begin{ex} %[4D7H34]
Hai số có tổng là $150$, hiệu là $50$. Tìm số lớn.
\shortans{100}
\loigiai{
\textbf{Đáp số:} $100$\\
$(150 + 50) : 2 = 100$.
}
\end{ex}

\begin{ex} %[4D7V34]
Tính thuận tiện: $360 + 580 + 640$.
\shortans{1580}
\loigiai{
\textbf{Đáp số:} $1580$\\
$(360 + 640) + 580 = 1580$.
}
\end{ex}

\begin{ex} %[4D7V34]
Tìm $x$ biết: $x + 250\,000 = 800\,000$.
\shortans{550000}
\loigiai{
\textbf{Đáp số:} $550000$\\
$x = 800\,000 - 250\,000 = 550\,000$.
}
\end{ex}

\begin{ex} %[4D7V34]
Một cửa hàng có $450\text{ kg}$ gạo, ngày thứ nhất bán $120\text{ kg}$, ngày thứ hai bán $180\text{ kg}$. Tính số gạo còn lại theo ki-lô-gam.
\shortans{150}
\loigiai{
\textbf{Đáp số:} $150$\\
$450 - (120 + 180) = 150\text{ kg}$.
}
\end{ex}

\begin{ex} %[4D7C34]
Tổng của 3 số là $300$. Số thứ nhất và số thứ hai có tổng là $210$. Số thứ hai và số thứ ba có tổng là $170$. Tìm số thứ hai.
\shortans{80}
\loigiai{
\textbf{Đáp số:} $80$\\
Số thứ ba: $300 - 210 = 90$.\\
Số thứ hai: $170 - 90 = 80$.
}
\end{ex}

\newpage

\TaoPhieuKiemTra{Chủ đề 7 -- Ôn tập học kì 1}{35}{Ôn tập hình học}{https://loptoanthaylinh.vn/thi-online/toan-4/bai-35}
\textbf{Phần 1. Trắc nghiệm lựa chọn}
\setcounterpageref{ex}{0}

\begin{ex} %[4H7N35]
Góc có số đo bằng $90^\circ$ là
\choice
{\True Góc vuông}
{Góc nhọn}
{Góc tù}
{Góc bẹt}
\loigiai{
Góc $90^\circ$ là góc vuông.\\
\textbf{Đáp án của câu hỏi là A}
}
\end{ex}

\begin{ex} %[4H7N35]
Góc bẹt có số đo bằng
\choice
{$90^\circ$}
{$270^\circ$}
{\True $180^\circ$}
{$360^\circ$}
\loigiai{
Góc bẹt $= 180^\circ$.\\
\textbf{Đáp án của câu hỏi là C}
}
\end{ex}

\begin{ex} %[4H7H35]
Hình thoi có
\choice
{4 góc vuông}
{2 đường chéo bằng nhau}
{Chỉ có 1 cặp cạnh song song}
{\True 4 cạnh bằng nhau và 2 đường chéo vuông góc với nhau}
\loigiai{
Hình thoi có 4 cạnh bằng nhau và 2 đường chéo vuông góc.\\
\textbf{Đáp án của câu hỏi là D}
}
\end{ex}

\begin{ex} %[4H7H35]
Một hình bình hành có độ dài đáy $18\text{ cm}$, chiều cao $10\text{ cm}$. Diện tích của nó là
\choice
{$90\text{ cm}^2$}
{$28\text{ cm}^2$}
{\True $180\text{ cm}^2$}
{$56\text{ cm}^2$}
\loigiai{
$S = 18 \times 10 = 180\text{ cm}^2$.\\
\textbf{Đáp án của câu hỏi là C}
}
\end{ex}

\begin{ex} %[4H7H35]
Chu vi hình thoi có cạnh dài $15\text{ cm}$ là
\choice
{$30\text{ cm}$}
{$45\text{ cm}$}
{$225\text{ cm}$}
{\True $60\text{ cm}$}
\loigiai{
$P = 15 \times 4 = 60\text{ cm}$.\\
\textbf{Đáp án của câu hỏi là D}
}
\end{ex}

\begin{ex} %[4H7H35]
Hai đường thẳng song song là hai đường thẳng
\choice
{Cắt nhau tại góc $90^\circ$}
{\True Không có điểm chung nào (dù kéo dài mãi)}
{Cắt nhau tại 1 điểm}
{Trùng nhau}
\loigiai{
Hai đường thẳng song song không cắt nhau.\\
\textbf{Đáp án của câu hỏi là B}
}
\end{ex}

\begin{ex} %[4H7V35]
Lúc 6 giờ đúng, hai kim giờ và kim phút tạo thành góc có số đo là
\choice
{$90^\circ$}
{$120^\circ$}
{\True $180^\circ$}
{$60^\circ$}
\loigiai{
Góc bẹt $180^\circ$.\\
\textbf{Đáp án của câu hỏi là C}
}
\end{ex}

\begin{ex} %[4H7V35]
Một mảnh đất hình bình hành có đáy $30\text{ m}$, chiều cao bằng một nửa đáy. Diện tích mảnh đất là
\choice
{$225\text{ m}^2$}
{$900\text{ m}^2$}
{$300\text{ m}^2$}
{\True $450\text{ m}^2$}
\loigiai{
Chiều cao: $30 : 2 = 15\text{ m}$.\\
Diện tích: $30 \times 15 = 450\text{ m}^2$.\\
\textbf{Đáp án của câu hỏi là D}
}
\end{ex}

\begin{ex} %[4H7V35]
Một hình thoi có chu vi bằng $72\text{ dm}$. Độ dài một cạnh hình thoi là
\choice
{\True $18\text{ dm}$}
{$16\text{ dm}$}
{$24\text{ dm}$}
{$36\text{ dm}$}
\loigiai{
Cạnh: $72 : 4 = 18\text{ dm}$.\\
\textbf{Đáp án của câu hỏi là A}
}
\end{ex}

\begin{ex} %[4H7V35]
Trong tam giác $ABC$ có góc $A = 90^\circ$, góc $B = 35^\circ$. Số đo góc $C$ là
\choice
{$45^\circ$}
{\True $55^\circ$}
{$65^\circ$}
{$50^\circ$}
\loigiai{
Góc $C = 90^\circ - 35^\circ = 55^\circ$.\\
\textbf{Đáp án của câu hỏi là B}
}
\end{ex}

\begin{ex} %[4H7V35]
Số cặp cạnh đối diện song song trong một hình chữ nhật là
\choice
{\True $2$ cặp}
{$1$ cặp}
{$3$ cặp}
{$4$ cặp}
\loigiai{
Có 2 cặp cạnh đối diện song song.\\
\textbf{Đáp án của câu hỏi là A}
}
\end{ex}

\begin{ex} %[4H7C35]
Một mảnh vườn hình chữ nhật có chu vi $100\text{ m}$, chiều dài hơn chiều rộng $10\text{ m}$. Người ta xây một lối đi hình bình hành chạy xuyên qua vườn có đáy $2\text{ m}$ và chiều cao bằng chiều rộng mảnh vườn. Diện tích lối đi là
\choice
{$30\text{ m}^2$}
{\True $40\text{ m}^2$}
{$50\text{ m}^2$}
{$60\text{ m}^2$}
\loigiai{
Nửa chu vi: $50\text{ m}$. Dài: $30\text{ m}$, Rộng: $20\text{ m}$.\\
Diện tích lối đi: $2 \times 20 = 40\text{ m}^2$.\\
\textbf{Đáp án của câu hỏi là B}
}
\end{ex}

\textbf{Phần 2. Trắc nghiệm đúng sai}
\setcounterpageref{ex}{0}

\begin{ex} %[4H7H35]
Xét tính đúng sai về các hình học đã học:
\choiceTF
{\True Góc nhọn $<$ Góc vuông $<$ Góc tù $<$ Góc bẹt.}
{\True Hai đường chéo hình thoi vuông góc với nhau.}
{Hình bình hành luôn có 4 cạnh bằng nhau.}
{\True Hình vuông vừa là hình chữ nhật đặc biệt vừa là hình thoi đặc biệt.}
\loigiai{
\begin{itemize}
    \item Thứ tự góc: Đúng.
    \item Hai đường chéo hình thoi vuông góc: Đúng.
    \item Hình bình hành chỉ có các cạnh đối bằng nhau. Sai.
    \item Hình vuông là hình chữ nhật và hình thoi đặc biệt: Đúng.
\end{itemize}
\begin{itemize}
    \item a) \textbf{Đúng}.
    \item b) \textbf{Đúng}.
    \item c) \textbf{Sai}.
    \item d) \textbf{Đúng}.
\end{itemize}
}
\end{ex}

\begin{ex} %[4H7V35]
Cho hình bình hành $MNPQ$ có đáy $PQ = 20\text{ cm}$, chiều cao tương ứng $MH = 12\text{ cm}$. Xét tính đúng sai:
\choiceTF
{\True Đoạn thẳng $MH$ vuông góc với cạnh đáy $PQ$.}
{\True Cạnh $MN$ song song và bằng cạnh $PQ = 20\text{ cm}$.}
{Diện tích hình bình hành là $120\text{ cm}^2$.}
{\True Diện tích hình bình hành là $240\text{ cm}^2$.}
\loigiai{
\begin{itemize}
    \item $MH \perp PQ$: Đúng.
    \item $MN // PQ$ và $MN = PQ = 20\text{ cm}$: Đúng.
    \item $S = 20 \times 12 = 240\text{ cm}^2 \neq 120\text{ cm}^2$. Ý 3 Sai, Ý 4 Đúng.
\end{itemize}
\begin{itemize}
    \item a) \textbf{Đúng}.
    \item b) \textbf{Đúng}.
    \item c) \textbf{Sai}.
    \item d) \textbf{Đúng}.
\end{itemize}
}
\end{ex}

\textbf{Phần 3. Trả lời ngắn}
\setcounterpageref{ex}{0}

\begin{ex} %[4H7H35]
Số đo của góc vuông là bao nhiêu độ?
\shortans{90}
\loigiai{
\textbf{Đáp số:} $90$\\
$90^\circ$.
}
\end{ex}

\begin{ex} %[4H7H35]
Tính chu vi hình thoi có cạnh dài $25\text{ cm}$.
\shortans{100}
\loigiai{
\textbf{Đáp số:} $100$\\
$25 \times 4 = 100\text{ cm}$.
}
\end{ex}

\begin{ex} %[4H7V35]
Một hình bình hành có đáy $15\text{ cm}$, diện tích $150\text{ cm}^2$. Tính chiều cao tương ứng theo xăng-ti-mét.
\shortans{10}
\loigiai{
\textbf{Đáp số:} $10$\\
$150 : 15 = 10\text{ cm}$.
}
\end{ex}

\begin{ex} %[4H7V35]
Lúc 4 giờ đúng, góc nhỏ tạo bởi hai kim đồng hồ có số đo bao nhiêu độ?
\shortans{120}
\loigiai{
\textbf{Đáp số:} $120$\\
$4 \times 30^\circ = 120^\circ$.
}
\end{ex}

\begin{ex} %[4H7V35]
Một hình thoi có chu vi $64\text{ cm}$. Tính độ dài cạnh hình thoi theo xăng-ti-mét.
\shortans{16}
\loigiai{
\textbf{Đáp số:} $16$\\
$64 : 4 = 16\text{ cm}$.
}
\end{ex}

\begin{ex} %[4H7C35]
Một hình vuông có cạnh $10\text{ cm}$. Người ta nối trung điểm các cạnh của hình vuông đó để được một hình thoi bên trong. Tính diện tích hình thoi đó theo xăng-ti-mét vuông.
\shortans{50}
\loigiai{
\textbf{Đáp số:} $50$\\
Diện tích hình thoi bằng một nửa diện tích hình vuông: $10 \times 10 : 2 = 50\text{ cm}^2$.
}
\end{ex}

\newpage

\TaoPhieuKiemTra{Chủ đề 7 -- Ôn tập học kì 1}{36}{Ôn tập đo lường}{https://loptoanthaylinh.vn/thi-online/toan-4/bai-36}
\textbf{Phần 1. Trắc nghiệm lựa chọn}
\setcounterpageref{ex}{0}

\begin{ex} %[4D7N36]
$1$ tấn bằng bao nhiêu tạ?
\choice
{$100\text{ tạ}$}
{\True $10\text{ tạ}$}
{$1000\text{ tạ}$}
{$5\text{ tạ}$}
\loigiai{
$1$ tấn $= 10\text{ tạ}$.\\
\textbf{Đáp án của câu hỏi là B}
}
\end{ex}

\begin{ex} %[4D7N36]
$1\text{ m}^2$ bằng bao nhiêu đề-xi-mét vuông?
\choice
{$10\text{ dm}^2$}
{$1000\text{ dm}^2$}
{\True $100\text{ dm}^2$}
{$10\,000\text{ dm}^2$}
\loigiai{
$1\text{ m}^2 = 100\text{ dm}^2$.\\
\textbf{Đáp án của câu hỏi là C}
}
\end{ex}

\begin{ex} %[4D7H36]
$1$ thế kỉ bằng bao nhiêu năm?
\choice
{\True $100$ năm}
{$10$ năm}
{$1000$ năm}
{$50$ năm}
\loigiai{
$1$ thế kỉ $= 100$ năm.\\
\textbf{Đáp án của câu hỏi là A}
}
\end{ex}

\begin{ex} %[4D7H36]
$3$ tạ $25\text{ kg}$ bằng bao nhiêu ki-lô-gam?
\choice
{$3250\text{ kg}$}
{$3025\text{ kg}$}
{$352\text{ kg}$}
{\True $325\text{ kg}$}
\loigiai{
$3 \times 100 + 25 = 325\text{ kg}$.\\
\textbf{Đáp án của câu hỏi là D}
}
\end{ex}

\begin{ex} %[4D7H36]
$2$ phút $45$ giây bằng bao nhiêu giây?
\choice
{$245$ giây}
{$175$ giây}
{\True $165$ giây}
{$185$ giây}
\loigiai{
$2 \times 60 + 45 = 165$ giây.\\
\textbf{Đáp án của câu hỏi là C}
}
\end{ex}

\begin{ex} %[4D7H36]
$5\text{ m}^2 6\text{ dm}^2$ bằng bao nhiêu đề-xi-mét vuông?
\choice
{$56\text{ dm}^2$}
{$560\text{ dm}^2$}
{$5060\text{ dm}^2$}
{\True $506\text{ dm}^2$}
\loigiai{
$5 \times 100 + 6 = 506\text{ dm}^2$.\\
\textbf{Đáp án của câu hỏi là D}
}
\end{ex}

\begin{ex} %[4D7V36]
Năm $1945$ Bác Hồ đọc Tuyên ngôn Độc lập. Năm đó thuộc thế kỉ nào?
\choice
{Thế kỉ XIX}
{\True Thế kỉ XX}
{Thế kỉ XXI}
{Thế kỉ XVIII}
\loigiai{
Năm $1945$ thuộc thế kỉ XX.\\
\textbf{Đáp án của câu hỏi là B}
}
\end{ex}

\begin{ex} %[4D7V36]
Một xe tải chở $2$ tấn $5$ tạ gạo. Xe đã dỡ xuống $8$ tạ gạo. Khối lượng gạo còn lại trên xe là
\choice
{\True $1\text{ tấn } 7\text{ tạ}$ ($17\text{ tạ}$)}
{$1\text{ tấn } 5\text{ tạ}$}
{$1\text{ tấn } 8\text{ tạ}$}
{$2\text{ tấn}$}
\loigiai{
$2$ tấn $5$ tạ $= 25$ tạ. Còn lại: $25 - 8 = 17$ tạ $= 1$ tấn $7$ tạ.\\
\textbf{Đáp án của câu hỏi là A}
}
\end{ex}

\begin{ex} %[4D7V36]
Một căn phòng hình chữ nhật dài $8\text{ m}$, rộng $5\text{ m}$. Diện tích căn phòng là
\choice
{$26\text{ m}^2$}
{\True $40\text{ m}^2$}
{$400\text{ dm}^2$}
{$80\text{ m}^2$}
\loigiai{
$8 \times 5 = 40\text{ m}^2$.\\
\textbf{Đáp án của câu hỏi là B}
}
\end{ex}

\begin{ex} %[4D7V36]
Tính giá trị biểu thức: $4$ tấn $- 15$ tạ $+ 500\text{ kg}$ theo đơn vị tạ.
\choice
{\True $30\text{ tạ}$}
{$25\text{ tạ}$}
{$35\text{ tạ}$}
{$20\text{ tạ}$}
\loigiai{
$40 - 15 + 5 = 30$ tạ.\\
\textbf{Đáp án của câu hỏi là A}
}
\end{ex}

\begin{ex} %[4D7V36]
Một con bò nặng $3$ tạ, một con lợn nặng $60\text{ kg}$. Con bò nặng gấp mấy lần con lợn?
\choice
{$4$ lần}
{$6$ lần}
{\True $5$ lần}
{$8$ lần}
\loigiai{
$3$ tạ $= 300\text{ kg} \Rightarrow 300 : 60 = 5$ lần.\\
\textbf{Đáp án của câu hỏi là C}
}
\end{ex}

\begin{ex} %[4D7C36]
Cần bao nhiêu viên gạch men hình vuông cạnh $5\text{ dm}$ để lát kín sàn một hội trường hình chữ nhật có chiều dài $20\text{ m}$ và chiều rộng $15\text{ m}$?
\choice
{$1000$ viên}
{$1500$ viên}
{$2000$ viên}
{\True $1200$ viên}
\loigiai{
Diện tích hội trường: $20 \times 15 = 300\text{ m}^2 = 30\,000\text{ dm}^2$.\\
Diện tích 1 viên gạch: $5 \times 5 = 25\text{ dm}^2$.\\
Số viên gạch: $30\,000 : 25 = 1200$ viên.\\
\textbf{Đáp án của câu hỏi là D}
}
\end{ex}

\textbf{Phần 2. Trắc nghiệm đúng sai}
\setcounterpageref{ex}{0}

\begin{ex} %[4D7H36]
Xét tính đúng sai của các biểu thức đổi đơn vị:
\choiceTF
{\True $2$ tấn $50\text{ kg} = 2050\text{ kg}$.}
{$3\text{ m}^2 4\text{ dm}^2 = 340\text{ dm}^2$.}
{\True $3\text{ m}^2 4\text{ dm}^2 = 304\text{ dm}^2$.}
{\True $1$ giờ $15$ phút $= 75$ phút.}
\loigiai{
\begin{itemize}
    \item $2$ tấn $50\text{ kg} = 2050\text{ kg}$. Đúng.
    \item $3\text{ m}^2 4\text{ dm}^2 = 304\text{ dm}^2 \neq 340\text{ dm}^2$. Ý 2 Sai, Ý 3 Đúng.
    \item $60 + 15 = 75$ phút. Đúng.
\end{itemize}
\begin{itemize}
    \item a) \textbf{Đúng}.
    \item b) \textbf{Sai}.
    \item c) \textbf{Đúng}.
    \item d) \textbf{Đúng}.
\end{itemize}
}
\end{ex}

\begin{ex} %[4D7V36]
Bác Bình thu hoạch được $3$ tấn $2$ tạ thóc: Bác bán đi $18$ tạ thóc, số thóc còn lại đóng đều vào các bao $50\text{ kg}$. Xét tính đúng sai:
\choiceTF
{\True Tổng số thóc thu hoạch là $32$ tạ ($3200\text{ kg}$).}
{\True Số thóc còn lại sau khi bán là $14$ tạ ($1400\text{ kg}$).}
{Bác Bình đóng được tất cả $20$ bao thóc.}
{\True Bác Bình đóng được tất cả $28$ bao thóc.}
\loigiai{
\begin{itemize}
    \item Tổng: $32$ tạ. Đúng.
    \item Còn lại: $32 - 18 = 14$ tạ $= 1400\text{ kg}$. Đúng.
    \item Số bao: $1400 : 50 = 28$ bao $\neq 20$ bao. Ý 3 Sai, Ý 4 Đúng.
\end{itemize}
\begin{itemize}
    \item a) \textbf{Đúng}.
    \item b) \textbf{Đúng}.
    \item c) \textbf{Sai}.
    \item d) \textbf{Đúng}.
\end{itemize}
}
\end{ex}

\textbf{Phần 3. Trả lời ngắn}
\setcounterpageref{ex}{0}

\begin{ex} %[4D7H36]
Đổi $5$ tấn $8$ tạ ra tạ.
\shortans{58}
\loigiai{
\textbf{Đáp số:} $58$\\
$58$ tạ.
}
\end{ex}

\begin{ex} %[4D7H36]
Đổi $7\text{ m}^2 25\text{ dm}^2$ ra đề-xi-mét vuông.
\shortans{725}
\loigiai{
\textbf{Đáp số:} $725$\\
$725\text{ dm}^2$.
}
\end{ex}

\begin{ex} %[4D7V36]
Đổi $300$ giây ra phút.
\shortans{5}
\loigiai{
\textbf{Đáp số:} $5$\\
$300 : 60 = 5$ phút.
}
\end{ex}

\begin{ex} %[4D7V36]
Năm $1789$ thuộc thế kỉ thứ mấy (ghi bằng số Ả Rập)?
\shortans{18}
\loigiai{
\textbf{Đáp số:} $18$\\
Thế kỉ 18 (XVIII).
}
\end{ex}

\begin{ex} %[4D7V36]
Một mảnh vườn hình chữ nhật dài $16\text{ m}$, rộng $10\text{ m}$. Tính diện tích mảnh vườn theo mét vuông.
\shortans{160}
\loigiai{
\textbf{Đáp số:} $160$\\
$16 \times 10 = 160\text{ m}^2$.
}
\end{ex}

\begin{ex} %[4D7C36]
Có hai xe tải chở hàng: Xe 1 chở $3$ tấn $5$ tạ, xe 2 chở ít hơn xe 1 là $8$ tạ. Tính tổng khối lượng hàng cả hai xe chở theo tạ.
\shortans{62}
\loigiai{
\textbf{Đáp số:} $62$\\
Xe 2: $35 - 8 = 27$ tạ. Tổng: $35 + 27 = 62$ tạ.
}
\end{ex}

\newpage

\TaoPhieuKiemTra{Chủ đề 7 -- Ôn tập học kì 1}{37}{Ôn tập chung học kì 1}{https://loptoanthaylinh.vn/thi-online/toan-4/bai-37}
\textbf{Phần 1. Trắc nghiệm lựa chọn}
\setcounterpageref{ex}{0}

\begin{ex} %[4D7N37]
Số gồm $8$ chục triệu, $5$ trăm nghìn và $3$ đơn vị viết là
\choice
{\True $80\,500\,003$}
{$85\,000\,003$}
{$80\,050\,003$}
{$850\,000\,003$}
\loigiai{
Số đó là $80\,500\,003$.\\
\textbf{Đáp án của câu hỏi là A}
}
\end{ex}

\begin{ex} %[4D7N37]
Kết quả của phép tính $560\,000 + 440\,000$ là
\choice
{$900\,000$}
{$1\,100\,000$}
{$10\,000\,000$}
{\True $1\,000\,000$}
\loigiai{
$560\,000 + 440\,000 = 1\,000\,000$.\\
\textbf{Đáp án của câu hỏi là D}
}
\end{ex}

\begin{ex} %[4D7H37]
Hai số có tổng bằng $150$ và hiệu bằng $50$. Số lớn là
\choice
{$90$}
{\True $100$}
{$110$}
{$120$}
\loigiai{
$(150 + 50) : 2 = 100$.\\
\textbf{Đáp án của câu hỏi là B}
}
\end{ex}

\begin{ex} %[4D7H37]
Một hình bình hành có độ dài đáy $25\text{ cm}$, chiều cao $14\text{ cm}$. Diện tích của nó là
\choice
{$175\text{ cm}^2$}
{$78\text{ cm}^2$}
{\True $350\text{ cm}^2$}
{$300\text{ cm}^2$}
\loigiai{
$S = 25 \times 14 = 350\text{ cm}^2$.\\
\textbf{Đáp án của câu hỏi là C}
}
\end{ex}

\begin{ex} %[4D7H37]
Làm tròn số $8\,650\,000$ đến hàng trăm nghìn ta được
\choice
{\True $8\,700\,000$}
{$8\,600\,000$}
{$9\,000\,000$}
{$8\,650\,000$}
\loigiai{
Hàng chục nghìn là 5 nên làm tròn lên $8\,700\,000$.\\
\textbf{Đáp án của câu hỏi là A}
}
\end{ex}

\begin{ex} %[4D7H37]
$4$ tấn $50\text{ kg}$ bằng bao nhiêu ki-lô-gam?
\choice
{$450\text{ kg}$}
{$4500\text{ kg}$}
{\True $4050\text{ kg}$}
{$40\,050\text{ kg}$}
\loigiai{
$4 \times 1000 + 50 = 4050\text{ kg}$.\\
\textbf{Đáp án của câu hỏi là C}
}
\end{ex}

\begin{ex} %[4D7V37]
Một khu đất hình chữ nhật có chu vi $140\text{ m}$, chiều dài hơn chiều rộng $10\text{ m}$. Diện tích khu đất đó là
\choice
{$1100\text{ m}^2$}
{$1300\text{ m}^2$}
{$1400\text{ m}^2$}
{\True $1200\text{ m}^2$}
\loigiai{
Nửa chu vi: $70\text{ m}$. Dài: $(70 + 10) : 2 = 40\text{ m}$. Rộng: $30\text{ m}$.\\
Diện tích: $40 \times 30 = 1200\text{ m}^2$.\\
\textbf{Đáp án của câu hỏi là D}
}
\end{ex}

\begin{ex} %[4D7V37]
Tính nhanh giá trị biểu thức: $125 + 348 + 875 + 652$.
\choice
{$1900$}
{\True $2000$}
{$2100$}
{$2200$}
\loigiai{
$(125 + 875) + (348 + 652) = 1000 + 1000 = 2000$.\\
\textbf{Đáp án của câu hỏi là B}
}
\end{ex}

\begin{ex} %[4D7V37]
Một hình thoi có chu vi $80\text{ cm}$. Độ dài cạnh của hình thoi là
\choice
{$15\text{ cm}$}
{$25\text{ cm}$}
{$40\text{ cm}$}
{\True $20\text{ cm}$}
\loigiai{
$80 : 4 = 20\text{ cm}$.\\
\textbf{Đáp án của câu hỏi là D}
}
\end{ex}

\begin{ex} %[4D7V37]
Lúc 5 giờ đúng, góc nhỏ tạo bởi kim giờ và kim phút của đồng hồ có số đo là
\choice
{$120^\circ$}
{\True $150^\circ$}
{$180^\circ$}
{$90^\circ$}
\loigiai{
$5 \times 30^\circ = 150^\circ$.\\
\textbf{Đáp án của câu hỏi là B}
}
\end{ex}

\begin{ex} %[4D7V37]
Một đội công nhân: Ngày 1 làm được $1200\text{ m}$ đường, ngày 2 làm gấp đôi ngày 1, ngày 3 làm ít hơn ngày 2 là $400\text{ m}$. Cả 3 ngày đội làm được số mét đường là
\choice
{\True $5600\text{ m}$}
{$5400\text{ m}$}
{$5800\text{ m}$}
{$6000\text{ m}$}
\loigiai{
Ngày 2: $2400\text{ m}$. Ngày 3: $2000\text{ m}$. Tổng: $1200 + 2400 + 2000 = 5600\text{ m}$.\\
\textbf{Đáp án của câu hỏi là A}
}
\end{ex}

\begin{ex} %[4D7C37]
Tìm số tự nhiên có hai chữ số biết rằng tổng hai chữ số bằng 11 và nếu đổi vị trí hai chữ số cho nhau thì số đó tăng thêm 27 đơn vị.
\choice
{$38$}
{$56$}
{\True $47$}
{$29$}
\loigiai{
Số đó là 47 vì $4 + 7 = 11$ và $74 - 47 = 27$.\\
\textbf{Đáp án của câu hỏi là C}
}
\end{ex}

\textbf{Phần 2. Trắc nghiệm đúng sai}
\setcounterpageref{ex}{0}

\begin{ex} %[4D7H37]
Xét tính đúng sai của các nhận định học kì 1:
\choiceTF
{\True Hai đường chéo của hình thoi vuông góc với nhau.}
{\True $1$ tấn $= 1000\text{ kg}$.}
{Góc $120^\circ$ là góc nhọn.}
{\True Số lớn nhất có 6 chữ số là $999\,999$.}
\loigiai{
\begin{itemize}
    \item Hai đường chéo hình thoi vuông góc: Đúng.
    \item 1 tấn $= 1000$ kg: Đúng.
    \item $120^\circ$ là góc tù (không phải góc nhọn). Sai.
    \item $999\,999$ lớn nhất có 6 chữ số: Đúng.
\end{itemize}
\begin{itemize}
    \item a) \textbf{Đúng}.
    \item b) \textbf{Đúng}.
    \item c) \textbf{Sai}.
    \item d) \textbf{Đúng}.
\end{itemize}
}
\end{ex}

\begin{ex} %[4D7V37]
Một mảnh đất hình chữ nhật có chu vi $120\text{ m}$, chiều dài hơn chiều rộng $16\text{ m}$. Xét tính đúng sai:
\choiceTF
{\True Nửa chu vi mảnh đất là $60\text{ m}$.}
{\True Chiều dài mảnh đất là $38\text{ m}$.}
{Chiều rộng mảnh đất là $24\text{ m}$.}
{\True Diện tích mảnh đất là $836\text{ m}^2$.}
\loigiai{
\begin{itemize}
    \item Nửa chu vi: $60\text{ m}$. Đúng.
    \item Chiều dài: $(60 + 16) : 2 = 38\text{ m}$. Đúng.
    \item Chiều rộng: $60 - 38 = 22\text{ m} \neq 24\text{ m}$. Sai.
    \item Diện tích: $38 \times 22 = 836\text{ m}^2$. Đúng.
\end{itemize}
\begin{itemize}
    \item a) \textbf{Đúng}.
    \item b) \textbf{Đúng}.
    \item c) \textbf{Sai}.
    \item d) \textbf{Đúng}.
\end{itemize}
}
\end{ex}

\textbf{Phần 3. Trả lời ngắn}
\setcounterpageref{ex}{0}

\begin{ex} %[4D7H37]
Tính tổng: $350\,000 + 650\,000$.
\shortans{1000000}
\loigiai{
\textbf{Đáp số:} $1000000$\\
$1\,000\,000$.
}
\end{ex}

\begin{ex} %[4D7H37]
Một hình thoi có cạnh $15\text{ cm}$. Tính chu vi hình thoi theo xăng-ti-mét.
\shortans{60}
\loigiai{
\textbf{Đáp số:} $60$\\
$15 \times 4 = 60\text{ cm}$.
}
\end{ex}

\begin{ex} %[4D7V37]
Đổi $3$ tấn $8$ tạ ra tạ.
\shortans{38}
\loigiai{
\textbf{Đáp số:} $38$\\
$38$ tạ.
}
\end{ex}

\begin{ex} %[4D7V37]
Hai số có tổng là $160$, hiệu là $40$. Tìm số lớn.
\shortans{100}
\loigiai{
\textbf{Đáp số:} $100$\\
$(160 + 40) : 2 = 100$.
}
\end{ex}

\begin{ex} %[4D7V37]
Một hình bình hành có đáy $16\text{ cm}$, chiều cao $12\text{ cm}$. Tính diện tích theo xăng-ti-mét vuông.
\shortans{192}
\loigiai{
\textbf{Đáp số:} $192$\\
$16 \times 12 = 192\text{ cm}^2$.
}
\end{ex}

\begin{ex} %[4D7C37]
Tính giá trị biểu thức: $1 + 2 + 3 + \dots + 50$.
\shortans{1275}
\loigiai{
\textbf{Đáp số:} $1275$\\
$(1 + 50) \times 50 : 2 = 1275$.
}
\end{ex}

\newpage

% =========================================================================
% CHỦ ĐỀ 8 - PHÉP NHÂN VÀ PHÉP CHIA
% =========================================================================
\section{Phép nhân và phép chia}

\TaoPhieuKiemTra{Chủ đề 8 -- Phép nhân và phép chia}{38}{Nhân với số có một chữ số}{https://loptoanthaylinh.vn/thi-online/toan-4/bai-38}
\textbf{Phần 1. Trắc nghiệm lựa chọn}
\setcounterpageref{ex}{0}

\begin{ex} %[4D8N38]
Kết quả của phép nhân $23\,142 \times 3$ là
\choice
{$69\,326$}
{$69\,416$}
{$68\,426$}
{\True $69\,426$}
\loigiai{
$23\,142 \times 3 = 69\,426$.\\
\textbf{Đáp án của câu hỏi là D}
}
\end{ex}

\begin{ex} %[4D8N38]
Kết quả của phép tính $120\,000 \times 4$ là
\choice
{\True $480\,000$}
{$460\,000$}
{$400\,000$}
{$500\,000$}
\loigiai{
$120\,000 \times 4 = 480\,000$.\\
\textbf{Đáp án của câu hỏi là A}
}
\end{ex}

\begin{ex} %[4D8H38]
Tìm $x$ biết: $x : 6 = 14\,250$.
\choice
{$84\,500$}
{$85\,200$}
{$86\,500$}
{\True $85\,500$}
\loigiai{
$x = 14\,250 \times 6 = 85\,500$.\\
\textbf{Đáp án của câu hỏi là D}
}
\end{ex}

\begin{ex} %[4D8H38]
Kết quả của phép tính $45\,678 \times 5$ là
\choice
{$228\,380$}
{\True $228\,390$}
{$225\,390$}
{$228\,490$}
\loigiai{
$45\,678 \times 5 = 228\,390$.\\
\textbf{Đáp án của câu hỏi là B}
}
\end{ex}

\begin{ex} %[4D8H38]
Một hộp có $125$ chiếc bút. Hỏi $8$ hộp như thế có tất cả bao nhiêu chiếc bút?
\choice
{$900$ chiếc}
{$1100$ chiếc}
{\True $1000$ chiếc}
{$800$ chiếc}
\loigiai{
$125 \times 8 = 1000$ chiếc.\\
\textbf{Đáp án của câu hỏi là C}
}
\end{ex}

\begin{ex} %[4D8H38]
Giá trị của biểu thức $102\,300 \times 3 + 45\,000$ là
\choice
{\True $351\,900$}
{$341\,900$}
{$350\,900$}
{$361\,900$}
\loigiai{
$306\,900 + 45\,000 = 351\,900$.\\
\textbf{Đáp án của câu hỏi là A}
}
\end{ex}

\begin{ex} %[4D8V38]
Một phân xưởng may mỗi ngày may được $1250$ bộ quần áo. Hỏi trong một tuần lễ làm việc 6 ngày, phân xưởng đó may được bao nhiêu bộ quần áo?
\choice
{$7200$ bộ}
{\True $7500$ bộ}
{$7800$ bộ}
{$8000$ bộ}
\loigiai{
$1250 \times 6 = 7500$ bộ.\\
\textbf{Đáp án của câu hỏi là B}
}
\end{ex}

\begin{ex} %[4D8V38]
Một ô tô chở hàng nặng $3$ tấn $500\text{ kg}$. Hỏi 4 ô tô chở như thế chở được tất cả bao nhiêu tấn hàng?
\choice
{$12\text{ tấn}$}
{$15\text{ tấn}$}
{\True $14\text{ tấn}$}
{$16\text{ tấn}$}
\loigiai{
$3500 \times 4 = 14\,000\text{ kg} = 14\text{ tấn}$.\\
\textbf{Đáp án của câu hỏi là C}
}
\end{ex}

\begin{ex} %[4D8V38]
Tính giá trị của biểu thức: $25\,000 \times 4 - 30\,000 : 3$.
\choice
{$80\,000$}
{$100\,000$}
{$70\,000$}
{\True $90\,000$}
\loigiai{
$100\,000 - 10\,000 = 90\,000$.\\
\textbf{Đáp án của câu hỏi là D}
}
\end{ex}

\begin{ex} %[4D8V38]
Một trường tiểu học mua 8 bộ bàn ghế cho phòng họp, mỗi bộ có giá $1\,450\,000$ đồng. Tổng số tiền mua bàn ghế là
\choice
{$11\,400\,000$ đồng}
{\True $11\,600\,000$ đồng}
{$11\,800\,000$ đồng}
{$12\,000\,000$ đồng}
\loigiai{
$1\,450\,000 \times 8 = 11\,600\,000$ đồng.\\
\textbf{Đáp án của câu hỏi là B}
}
\end{ex}

\begin{ex} %[4D8V38]
Tìm chữ số $a$ biết: $12a,400 \times 3 = 373\,200$.
\choice
{\True $a = 4$}
{$a = 3$}
{$a = 5$}
{$a = 2$}
\loigiai{
$373\,200 : 3 = 124\,400 \Rightarrow a = 4$.\\
\textbf{Đáp án của câu hỏi là A}
}
\end{ex}

\begin{ex} %[4D8C38]
Khi nhân một số có 5 chữ số với 4, một học sinh đã viết nhầm chữ số hàng trăm của số đó từ 3 thành 8 nên tích tăng thêm bao nhiêu đơn vị?
\choice
{$500$ đơn vị}
{$1500$ đơn vị}
{\True $2000$ đơn vị}
{$5000$ đơn vị}
\loigiai{
Thừa số thứ nhất tăng thêm: $800 - 300 = 500$.\\
Tích tăng thêm: $500 \times 4 = 2000$ đơn vị.\\
\textbf{Đáp án của câu hỏi là C}
}
\end{ex}

\textbf{Phần 2. Trắc nghiệm đúng sai}
\setcounterpageref{ex}{0}

\begin{ex} %[4D8H38]
Xét tính đúng sai của các phép nhân:
\choiceTF
{\True $12\,345 \times 2 = 24\,690$.}
{$25\,000 \times 4 = 1\,000\,000$.}
{\True $25\,000 \times 4 = 100\,000$.}
{\True $150\,000 \times 6 = 900\,000$.}
\loigiai{
\begin{itemize}
    \item $12\,345 \times 2 = 24\,690$. Đúng.
    \item $25\,000 \times 4 = 100\,000 \neq 1\,000\,000$. Ý 2 Sai, Ý 3 Đúng.
    \item $150\,000 \times 6 = 900\,000$. Đúng.
\end{itemize}
\begin{itemize}
    \item a) \textbf{Đúng}.
    \item b) \textbf{Sai}.
    \item c) \textbf{Đúng}.
    \item d) \textbf{Đúng}.
\end{itemize}
}
\end{ex}

\begin{ex} %[4D8V38]
Một cửa hàng nhập về 5 bao đường, mỗi bao nặng $50\text{ kg}$. Giá mỗi ki-lô-gam đường là $22\,000$ đồng. Xét tính đúng sai:
\choiceTF
{\True Cửa hàng đã nhập về tất cả $250\text{ kg}$ đường.}
{\True Số đường đó tương đương $2{,}5$ tạ đường.}
{Tổng số tiền mua đường là $5\,000\,000$ đồng.}
{\True Tổng số tiền mua đường là $5\,500\,000$ đồng.}
\loigiai{
\begin{itemize}
    \item $5 \times 50 = 250\text{ kg}$. Đúng.
    \item $250\text{ kg} = 2{,}5$ tạ. Đúng.
    \item Tiền: $250 \times 22\,000 = 5\,500\,000$ đồng $\neq 5\,000\,000$ đồng. Ý 3 Sai, Ý 4 Đúng.
\end{itemize}
\begin{itemize}
    \item a) \textbf{Đúng}.
    \item b) \textbf{Đúng}.
    \item c) \textbf{Sai}.
    \item d) \textbf{Đúng}.
\end{itemize}
}
\end{ex}

\textbf{Phần 3. Trả lời ngắn}
\setcounterpageref{ex}{0}

\begin{ex} %[4D8H38]
Tính tích: $15\,000 \times 6$.
\shortans{90000}
\loigiai{
\textbf{Đáp số:} $90000$\\
$90\,000$.
}
\end{ex}

\begin{ex} %[4D8H38]
Tính tích: $125\,000 \times 8$.
\shortans{1000000}
\loigiai{
\textbf{Đáp số:} $1000000$\\
$1\,000\,000$.
}
\end{ex}

\begin{ex} %[4D8V38]
Tìm $x$ biết: $x : 7 = 35\,000$.
\shortans{245000}
\loigiai{
\textbf{Đáp số:} $245000$\\
$x = 35\,000 \times 7 = 245\,000$.
}
\end{ex}

\begin{ex} %[4D8V38]
Một ngày có $24$ giờ, mỗi giờ có $60$ phút. Tính số phút trong 3 ngày.
\shortans{4320}
\loigiai{
\textbf{Đáp số:} $4320$\\
$24 \times 60 \times 3 = 4320$ phút.
}
\end{ex}

\begin{ex} %[4D8V38]
Một cửa hàng bán được 8 thùng dầu, mỗi thùng chứa $25\text{ l}$ dầu. Giá mỗi lít dầu là $20\,000$ đồng. Tính tổng số tiền thu được theo nghìn đồng (chỉ ghi số nghìn đồng).
\shortans{4000}
\loigiai{
\textbf{Đáp số:} $4000$\\
$8 \times 25 \times 20\,000 = 4\,000\,000$ đồng $= 4000$ nghìn đồng.
}
\end{ex}

\begin{ex} %[4D8C38]
Tích của hai số là $45\,000$. Nếu một thừa số gấp lên 3 lần và thừa số kia giữ nguyên thì tích mới là bao nhiêu?
\shortans{135000}
\loigiai{
\textbf{Đáp số:} $135000$\\
$45\,000 \times 3 = 135\,000$.
}
\end{ex}

\newpage

\TaoPhieuKiemTra{Chủ đề 8 -- Phép nhân và phép chia}{39}{Chia cho số có một chữ số}{https://loptoanthaylinh.vn/thi-online/toan-4/bai-39}
\textbf{Phần 1. Trắc nghiệm lựa chọn}
\setcounterpageref{ex}{0}

\begin{ex} %[4D8N39]
Kết quả của phép chia $84\,628 : 2$ là
\choice
{\True $42\,314$}
{$42\,313$}
{$42\,304$}
{$41\,314$}
\loigiai{
$84\,628 : 2 = 42\,314$.\\
\textbf{Đáp án của câu hỏi là A}
}
\end{ex}

\begin{ex} %[4D8N39]
Số dư của phép chia $245\,678 : 5$ là
\choice
{$1$}
{$2$}
{\True $3$}
{$4$}
\loigiai{
$8 : 5 = 1$ dư 3.\\
\textbf{Đáp án của câu hỏi là C}
}
\end{ex}

\begin{ex} %[4D8H39]
Kết quả của phép chia $360\,000 : 9$ là
\choice
{$4000$}
{\True $40\,000$}
{$400\,000$}
{$30\,000$}
\loigiai{
$360\,000 : 9 = 40\,000$.\\
\textbf{Đáp án của câu hỏi là B}
}
\end{ex}

\begin{ex} %[4D8H39]
Tìm $x$ biết: $x \times 4 = 184\,000$.
\choice
{$45\,000$}
{$47\,000$}
{$48\,000$}
{\True $46\,000$}
\loigiai{
$x = 184\,000 : 4 = 46\,000$.\\
\textbf{Đáp án của câu hỏi là D}
}
\end{ex}

\begin{ex} %[4D8H39]
Một hình vuông có chu vi $96\text{ cm}$. Độ dài cạnh hình vuông là
\choice
{$22\text{ cm}$}
{\True $24\text{ cm}$}
{$26\text{ cm}$}
{$48\text{ cm}$}
\loigiai{
$96 : 4 = 24\text{ cm}$.\\
\textbf{Đáp án của câu hỏi là B}
}
\end{ex}

\begin{ex} %[4D8H39]
Kết quả của phép chia $125\,000 : 5$ là
\choice
{\True $25\,000$}
{$20\,000$}
{$30\,000$}
{$250\,000$}
\loigiai{
$125\,000 : 5 = 25\,000$.\\
\textbf{Đáp án của câu hỏi là A}
}
\end{ex}

\begin{ex} %[4D8V39]
Người ta đóng $12\,450\text{ kg}$ gạo vào các bao loại $5\text{ kg}$. Số bao gạo đóng được là
\choice
{$2480$ bao}
{$2500$ bao}
{$2450$ bao}
{\True $2490$ bao}
\loigiai{
$12\,450 : 5 = 2490$ bao.\\
\textbf{Đáp án của câu hỏi là D}
}
\end{ex}

\begin{ex} %[4D8V39]
Có $350$ học sinh đi tham quan bằng xe ô tô loại 45 chỗ ngồi (không tính ghế lái xe). Cần ít nhất bao nhiêu xe ô tô như thế để chở hết số học sinh?
\choice
{$7$ xe}
{$9$ xe}
{\True $8$ xe}
{$10$ xe}
\loigiai{
$350 : 45 = 7$ dư 35 em. Cần ít nhất $7 + 1 = 8$ xe.\\
\textbf{Đáp án của câu hỏi là C}
}
\end{ex}

\begin{ex} %[4D8V39]
Một xí nghiệp sản xuất được $48\,600$ hộp bánh trong 6 ngày. Hỏi trung bình mỗi ngày xí nghiệp sản xuất được bao nhiêu hộp bánh?
\choice
{\True $8100$ hộp}
{$8000$ hộp}
{$8200$ hộp}
{$8300$ hộp}
\loigiai{
$48\,600 : 6 = 8100$ hộp.\\
\textbf{Đáp án của câu hỏi là A}
}
\end{ex}

\begin{ex} %[4D8V39]
Tính giá trị biểu thức: $450\,000 - 360\,000 : 4$.
\choice
{$340\,000$}
{$380\,000$}
{$22\,500$}
{\True $360\,000$}
\loigiai{
$450\,000 - 90\,000 = 360\,000$.\\
\textbf{Đáp án của câu hỏi là D}
}
\end{ex}

\begin{ex} %[4D8V39]
Tìm số bị chia trong phép chia cho 7 có thương là $12\,340$ và số dư lớn nhất có thể.
\choice
{$86\,380$}
{$86\,385$}
{\True $86\,386$}
{$86\,387$}
\loigiai{
Số dư lớn nhất là 6.\\
Số bị chia: $12\,340 \times 7 + 6 = 86\,380 + 6 = 86\,386$.\\
\textbf{Đáp án của câu hỏi là C}
}
\end{ex}

\begin{ex} %[4D8C39]
Khi chia một số cho 6 được thương là $15\,200$ và không còn dư. Nếu chia số đó cho 8 thì được thương là bao nhiêu?
\choice
{$11\,300$}
{\True $11\,400$}
{$11\,500$}
{$12\,000$}
\loigiai{
Số đó là: $15\,200 \times 6 = 91\,200$.\\
Chia cho 8: $91\,200 : 8 = 11\,400$.\\
\textbf{Đáp án của câu hỏi là B}
}
\end{ex}

\textbf{Phần 2. Trắc nghiệm đúng sai}
\setcounterpageref{ex}{0}

\begin{ex} %[4D8H39]
Xét tính đúng sai của các phép chia:
\choiceTF
{\True $150\,000 : 3 = 50\,000$.}
{$45\,678 : 9 = 5075$ dư 5.}
{\True $45\,678 : 9 = 5075$ dư 3.}
{\True Trong phép chia có dư, số dư luôn nhỏ hơn số chia.}
\loigiai{
\begin{itemize}
    \item $150\,000 : 3 = 50\,000$. Đúng.
    \item $45\,678 : 9 = 5075$ dư 3 ($5075 \times 9 + 3 = 45678$). Ý 2 Sai, Ý 3 Đúng.
    \item Số dư $<$ số chia: Đúng.
\end{itemize}
\begin{itemize}
    \item a) \textbf{Đúng}.
    \item b) \textbf{Sai}.
    \item c) \textbf{Đúng}.
    \item d) \textbf{Đúng}.
\end{itemize}
}
\end{ex}

\begin{ex} %[4D8V39]
Có $120$ lít mật ong đóng đều vào các can $5$ lít. Mỗi can bán với giá $450\,000$ đồng. Xét tính đúng sai:
\choiceTF
{\True Đóng được tất cả $24$ can mật ong.}
{\True Mỗi lít mật ong có giá là $90\,000$ đồng.}
{Tổng số tiền bán mật ong là $10\,000\,000$ đồng.}
{\True Tổng số tiền bán mật ong là $10\,800\,000$ đồng.}
\loigiai{
\begin{itemize}
    \item $120 : 5 = 24$ can. Đúng.
    \item $450\,000 : 5 = 90\,000$ đồng/lít. Đúng.
    \item Tổng tiền: $24 \times 450\,000 = 10\,800\,000$ đồng $\neq 10\,000\,000$ đồng. Ý 3 Sai, Ý 4 Đúng.
\end{itemize}
\begin{itemize}
    \item a) \textbf{Đúng}.
    \item b) \textbf{Đúng}.
    \item c) \textbf{Sai}.
    \item d) \textbf{Đúng}.
\end{itemize}
}
\end{ex}

\textbf{Phần 3. Trả lời ngắn}
\setcounterpageref{ex}{0}

\begin{ex} %[4D8H39]
Tính thương: $240\,000 : 6$.
\shortans{40000}
\loigiai{
\textbf{Đáp số:} $40000$\\
$40\,000$.
}
\end{ex}

\begin{ex} %[4D8H39]
Tính thương: $500\,000 : 5$.
\shortans{100000}
\loigiai{
\textbf{Đáp số:} $100000$\\
$100\,000$.
}
\end{ex}

\begin{ex} %[4D8V39]
Tìm $x$ biết: $x \times 8 = 320\,000$.
\shortans{40000}
\loigiai{
\textbf{Đáp số:} $40000$\\
$x = 320\,000 : 8 = 40\,000$.
}
\end{ex}

\begin{ex} %[4D8V39]
Một cuộn dây dài $450\text{ m}$, người ta cắt thành các đoạn dài $9\text{ m}$. Hỏi cắt được bao nhiêu đoạn dây?
\shortans{50}
\loigiai{
\textbf{Đáp số:} $50$\\
$450 : 9 = 50$ đoạn.
}
\end{ex}

\begin{ex} %[4D8V39]
Có $155\text{ kg}$ gạo đóng vào các túi $5\text{ kg}$. Hỏi đóng được bao nhiêu túi?
\shortans{31}
\loigiai{
\textbf{Đáp số:} $31$\\
$155 : 5 = 31$ túi.
}
\end{ex}

\begin{ex} %[4D8C39]
Một số chia cho 5 dư 3. Hỏi số đó cộng thêm ít nhất bao nhiêu đơn vị để được số chia hết cho 5?
\shortans{2}
\loigiai{
\textbf{Đáp số:} $2$\\
$5 - 3 = 2$ đơn vị.
}
\end{ex}

\newpage

\TaoPhieuKiemTra{Chủ đề 8 -- Phép nhân và phép chia}{40}{Tính chất giao hoán, kết hợp của phép nhân}{https://loptoanthaylinh.vn/thi-online/toan-4/bai-40}
\textbf{Phần 1. Trắc nghiệm lựa chọn}
\setcounterpageref{ex}{0}

\begin{ex} %[4D8N40]
Tính chất giao hoán của phép nhân là
\choice
{$a \times b = a + b$}
{$(a \times b) \times c = a \times (b \times c)$}
{\True $a \times b = b \times a$}
{$a \times 1 = a$}
\loigiai{
$a \times b = b \times a$.\\
\textbf{Đáp án của câu hỏi là C}
}
\end{ex}

\begin{ex} %[4D8N40]
Tính chất kết hợp của phép nhân là
\choice
{$a \times b = b \times a$}
{\True $(a \times b) \times c = a \times (b \times c)$}
{$a \times (b + c) = a \times b + a \times c$}
{$a \times 0 = 0$}
\loigiai{
$(a \times b) \times c = a \times (b \times c)$.\\
\textbf{Đáp án của câu hỏi là B}
}
\end{ex}

\begin{ex} %[4D8H40]
Tính thuận tiện: $25 \times 37 \times 4$.
\choice
{$3600$}
{$3800$}
{\True $3700$}
{$3500$}
\loigiai{
$(25 \times 4) \times 37 = 100 \times 37 = 3700$.\\
\textbf{Đáp án của câu hỏi là C}
}
\end{ex}

\begin{ex} %[4D8H40]
Tính thuận tiện: $125 \times 48 \times 8$.
\choice
{$46\,000$}
{\True $48\,000$}
{$50\,000$}
{$480\,000$}
\loigiai{
$(125 \times 8) \times 48 = 1000 \times 48 = 48\,000$.\\
\textbf{Đáp án của câu hỏi là B}
}
\end{ex}

\begin{ex} %[4D8H40]
Tính thuận tiện: $50 \times 19 \times 2$.
\choice
{\True $1900$}
{$1800$}
{$2000$}
{$1950$}
\loigiai{
$(50 \times 2) \times 19 = 100 \times 19 = 1900$.\\
\textbf{Đáp án của câu hỏi là A}
}
\end{ex}

\begin{ex} %[4D8H40]
Tính thuận tiện: $2 \times 456 \times 5$.
\choice
{$4460$}
{\True $4560$}
{$4660$}
{$4500$}
\loigiai{
$(2 \times 5) \times 456 = 10 \times 456 = 4560$.\\
\textbf{Đáp án của câu hỏi là B}
}
\end{ex}

\begin{ex} %[4D8V40]
Tính nhanh giá trị biểu thức: $4 \times 18 \times 25$.
\choice
{$1600$}
{$2000$}
{$1700$}
{\True $1800$}
\loigiai{
$(4 \times 25) \times 18 = 100 \times 18 = 1800$.\\
\textbf{Đáp án của câu hỏi là D}
}
\end{ex}

\begin{ex} %[4D8V40]
Tính nhanh: $8 \times 35 \times 125$.
\choice
{\True $35\,000$}
{$34\,000$}
{$36\,000$}
{$350\,000$}
\loigiai{
$(8 \times 125) \times 35 = 1000 \times 35 = 35\,000$.\\
\textbf{Đáp án của câu hỏi là A}
}
\end{ex}

\begin{ex} %[4D8V40]
Một phòng học có 8 dãy bàn, mỗi dãy có 5 bàn, mỗi bàn có 2 học sinh ngồi. Số học sinh trong phòng học đó tính thuận tiện là
\choice
{$70$ học sinh}
{$90$ học sinh}
{$100$ học sinh}
{\True $80$ học sinh}
\loigiai{
$8 \times (5 \times 2) = 8 \times 10 = 80$ học sinh.\\
\textbf{Đáp án của câu hỏi là D}
}
\end{ex}

\begin{ex} %[4D8V40]
Một cửa hàng xếp 4 thùng bánh, mỗi thùng có 25 hộp bánh, mỗi hộp bánh giá $15\,000$ đồng. Tổng giá trị của 4 thùng bánh là
\choice
{\True $1\,500\,000$ đồng}
{$1\,400\,000$ đồng}
{$1\,600\,000$ đồng}
{$1\,800\,000$ đồng}
\loigiai{
$(4 \times 25) \times 15\,000 = 100 \times 15\,000 = 1\,500\,000$ đồng.\\
\textbf{Đáp án của câu hỏi là A}
}
\end{ex}

\begin{ex} %[4D8V40]
Tìm giá trị của $m$ biết: $m \times 25 \times 4 = 8400$.
\choice
{$82$}
{$86$}
{$88$}
{\True $84$}
\loigiai{
$m \times 100 = 8400 \Rightarrow m = 84$.\\
\textbf{Đáp án của câu hỏi là D}
}
\end{ex}

\begin{ex} %[4D8C40]
Tính tích: $A = 2 \times 4 \times 5 \times 25 \times 8 \times 125$.
\choice
{$1\,000\,000$}
{$10\,000\,000$}
{\True $1\,000\,000$}
{$100\,000$}
\loigiai{
$(2 \times 5) \times (4 \times 25) \times (8 \times 125) = 10 \times 100 \times 1000 = 1\,000\,000$.\\
\textbf{Đáp án của câu hỏi là C}
}
\end{ex}

\textbf{Phần 2. Trắc nghiệm đúng sai}
\setcounterpageref{ex}{0}

\begin{ex} %[4D8H40]
Xét tính đúng sai về tính chất phép nhân:
\choiceTF
{\True $a \times b = b \times a$.}
{$(a \times b) \times c = a \times b + a \times c$.}
{\True $(a \times b) \times c = a \times (b \times c)$.}
{\True $25 \times 18 \times 4 = (25 \times 4) \times 18 = 1800$.}
\loigiai{
\begin{itemize}
    \item Giao hoán: Đúng.
    \item Đây không phải tính chất kết hợp. Sai.
    \item Kết hợp: Đúng.
    \item Tính thuận tiện đúng.
\end{itemize}
\begin{itemize}
    \item a) \textbf{Đúng}.
    \item b) \textbf{Sai}.
    \item c) \textbf{Đúng}.
    \item d) \textbf{Đúng}.
\end{itemize}
}
\end{ex}

\begin{ex} %[4D8V40]
Cho hai biểu thức: $M = 125 \times 36 \times 8$ và $N = 25 \times 36 \times 40$. Xét tính đúng sai:
\choiceTF
{\True $M = 36\,000$.}
{\True $N = 36\,000$.}
{Giá trị biểu thức $M$ lớn hơn giá trị biểu thức $N$.}
{\True $M = N$.}
\loigiai{
\begin{itemize}
    \item $M = (125 \times 8) \times 36 = 1000 \times 36 = 36\,000$. Đúng.
    \item $N = (25 \times 40) \times 36 = 1000 \times 36 = 36\,000$. Đúng.
    \item $M = N$ nên $M > N$ là Sai.
    \item $M = N$. Đúng.
\end{itemize}
\begin{itemize}
    \item a) \textbf{Đúng}.
    \item b) \textbf{Đúng}.
    \item c) \textbf{Sai}.
    \item d) \textbf{Đúng}.
\end{itemize}
}
\end{ex}

\textbf{Phần 3. Trả lời ngắn}
\setcounterpageref{ex}{0}

\begin{ex} %[4D8H40]
Tính thuận tiện: $25 \times 49 \times 4$.
\shortans{4900}
\loigiai{
\textbf{Đáp số:} $4900$\\
$(25 \times 4) \times 49 = 4900$.
}
\end{ex}

\begin{ex} %[4D8H40]
Tính thuận tiện: $125 \times 73 \times 8$.
\shortans{73000}
\loigiai{
\textbf{Đáp số:} $73000$\\
$(125 \times 8) \times 73 = 73\,000$.
}
\end{ex}

\begin{ex} %[4D8V40]
Tính nhanh: $50 \times 68 \times 2$.
\shortans{6800}
\loigiai{
\textbf{Đáp số:} $6800$\\
$(50 \times 2) \times 68 = 6800$.
}
\end{ex}

\begin{ex} %[4D8V40]
Tìm $x$ biết: $x : (25 \times 4) = 45$.
\shortans{4500}
\loigiai{
\textbf{Đáp số:} $4500$\\
$x : 100 = 45 \Rightarrow x = 4500$.
}
\end{ex}

\begin{ex} %[4D8V40]
Một đoàn xe có 5 xe, mỗi xe chở 20 bao gạo, mỗi bao gạo nặng $50\text{ kg}$. Tính tổng khối lượng gạo đoàn xe chở theo tấn.
\shortans{5}
\loigiai{
\textbf{Đáp số:} $5$\\
$5 \times 20 \times 50 = 5000\text{ kg} = 5$ tấn.
}
\end{ex}

\begin{ex} %[4D8C40]
Tính tích: $2 \times 5 \times 4 \times 25 \times 8 \times 125$.
\shortans{1000000}
\loigiai{
\textbf{Đáp số:} $1000000$\\
$10 \times 100 \times 1000 = 1\,000\,000$.
}
\end{ex}

\newpage

\TaoPhieuKiemTra{Chủ đề 8 -- Phép nhân và phép chia}{41}{Nhân, chia với 10, 100, 1000...}{https://loptoanthaylinh.vn/thi-online/toan-4/bai-41}
\textbf{Phần 1. Trắc nghiệm lựa chọn}
\setcounterpageref{ex}{0}

\begin{ex} %[4D8N41]
Khi nhân một số tự nhiên với 10, ta chỉ việc
\choice
{Thêm hai chữ số 0 vào bên phải số đó}
{\True Thêm một chữ số 0 vào bên phải số đó}
{Thêm ba chữ số 0 vào bên phải số đó}
{Nhân số đó với 1}
\loigiai{
Nhân với 10 viết thêm một chữ số 0 vào bên phải.\\
\textbf{Đáp án của câu hỏi là B}
}
\end{ex}

\begin{ex} %[4D8N41]
Khi chia một số tròn trăm cho 100, ta chỉ việc
\choice
{Bớt một chữ số 0 ở bên phải số đó}
{Bớt ba chữ số 0 ở bên phải số đó}
{\True Bớt hai chữ số 0 ở bên phải số đó}
{Không thay đổi}
\loigiai{
Chia cho 100 bớt đi hai chữ số 0 ở bên phải.\\
\textbf{Đáp án của câu hỏi là C}
}
\end{ex}

\begin{ex} %[4D8H41]
Kết quả của phép nhân $345 \times 100$ là
\choice
{$3450$}
{$345\,000$}
{$3450$}
{\True $34\,500$}
\loigiai{
$345 \times 100 = 34\,500$.\\
\textbf{Đáp án của câu hỏi là D}
}
\end{ex}

\begin{ex} %[4D8H41]
Kết quả của phép chia $280\,000 : 1000$ là
\choice
{\True $280$}
{$28$}
{$2800$}
{$28\,000$}
\loigiai{
$280\,000 : 1000 = 280$.\\
\textbf{Đáp án của câu hỏi là A}
}
\end{ex}

\begin{ex} %[4D8H41]
Điền số thích hợp vào chỗ chấm: $450 \times \dots = 45\,000$.
\choice
{$10$}
{$1000$}
{$10\,000$}
{\True $100$}
\loigiai{
$45\,000 : 450 = 100$.\\
\textbf{Đáp án của câu hỏi là D}
}
\end{ex}

\begin{ex} %[4D8H41]
Điền số thích hợp vào chỗ chấm: $670\,000 : \dots = 670$.
\choice
{\True $1000$}
{$100$}
{$10$}
{$10\,000$}
\loigiai{
$670\,000 : 670 = 1000$.\\
\textbf{Đáp án của câu hỏi là A}
}
\end{ex}

\begin{ex} %[4D8V41]
Một bao gạo nặng $50\text{ kg}$. Hỏi $100$ bao gạo như thế nặng bao nhiêu tấn?
\choice
{$0{,}5\text{ tấn}$}
{$50\text{ tấn}$}
{\True $5\text{ tấn}$}
{$500\text{ tấn}$}
\loigiai{
$50 \times 100 = 5000\text{ kg} = 5\text{ tấn}$.\\
\textbf{Đáp án của câu hỏi là C}
}
\end{ex}

\begin{ex} %[4D8V41]
Một xí nghiệp đóng $150\,000$ gói kẹo vào các thùng, mỗi thùng $100$ gói. Số thùng kẹo đóng được là
\choice
{$150$ thùng}
{\True $1500$ thùng}
{$15\,000$ thùng}
{$15$ thùng}
\loigiai{
$150\,000 : 100 = 1500$ thùng.\\
\textbf{Đáp án của câu hỏi là B}
}
\end{ex}

\begin{ex} %[4D8V41]
Tính giá trị biểu thức: $125 \times 100 + 450 \times 10$.
\choice
{$16\,000$}
{$18\,000$}
{$17\,500$}
{\True $17\,000$}
\loigiai{
$12\,500 + 4500 = 17\,000$.\\
\textbf{Đáp án của câu hỏi là D}
}
\end{ex}

\begin{ex} %[4D8V41]
Tìm $x$ biết: $x \times 100 = 24\,000 : 10$.
\choice
{$240$}
{\True $24$}
{$2400$}
{$2{,}4$}
\loigiai{
$x \times 100 = 2400 \Rightarrow x = 24$.\\
\textbf{Đáp án của câu hỏi là B}
}
\end{ex}

\begin{ex} %[4D8V41]
Một mảnh đất hình chữ nhật có chiều dài $40\text{ m}$, chiều rộng $25\text{ m}$. Diện tích mảnh đất theo mét vuông là
\choice
{\True $1000\text{ m}^2$}
{$100\text{ m}^2$}
{$10\,000\text{ m}^2$}
{$65\text{ m}^2$}
\loigiai{
$40 \times 25 = 1000\text{ m}^2$.\\
\textbf{Đáp án của câu hỏi là A}
}
\end{ex}

\begin{ex} %[4D8C41]
Nếu gấp một số lên 100 lần rồi giảm đi 10 lần thì được số $450$. Số ban đầu là
\choice
{$4500$}
{$450$}
{\True $45$}
{$4{,}5$}
\loigiai{
Gấp 100 lần rồi giảm 10 lần tức là gấp 10 lần. Số ban đầu là $450 : 10 = 45$.\\
\textbf{Đáp án của câu hỏi là C}
}
\end{ex}

\textbf{Phần 2. Trắc nghiệm đúng sai}
\setcounterpageref{ex}{0}

\begin{ex} %[4D8H41]
Xét tính đúng sai của các phép nhân, chia:
\choiceTF
{\True $78 \times 1000 = 78\,000$.}
{$450\,000 : 100 = 450$.}
{\True $450\,000 : 100 = 4500$.}
{\True $120 \times 10 = 1200$.}
\loigiai{
\begin{itemize}
    \item $78 \times 1000 = 78\,000$. Đúng.
    \item $450\,000 : 100 = 4500 \neq 450$. Ý 2 Sai, Ý 3 Đúng.
    \item $120 \times 10 = 1200$. Đúng.
\end{itemize}
\begin{itemize}
    \item a) \textbf{Đúng}.
    \item b) \textbf{Sai}.
    \item c) \textbf{Đúng}.
    \item d) \textbf{Đúng}.
\end{itemize}
}
\end{ex}

\begin{ex} %[4D8V41]
Một trang trại nuôi bò sữa mỗi ngày vắt được $1200\text{ l}$ sữa. Sữa được đóng vào các chai $1\text{ l}$ rồi đóng vào các thùng, mỗi thùng $10$ chai. Xét tính đúng sai:
\choiceTF
{\True Mỗi ngày trang trại đóng được $1200$ chai sữa.}
{\True Mỗi ngày trang trại đóng được $120$ thùng sữa.}
{Trong 10 ngày trang trại đóng được $1000$ thùng sữa.}
{\True Trong 10 ngày trang trại đóng được $1200$ thùng sữa.}
\loigiai{
\begin{itemize}
    \item 1200 chai: Đúng.
    \item $1200 : 10 = 120$ thùng: Đúng.
    \item 10 ngày: $120 \times 10 = 1200$ thùng $\neq 1000$ thùng. Ý 3 Sai, Ý 4 Đúng.
\end{itemize}
\begin{itemize}
    \item a) \textbf{Đúng}.
    \item b) \textbf{Đúng}.
    \item c) \textbf{Sai}.
    \item d) \textbf{Đúng}.
\end{itemize}
}
\end{ex}

\textbf{Phần 3. Trả lời ngắn}
\setcounterpageref{ex}{0}

\begin{ex} %[4D8H41]
Tính tích: $560 \times 100$.
\shortans{56000}
\loigiai{
\textbf{Đáp số:} $56000$\\
$56\,000$.
}
\end{ex}

\begin{ex} %[4D8H41]
Tính thương: $890\,000 : 1000$.
\shortans{890}
\loigiai{
\textbf{Đáp số:} $890$\\
$890$.
}
\end{ex}

\begin{ex} %[4D8V41]
Tìm $x$ biết: $x \times 100 = 75\,000$.
\shortans{750}
\loigiai{
\textbf{Đáp số:} $750$\\
$x = 750$.
}
\end{ex}

\begin{ex} %[4D8V41]
Một trường tiểu học mua 100 quyển sách tham khảo giá $35\,000$ đồng một quyển. Tổng số tiền mua sách là bao nhiêu nghìn đồng (chỉ ghi số nghìn đồng)?
\shortans{3500}
\loigiai{
\textbf{Đáp số:} $3500$\\
$35\,000 \times 100 = 3\,500\,000$ đồng $= 3500$ nghìn đồng.
}
\end{ex}

\begin{ex} %[4D8V41]
Tính giá trị biểu thức: $450 \times 100 - 35\,000$.
\shortans{10000}
\loigiai{
\textbf{Đáp số:} $10000$\\
$45\,000 - 35\,000 = 10\,000$.
}
\end{ex}

\begin{ex} %[4D8C41]
Một số tròn nghìn chia cho 1000 được số bé nhất có 3 chữ số khác nhau. Tìm số tròn nghìn đó.
\shortans{102000}
\loigiai{
\textbf{Đáp số:} $102000$\\
Số bé nhất có 3 chữ số khác nhau là 102. Số cần tìm: $102 \times 1000 = 102\,000$.
}
\end{ex}

\newpage

\TaoPhieuKiemTra{Chủ đề 8 -- Phép nhân và phép chia}{42}{Tính chất phân phối của phép nhân đối với phép cộng}{https://loptoanthaylinh.vn/thi-online/toan-4/bai-42}
\textbf{Phần 1. Trắc nghiệm lựa chọn}
\setcounterpageref{ex}{0}

\begin{ex} %[4D8N42]
Tính chất nhân một số với một tổng được viết dưới dạng công thức là
\choice
{\True $a \times (b + c) = a \times b + a \times c$}
{$a \times (b + c) = a \times b + c$}
{$a \times (b + c) = a + b \times c$}
{$a \times (b + c) = a \times b \times c$}
\loigiai{
$a \times (b + c) = a \times b + a \times c$.\\
\textbf{Đáp án của câu hỏi là A}
}
\end{ex}

\begin{ex} %[4D8N42]
Tính chất nhân một số với một hiệu được viết dưới dạng công thức là
\choice
{$a \times (b - c) = a \times b - c$}
{$a \times (b - c) = a - b \times c$}
{$a \times (b - c) = a \times b : c$}
{\True $a \times (b - c) = a \times b - a \times c$}
\loigiai{
$a \times (b - c) = a \times b - a \times c$.\\
\textbf{Đáp án của câu hỏi là D}
}
\end{ex}

\begin{ex} %[4D8H42]
Tính thuận tiện: $35 \times 18 + 35 \times 82$.
\choice
{$3400$}
{\True $3500$}
{$3600$}
{$3700$}
\loigiai{
$35 \times (18 + 82) = 35 \times 100 = 3500$.\\
\textbf{Đáp án của câu hỏi là B}
}
\end{ex}

\begin{ex} %[4D8H42]
Tính thuận tiện: $48 \times 115 - 48 \times 15$.
\choice
{$4600$}
{$5000$}
{\True $4800$}
{$4700$}
\loigiai{
$48 \times (115 - 15) = 48 \times 100 = 4800$.\\
\textbf{Đáp án của câu hỏi là C}
}
\end{ex}

\begin{ex} %[4D8H42]
Biểu thức $125 \times 99 + 125$ có giá trị bằng
\choice
{\True $12\,500$}
{$12\,400$}
{$12\,600$}
{$13\,000$}
\loigiai{
$125 \times (99 + 1) = 125 \times 100 = 12\,500$.\\
\textbf{Đáp án của câu hỏi là A}
}
\end{ex}

\begin{ex} %[4D8H42]
Tính nhanh: $27 \times 11$.
\choice
{$287$}
{$307$}
{\True $297$}
{$277$}
\loigiai{
$27 \times (10 + 1) = 270 + 27 = 297$.\\
\textbf{Đáp án của câu hỏi là C}
}
\end{ex}

\begin{ex} %[4D8V42]
Tính nhanh giá trị biểu thức: $135 \times 45 + 135 \times 55$.
\choice
{$13\,400$}
{$13\,600$}
{$14\,000$}
{\True $13\,500$}
\loigiai{
$135 \times (45 + 55) = 135 \times 100 = 13\,500$.\\
\textbf{Đáp án của câu hỏi là D}
}
\end{ex}

\begin{ex} %[4D8V42]
Một cửa hàng ngày thứ nhất bán được 25 bao gạo, ngày thứ hai bán được 35 bao gạo. Mỗi bao gạo nặng $50\text{ kg}$. Tính tổng số ki-lô-gam gạo bán trong 2 ngày bằng cách thuận tiện nhất.
\choice
{$2800\text{ kg}$}
{\True $3000\text{ kg}$}
{$3200\text{ kg}$}
{$3500\text{ kg}$}
\loigiai{
$(25 + 35) \times 50 = 60 \times 50 = 3000\text{ kg}$.\\
\textbf{Đáp án của câu hỏi là B}
}
\end{ex}

\begin{ex} %[4D8V42]
Tính nhanh: $64 \times 101 - 64$.
\choice
{$6300$}
{$6500$}
{\True $6400$}
{$6464$}
\loigiai{
$64 \times (101 - 1) = 64 \times 100 = 6400$.\\
\textbf{Đáp án của câu hỏi là C}
}
\end{ex}

\begin{ex} %[4D8V42]
Một mảnh đất hình chữ nhật có chiều dài $28\text{ m}$, chiều rộng $12\text{ m}$. Nếu tăng chiều dài thêm $5\text{ m}$ thì diện tích mảnh đất tăng thêm bao nhiêu mét vuông?
\choice
{$50\text{ m}^2$}
{$70\text{ m}^2$}
{$80\text{ m}^2$}
{\True $60\text{ m}^2$}
\loigiai{
Diện tích tăng: $5 \times 12 = 60\text{ m}^2$.\\
\textbf{Đáp án của câu hỏi là D}
}
\end{ex}

\begin{ex} %[4D8V42]
Tính giá trị biểu thức: $45 \times 12 + 45 \times 18 + 45 \times 70$.
\choice
{\True $4500$}
{$4400$}
{$4600$}
{$5000$}
\loigiai{
$45 \times (12 + 18 + 70) = 45 \times 100 = 4500$.\\
\textbf{Đáp án của câu hỏi là A}
}
\end{ex}

\begin{ex} %[4D8C42]
Tính nhanh giá trị biểu thức: $A = 99 \times 45 + 45$.
\choice
{$4400$}
{\True $4500$}
{$4600$}
{$4455$}
\loigiai{
$A = 45 \times (99 + 1) = 45 \times 100 = 4500$.\\
\textbf{Đáp án của câu hỏi là B}
}
\end{ex}

\textbf{Phần 2. Trắc nghiệm đúng sai}
\setcounterpageref{ex}{0}

\begin{ex} %[4D8H42]
Xét tính đúng sai của các công thức phân phối:
\choiceTF
{\True $a \times (b + c) = a \times b + a \times c$.}
{\True $a \times (b - c) = a \times b - a \times c$.}
{$(a + b) \times c = a \times c + b$.}
{\True $(a + b) \times c = a \times c + b \times c$.}
\loigiai{
\begin{itemize}
    \item $a \times (b + c) = a \times b + a \times c$. Đúng.
    \item $a \times (b - c) = a \times b - a \times c$. Đúng.
    \item Phải nhân cả $b$ với $c$. Ý 3 Sai, Ý 4 Đúng.
\end{itemize}
\begin{itemize}
    \item a) \textbf{Đúng}.
    \item b) \textbf{Đúng}.
    \item c) \textbf{Sai}.
    \item d) \textbf{Đúng}.
\end{itemize}
}
\end{ex}

\begin{ex} %[4D8V42]
Hai tổ may khẩu trang: Tổ 1 may được 140 hộp, tổ 2 may được 160 hộp. Mỗi hộp chứa 50 chiếc khẩu trang. Xét tính đúng sai:
\choiceTF
{\True Cả hai tổ may được tất cả 300 hộp khẩu trang.}
{\True Tổ 2 may nhiều hơn tổ 1 là 1000 chiếc khẩu trang.}
{Tổng số khẩu trang cả hai tổ may được là 12000 chiếc.}
{\True Tổng số khẩu trang cả hai tổ may được là 15000 chiếc.}
\loigiai{
\begin{itemize}
    \item $140 + 160 = 300$ hộp: Đúng.
    \item $(160 - 140) \times 50 = 20 \times 50 = 1000$ chiếc: Đúng.
    \item Tổng: $300 \times 50 = 15\,000$ chiếc $\neq 12\,000$ chiếc. Ý 3 Sai, Ý 4 Đúng.
\end{itemize}
\begin{itemize}
    \item a) \textbf{Đúng}.
    \item b) \textbf{Đúng}.
    \item c) \textbf{Sai}.
    \item d) \textbf{Đúng}.
\end{itemize}
}
\end{ex}

\textbf{Phần 3. Trả lời ngắn}
\setcounterpageref{ex}{0}

\begin{ex} %[4D8H42]
Tính thuận tiện: $45 \times 38 + 45 \times 62$.
\shortans{4500}
\loigiai{
\textbf{Đáp số:} $4500$\\
$45 \times 100 = 4500$.
}
\end{ex}

\begin{ex} %[4D8H42]
Tính thuận tiện: $78 \times 125 - 78 \times 25$.
\shortans{7800}
\loigiai{
\textbf{Đáp số:} $7800$\\
$78 \times 100 = 7800$.
}
\end{ex}

\begin{ex} %[4D8V42]
Tính nhanh: $56 \times 99 + 56$.
\shortans{5600}
\loigiai{
\textbf{Đáp số:} $5600$\\
$56 \times 100 = 5600$.
}
\end{ex}

\begin{ex} %[4D8V42]
Tính nhanh: $84 \times 101 - 84$.
\shortans{8400}
\loigiai{
\textbf{Đáp số:} $8400$\\
$84 \times 100 = 8400$.
}
\end{ex}

\begin{ex} %[4D8V42]
Một lớp học mua 15 quyển vở Toán và 15 quyển vở Tiếng Việt cùng giá $8000$ đồng một quyển. Tính tổng số tiền mua vở theo nghìn đồng (chỉ ghi số nghìn đồng).
\shortans{240}
\loigiai{
\textbf{Đáp số:} $240$\\
$15 \times (8 + 8) = 240$ nghìn đồng.
}
\end{ex}

\begin{ex} %[4D8C42]
Tính nhanh: $12 \times 35 + 12 \times 65 + 24 \times 50$.
\shortans{2400}
\loigiai{
\textbf{Đáp số:} $2400$\\
$12 \times (35 + 65) + 12 \times 100 = 12 \times 100 + 12 \times 100 = 2400$.
}
\end{ex}

\newpage

\TaoPhieuKiemTra{Chủ đề 8 -- Phép nhân và phép chia}{43}{Nhân với số có hai chữ số}{https://loptoanthaylinh.vn/thi-online/toan-4/bai-43}
\textbf{Phần 1. Trắc nghiệm lựa chọn}
\setcounterpageref{ex}{0}

\begin{ex} %[4D8N43]
Kết quả của phép nhân $34 \times 12$ là
\choice
{$406$}
{$418$}
{$398$}
{\True $408$}
\loigiai{
$34 \times 12 = 408$.\\
\textbf{Đáp án của câu hỏi là D}
}
\end{ex}

\begin{ex} %[4D8N43]
Trong phép nhân với số có hai chữ số, tích riêng thứ hai được viết
\choice
{Thẳng cột với tích riêng thứ nhất}
{Lùi sang bên phải một cột}
{\True Lùi sang bên trái một cột so với tích riêng thứ nhất}
{Lùi sang bên trái hai cột}
\loigiai{
Tích riêng thứ hai lùi sang bên trái một cột.\\
\textbf{Đáp án của câu hỏi là C}
}
\end{ex}

\begin{ex} %[4D8H43]
Kết quả của phép nhân $125 \times 24$ là
\choice
{\True $3000$}
{$2800$}
{$3200$}
{$3500$}
\loigiai{
$125 \times 24 = 3000$.\\
\textbf{Đáp án của câu hỏi là A}
}
\end{ex}

\begin{ex} %[4D8H43]
Kết quả của phép tính $245 \times 36$ là
\choice
{$8720$}
{$8920$}
{\True $8820$}
{$8620$}
\loigiai{
$245 \times 36 = 8820$.\\
\textbf{Đáp án của câu hỏi là C}
}
\end{ex}

\begin{ex} %[4D8H43]
Một quyển sách dày 148 trang. Hỏi 25 quyển sách như thế có tất cả bao nhiêu trang?
\choice
{$3600$ trang}
{\True $3700$ trang}
{$3800$ trang}
{$3500$ trang}
\loigiai{
$148 \times 25 = 3700$ trang.\\
\textbf{Đáp án của câu hỏi là B}
}
\end{ex}

\begin{ex} %[4D8H43]
Một hình chữ nhật có chiều dài $45\text{ m}$, chiều rộng $28\text{ m}$. Diện tích hình chữ nhật là
\choice
{$1240\text{ m}^2$}
{$1280\text{ m}^2$}
{\True $1260\text{ m}^2$}
{$1300\text{ m}^2$}
\loigiai{
$45 \times 28 = 1260\text{ m}^2$.\\
\textbf{Đáp án của câu hỏi là C}
}
\end{ex}

\begin{ex} %[4D8V43]
Một người thợ làm việc trong 24 ngày, mỗi ngày làm được 15 sản phẩm. Tổng số sản phẩm người đó làm được là
\choice
{$340$ sản phẩm}
{\True $360$ sản phẩm}
{$380$ sản phẩm}
{$400$ sản phẩm}
\loigiai{
$24 \times 15 = 360$ sản phẩm.\\
\textbf{Đáp án của câu hỏi là B}
}
\end{ex}

\begin{ex} %[4D8V43]
Một đoàn tàu hỏa có 12 toa xe, mỗi toa xe chở $35$ tấn hàng. Tổng khối lượng hàng đoàn tàu chở là
\choice
{$400\text{ tấn}$}
{$440\text{ tấn}$}
{$450\text{ tấn}$}
{\True $420\text{ tấn}$}
\loigiai{
$12 \times 35 = 420$ tấn.\\
\textbf{Đáp án của câu hỏi là D}
}
\end{ex}

\begin{ex} %[4D8V43]
Tính giá trị biểu thức: $245 \times 18 + 120 \times 25$.
\choice
{\True $7410$}
{$7310$}
{$7510$}
{$7210$}
\loigiai{
$4410 + 3000 = 7410$.\\
\textbf{Đáp án của câu hỏi là A}
}
\end{ex}

\begin{ex} %[4D8V43]
Một cửa hàng bán được 35 chiếc xe đạp với giá $1\,250\,000$ đồng một chiếc. Tổng số tiền bán xe đạp là
\choice
{$42\,750\,000$ đồng}
{\True $43\,750\,000$ đồng}
{$44\,750\,000$ đồng}
{$45\,000\,000$ đồng}
\loigiai{
$1\,250\,000 \times 35 = 43\,750\,000$ đồng.\\
\textbf{Đáp án của câu hỏi là B}
}
\end{ex}

\begin{ex} %[4D8V43]
Tìm $x$ biết: $x : 24 = 135$.
\choice
{\True $3240$}
{$3140$}
{$3340$}
{$3440$}
\loigiai{
$x = 135 \times 24 = 3240$.\\
\textbf{Đáp án của câu hỏi là A}
}
\end{ex}

\begin{ex} %[4D8C43]
Khi nhân một số tự nhiên với 35, một bạn học sinh đã đặt các tích riêng thẳng cột như trong phép cộng nên được kết quả là $680$. Tích đúng của phép nhân đó là
\choice
{$2875$}
{$3075$}
{$3175$}
{\True $2975$}
\loigiai{
Đặt tích riêng thẳng cột tức là nhân số đó với $3 + 5 = 8$.\\
Số đó là: $680 : 8 = 85$.\\
Tích đúng: $85 \times 35 = 2975$.\\
\textbf{Đáp án của câu hỏi là D}
}
\end{ex}

\textbf{Phần 2. Trắc nghiệm đúng sai}
\setcounterpageref{ex}{0}

\begin{ex} %[4D8H43]
Xét tính đúng sai của các phép tính:
\choiceTF
{\True $25 \times 14 = 350$.}
{$42 \times 23 = 946$.}
{\True $42 \times 23 = 966$.}
{\True Tích riêng thứ hai lùi sang trái một cột so với tích riêng thứ nhất.}
\loigiai{
\begin{itemize}
    \item $25 \times 14 = 350$. Đúng.
    \item $42 \times 23 = 966 \neq 946$. Ý 2 Sai, Ý 3 Đúng.
    \item Quy tắc đặt tính đúng.
\end{itemize}
\begin{itemize}
    \item a) \textbf{Đúng}.
    \item b) \textbf{Sai}.
    \item c) \textbf{Đúng}.
    \item d) \textbf{Đúng}.
\end{itemize}
}
\end{ex}

\begin{ex} %[4D8V43]
Một xí nghiệp may có 25 công nhân, mỗi ngày mỗi công nhân may được 12 chiếc áo. Mỗi chiếc áo bán với giá $150\,000$ đồng. Xét tính đúng sai:
\choiceTF
{\True Mỗi ngày xí nghiệp may được 300 chiếc áo.}
{\True Trong 10 ngày xí nghiệp may được 3000 chiếc áo.}
{Doanh thu mỗi ngày của xí nghiệp là 40 triệu đồng.}
{\True Doanh thu mỗi ngày của xí nghiệp là 45 triệu đồng.}
\loigiai{
\begin{itemize}
    \item $25 \times 12 = 300$ áo: Đúng.
    \item 10 ngày: $300 \times 10 = 3000$ áo: Đúng.
    \item Doanh thu: $300 \times 150\,000 = 45\,000\,000$ đồng $\neq 40$ triệu. Ý 3 Sai, Ý 4 Đúng.
\end{itemize}
\begin{itemize}
    \item a) \textbf{Đúng}.
    \item b) \textbf{Đúng}.
    \item c) \textbf{Sai}.
    \item d) \textbf{Đúng}.
\end{itemize}
}
\end{ex}

\textbf{Phần 3. Trả lời ngắn}
\setcounterpageref{ex}{0}

\begin{ex} %[4D8H43]
Tính tích: $36 \times 25$.
\shortans{900}
\loigiai{
\textbf{Đáp số:} $900$\\
$900$.
}
\end{ex}

\begin{ex} %[4D8H43]
Tính tích: $124 \times 15$.
\shortans{1860}
\loigiai{
\textbf{Đáp số:} $1860$\\
$1860$.
}
\end{ex}

\begin{ex} %[4D8V43]
Một mảnh vườn hình chữ nhật dài $35\text{ m}$, rộng $18\text{ m}$. Tính diện tích mảnh vườn theo mét vuông.
\shortans{630}
\loigiai{
\textbf{Đáp số:} $630$\\
$35 \times 18 = 630\text{ m}^2$.
}
\end{ex}

\begin{ex} %[4D8V43]
Một thùng sách có 45 cuốn, mỗi cuốn giá $24\,000$ đồng. Tính tổng giá trị thùng sách theo nghìn đồng (chỉ ghi số nghìn đồng).
\shortans{1080}
\loigiai{
\textbf{Đáp số:} $1080$\\
$45 \times 24 = 1080$ nghìn đồng.
}
\end{ex}

\begin{ex} %[4D8V43]
Tìm $x$ biết: $x : 16 = 250$.
\shortans{4000}
\loigiai{
\textbf{Đáp số:} $4000$\\
$x = 250 \times 16 = 4000$.
}
\end{ex}

\begin{ex} %[4D8C43]
Khi nhân một số với 24, bạn Bình quên lùi tích riêng thứ hai nên kết quả bị giảm đi $432$ đơn vị. Tìm số đem nhân với 24.
\shortans{24}
\loigiai{
\textbf{Đáp số:} $24$\\
Quên lùi tích riêng tức là đã nhân với $2 + 4 = 6$.\\
Tích giảm đi: $24 - 6 = 18$ lần số đó.\\
Số đó là: $432 : 18 = 24$.
}
\end{ex}

\newpage

\TaoPhieuKiemTra{Chủ đề 8 -- Phép nhân và phép chia}{44}{Chia cho số có hai chữ số}{https://loptoanthaylinh.vn/thi-online/toan-4/bai-44}
\textbf{Phần 1. Trắc nghiệm lựa chọn}
\setcounterpageref{ex}{0}

\begin{ex} %[4D8N44]
Kết quả của phép chia $288 : 24$ là
\choice
{$11$}
{\True $12$}
{$13$}
{$14$}
\loigiai{
$288 : 24 = 12$.\\
\textbf{Đáp án của câu hỏi là B}
}
\end{ex}

\begin{ex} %[4D8N44]
Số dư trong phép chia $150 : 20$ là
\choice
{$0$}
{$5$}
{\True $10$}
{$15$}
\loigiai{
$150 : 20 = 7$ dư 10.\\
\textbf{Đáp án của câu hỏi là C}
}
\end{ex}

\begin{ex} %[4D8H44]
Kết quả của phép chia $1440 : 36$ là
\choice
{$30$}
{$50$}
{$45$}
{\True $40$}
\loigiai{
$1440 : 36 = 40$.\\
\textbf{Đáp án của câu hỏi là D}
}
\end{ex}

\begin{ex} %[4D8H44]
Tìm $x$ biết: $x \times 25 = 1250$.
\choice
{$40$}
{\True $50$}
{$60$}
{$55$}
\loigiai{
$x = 1250 : 25 = 50$.\\
\textbf{Đáp án của câu hỏi là B}
}
\end{ex}

\begin{ex} %[4D8H44]
Một hình chữ nhật có diện tích $360\text{ m}^2$, chiều rộng $15\text{ m}$. Chiều dài của hình chữ nhật là
\choice
{\True $24\text{ m}$}
{$22\text{ m}$}
{$26\text{ m}$}
{$28\text{ m}$}
\loigiai{
$360 : 15 = 24\text{ m}$.\\
\textbf{Đáp án của câu hỏi là A}
}
\end{ex}

\begin{ex} %[4D8H44]
Kết quả của phép chia $7200 : 45$ là
\choice
{$150$}
{\True $160$}
{$170$}
{$180$}
\loigiai{
$7200 : 45 = 160$.\\
\textbf{Đáp án của câu hỏi là B}
}
\end{ex}

\begin{ex} %[4D8V44]
Người ta xếp $1800$ quyển vở vào các thùng, mỗi thùng chứa được $40$ quyển. Số thùng cần dùng là
\choice
{\True $45$ thùng}
{$40$ thùng}
{$50$ thùng}
{$55$ thùng}
\loigiai{
$1800 : 40 = 45$ thùng.\\
\textbf{Đáp án của câu hỏi là A}
}
\end{ex}

\begin{ex} %[4D8V44]
Một đoàn khách du lịch gồm $265$ người. Người ta thuê các xe ô tô loại $29$ chỗ ngồi để đi tham quan. Cần thuê ít nhất bao nhiêu xe như thế?
\choice
{$9$ xe}
{$11$ xe}
{$12$ xe}
{\True $10$ xe}
\loigiai{
$265 : 29 = 9$ dư 4 người. Cần ít nhất $9 + 1 = 10$ xe.\\
\textbf{Đáp án của câu hỏi là D}
}
\end{ex}

\begin{ex} %[4D8V44]
Tính giá trị biểu thức: $4500 : 15 + 3600 : 18$.
\choice
{\True $500$}
{$400$}
{$600$}
{$550$}
\loigiai{
$300 + 200 = 500$.\\
\textbf{Đáp án của câu hỏi là A}
}
\end{ex}

\begin{ex} %[4D8V44]
Một xí nghiệp có $3600\text{ kg}$ đường đóng vào các bao loại $50\text{ kg}$. Người ta đã bán đi $40$ bao đường đó. Số bao đường còn lại là
\choice
{$30$ bao}
{$35$ bao}
{\True $32$ bao}
{$40$ bao}
\loigiai{
Tổng số bao: $3600 : 50 = 72$ bao.\\
Còn lại: $72 - 40 = 32$ bao.\\
\textbf{Đáp án của câu hỏi là C}
}
\end{ex}

\begin{ex} %[4D8V44]
Một mảnh đất hình chữ nhật có chu vi $140\text{ m}$, chiều rộng $25\text{ m}$. Diện tích mảnh đất là
\choice
{$1000\text{ m}^2$}
{$1200\text{ m}^2$}
{$1250\text{ m}^2$}
{\True $1125\text{ m}^2$}
\loigiai{
Nửa chu vi: $70\text{ m}$. Chiều dài: $70 - 25 = 45\text{ m}$.\\
Diện tích: $45 \times 25 = 1125\text{ m}^2$.\\
\textbf{Đáp án của câu hỏi là D}
}
\end{ex}

\begin{ex} %[4D8C44]
Tìm số tự nhiên $x$ lớn nhất có hai chữ số sao cho $2500 : x > 50$.
\choice
{$48$}
{$50$}
{\True $49$}
{$51$}
\loigiai{
$2500 : 50 = 50 \Rightarrow x < 50$. Số tự nhiên $x$ lớn nhất có hai chữ số là 49.\\
\textbf{Đáp án của câu hỏi là C}
}
\end{ex}

\textbf{Phần 2. Trắc nghiệm đúng sai}
\setcounterpageref{ex}{0}

\begin{ex} %[4D8H44]
Xét tính đúng sai của các phép chia cho số có hai chữ số:
\choiceTF
{\True $480 : 12 = 40$.}
{$750 : 25 = 25$.}
{\True $750 : 25 = 30$.}
{\True Trong phép chia có dư, số dư luôn nhỏ hơn số chia.}
\loigiai{
\begin{itemize}
    \item $480 : 12 = 40$. Đúng.
    \item $750 : 25 = 30 \neq 25$. Ý 2 Sai, Ý 3 Đúng.
    \item Số dư $<$ số chia: Đúng.
\end{itemize}
\begin{itemize}
    \item a) \textbf{Đúng}.
    \item b) \textbf{Sai}.
    \item c) \textbf{Đúng}.
    \item d) \textbf{Đúng}.
\end{itemize}
}
\end{ex}

\begin{ex} %[4D8V44]
Một cửa hàng nhập về $1200\text{ kg}$ gạo tẻ và $800\text{ kg}$ gạo nếp đóng vào các bao loại $40\text{ kg}$. Xét tính đúng sai:
\choiceTF
{\True Số bao gạo tẻ là $30$ bao.}
{\True Số bao gạo nếp là $20$ bao.}
{Tổng số bao gạo là $60$ bao.}
{\True Tổng số bao gạo là $50$ bao.}
\loigiai{
\begin{itemize}
    \item Gạo tẻ: $1200 : 40 = 30$ bao. Đúng.
    \item Gạo nếp: $800 : 40 = 20$ bao. Đúng.
    \item Tổng: $30 + 20 = 50$ bao $\neq 60$ bao. Ý 3 Sai, Ý 4 Đúng.
\end{itemize}
\begin{itemize}
    \item a) \textbf{Đúng}.
    \item b) \textbf{Đúng}.
    \item c) \textbf{Sai}.
    \item d) \textbf{Đúng}.
\end{itemize}
}
\end{ex}

\textbf{Phần 3. Trả lời ngắn}
\setcounterpageref{ex}{0}

\begin{ex} %[4D8H44]
Tính thương: $720 : 18$.
\shortans{40}
\loigiai{
\textbf{Đáp số:} $40$\\
$40$.
}
\end{ex}

\begin{ex} %[4D8H44]
Tính thương: $1250 : 25$.
\shortans{50}
\loigiai{
\textbf{Đáp số:} $50$\\
$50$.
}
\end{ex}

\begin{ex} %[4D8V44]
Tìm $x$ biết: $x \times 35 = 1400$.
\shortans{40}
\loigiai{
\textbf{Đáp số:} $40$\\
$x = 1400 : 35 = 40$.
}
\end{ex}

\begin{ex} %[4D8V44]
Một mảnh vườn hình chữ nhật có diện tích $480\text{ m}^2$, chiều dài $24\text{ m}$. Tính chu vi mảnh vườn theo mét.
\shortans{88}
\loigiai{
\textbf{Đáp số:} $88$\\
Rộng: $480 : 24 = 20\text{ m}$. Chu vi: $(24 + 20) \times 2 = 88\text{ m}$.
}
\end{ex}

\begin{ex} %[4D8V44]
Có $960$ quyển sách xếp đều vào 24 ngăn tủ. Hỏi mỗi ngăn tủ có bao nhiêu quyển sách?
\shortans{40}
\loigiai{
\textbf{Đáp số:} $40$\\
$960 : 24 = 40$.
}
\end{ex}

\begin{ex} %[4D8C44]
Tìm số bị chia trong phép chia cho 32 có thương là 45 và số dư là số dư lớn nhất có thể.
\shortans{1471}
\loigiai{
\textbf{Đáp số:} $1471$\\
Số dư lớn nhất: 31. Số bị chia: $45 \times 32 + 31 = 1440 + 31 = 1471$.
}
\end{ex}

\newpage

\TaoPhieuKiemTra{Chủ đề 8 -- Phép nhân và phép chia}{45}{Thực hành ước lượng trong tính toán}{https://loptoanthaylinh.vn/thi-online/toan-4/bai-45}
\textbf{Phần 1. Trắc nghiệm lựa chọn}
\setcounterpageref{ex}{0}

\begin{ex} %[4D8N45]
Ước lượng kết quả phép tính $49 \times 21$ bằng cách làm tròn các thừa số đến hàng chục ta được
\choice
{$800$}
{\True $1000$}
{$1200$}
{$900$}
\loigiai{
$50 \times 20 = 1000$.\\
\textbf{Đáp án của câu hỏi là B}
}
\end{ex}

\begin{ex} %[4D8N45]
Ước lượng kết quả của phép chia $398 : 19$ ta được khoảng
\choice
{\True $20$}
{$10$}
{$30$}
{$40$}
\loigiai{
$400 : 20 = 20$.\\
\textbf{Đáp án của câu hỏi là A}
}
\end{ex}

\begin{ex} %[4D8H45]
Ước lượng số tiền mua 19 quyển vở giá $9800$ đồng mỗi quyển bằng cách làm tròn ta được khoảng
\choice
{$180\,000$ đồng}
{$220\,000$ đồng}
{$150\,000$ đồng}
{\True $200\,000$ đồng}
\loigiai{
$20 \times 10\,000 = 200\,000$ đồng.\\
\textbf{Đáp án của câu hỏi là D}
}
\end{ex}

\begin{ex} %[4D8H45]
Một rạp chiếu phim có 28 hàng ghế, mỗi hàng có 32 ghế. Ước lượng tổng số ghế trong rạp là khoảng
\choice
{\True $900$ ghế}
{$600$ ghế}
{$1200$ ghế}
{$800$ ghế}
\loigiai{
$30 \times 30 = 900$ ghế.\\
\textbf{Đáp án của câu hỏi là A}
}
\end{ex}

\begin{ex} %[4D8H45]
Ước lượng chiều dài sân trường khoảng bao nhiêu mét?
\choice
{$5\text{ m}$}
{$500\text{ m}$}
{\True $50\text{ m}$}
{$5\text{ km}$}
\loigiai{
Chiều dài sân trường khoảng $40 - 60\text{ m}$.\\
\textbf{Đáp án của câu hỏi là C}
}
\end{ex}

\begin{ex} %[4D8H45]
Ước lượng kết quả tổng $4980 + 3020 + 2010$ khoảng
\choice
{$9000$}
{\True $10\,000$}
{$11\,000$}
{$8000$}
\loigiai{
$5000 + 3000 + 2000 = 10\,000$.\\
\textbf{Đáp án của câu hỏi là B}
}
\end{ex}

\begin{ex} %[4D8V45]
Mẹ mang $500\,000$ đồng đi siêu thị. Mẹ dự định mua: $1$ thùng sữa giá $295\,000$ đồng, $2\text{ kg}$ thịt lợn giá $110\,000$ đồng/kg. Mẹ có đủ tiền mua không?
\choice
{Vừa đủ}
{Thừa tiền rất nhiều}
{Không ước lượng được}
{\True Thiếu tiền (cần khoảng $515\,000$ đồng)}
\loigiai{
$295\,000 + 2 \times 110\,000 = 515\,000 > 500\,000$ đồng nên thiếu tiền.\\
\textbf{Đáp án của câu hỏi là D}
}
\end{ex}

\begin{ex} %[4D8V45]
Một người đi xe máy mỗi giờ đi được khoảng $38\text{ km}$. Ước lượng quãng đường người đó đi trong 5 giờ khoảng
\choice
{$150\text{ km}$}
{$250\text{ km}$}
{\True $200\text{ km}$}
{$300\text{ km}$}
\loigiai{
$40 \times 5 = 200\text{ km}$.\\
\textbf{Đáp án của câu hỏi là C}
}
\end{ex}

\begin{ex} %[4D8V45]
Một cửa hàng bán được $48$ chiếc áo thun với giá $195\,000$ đồng một chiếc. Ước lượng doanh thu bán áo khoảng
\choice
{$8\,000\,000$ đồng}
{$12\,000\,000$ đồng}
{$15\,000\,000$ đồng}
{\True $10\,000\,000$ đồng}
\loigiai{
$50 \times 200\,000 = 10\,000\,000$ đồng.\\
\textbf{Đáp án của câu hỏi là D}
}
\end{ex}

\begin{ex} %[4D8V45]
Một khối lớp 4 có 5 lớp, mỗi lớp có khoảng 32 học sinh. Ước lượng tổng số học sinh khối 4 là khoảng
\choice
{$120$ học sinh}
{$200$ học sinh}
{\True $150$ học sinh (hoặc 160 học sinh)}
{$180$ học sinh}
\loigiai{
$5 \times 30 = 150$ học sinh (chính xác là 160).\\
\textbf{Đáp án của câu hỏi là C}
}
\end{ex}

\begin{ex} %[4D8V45]
Một bể nước hình hộp chữ nhật chứa được khoảng $2950\text{ l}$ nước. Khi đã dùng hết một nửa, lượng nước còn lại khoảng
\choice
{\True $1500\text{ l}$}
{$1000\text{ l}$}
{$2000\text{ l}$}
{$2500\text{ l}$}
\loigiai{
$3000 : 2 = 1500\text{ l}$.\\
\textbf{Đáp án của câu hỏi là A}
}
\end{ex}

\begin{ex} %[4D8C45]
Để chuẩn bị nước uống cho chuyến dã ngoại 150 học sinh, biết mỗi học sinh uống khoảng $450\text{ ml}$ nước. Cần chuẩn bị khoảng bao nhiêu bình nước loại $20$ lít?
\choice
{$2$ bình}
{\True $4$ bình (khoảng $67{,}5\text{ l}$)}
{$6$ bình}
{$8$ bình}
\loigiai{
Tổng lượng nước: $150 \times 450 = 67\,500\text{ ml} = 67{,}5\text{ l}$.\\
Số bình 20 lít: $67{,}5 : 20 \approx 3{,}375 \Rightarrow$ cần 4 bình.\\
\textbf{Đáp án của câu hỏi là B}
}
\end{ex}

\textbf{Phần 2. Trắc nghiệm đúng sai}
\setcounterpageref{ex}{0}

\begin{ex} %[4D8H45]
Xét tính hợp lí của các ước lượng:
\choiceTF
{\True Cân nặng của một chiếc xe máy khoảng $100\text{ kg}$.}
{Thời gian ngủ mỗi đêm của học sinh tiểu học khoảng $15$ giờ.}
{\True Thời gian ngủ mỗi đêm của học sinh tiểu học khoảng $8 - 9$ giờ.}
{\True Ước lượng tích $98 \times 49$ khoảng $100 \times 50 = 5000$.}
\loigiai{
\begin{itemize}
    \item Xe máy ~100 kg: Đúng.
    \item Học sinh ngủ 8-9h (15h là sai). Ý 2 Sai, Ý 3 Đúng.
    \item Ước lượng tích đúng.
\end{itemize}
\begin{itemize}
    \item a) \textbf{Đúng}.
    \item b) \textbf{Sai}.
    \item c) \textbf{Đúng}.
    \item d) \textbf{Đúng}.
\end{itemize}
}
\end{ex}

\begin{ex} %[4D8V45]
Bạn An có $100\,000$ đồng đi mua đồ dùng học tập: mua 1 hộp bút giá $38\,000$ đồng, 5 quyển vở giá $9500$ đồng/quyển và 1 cây thước kẻ giá $5000$ đồng. Xét tính đúng sai:
\choiceTF
{\True Tiền mua 5 quyển vở hết $47\,500$ đồng.}
{\True Tổng số tiền mua 3 món đồ là $90\,500$ đồng.}
{Bạn An không đủ tiền để thanh toán.}
{\True Bạn An còn thừa khoảng $9500$ đồng.}
\loigiai{
\begin{itemize}
    \item Vở: $5 \times 9500 = 47\,500$ đồng: Đúng.
    \item Tổng: $38\,000 + 47\,500 + 5000 = 90\,500$ đồng: Đúng.
    \item Đủ tiền vì $90\,500 < 100\,000$. Sai.
    \item Thừa: $100\,000 - 90\,500 = 9500$ đồng. Đúng.
\end{itemize}
\begin{itemize}
    \item a) \textbf{Đúng}.
    \item b) \textbf{Đúng}.
    \item c) \textbf{Sai}.
    \item d) \textbf{Đúng}.
\end{itemize}
}
\end{ex}

\textbf{Phần 3. Trả lời ngắn}
\setcounterpageref{ex}{0}

\begin{ex} %[4D8H45]
Ước lượng tích $51 \times 39$ bằng cách làm tròn đến hàng chục.
\shortans{2000}
\loigiai{
\textbf{Đáp số:} $2000$\\
$50 \times 40 = 2000$.
}
\end{ex}

\begin{ex} %[4D8H45]
Ước lượng thương $595 : 31$ bằng cách làm tròn.
\shortans{20}
\loigiai{
\textbf{Đáp số:} $20$\\
$600 : 30 = 20$.
}
\end{ex}

\begin{ex} %[4D8V45]
Một cửa hàng bán được 29 cái cặp sách giá $190\,000$ đồng một cái. Ước lượng doanh thu theo triệu đồng (làm tròn đến hàng triệu).
\shortans{6}
\loigiai{
\textbf{Đáp số:} $6$\\
$30 \times 200\,000 = 6\,000\,000$ đồng $= 6$ triệu đồng.
}
\end{ex}

\begin{ex} %[4D8V45]
Một xe tải chở 48 thùng hàng, mỗi thùng nặng $21\text{ kg}$. Ước lượng tổng khối lượng hàng theo ki-lô-gam (làm tròn đến hàng trăm).
\shortans{1000}
\loigiai{
\textbf{Đáp số:} $1000$\\
$50 \times 20 = 1000\text{ kg}$.
}
\end{ex}

\begin{ex} %[4D8V45]
Ước lượng tổng $1980 + 2030 + 4990$ (làm tròn đến hàng nghìn).
\shortans{9000}
\loigiai{
\textbf{Đáp số:} $9000$\\
$2000 + 2000 + 5000 = 9000$.
}
\end{ex}

\begin{ex} %[4D8C45]
Một trường học có 980 học sinh, mỗi ngày mỗi học sinh uống khoảng $250\text{ ml}$ sữa tươi. Ước lượng số lít sữa toàn trường tiêu thụ mỗi ngày (làm tròn đến hàng chục lít).
\shortans{250}
\loigiai{
\textbf{Đáp số:} $250$\\
$1000 \times 250\text{ ml} = 250\,000\text{ ml} = 250$ lít.
}
\end{ex}

\newpage

\TaoPhieuKiemTra{Chủ đề 8 -- Phép nhân và phép chia}{46}{Tìm số trung bình cộng}{https://loptoanthaylinh.vn/thi-online/toan-4/bai-46}
\textbf{Phần 1. Trắc nghiệm lựa chọn}
\setcounterpageref{ex}{0}

\begin{ex} %[4D8N46]
Muốn tìm số trung bình cộng của nhiều số, ta
\choice
{\True Tính tổng của các số đó rồi chia tổng đó cho số các số hạng}
{Lấy tổng của các số đó nhân với số các số hạng}
{Lấy số lớn nhất cộng số bé nhất chia đôi}
{Lấy hiệu của các số chia cho số các số hạng}
\loigiai{
Trung bình cộng $=$ Tổng các số $:$ Số các số hạng.\\
\textbf{Đáp án của câu hỏi là A}
}
\end{ex}

\begin{ex} %[4D8N46]
Trung bình cộng của hai số $40$ và $60$ là
\choice
{$45$}
{$55$}
{\True $50$}
{$100$}
\loigiai{
$(40 + 60) : 2 = 50$.\\
\textbf{Đáp án của câu hỏi là C}
}
\end{ex}

\begin{ex} %[4D8H46]
Số trung bình cộng của 3 số: $15; 25; 35$ là
\choice
{$20$}
{\True $25$}
{$30$}
{$75$}
\loigiai{
$(15 + 25 + 35) : 3 = 25$.\\
\textbf{Đáp án của câu hỏi là B}
}
\end{ex}

\begin{ex} %[4D8H46]
Bốn bạn có số bi lần lượt là $12, 14, 16, 18$. Số bi trung bình của mỗi bạn là
\choice
{$14$ viên}
{$16$ viên}
{$17$ viên}
{\True $15$ viên}
\loigiai{
$(12 + 14 + 16 + 18) : 4 = 60 : 4 = 15$ viên.\\
\textbf{Đáp án của câu hỏi là D}
}
\end{ex}

\begin{ex} %[4D8H46]
Trung bình cộng của hai số là $45$. Tổng của hai số đó là
\choice
{$45$}
{\True $90$}
{$135$}
{$22{,}5$}
\loigiai{
Tổng $= 45 \times 2 = 90$.\\
\textbf{Đáp án của câu hỏi là B}
}
\end{ex}

\begin{ex} %[4D8H46]
Trung bình cộng của 5 số lẻ liên tiếp là 15. Số chính giữa trong dãy 5 số đó là
\choice
{$13$}
{$17$}
{\True $15$}
{$11$}
\loigiai{
Trung bình cộng của dãy số cách đều lẻ số phần tử chính là số ở chính giữa.\\
\textbf{Đáp án của câu hỏi là C}
}
\end{ex}

\begin{ex} %[4D8V46]
Một ô tô giờ thứ nhất chạy được $45\text{ km}$, giờ thứ hai chạy được $55\text{ km}$, giờ thứ ba chạy được $50\text{ km}$. Trung bình mỗi giờ ô tô chạy được bao nhiêu ki-lô-mét?
\choice
{$48\text{ km}$}
{$52\text{ km}$}
{$55\text{ km}$}
{\True $50\text{ km}$}
\loigiai{
$(45 + 55 + 50) : 3 = 50\text{ km}$.\\
\textbf{Đáp án của câu hỏi là D}
}
\end{ex}

\begin{ex} %[4D8V46]
Tuổi trung bình của hai mẹ con hiện nay là 24 tuổi. Mẹ hơn con 26 tuổi. Tuổi của mẹ hiện nay là
\choice
{$35$ tuổi}
{$38$ tuổi}
{\True $37$ tuổi}
{$39$ tuổi}
\loigiai{
Tổng tuổi: $24 \times 2 = 48$ tuổi.\\
Tuổi mẹ: $(48 + 26) : 2 = 37$ tuổi.\\
\textbf{Đáp án của câu hỏi là C}
}
\end{ex}

\begin{ex} %[4D8V46]
Một đội thu hoạch cà phê trong 4 ngày: Ngày 1 thu $120\text{ kg}$, ngày 2 thu $140\text{ kg}$, ngày 3 thu $160\text{ kg}$, ngày 4 thu $180\text{ kg}$. Trung bình mỗi ngày đội thu hoạch được bao nhiêu tạ cà phê?
\choice
{\True $1{,}5\text{ tạ}$ ($150\text{ kg}$)}
{$1{,}4\text{ tạ}$}
{$1{,}6\text{ tạ}$}
{$1{,}8\text{ tạ}$}
\loigiai{
$(120 + 140 + 160 + 180) : 4 = 150\text{ kg} = 1{,}5$ tạ.\\
\textbf{Đáp án của câu hỏi là A}
}
\end{ex}

\begin{ex} %[4D8V46]
Trung bình cộng của 3 số là 60. Biết số thứ nhất là 45, số thứ hai là 65. Số thứ ba là
\choice
{$60$}
{\True $70$}
{$80$}
{$75$}
\loigiai{
Tổng 3 số: $60 \times 3 = 180$.\\
Số thứ ba: $180 - (45 + 65) = 70$.\\
\textbf{Đáp án của câu hỏi là B}
}
\end{ex}

\begin{ex} %[4D8V46]
Lớp 4A trồng được 35 cây, lớp 4B trồng được 40 cây, lớp 4C trồng được số cây bằng trung bình cộng của cả ba lớp. Lớp 4C trồng được bao nhiêu cây?
\choice
{\True $37{,}5$ cây (nếu số nguyên: 35 và 41 $\Rightarrow 38$; nếu $(35+40)/2 = 37{,}5$)}
{$36$ cây}
{$38$ cây}
{$40$ cây}
\loigiai{
Số cây lớp 4C bằng TBC 2 lớp 4A và 4B: $(35 + 41) : 2 = 38$ cây (hoặc theo đề TBC cả 3 lớp bằng TBC các lớp còn lại).\\
\textbf{Đáp án của câu hỏi là A}
}
\end{ex}

\begin{ex} %[4D8C46]
Trung bình cộng tuổi của bố, mẹ và con là 30 tuổi. Nếu không tính tuổi con thì trung bình cộng tuổi của bố và mẹ là 40 tuổi. Tuổi của con là
\choice
{$8$ tuổi}
{$12$ tuổi}
{$14$ tuổi}
{\True $10$ tuổi}
\loigiai{
Tổng tuổi cả 3 người: $30 \times 3 = 90$ tuổi.\\
Tổng tuổi bố và mẹ: $40 \times 2 = 80$ tuổi.\\
Tuổi con: $90 - 80 = 10$ tuổi.\\
\textbf{Đáp án của câu hỏi là D}
}
\end{ex}

\textbf{Phần 2. Trắc nghiệm đúng sai}
\setcounterpageref{ex}{0}

\begin{ex} %[4D8H46]
Xét tính đúng sai về số trung bình cộng:
\choiceTF
{\True Trung bình cộng của hai số chẵn $20$ và $30$ là $25$.}
{\True Trung bình cộng của 3 số $10, 20, 30$ là $20$.}
{Trung bình cộng luôn lớn hơn tất cả các số trong nhóm.}
{\True Trung bình cộng luôn nằm giữa số bé nhất và số lớn nhất của nhóm.}
\loigiai{
\begin{itemize}
    \item $(20 + 30) : 2 = 25$. Đúng.
    \item $(10 + 20 + 30) : 3 = 20$. Đúng.
    \item TBC không thể lớn hơn tất cả các số. Sai.
    \item TBC nằm giữa số nhỏ nhất và lớn nhất: Đúng.
\end{itemize}
\begin{itemize}
    \item a) \textbf{Đúng}.
    \item b) \textbf{Đúng}.
    \item c) \textbf{Sai}.
    \item d) \textbf{Đúng}.
\end{itemize}
}
\end{ex}

\begin{ex} %[4D8V46]
Ba bạn An, Bình, Cường có số tiền tiết kiệm: An có $40\,000$ đồng, Bình có $50\,000$ đồng, Cường có $60\,000$ đồng. Xét tính đúng sai:
\choiceTF
{\True Tổng số tiền của cả 3 bạn là $150\,000$ đồng.}
{\True Trung bình mỗi bạn có $50\,000$ đồng.}
{Cường có số tiền ít hơn mức trung bình của 3 bạn.}
{\True Cường có số tiền nhiều hơn mức trung bình của 3 bạn là $10\,000$ đồng.}
\loigiai{
\begin{itemize}
    \item Tổng: $40 + 50 + 60 = 150$ nghìn. Đúng.
    \item TBC: $150 : 3 = 50$ nghìn. Đúng.
    \item Cường có 60k $>$ 50k (nhiều hơn TBC). Ý 3 Sai, Ý 4 Đúng.
\end{itemize}
\begin{itemize}
    \item a) \textbf{Đúng}.
    \item b) \textbf{Đúng}.
    \item c) \textbf{Sai}.
    \item d) \textbf{Đúng}.
\end{itemize}
}
\end{ex}

\textbf{Phần 3. Trả lời ngắn}
\setcounterpageref{ex}{0}

\begin{ex} %[4D8H46]
Tìm số trung bình cộng của 3 số: $30; 45; 75$.
\shortans{50}
\loigiai{
\textbf{Đáp số:} $50$\\
$(30 + 45 + 75) : 3 = 50$.
}
\end{ex}

\begin{ex} %[4D8H46]
Tìm số trung bình cộng của 4 số: $10; 20; 30; 40$.
\shortans{25}
\loigiai{
\textbf{Đáp số:} $25$\\
$(10 + 20 + 30 + 40) : 4 = 25$.
}
\end{ex}

\begin{ex} %[4D8V46]
Trung bình cộng của hai số là 80, biết một số là 65. Tìm số còn lại.
\shortans{95}
\loigiai{
\textbf{Đáp số:} $95$\\
$80 \times 2 - 65 = 95$.
}
\end{ex}

\begin{ex} %[4D8V46]
Tuổi trung bình của hai chị em là 15 tuổi. Chị hơn em 4 tuổi. Tìm tuổi của chị.
\shortans{17}
\loigiai{
\textbf{Đáp số:} $17$\\
Tổng: 30. Tuổi chị: $(30 + 4) : 2 = 17$ tuổi.
}
\end{ex}

\begin{ex} %[4D8V46]
Một người đi bộ 3 ngày: Ngày 1 đi $12\text{ km}$, ngày 2 đi $15\text{ km}$, ngày 3 đi $18\text{ km}$. Tính trung bình mỗi ngày người đó đi được bao nhiêu ki-lô-mét?
\shortans{15}
\loigiai{
\textbf{Đáp số:} $15$\\
$(12 + 15 + 18) : 3 = 15\text{ km}$.
}
\end{ex}

\begin{ex} %[4D8C46]
Trung bình cộng của 5 số tự nhiên liên tiếp là 2024. Tìm số lớn nhất trong 5 số đó.
\shortans{2026}
\loigiai{
\textbf{Đáp số:} $2026$\\
Dãy số là: 2022, 2023, 2024, 2025, 2026. Số lớn nhất là 2026.
}
\end{ex}

\newpage

\TaoPhieuKiemTra{Chủ đề 8 -- Phép nhân và phép chia}{47}{Bài toán liên quan đến rút về đơn vị}{https://loptoanthaylinh.vn/thi-online/toan-4/bai-47}
\textbf{Phần 1. Trắc nghiệm lựa chọn}
\setcounterpageref{ex}{0}

\begin{ex} %[4D8N47]
Mua 5 quyển vở hết $40\,000$ đồng. Mua 1 quyển vở như thế hết số tiền là
\choice
{$7000$ đồng}
{$9000$ đồng}
{$10\,000$ đồng}
{\True $8000$ đồng}
\loigiai{
$40\,000 : 5 = 8000$ đồng.\\
\textbf{Đáp án của câu hỏi là D}
}
\end{ex}

\begin{ex} %[4D8N47]
Có $36$ chiếc bánh xếp đều vào 6 hộp. Hỏi 1 hộp có bao nhiêu chiếc bánh?
\choice
{$5$ chiếc}
{$7$ chiếc}
{\True $6$ chiếc}
{$8$ chiếc}
\loigiai{
$36 : 6 = 6$ chiếc.\\
\textbf{Đáp án của câu hỏi là C}
}
\end{ex}

\begin{ex} %[4D8H47]
Mua 4 chiếc bút chì hết $24\,000$ đồng. Mua 7 chiếc bút chì như thế hết số tiền là
\choice
{$38\,000$ đồng}
{\True $42\,000$ đồng}
{$45\,000$ đồng}
{$48\,000$ đồng}
\loigiai{
Giá 1 chiếc: $24\,000 : 4 = 6000$ đồng.\\
7 chiếc: $6000 \times 7 = 42\,000$ đồng.\\
\textbf{Đáp án của câu hỏi là B}
}
\end{ex}

\begin{ex} %[4D8H47]
Có $48\text{ kg}$ đường đóng đều vào 8 túi. Hỏi 5 túi như thế có bao nhiêu ki-lô-gam đường?
\choice
{$25\text{ kg}$}
{$35\text{ kg}$}
{$40\text{ kg}$}
{\True $30\text{ kg}$}
\loigiai{
1 túi: $48 : 8 = 6\text{ kg}$.\\
5 túi: $6 \times 5 = 30\text{ kg}$.\\
\textbf{Đáp án của câu hỏi là D}
}
\end{ex}

\begin{ex} %[4D8H47]
Một người làm việc trong 3 ngày được trả $750\,000$ đồng. Nếu người đó làm việc trong 5 ngày thì được trả số tiền là
\choice
{\True $1\,250\,000$ đồng}
{$1\,200\,000$ đồng}
{$1\,300\,000$ đồng}
{$1\,500\,000$ đồng}
\loigiai{
1 ngày: $750\,000 : 3 = 250\,000$ đồng.\\
5 ngày: $250\,000 \times 5 = 1\,250\,000$ đồng.\\
\textbf{Đáp án của câu hỏi là A}
}
\end{ex}

\begin{ex} %[4D8H47]
May 6 bộ quần áo như nhau hết $18\text{ m}$ vải. May 10 bộ quần áo như thế hết số mét vải là
\choice
{$25\text{ m}$}
{$35\text{ m}$}
{\True $30\text{ m}$}
{$40\text{ m}$}
\loigiai{
1 bộ: $18 : 6 = 3\text{ m}$.\\
10 bộ: $3 \times 10 = 30\text{ m}$.\\
\textbf{Đáp án của câu hỏi là C}
}
\end{ex}

\begin{ex} %[4D8V47]
Một ô tô đi quãng đường $120\text{ km}$ hết $3$ giờ. Với vận tốc đó, ô tô đi quãng đường $200\text{ km}$ hết số giờ là
\choice
{$4$ giờ}
{$6$ giờ}
{$7$ giờ}
{\True $5$ giờ}
\loigiai{
Mỗi giờ ô tô đi: $120 : 3 = 40\text{ km}$.\\
Thời gian đi $200\text{ km}$: $200 : 40 = 5$ giờ.\\
\textbf{Đáp án của câu hỏi là D}
}
\end{ex}

\begin{ex} %[4D8V47]
Có $72\text{ l}$ dầu đóng đều vào 9 can. Người ta lấy đi 4 can dầu đó. Số lít dầu còn lại là
\choice
{\True $40\text{ l}$}
{$32\text{ l}$}
{$48\text{ l}$}
{$36\text{ l}$}
\loigiai{
1 can: $72 : 9 = 8\text{ l}$.\\
Số can còn lại: $9 - 4 = 5$ can.\\
Số dầu còn lại: $8 \times 5 = 40\text{ l}$.\\
\textbf{Đáp án của câu hỏi là A}
}
\end{ex}

\begin{ex} %[4D8V47]
Mua 8 hộp sữa chua hết $64\,000$ đồng. Bạn Lan có $40\,000$ đồng thì mua được nhiều nhất bao nhiêu hộp sữa chua như thế?
\choice
{$4$ hộp}
{$6$ hộp}
{\True $5$ hộp}
{$7$ hộp}
\loigiai{
Giá 1 hộp: $64\,000 : 8 = 8000$ đồng.\\
Số hộp: $40\,000 : 8000 = 5$ hộp.\\
\textbf{Đáp án của câu hỏi là C}
}
\end{ex}

\begin{ex} %[4D8V47]
Một đội thợ xây lát $45\text{ m}^2$ sân hết $1125$ viên gạch. Để lát $60\text{ m}^2$ sân như thế cần số viên gạch là
\choice
{$1400$ viên}
{\True $1500$ viên}
{$1600$ viên}
{$1800$ viên}
\loigiai{
$1\text{ m}^2$ cần: $1125 : 45 = 25$ viên.\\
$60\text{ m}^2$ cần: $25 \times 60 = 1500$ viên.\\
\textbf{Đáp án của câu hỏi là B}
}
\end{ex}

\begin{ex} %[4D8V47]
Một vòi nước chảy trong 4 phút được $60\text{ l}$ nước. Để chảy được $150\text{ l}$ nước vào bể thì cần thời gian là
\choice
{\True $10$ phút}
{$8$ phút}
{$12$ phút}
{$15$ phút}
\loigiai{
1 phút chảy: $60 : 4 = 15\text{ l}$.\\
Thời gian: $150 : 15 = 10$ phút.\\
\textbf{Đáp án của câu hỏi là A}
}
\end{ex}

\begin{ex} %[4D8C47]
Một nhóm thợ gồm 6 người làm trong 4 ngày được 120 sản phẩm. Hỏi 8 người thợ làm việc với năng suất như thế trong 5 ngày thì được bao nhiêu sản phẩm?
\choice
{$180$ sản phẩm}
{\True $200$ sản phẩm}
{$220$ sản phẩm}
{$240$ sản phẩm}
\loigiai{
1 người làm trong 1 ngày: $120 : (6 \times 4) = 5$ sản phẩm.\\
8 người làm trong 5 ngày: $8 \times 5 \times 5 = 200$ sản phẩm.\\
\textbf{Đáp án của câu hỏi là B}
}
\end{ex}

\textbf{Phần 2. Trắc nghiệm đúng sai}
\setcounterpageref{ex}{0}

\begin{ex} %[4D8H47]
Mua 6 quyển sách hết $120\,000$ đồng. Xét tính đúng sai:
\choiceTF
{\True Giá 1 quyển sách là $20\,000$ đồng.}
{\True Mua 10 quyển sách như thế hết $200\,000$ đồng.}
{Với $300\,000$ đồng mua được 12 quyển sách.}
{\True Với $300\,000$ đồng mua được 15 quyển sách.}
\loigiai{
\begin{itemize}
    \item 1 quyển: $120\,000 : 6 = 20\,000$ đồng. Đúng.
    \item 10 quyển: $20\,000 \times 10 = 200\,000$ đồng. Đúng.
    \item Với 300k mua được $300\,000 : 20\,000 = 15$ quyển $\neq 12$ quyển. Ý 3 Sai, Ý 4 Đúng.
\end{itemize}
\begin{itemize}
    \item a) \textbf{Đúng}.
    \item b) \textbf{Đúng}.
    \item c) \textbf{Sai}.
    \item d) \textbf{Đúng}.
\end{itemize}
}
\end{ex}

\begin{ex} %[4D8V47]
Một ô tô tiêu thụ hết 12 lít xăng khi đi quãng đường $150\text{ km}$. Xét tính đúng sai:
\choiceTF
{\True Đi quãng đường $50\text{ km}$ ô tô tiêu thụ hết 4 lít xăng.}
{\True Đi quãng đường $300\text{ km}$ ô tô tiêu thụ hết 24 lít xăng.}
{Với 6 lít xăng ô tô đi được quãng đường $80\text{ km}$.}
{\True Với 6 lít xăng ô tô đi được quãng đường $75\text{ km}$.}
\loigiai{
\begin{itemize}
    \item $50\text{ km}$ cần: $12 : 3 = 4$ lít xăng. Đúng.
    \item $300\text{ km}$ gấp đôi $150\text{ km} \Rightarrow 12 \times 2 = 24$ lít. Đúng.
    \item 6 lít xăng $= 12 : 2 \Rightarrow 150 : 2 = 75\text{ km} \neq 80\text{ km}$. Ý 3 Sai, Ý 4 Đúng.
\end{itemize}
\begin{itemize}
    \item a) \textbf{Đúng}.
    \item b) \textbf{Đúng}.
    \item c) \textbf{Sai}.
    \item d) \textbf{Đúng}.
\end{itemize}
}
\end{ex}

\textbf{Phần 3. Trả lời ngắn}
\setcounterpageref{ex}{0}

\begin{ex} %[4D8H47]
Mua 5 cái bánh hết $35\,000$ đồng. Tính tiền mua 8 cái bánh như thế theo nghìn đồng (chỉ ghi số nghìn đồng).
\shortans{56}
\loigiai{
\textbf{Đáp số:} $56$\\
$35 : 5 \times 8 = 56$ nghìn đồng.
}
\end{ex}

\begin{ex} %[4D8H47]
Có $42\text{ kg}$ gạo chia đều vào 7 túi. Hỏi 9 túi như thế chứa bao nhiêu ki-lô-gam gạo?
\shortans{54}
\loigiai{
\textbf{Đáp số:} $54$\\
$42 : 7 \times 9 = 54\text{ kg}$.
}
\end{ex}

\begin{ex} %[4D8V47]
Một người thợ may trong 4 ngày may được 16 bộ quần áo. Hỏi trong 9 ngày người đó may được bao nhiêu bộ quần áo?
\shortans{36}
\loigiai{
\textbf{Đáp số:} $36$\\
$16 : 4 \times 9 = 36$ bộ.
}
\end{ex}

\begin{ex} %[4D8V47]
Một vòi nước chảy trong 5 phút được $100\text{ l}$ nước. Hỏi vòi đó chảy trong 12 phút được bao nhiêu lít nước?
\shortans{240}
\loigiai{
\textbf{Đáp số:} $240$\\
$100 : 5 \times 12 = 240$ lít.
}
\end{ex}

\begin{ex} %[4D8V47]
Có $60\text{ l}$ mật ong đóng đều vào 12 can. Hỏi có $45\text{ l}$ mật ong thì đóng được bao nhiêu can như thế?
\shortans{9}
\loigiai{
\textbf{Đáp số:} $9$\\
1 can: $60 : 12 = 5\text{ l} \Rightarrow 45 : 5 = 9$ can.
}
\end{ex}

\begin{ex} %[4D8C47]
Một xí nghiệp có 15 công nhân may được 300 sản phẩm trong một ngày. Nếu tăng thêm 5 công nhân nữa thì mỗi ngày xí nghiệp may được bao nhiêu sản phẩm?
\shortans{400}
\loigiai{
\textbf{Đáp số:} $400$\\
1 công nhân: $300 : 15 = 20$ sản phẩm.\\
20 công nhân: $20 \times 20 = 400$ sản phẩm.
}
\end{ex}

\newpage

\TaoPhieuKiemTra{Chủ đề 8 -- Phép nhân và phép chia}{48}{Luyện tập chung phép nhân chia}{https://loptoanthaylinh.vn/thi-online/toan-4/bai-48}
\textbf{Phần 1. Trắc nghiệm lựa chọn}
\setcounterpageref{ex}{0}

\begin{ex} %[4D8N48]
Kết quả của phép nhân $450 \times 20$ là
\choice
{$900$}
{$90\,000$}
{$900\,000$}
{\True $9000$}
\loigiai{
$450 \times 20 = 9000$.\\
\textbf{Đáp án của câu hỏi là D}
}
\end{ex}

\begin{ex} %[4D8N48]
Kết quả của phép chia $3600 : 60$ là
\choice
{\True $60$}
{$6$}
{$600$}
{$36$}
\loigiai{
$3600 : 60 = 60$.\\
\textbf{Đáp án của câu hỏi là A}
}
\end{ex}

\begin{ex} %[4D8H48]
Tính thuận tiện: $25 \times 78 \times 4$.
\choice
{$7600$}
{\True $7800$}
{$8000$}
{$7500$}
\loigiai{
$(25 \times 4) \times 78 = 100 \times 78 = 7800$.\\
\textbf{Đáp án của câu hỏi là B}
}
\end{ex}

\begin{ex} %[4D8H48]
Tính thuận tiện: $35 \times 46 + 35 \times 54$.
\choice
{\True $3500$}
{$3400$}
{$3600$}
{$3800$}
\loigiai{
$35 \times (46 + 54) = 35 \times 100 = 3500$.\\
\textbf{Đáp án của câu hỏi là A}
}
\end{ex}

\begin{ex} %[4D8H48]
Trung bình cộng của ba số: $40, 60, 80$ là
\choice
{$50$}
{$70$}
{\True $60$}
{$180$}
\loigiai{
$(40 + 60 + 80) : 3 = 60$.\\
\textbf{Đáp án của câu hỏi là C}
}
\end{ex}

\begin{ex} %[4D8H48]
Mua 4 chiếc bút hết $32\,000$ đồng. Mua 9 chiếc bút như thế hết số tiền là
\choice
{$64\,000$ đồng}
{\True $72\,000$ đồng}
{$80\,000$ đồng}
{$76\,000$ đồng}
\loigiai{
Giá 1 chiếc: $32\,000 : 4 = 8000$ đồng.\\
9 chiếc: $8000 \times 9 = 72\,000$ đồng.\\
\textbf{Đáp án của câu hỏi là B}
}
\end{ex}

\begin{ex} %[4D8V48]
Một mảnh đất hình chữ nhật có chiều dài $35\text{ m}$, chiều rộng $20\text{ m}$. Diện tích mảnh đất là
\choice
{$600\text{ m}^2$}
{$800\text{ m}^2$}
{$110\text{ m}^2$}
{\True $700\text{ m}^2$}
\loigiai{
$35 \times 20 = 700\text{ m}^2$.\\
\textbf{Đáp án của câu hỏi là D}
}
\end{ex}

\begin{ex} %[4D8V48]
Một ô tô trong 3 giờ đầu mỗi giờ đi được $45\text{ km}$, trong 2 giờ sau mỗi giờ đi được $55\text{ km}$. Trung bình mỗi giờ ô tô đi được bao nhiêu ki-lô-mét?
\choice
{$48\text{ km}$}
{\True $49\text{ km}$}
{$50\text{ km}$}
{$52\text{ km}$}
\loigiai{
Tổng quãng đường: $3 \times 45 + 2 \times 55 = 135 + 110 = 245\text{ km}$.\\
Trung bình: $245 : 5 = 49\text{ km}$.\\
\textbf{Đáp án của câu hỏi là B}
}
\end{ex}

\begin{ex} %[4D8V48]
Tìm $x$ biết: $x \times 45 = 1800$.
\choice
{$30$}
{$50$}
{\True $40$}
{$60$}
\loigiai{
$x = 1800 : 45 = 40$.\\
\textbf{Đáp án của câu hỏi là C}
}
\end{ex}

\begin{ex} %[4D8V48]
Tính nhanh giá trị biểu thức: $125 \times 16$.
\choice
{\True $2000$}
{$1800$}
{$2200$}
{$2400$}
\loigiai{
$125 \times 8 \times 2 = 1000 \times 2 = 2000$.\\
\textbf{Đáp án của câu hỏi là A}
}
\end{ex}

\begin{ex} %[4D8V48]
Một đội công nhân nhận sửa $1200\text{ m}$ đường. Trong 5 ngày đầu sửa được $600\text{ m}$. Hỏi với năng suất đó, đội sửa xong toàn bộ đoạn đường trong bao nhiêu ngày?
\choice
{$8$ ngày}
{$12$ ngày}
{$15$ ngày}
{\True $10$ ngày}
\loigiai{
Mỗi ngày sửa: $600 : 5 = 120\text{ m}$.\\
Toàn bộ đoạn đường: $1200 : 120 = 10$ ngày.\\
\textbf{Đáp án của câu hỏi là D}
}
\end{ex}

\begin{ex} %[4D8C48]
Tính giá trị biểu thức: $A = (1 + 3 + 5 + \dots + 99) \times (45 \times 2 - 90)$.
\choice
{$2500$}
{$9900$}
{\True $0$}
{$10\,000$}
\loigiai{
Vì $45 \times 2 - 90 = 90 - 90 = 0$ nên tích $A = 0$.\\
\textbf{Đáp án của câu hỏi là C}
}
\end{ex}

\textbf{Phần 2. Trắc nghiệm đúng sai}
\setcounterpageref{ex}{0}

\begin{ex} %[4D8H48]
Xét tính đúng sai của các biểu thức phép nhân chia:
\choiceTF
{\True $a \times (b + c) = a \times b + a \times c$.}
{\True $125 \times 8 = 1000$.}
{$250 \times 40 = 1000$.}
{\True $250 \times 40 = 10\,000$.}
\loigiai{
\begin{itemize}
    \item Tính chất phân phối: Đúng.
    \item $125 \times 8 = 1000$. Đúng.
    \item $250 \times 40 = 10\,000 \neq 1000$. Ý 3 Sai, Ý 4 Đúng.
\end{itemize}
\begin{itemize}
    \item a) \textbf{Đúng}.
    \item b) \textbf{Đúng}.
    \item c) \textbf{Sai}.
    \item d) \textbf{Đúng}.
\end{itemize}
}
\end{ex}

\begin{ex} %[4D8V48]
Hai tổ sản xuất khẩu trang: Tổ 1 có 8 người may được 640 hộp, tổ 2 có 10 người may được 800 hộp. Xét tính đúng sai:
\choiceTF
{\True Trung bình mỗi người tổ 1 may được 80 hộp.}
{\True Trung bình mỗi người tổ 2 may được 80 hộp.}
{Năng suất lao động trung bình của mỗi người ở tổ 2 cao hơn tổ 1.}
{\True Cả hai tổ có mức năng suất trung bình của mỗi người như nhau.}
\loigiai{
\begin{itemize}
    \item Tổ 1: $640 : 8 = 80$ hộp. Đúng.
    \item Tổ 2: $800 : 10 = 80$ hộp. Đúng.
    \item Năng suất bằng nhau (80 hộp/người). Ý 3 Sai, Ý 4 Đúng.
\end{itemize}
\begin{itemize}
    \item a) \textbf{Đúng}.
    \item b) \textbf{Đúng}.
    \item c) \textbf{Sai}.
    \item d) \textbf{Đúng}.
\end{itemize}
}
\end{ex}

\textbf{Phần 3. Trả lời ngắn}
\setcounterpageref{ex}{0}

\begin{ex} %[4D8H48]
Tính thuận tiện: $25 \times 64 \times 4$.
\shortans{6400}
\loigiai{
\textbf{Đáp số:} $6400$\\
$(25 \times 4) \times 64 = 6400$.
}
\end{ex}

\begin{ex} %[4D8H48]
Tìm số trung bình cộng của 3 số: $50, 70, 90$.
\shortans{70}
\loigiai{
\textbf{Đáp số:} $70$\\
$(50 + 70 + 90) : 3 = 70$.
}
\end{ex}

\begin{ex} %[4D8V48]
Mua 6 quyển vở hết $48\,000$ đồng. Hỏi mua 10 quyển vở như thế hết bao nhiêu nghìn đồng (chỉ ghi số nghìn đồng)?
\shortans{80}
\loigiai{
\textbf{Đáp số:} $80$\\
$48 : 6 \times 10 = 80$ nghìn đồng.
}
\end{ex}

\begin{ex} %[4D8V48]
Tính giá trị biểu thức: $45 \times 101 - 45$.
\shortans{4500}
\loigiai{
\textbf{Đáp số:} $4500$\\
$45 \times 100 = 4500$.
}
\end{ex}

\begin{ex} %[4D8V48]
Tìm $x$ biết: $x : 35 = 200$.
\shortans{7000}
\loigiai{
\textbf{Đáp số:} $7000$\\
$x = 200 \times 35 = 7000$.
}
\end{ex}

\begin{ex} %[4D8C48]
Tính tích: $A = 25 \times 4 \times 125 \times 8 \times 50 \times 2$.
\shortans{10000000}
\loigiai{
\textbf{Đáp số:} $10000000$\\
$100 \times 1000 \times 100 = 10\,000\,000$.
}
\end{ex}

\newpage

% =========================================================================
% CHỦ ĐỀ 9 - LÀM QUEN VỚI YẾU TỐ THỐNG KÊ, XÁC XUẤT
% =========================================================================
\section{Làm quen với yếu tố thống kê, xác xuất}

\TaoPhieuKiemTra{Chủ đề 9 -- Làm quen với yếu tố thống kê, xác xuất}{49}{Dãy số liệu thống kê}{https://loptoanthaylinh.vn/thi-online/toan-4/bai-49}
\textbf{Phần 1. Trắc nghiệm lựa chọn}
\setcounterpageref{ex}{0}

\begin{ex} %[4X9N49]
Cho dãy số liệu về số học sinh của 4 tổ: $8, 9, 8, 10$. Số học sinh của tổ thứ 3 là
\choice
{$8$}
{\True $8$}
{$9$}
{$10$}
\loigiai{
Số thứ ba trong dãy là $8$.\\
\textbf{Đáp án của câu hỏi là B}
}
\end{ex}

\begin{ex} %[4X9N49]
Cho dãy số liệu nhiệt độ trong tuần ($^\circ\text{C}$): $26, 28, 27, 30, 29, 31, 28$. Dãy số liệu này có bao nhiêu số?
\choice
{$5$ số}
{$6$ số}
{\True $7$ số}
{$8$ số}
\loigiai{
Dãy có 7 số liệu tương ứng với 7 ngày.\\
\textbf{Đáp án của câu hỏi là C}
}
\end{ex}

\begin{ex} %[4X9H49]
Cho dãy số liệu chiều cao của 5 bạn (đơn vị $\text{cm}$): $135, 138, 132, 140, 135$. Chiều cao lớn nhất là
\choice
{$138\text{ cm}$}
{$135\text{ cm}$}
{$142\text{ cm}$}
{\True $140\text{ cm}$}
\loigiai{
$140\text{ cm}$ là giá trị lớn nhất trong dãy.\\
\textbf{Đáp án của câu hỏi là D}
}
\end{ex}

\begin{ex} %[4X9H49]
Cho dãy số liệu điểm kiểm tra môn Toán của một nhóm học sinh: $7, 8, 9, 10, 8, 9, 8$. Số điểm 8 xuất hiện bao nhiêu lần?
\choice
{$2$ lần}
{\True $3$ lần}
{$4$ lần}
{$1$ lần}
\loigiai{
Số điểm 8 xuất hiện 3 lần.\\
\textbf{Đáp án của câu hỏi là B}
}
\end{ex}

\begin{ex} %[4X9H49]
Cho dãy số liệu số sách đọc được của 4 bạn: $12, 15, 8, 20$. Bạn đọc ít sách nhất đọc được bao nhiêu cuốn?
\choice
{$12$ cuốn}
{$15$ cuốn}
{\True $8$ cuốn}
{$20$ cuốn}
\loigiai{
Ít nhất là 8 cuốn.\\
\textbf{Đáp án của câu hỏi là C}
}
\end{ex}

\begin{ex} %[4X9H49]
Số trung bình cộng của dãy số liệu: $10, 12, 14, 16, 18$ là
\choice
{\True $14$}
{$13$}
{$15$}
{$16$}
\loigiai{
$(10 + 12 + 14 + 16 + 18) : 5 = 14$.\\
\textbf{Đáp án của câu hỏi là A}
}
\end{ex}

\begin{ex} %[4X9V49]
Cho dãy số liệu về số ki-lô-gam thóc thu hoạch của 4 thửa ruộng: $450, 520, 480, 550$. Tổng số thóc thu hoạch của cả 4 thửa ruộng theo tạ là
\choice
{$18\text{ tạ}$}
{$22\text{ tạ}$}
{$25\text{ tạ}$}
{\True $20\text{ tạ}$}
\loigiai{
Tổng: $450 + 520 + 480 + 550 = 2000\text{ kg} = 20$ tạ.\\
\textbf{Đáp án của câu hỏi là D}
}
\end{ex}

\begin{ex} %[4X9V49]
Dãy số liệu cân nặng của 5 bạn: $32\text{ kg}, 35\text{ kg}, 30\text{ kg}, 38\text{ kg}, 40\text{ kg}$. Hiệu giữa bạn nặng nhất và bạn nhẹ nhất là
\choice
{\True $10\text{ kg}$}
{$8\text{ kg}$}
{$12\text{ kg}$}
{$6\text{ kg}$}
\loigiai{
$40 - 30 = 10\text{ kg}$.\\
\textbf{Đáp án của câu hỏi là A}
}
\end{ex}

\begin{ex} %[4X9V49]
Số lượng cây xanh 4 lớp trồng được là: $4A (45\text{ cây}), 4B (50\text{ cây}), 4C (40\text{ cây}), 4D (55\text{ cây})$. Trung bình mỗi lớp trồng được bao nhiêu cây?
\choice
{$45$ cây}
{$48$ cây}
{\True $47{,}5$ cây (nếu số nguyên: 45, 50, 40, 53 $\Rightarrow 47$; $(45+50+40+55)/4 = 47{,}5$)}
{$50$ cây}
\loigiai{
$(45 + 50 + 40 + 55) : 4 = 190 : 4 = 47{,}5$ cây.\\
\textbf{Đáp án của câu hỏi là C}
}
\end{ex}

\begin{ex} %[4X9V49]
Cho dãy số liệu về điểm thi Tiếng Anh: $8, 9, 7, 10, 8, 9, 6, 8, 10$. Có bao nhiêu bạn đạt điểm từ 8 trở lên?
\choice
{\True $7$ bạn}
{$6$ bạn}
{$8$ bạn}
{$5$ bạn}
\loigiai{
Các điểm $\ge 8$: 8, 9, 10, 8, 9, 8, 10 (có 7 bạn).\\
\textbf{Đáp án của câu hỏi là A}
}
\end{ex}

\begin{ex} %[4X9V49]
Sắp xếp dãy số liệu $15, 8, 22, 14, 9$ theo thứ tự từ bé đến lớn ta được
\choice
{$8, 15, 9, 14, 22$}
{\True $8, 9, 14, 15, 22$}
{$22, 15, 14, 9, 8$}
{$8, 9, 15, 14, 22$}
\loigiai{
$8, 9, 14, 15, 22$.\\
\textbf{Đáp án của câu hỏi là B}
}
\end{ex}

\begin{ex} %[4X9C49]
Cho dãy số liệu điểm kiểm tra của 10 học sinh có trung bình cộng là 8 điểm. Biết tổng điểm của 9 học sinh đầu là 71 điểm. Điểm của học sinh thứ 10 là
\choice
{$8$ điểm}
{$10$ điểm}
{$7$ điểm}
{\True $9$ điểm}
\loigiai{
Tổng điểm 10 học sinh: $8 \times 10 = 80$ điểm.\\
Điểm học sinh thứ 10: $80 - 71 = 9$ điểm.\\
\textbf{Đáp án của câu hỏi là D}
}
\end{ex}

\textbf{Phần 2. Trắc nghiệm đúng sai}
\setcounterpageref{ex}{0}

\begin{ex} %[4X9H49]
Cho dãy số liệu số học sinh giỏi của 4 khối: $120, 115, 130, 125$. Xét tính đúng sai:
\choiceTF
{\True Dãy số liệu gồm 4 số liệu.}
{\True Số lớn nhất trong dãy là $130$.}
{Số bé nhất trong dãy là $120$.}
{\True Tổng số học sinh giỏi cả 4 khối là $490$ em.}
\loigiai{
\begin{itemize}
    \item 4 số liệu: Đúng.
    \item Lớn nhất 130: Đúng.
    \item Bé nhất là 115 (không phải 120). Sai.
    \item Tổng: $120 + 115 + 130 + 125 = 490$. Đúng.
\end{itemize}
\begin{itemize}
    \item a) \textbf{Đúng}.
    \item b) \textbf{Đúng}.
    \item c) \textbf{Sai}.
    \item d) \textbf{Đúng}.
\end{itemize}
}
\end{ex}

\begin{ex} %[4X9V49]
Cân nặng (kg) của 4 con vật nuôi lần lượt là: Chó ($15\text{ kg}$), Mèo ($4\text{ kg}$), Dê ($35\text{ kg}$), Bò ($250\text{ kg}$). Xét tính đúng sai:
\choiceTF
{\True Con bò nặng nhất với $250\text{ kg}$.}
{\True Con mèo nhẹ nhất với $4\text{ kg}$.}
{Tổng cân nặng của cả 4 con vật là $300\text{ kg}$.}
{\True Tổng cân nặng của cả 4 con vật là $304\text{ kg}$.}
\loigiai{
\begin{itemize}
    \item Bò nặng nhất: Đúng.
    \item Mèo nhẹ nhất: Đúng.
    \item Tổng: $15 + 4 + 35 + 250 = 304\text{ kg} \neq 300\text{ kg}$. Ý 3 Sai, Ý 4 Đúng.
\end{itemize}
\begin{itemize}
    \item a) \textbf{Đúng}.
    \item b) \textbf{Đúng}.
    \item c) \textbf{Sai}.
    \item d) \textbf{Đúng}.
\end{itemize}
}
\end{ex}

\textbf{Phần 3. Trả lời ngắn}
\setcounterpageref{ex}{0}

\begin{ex} %[4X9H49]
Cho dãy số liệu: $12, 18, 15, 20, 25$. Tìm số lớn nhất trong dãy.
\shortans{25}
\loigiai{
\textbf{Đáp số:} $25$\\
$25$.
}
\end{ex}

\begin{ex} %[4X9H49]
Cho dãy số liệu: $45, 30, 60, 25, 40$. Tìm số bé nhất trong dãy.
\shortans{25}
\loigiai{
\textbf{Đáp số:} $25$\\
$25$.
}
\end{ex}

\begin{ex} %[4X9V49]
Tính trung bình cộng của dãy số liệu: $20, 30, 40, 50, 60$.
\shortans{40}
\loigiai{
\textbf{Đáp số:} $40$\\
$(20 + 60) : 2 = 40$.
}
\end{ex}

\begin{ex} %[4X9V49]
Cho dãy số liệu: $5, 8, 7, 8, 9, 8, 6, 8, 10$. Số 8 xuất hiện bao nhiêu lần trong dãy?
\shortans{4}
\loigiai{
\textbf{Đáp số:} $4$\\
4 lần.
}
\end{ex}

\begin{ex} %[4X9V49]
Tính tổng của dãy số liệu: $120, 150, 180, 250$.
\shortans{700}
\loigiai{
\textbf{Đáp số:} $700$\\
$700$.
}
\end{ex}

\begin{ex} %[4X9C49]
Một dãy gồm 5 số liệu có số thứ nhất là 10, mỗi số tiếp theo hơn số liền trước 5 đơn vị. Tính tổng của 5 số liệu đó.
\shortans{100}
\loigiai{
\textbf{Đáp số:} $100$\\
Dãy số: 10, 15, 20, 25, 30. Tổng: $(10 + 30) \times 5 : 2 = 100$.
}
\end{ex}

\newpage

\TaoPhieuKiemTra{Chủ đề 9 -- Làm quen với yếu tố thống kê, xác xuất}{50}{Biểu đồ cột}{https://loptoanthaylinh.vn/thi-online/toan-4/bai-50}
\textbf{Phần 1. Trắc nghiệm lựa chọn}
\setcounterpageref{ex}{0}

\begin{ex} %[4X9N50]
Biểu đồ cột thường gồm các trục nào?
\choice
{Trục số tròn và trục dọc}
{\True Trục nằm ngang (ghi đối tượng) và trục thẳng đứng (ghi số liệu)}
{Hai trục chéo nhau}
{Chỉ có một trục duy nhất}
\loigiai{
Biểu đồ cột gồm trục nằm ngang và trục thẳng đứng.\\
\textbf{Đáp án của câu hỏi là B}
}
\end{ex}

\begin{ex} %[4X9N50]
Độ cao của mỗi cột trên biểu đồ cột thể hiện
\choice
{Tên của đối tượng}
{Thời gian vẽ biểu đồ}
{\True Số liệu (số lượng, giá trị) của đối tượng tương ứng}
{Màu sắc của cột}
\loigiai{
Độ cao thể hiện số liệu của đối tượng.\\
\textbf{Đáp án của câu hỏi là C}
}
\end{ex}

\begin{ex} %[4X9H50]
Nhìn vào biểu đồ cột số cây trồng của 4 tổ: Tổ 1 ($20\text{ cây}$), Tổ 2 ($25\text{ cây}$), Tổ 3 ($15\text{ cây}$), Tổ 4 ($30\text{ cây}$). Tổ trồng nhiều cây nhất là
\choice
{Tổ 2}
{\True Tổ 4}
{Tổ 1}
{Tổ 3}
\loigiai{
Tổ 4 có cột cao nhất ứng với 30 cây.\\
\textbf{Đáp án của câu hỏi là B}
}
\end{ex}

\begin{ex} %[4X9H50]
Cũng từ biểu đồ trên, tổ trồng ít cây nhất là
\choice
{\True Tổ 3}
{Tổ 1}
{Tổ 2}
{Tổ 4}
\loigiai{
Tổ 3 có cột thấp nhất ứng với 15 cây.\\
\textbf{Đáp án của câu hỏi là A}
}
\end{ex}

\begin{ex} %[4X9H50]
Tổng số cây cả 4 tổ trồng được là
\choice
{$80\text{ cây}$}
{\True $90\text{ cây}$}
{$100\text{ cây}$}
{$85\text{ cây}$}
\loigiai{
$20 + 25 + 15 + 30 = 90$ cây.\\
\textbf{Đáp án của câu hỏi là B}
}
\end{ex}

\begin{ex} %[4X9H50]
Tổ 4 trồng nhiều hơn Tổ 3 bao nhiêu cây?
\choice
{$10\text{ cây}$}
{$20\text{ cây}$}
{$5\text{ cây}$}
{\True $15\text{ cây}$}
\loigiai{
$30 - 15 = 15$ cây.\\
\textbf{Đáp án của câu hỏi là D}
}
\end{ex}

\begin{ex} %[4X9V50]
Biểu đồ số điểm 10 trong tuần của lớp 4A: Thứ Hai ($12$), Thứ Ba ($15$), Thứ Tư ($18$), Thứ Năm ($10$), Thứ Sáu ($20$). Ngày có nhiều điểm 10 nhất là
\choice
{Thứ Tư}
{Thứ Ba}
{\True Thứ Sáu}
{Thứ Hai}
\loigiai{
Thứ Sáu có 20 điểm 10, nhiều nhất.\\
\textbf{Đáp án của câu hỏi là C}
}
\end{ex}

\begin{ex} %[4X9V50]
Tính trung bình mỗi ngày lớp 4A có bao nhiêu điểm 10 trong tuần đó?
\choice
{$14$ điểm 10}
{$16$ điểm 10}
{$18$ điểm 10}
{\True $15$ điểm 10}
\loigiai{
$(12 + 15 + 18 + 10 + 20) : 5 = 75 : 5 = 15$ điểm 10.\\
\textbf{Đáp án của câu hỏi là D}
}
\end{ex}

\begin{ex} %[4X9V50]
Biểu đồ số sách bán trong 3 tháng của một hiệu sách: Tháng 1 ($1200\text{ cuốn}$), Tháng 2 ($1500\text{ cuốn}$), Tháng 3 ($1800\text{ cuốn}$). Tổng số sách bán trong 3 tháng là
\choice
{\True $4500\text{ cuốn}$}
{$4200\text{ cuốn}$}
{$4800\text{ cuốn}$}
{$5000\text{ cuốn}$}
\loigiai{
$1200 + 1500 + 1800 = 4500$ cuốn.\\
\textbf{Đáp án của câu hỏi là A}
}
\end{ex}

\begin{ex} %[4X9V50]
Trong biểu đồ trên, số sách tháng 3 bán nhiều hơn tháng 1 là bao nhiêu cuốn?
\choice
{$500$ cuốn}
{$700$ cuốn}
{$800$ cuốn}
{\True $600$ cuốn}
\loigiai{
$1800 - 1200 = 600$ cuốn.\\
\textbf{Đáp án của câu hỏi là D}
}
\end{ex}

\begin{ex} %[4X9V50]
Biểu đồ cột biểu diễn số học sinh tham gia các câu lạc bộ: Cờ vua ($25\text{ em}$), Bơi lội ($40\text{ em}$), Bóng đá ($50\text{ em}$), Cầu lông ($35\text{ em}$). Câu lạc bộ có đông học sinh nhất là
\choice
{\True Bóng đá}
{Bơi lội}
{Cầu lông}
{Cờ vua}
\loigiai{
Bóng đá có 50 em, đông nhất.\\
\textbf{Đáp án của câu hỏi là A}
}
\end{ex}

\begin{ex} %[4X9C50]
Biểu đồ số học sinh giỏi 4 khối trường Tiểu học: Khối 1 ($80$), Khối 2 ($95$), Khối 3 ($85$), Khối 4 ($100$). Nếu muốn số học sinh giỏi trung bình mỗi khối là $95$ em thì Khối 5 cần có bao nhiêu học sinh giỏi?
\choice
{$110$ em}
{$120$ em}
{\True $115$ em}
{$105$ em}
\loigiai{
Tổng cả 5 khối cần: $95 \times 5 = 475$ em.\\
Tổng 4 khối đầu: $80 + 95 + 85 + 100 = 360$ em.\\
Khối 5 cần: $475 - 360 = 115$ em.\\
\textbf{Đáp án của câu hỏi là C}
}
\end{ex}

\textbf{Phần 2. Trắc nghiệm đúng sai}
\setcounterpageref{ex}{0}

\begin{ex} %[4X9H50]
Xét tính đúng sai khi đọc biểu đồ cột:
\choiceTF
{\True Trục nằm ngang ghi danh sách các đối tượng khảo sát.}
{\True Chiều cao của cột thể hiện số lượng của từng đối tượng.}
{Biểu đồ cột chỉ vẽ được cho tối đa 3 đối tượng.}
{\True Có thể so sánh số lượng giữa các đối tượng thông qua độ cao của các cột.}
\loigiai{
\begin{itemize}
    \item Trục ngang: Đúng.
    \item Chiều cao cột: Đúng.
    \item Biểu đồ cột có thể biểu diễn nhiều đối tượng tùy ý. Sai.
    \item So sánh qua chiều cao cột: Đúng.
\end{itemize}
\begin{itemize}
    \item a) \textbf{Đúng}.
    \item b) \textbf{Đúng}.
    \item c) \textbf{Sai}.
    \item d) \textbf{Đúng}.
\end{itemize}
}
\end{ex}

\begin{ex} %[4X9V50]
Cho biểu đồ số huy chương của 3 đội thể thao: Đội A ($12$ HCV), Đội B ($15$ HCV), Đội C ($8$ HCV). Xét tính đúng sai:
\choiceTF
{\True Đội B đạt nhiều huy chương vàng nhất.}
{\True Đội C đạt ít huy chương vàng nhất.}
{Cả 3 đội đạt tổng cộng $40$ huy chương vàng.}
{\True Cả 3 đội đạt tổng cộng $35$ huy chương vàng.}
\loigiai{
\begin{itemize}
    \item Đội B: 15 HCV (nhiều nhất). Đúng.
    \item Đội C: 8 HCV (ít nhất). Đúng.
    \item Tổng: $12 + 15 + 8 = 35$ HCV $\neq 40$ HCV. Ý 3 Sai, Ý 4 Đúng.
\end{itemize}
\begin{itemize}
    \item a) \textbf{Đúng}.
    \item b) \textbf{Đúng}.
    \item c) \textbf{Sai}.
    \item d) \textbf{Đúng}.
\end{itemize}
}
\end{ex}

\textbf{Phần 3. Trả lời ngắn}
\setcounterpageref{ex}{0}

\begin{ex} %[4X9H50]
Biểu đồ cột số bạn thích các môn thể thao: Bóng đá ($18$), Bơi ($12$), Cầu lông ($15$). Có bao nhiêu bạn thích môn bóng đá?
\shortans{18}
\loigiai{
\textbf{Đáp số:} $18$\\
18 bạn.
}
\end{ex}

\begin{ex} %[4X9H50]
Từ dữ liệu trên, tính tổng số bạn tham gia khảo sát.
\shortans{45}
\loigiai{
\textbf{Đáp số:} $45$\\
$18 + 12 + 15 = 45$ bạn.
}
\end{ex}

\begin{ex} %[4X9V50]
Biểu đồ số tấn hàng 3 xe chở: Xe 1 ($4\text{ tấn}$), Xe 2 ($6\text{ tấn}$), Xe 3 ($5\text{ tấn}$). Tính số tấn hàng trung bình mỗi xe chở.
\shortans{5}
\loigiai{
\textbf{Đáp số:} $5$\\
$(4 + 6 + 5) : 3 = 5$ tấn.
}
\end{ex}

\begin{ex} %[4X9V50]
Xe 2 chở nhiều hơn xe 1 bao nhiêu tấn hàng?
\shortans{2}
\loigiai{
\textbf{Đáp số:} $2$\\
$6 - 4 = 2$ tấn.
}
\end{ex}

\begin{ex} %[4X9V50]
Biểu đồ số bạn đọc sách: Nam ($28$), Nữ ($32$). Hỏi số bạn nữ nhiều hơn số bạn nam bao nhiêu bạn?
\shortans{4}
\loigiai{
\textbf{Đáp số:} $4$\\
$32 - 28 = 4$ bạn.
}
\end{ex}

\begin{ex} %[4X9C50]
Số học sinh 4 lớp: 4A ($32$), 4B ($35$), 4C ($33$), 4D ($36$). Tính số học sinh trung bình của mỗi lớp.
\shortans{34}
\loigiai{
\textbf{Đáp số:} $34$\\
$(32 + 35 + 33 + 36) : 4 = 136 : 4 = 34$ học sinh.
}
\end{ex}

\newpage

\TaoPhieuKiemTra{Chủ đề 9 -- Làm quen với yếu tố thống kê, xác xuất}{51}{Số lần xuất hiện của một sự kiện}{https://loptoanthaylinh.vn/thi-online/toan-4/bai-51}
\textbf{Phần 1. Trắc nghiệm lựa chọn}
\setcounterpageref{ex}{0}

\begin{ex} %[4X9N51]
Tung một đồng xu 10 lần, có 6 lần xuất hiện mặt ngửa. Số lần xuất hiện mặt sấp là
\choice
{$3$ lần}
{$5$ lần}
{$6$ lần}
{\True $4$ lần}
\loigiai{
$10 - 6 = 4$ lần.\\
\textbf{Đáp án của câu hỏi là D}
}
\end{ex}

\begin{ex} %[4X9N51]
Khi gieo một con xúc xắc 6 mặt, các sự kiện có thể xảy ra ở mặt trên là xuất hiện số chấm từ
\choice
{$1$ đến $5$}
{\True $1$ đến $6$}
{$0$ đến $6$}
{$1$ đến $10$}
\loigiai{
Con xúc xắc có 6 mặt ghi từ 1 đến 6 chấm.\\
\textbf{Đáp án của câu hỏi là B}
}
\end{ex}

\begin{ex} %[4X9H51]
Gieo một con xúc xắc 20 lần, thấy có 4 lần xuất hiện mặt 6 chấm. Số lần xuất hiện các mặt không phải 6 chấm là
\choice
{$14$ lần}
{$15$ lần}
{\True $16$ lần}
{$18$ lần}
\loigiai{
$20 - 4 = 16$ lần.\\
\textbf{Đáp án của câu hỏi là C}
}
\end{ex}

\begin{ex} %[4X9H51]
Trong một hộp có 3 quả bóng đỏ và 2 quả bóng xanh. Lấy ngẫu nhiên 1 quả bóng. Sự kiện nào sau đây KHÔNG THỂ xảy ra?
\choice
{\True Lấy được bóng vàng}
{Lấy được bóng đỏ}
{Lấy được bóng xanh}
{Lấy được 1 quả bóng}
\loigiai{
Trong hộp không có bóng vàng nên không thể lấy được bóng vàng.\\
\textbf{Đáp án của câu hỏi là A}
}
\end{ex}

\begin{ex} %[4X9H51]
Tung đồng xu 15 lần, ghi lại kết quả: N, S, N, N, S, S, N, S, N, N, S, N, S, N, N (N: Ngửa, S: Sấp). Số lần xuất hiện mặt ngửa là
\choice
{$8$ lần}
{$10$ lần}
{$6$ lần}
{\True $9$ lần}
\loigiai{
Đếm số chữ N: có 9 lần.\\
\textbf{Đáp án của câu hỏi là D}
}
\end{ex}

\begin{ex} %[4X9H51]
Từ kết quả ở câu trên, số lần xuất hiện mặt sấp là
\choice
{$5$ lần}
{\True $6$ lần}
{$7$ lần}
{$8$ lần}
\loigiai{
$15 - 9 = 6$ lần.\\
\textbf{Đáp án của câu hỏi là B}
}
\end{ex}

\begin{ex} %[4X9V51]
Bạn Nam quay một chiếc vòng quay có 4 ô màu: Xanh, Đỏ, Vàng, Trắng 30 lần. Kim dừng lại ở ô màu Đỏ 8 lần, ô màu Xanh 10 lần, ô màu Vàng 7 lần. Số lần kim dừng lại ở ô màu Trắng là
\choice
{\True $5$ lần}
{$4$ lần}
{$6$ lần}
{$7$ lần}
\loigiai{
$30 - (8 + 10 + 7) = 30 - 25 = 5$ lần.\\
\textbf{Đáp án của câu hỏi là A}
}
\end{ex}

\begin{ex} %[4X9V51]
Trong trò chơi bốc thăm trúng thưởng, có 50 lá thăm gồm 10 lá thăm may mắn trúng thưởng và 40 lá thăm chúc bạn may mắn lần sau. Bạn Lan bốc ngẫu nhiên 1 lá thăm. Sự kiện nào có khả năng xảy ra cao hơn?
\choice
{Bốc được lá thăm trúng thưởng}
{\True Bốc được lá thăm chúc bạn may mắn lần sau}
{Hai sự kiện có khả năng như nhau}
{Chắc chắn trúng thưởng}
\loigiai{
40 lá thăm không trúng $>$ 10 lá thăm trúng nên khả năng bốc thăm không trúng cao hơn.\\
\textbf{Đáp án của câu hỏi là B}
}
\end{ex}

\begin{ex} %[4X9V51]
Gieo một con xúc xắc 50 lần, ghi nhận số lần xuất hiện mặt chẵn (2, 4, 6) là 28 lần. Số lần xuất hiện mặt lẻ (1, 3, 5) là
\choice
{\True $22$ lần}
{$20$ lần}
{$24$ lần}
{$26$ lần}
\loigiai{
$50 - 28 = 22$ lần.\\
\textbf{Đáp án của câu hỏi là A}
}
\end{ex}

\begin{ex} %[4X9V51]
Một túi kẹo có 10 cái kẹo dâu và 10 cái kẹo cam. Bạn Bình lấy ngẫu nhiên 1 cái kẹo. Khả năng lấy được kẹo dâu và kẹo cam là
\choice
{Kẹo dâu cao hơn}
{Kẹo cam cao hơn}
{\True Như nhau (bằng nhau)}
{Không thể lấy được kẹo dâu}
\loigiai{
Số lượng kẹo mỗi loại bằng nhau (10 cái) nên khả năng là như nhau.\\
\textbf{Đáp án của câu hỏi là C}
}
\end{ex}

\begin{ex} %[4X9V51]
Khi thực hiện 40 lần tung đồng xu, số lần xuất hiện mặt ngửa gấp 3 lần số lần xuất hiện mặt sấp. Số lần xuất hiện mặt ngửa là
\choice
{$20$ lần}
{$25$ lần}
{$35$ lần}
{\True $30$ lần}
\loigiai{
Tổng số phần: $1 + 3 = 4$ phần.\\
Mặt ngửa: $40 : 4 \times 3 = 30$ lần.\\
\textbf{Đáp án của câu hỏi là D}
}
\end{ex}

\begin{ex} %[4X9C51]
Trong hộp kín có các viên bi xanh và đỏ. Bạn An lấy ngẫu nhiên 1 viên bi, ghi lại màu rồi trả lại hộp, lặp lại 100 lần. Kết quả: 70 lần lấy được bi đỏ, 30 lần lấy được bi xanh. Dự đoán nào sau đây là hợp lí nhất?
\choice
{Số bi xanh trong hộp nhiều hơn số bi đỏ}
{Trong hộp chỉ có toàn bi đỏ}
{\True Số bi đỏ trong hộp nhiều hơn số bi xanh}
{Trong hộp có số bi xanh bằng bi đỏ}
\loigiai{
Tần số xuất hiện bi đỏ cao hơn nhiều (70 lần so với 30 lần) nên số bi đỏ trong hộp nhiều hơn.\\
\textbf{Đáp án của câu hỏi là C}
}
\end{ex}

\textbf{Phần 2. Trắc nghiệm đúng sai}
\setcounterpageref{ex}{0}

\begin{ex} %[4X9H51]
Gieo một con xúc xắc 30 lần: mặt 1 chấm xuất hiện 5 lần, mặt 2 chấm xuất hiện 6 lần, mặt 3 chấm xuất hiện 4 lần, mặt 4 chấm xuất hiện 5 lần, mặt 5 chấm xuất hiện 7 lần, mặt 6 chấm xuất hiện 3 lần. Xét tính đúng sai:
\choiceTF
{\True Mặt 5 chấm xuất hiện nhiều lần nhất ($7$ lần).}
{\True Mặt 6 chấm xuất hiện ít lần nhất ($3$ lần).}
{Tổng số lần xuất hiện mặt có số chấm là số chẵn là $15$ lần.}
{\True Tổng số lần xuất hiện mặt có số chấm là số chẵn là $14$ lần ($6 + 5 + 3 = 14$).}
\loigiai{
\begin{itemize}
    \item Mặt 5 chấm: 7 lần (nhiều nhất). Đúng.
    \item Mặt 6 chấm: 3 lần (ít nhất). Đúng.
    \item Số chẵn (2, 4, 6): $6 + 5 + 3 = 14$ lần $\neq 15$ lần. Ý 3 Sai, Ý 4 Đúng.
\end{itemize}
\begin{itemize}
    \item a) \textbf{Đúng}.
    \item b) \textbf{Đúng}.
    \item c) \textbf{Sai}.
    \item d) \textbf{Đúng}.
\end{itemize}
}
\end{ex}

\begin{ex} %[4X9V51]
Trong hộp có 5 viên bi xanh, 3 viên bi đỏ và 2 viên bi vàng. Lấy ngẫu nhiên 1 viên bi. Xét tính đúng sai:
\choiceTF
{\True Khả năng lấy được bi xanh là cao nhất.}
{\True Khả năng lấy được bi vàng là thấp nhất.}
{Chắc chắn sẽ lấy được viên bi màu đỏ.}
{\True Có thể lấy được viên bi màu xanh, đỏ hoặc vàng.}
\loigiai{
\begin{itemize}
    \item 5 xanh (nhiều nhất): Đúng.
    \item 2 vàng (ít nhất): Đúng.
    \item Không chắc chắn lấy được bi đỏ (vì có thể lấy bi khác). Sai.
    \item Có thể ra xanh, đỏ, vàng: Đúng.
\end{itemize}
\begin{itemize}
    \item a) \textbf{Đúng}.
    \item b) \textbf{Đúng}.
    \item c) \textbf{Sai}.
    \item d) \textbf{Đúng}.
\end{itemize}
}
\end{ex}

\textbf{Phần 3. Trả lời ngắn}
\setcounterpageref{ex}{0}

\begin{ex} %[4X9H51]
Tung một đồng xu 25 lần, có 13 lần xuất hiện mặt ngửa. Hỏi mặt sấp xuất hiện bao nhiêu lần?
\shortans{12}
\loigiai{
\textbf{Đáp số:} $12$\\
$25 - 13 = 12$ lần.
}
\end{ex}

\begin{ex} %[4X9H51]
Gieo một con xúc xắc 40 lần, thấy mặt 1 chấm xuất hiện 8 lần. Hỏi các mặt còn lại xuất hiện tất cả bao nhiêu lần?
\shortans{32}
\loigiai{
\textbf{Đáp số:} $32$\\
$40 - 8 = 32$ lần.
}
\end{ex}

\begin{ex} %[4X9V51]
Quay vòng quay 50 lần, mũi tên chỉ vào ô màu đỏ 18 lần, ô màu xanh 22 lần, còn lại là ô màu vàng. Hỏi mũi tên chỉ vào ô màu vàng bao nhiêu lần?
\shortans{10}
\loigiai{
\textbf{Đáp số:} $10$\\
$50 - (18 + 22) = 10$ lần.
}
\end{ex}

\begin{ex} %[4X9V51]
Tung đồng xu 60 lần, số lần xuất hiện mặt ngửa nhiều hơn mặt sấp là 6 lần. Tính số lần xuất hiện mặt ngửa.
\shortans{33}
\loigiai{
\textbf{Đáp số:} $33$\\
$(60 + 6) : 2 = 33$ lần.
}
\end{ex}

\begin{ex} %[4X9V51]
Tung đồng xu 60 lần như trên, tính số lần xuất hiện mặt sấp.
\shortans{27}
\loigiai{
\textbf{Đáp số:} $27$\\
$60 - 33 = 27$ lần.
}
\end{ex}

\begin{ex} %[4X9C51]
Một hộp có 20 thẻ đánh số từ 1 đến 20. Rút ngẫu nhiên 1 thẻ. Có bao nhiêu khả năng rút được thẻ mang số chia hết cho 5?
\shortans{4}
\loigiai{
\textbf{Đáp số:} $4$\\
Các số: 5, 10, 15, 20 (có 4 thẻ).
}
\end{ex}

\newpage

\TaoPhieuKiemTra{Chủ đề 9 -- Làm quen với yếu tố thống kê, xác xuất}{52}{Luyện tập chung thống kê, xác suất}{https://loptoanthaylinh.vn/thi-online/toan-4/bai-52}
\textbf{Phần 1. Trắc nghiệm lựa chọn}
\setcounterpageref{ex}{0}

\begin{ex} %[4X9N52]
Dãy số liệu về chiều cao 4 bạn ($132\text{ cm}, 135\text{ cm}, 130\text{ cm}, 139\text{ cm}$) có giá trị lớn nhất là
\choice
{\True $139\text{ cm}$}
{$135\text{ cm}$}
{$132\text{ cm}$}
{$140\text{ cm}$}
\loigiai{
$139\text{ cm}$ lớn nhất.\\
\textbf{Đáp án của câu hỏi là A}
}
\end{ex}

\begin{ex} %[4X9N52]
Biểu đồ cột dùng để
\choice
{Vẽ tranh phong cảnh}
{Tính nhẩm phép nhân}
{Đo độ dài đoạn thẳng}
{\True Biểu diễn số liệu thống kê một cách trực quan}
\loigiai{
Biểu diễn số liệu thống kê trực quan.\\
\textbf{Đáp án của câu hỏi là D}
}
\end{ex}

\begin{ex} %[4X9H52]
Số trung bình cộng của dãy số liệu $12, 14, 16, 18$ là
\choice
{$14$}
{\True $15$}
{$16$}
{$17$}
\loigiai{
$(12 + 14 + 16 + 18) : 4 = 60 : 4 = 15$.\\
\textbf{Đáp án của câu hỏi là B}
}
\end{ex}

\begin{ex} %[4X9H52]
Tung đồng xu 20 lần có 12 lần ngửa. Số lần xuất hiện mặt sấp là
\choice
{$6$ lần}
{$10$ lần}
{$12$ lần}
{\True $8$ lần}
\loigiai{
$20 - 12 = 8$ lần.\\
\textbf{Đáp án của câu hỏi là D}
}
\end{ex}

\begin{ex} %[4X9H52]
Biểu đồ cột số điểm tốt của 3 tổ: Tổ 1 ($35$), Tổ 2 ($40$), Tổ 3 ($45$). Tổng số điểm tốt của cả 3 tổ là
\choice
{$110$}
{$130$}
{\True $120$}
{$115$}
\loigiai{
$35 + 40 + 45 = 120$.\\
\textbf{Đáp án của câu hỏi là C}
}
\end{ex}

\begin{ex} %[4X9H52]
Từ dữ liệu trên, trung bình mỗi tổ có bao nhiêu điểm tốt?
\choice
{$38$ điểm tốt}
{\True $40$ điểm tốt}
{$42$ điểm tốt}
{$45$ điểm tốt}
\loigiai{
$120 : 3 = 40$ điểm tốt.\\
\textbf{Đáp án của câu hỏi là B}
}
\end{ex}

\begin{ex} %[4X9V52]
Dãy số liệu số thóc 4 kho chứa ($25\text{ tấn}, 30\text{ tấn}, 35\text{ tấn}, 40\text{ tấn}$). Kho chứa nhiều nhất hơn kho chứa ít nhất bao nhiêu tấn thóc?
\choice
{\True $15\text{ tấn}$}
{$10\text{ tấn}$}
{$20\text{ tấn}$}
{$5\text{ tấn}$}
\loigiai{
$40 - 25 = 15$ tấn.\\
\textbf{Đáp án của câu hỏi là A}
}
\end{ex}

\begin{ex} %[4X9V52]
Trong một hộp có 4 bút bi xanh và 4 bút bi đỏ. Bạn Nam nhắm mắt lấy ra 1 chiếc bút. Khả năng lấy được bút bi xanh là
\choice
{Chắc chắn xảy ra}
{Không thể xảy ra}
{\True Bằng khả năng lấy được bút bi đỏ}
{Lớn hơn khả năng lấy bút bi đỏ}
\loigiai{
Số lượng 2 màu bằng nhau (4 chiếc) nên khả năng bằng nhau.\\
\textbf{Đáp án của câu hỏi là C}
}
\end{ex}

\begin{ex} %[4X9V52]
Gieo con xúc xắc 60 lần, mặt 6 chấm xuất hiện 12 lần. Số lần xuất hiện các mặt còn lại là
\choice
{$46$ lần}
{$50$ lần}
{$52$ lần}
{\True $48$ lần}
\loigiai{
$60 - 12 = 48$ lần.\\
\textbf{Đáp án của câu hỏi là D}
}
\end{ex}

\begin{ex} %[4X9V52]
Biểu đồ số học sinh 4 lớp: 4A ($35$), 4B ($38$), 4C ($34$), 4D ($37$). Tổng số học sinh của cả khối 4 là
\choice
{$140$ học sinh}
{$145$ học sinh}
{\True $144$ học sinh}
{$150$ học sinh}
\loigiai{
$35 + 38 + 34 + 37 = 144$ học sinh.\\
\textbf{Đáp án của câu hỏi là C}
}
\end{ex}

\begin{ex} %[4X9V52]
Trung bình mỗi lớp khối 4 có bao nhiêu học sinh?
\choice
{$35$ học sinh}
{\True $36$ học sinh}
{$37$ học sinh}
{$38$ học sinh}
\loigiai{
$144 : 4 = 36$ học sinh.\\
\textbf{Đáp án của câu hỏi là B}
}
\end{ex}

\begin{ex} %[4X9C52]
Một hộp có 10 quả bóng bàn gồm 3 màu: trắng, vàng, cam. Biết số bóng trắng gấp đôi số bóng vàng và số bóng cam là 1 quả. Bạn lấy ngẫu nhiên 1 quả bóng. Màu bóng có khả năng lấy được cao nhất là
\choice
{\True Bóng trắng (6 quả)}
{Bóng vàng}
{Bóng cam}
{Cả ba màu như nhau}
\loigiai{
Số bóng trắng $+$ vàng $= 10 - 1 = 9$ quả. Vàng: $9 : 3 = 3$ quả, Trắng: $3 \times 2 = 6$ quả.\\
Bóng trắng có 6 quả (nhiều nhất) nên khả năng lấy được cao nhất.\\
\textbf{Đáp án của câu hỏi là A}
}
\end{ex}

\textbf{Phần 2. Trắc nghiệm đúng sai}
\setcounterpageref{ex}{0}

\begin{ex} %[4X9H52]
Xét tính đúng sai về thống kê và xác suất:
\choiceTF
{\True Dãy số liệu thống kê giúp ta thu thập và xử lí thông tin số lượng.}
{\True Biểu đồ cột giúp so sánh các số liệu một cách trực quan.}
{Tung đồng xu 10 lần thì chắc chắn phải có 5 lần ngửa và 5 lần sấp.}
{\True Tung đồng xu 10 lần thì số lần ngửa và sấp có thể ngẫu nhiên khác nhau.}
\loigiai{
\begin{itemize}
    \item Dãy số liệu: Đúng.
    \item Biểu đồ cột trực quan: Đúng.
    \item Không chắc chắn đúng tỉ lệ 5-5 (xác suất là ngẫu nhiên). Sai.
    \item Số lần ngửa sấp là ngẫu nhiên: Đúng.
\end{itemize}
\begin{itemize}
    \item a) \textbf{Đúng}.
    \item b) \textbf{Đúng}.
    \item c) \textbf{Sai}.
    \item d) \textbf{Đúng}.
\end{itemize}
}
\end{ex}

\begin{ex} %[4X9V52]
Cho biểu đồ số ki-lô-gam giấy vụn thu gom của 3 lớp: 4A ($45\text{ kg}$), 4B ($50\text{ kg}$), 4C ($55\text{ kg}$). Xét tính đúng sai:
\choiceTF
{\True Lớp 4C thu gom được nhiều giấy vụn nhất.}
{\True Lớp 4A thu gom được ít giấy vụn nhất.}
{Cả 3 lớp thu gom được tất cả $160\text{ kg}$ giấy vụn.}
{\True Cả 3 lớp thu gom được tất cả $150\text{ kg}$ giấy vụn ($1{,}5$ tạ).}
\loigiai{
\begin{itemize}
    \item 4C: 55 kg (nhiều nhất). Đúng.
    \item 4A: 45 kg (ít nhất). Đúng.
    \item Tổng: $45 + 50 + 55 = 150\text{ kg} \neq 160\text{ kg}$. Ý 3 Sai, Ý 4 Đúng.
\end{itemize}
\begin{itemize}
    \item a) \textbf{Đúng}.
    \item b) \textbf{Đúng}.
    \item c) \textbf{Sai}.
    \item d) \textbf{Đúng}.
\end{itemize}
}
\end{ex}

\textbf{Phần 3. Trả lời ngắn}
\setcounterpageref{ex}{0}

\begin{ex} %[4X9H52]
Tìm trung bình cộng của 3 số liệu: $25, 35, 45$.
\shortans{35}
\loigiai{
\textbf{Đáp số:} $35$\\
$(25 + 35 + 45) : 3 = 35$.
}
\end{ex}

\begin{ex} %[4X9H52]
Tung đồng xu 30 lần có 16 lần ngửa. Số lần sấp là bao nhiêu?
\shortans{14}
\loigiai{
\textbf{Đáp số:} $14$\\
$30 - 16 = 14$ lần.
}
\end{ex}

\begin{ex} %[4X9V52]
Cho dãy số liệu điểm kiểm tra: $7, 8, 9, 8, 10, 8, 9$. Số điểm 8 xuất hiện bao nhiêu lần?
\shortans{3}
\loigiai{
\textbf{Đáp số:} $3$\\
3 lần.
}
\end{ex}

\begin{ex} %[4X9V52]
Biểu đồ số cây trồng của 3 tổ: $30, 40, 50$. Tính số cây trung bình mỗi tổ trồng được.
\shortans{40}
\loigiai{
\textbf{Đáp số:} $40$\\
$(30 + 40 + 50) : 3 = 40$ cây.
}
\end{ex}

\begin{ex} %[4X9V52]
Hiệu số cây giữa tổ trồng nhiều nhất (50) và tổ ít nhất (30) là bao nhiêu cây?
\shortans{20}
\loigiai{
\textbf{Đáp số:} $20$\\
$50 - 30 = 20$ cây.
}
\end{ex}

\begin{ex} %[4X9C52]
Một nhóm 4 học sinh có số đo chiều cao trung bình là $136\text{ cm}$. Khi có thêm bạn thứ 5 cao $141\text{ cm}$ tham gia nhóm thì chiều cao trung bình của cả 5 bạn là bao nhiêu xăng-ti-mét?
\shortans{137}
\loigiai{
\textbf{Đáp số:} $137$\\
Tổng 4 bạn: $136 \times 4 = 544\text{ cm}$.\\
Tổng 5 bạn: $544 + 141 = 685\text{ cm}$.\\
Trung bình 5 bạn: $685 : 5 = 137\text{ cm}$.
}
\end{ex}

\newpage

% =========================================================================
% CHỦ ĐỀ 10 - PHÂN SỐ
% =========================================================================
\section{Phân số}

\TaoPhieuKiemTra{Chủ đề 10 -- Phân số}{53}{Khái niệm phân số}{https://loptoanthaylinh.vn/thi-online/toan-4/bai-53}
\textbf{Phần 1. Trắc nghiệm lựa chọn}
\setcounterpageref{ex}{0}

\begin{ex} %[4D10N53]
Phân số chỉ phần đã tô màu gồm 3 phần bằng nhau trong tổng số 4 phần bằng nhau của hình tròn là
\choice
{$\dfrac{4}{3}$}
{$\dfrac{1}{4}$}
{\True $\dfrac{3}{4}$}
{$\dfrac{3}{7}$}
\loigiai{
Phân số chỉ phần đã tô màu là $\dfrac{3}{4}$.\\
\textbf{Đáp án của câu hỏi là C}
}
\end{ex}

\begin{ex} %[4D10N53]
Trong phân số $\dfrac{5}{8}$, tử số là
\choice
{$8$}
{$13$}
{$3$}
{\True $5$}
\loigiai{
Tử số là số viết trên gạch ngang: 5.\\
\textbf{Đáp án của câu hỏi là D}
}
\end{ex}

\begin{ex} %[4D10H53]
Trong phân số $\dfrac{7}{12}$, mẫu số là
\choice
{$7$}
{$19$}
{\True $12$}
{$5$}
\loigiai{
Mẫu số là số viết dưới gạch ngang: 12.\\
\textbf{Đáp án của câu hỏi là C}
}
\end{ex}

\begin{ex} %[4D10H53]
Phân số \textit{Bảy phần mười một} được viết là
\choice
{$\dfrac{11}{7}$}
{\True $\dfrac{7}{11}$}
{$7{,}11$}
{$11{,}7$}
\loigiai{
Viết là $\dfrac{7}{11}$.\\
\textbf{Đáp án của câu hỏi là B}
}
\end{ex}

\begin{ex} %[4D10H53]
Mẫu số của một phân số luôn phải thỏa mãn điều kiện nào?
\choice
{Là số tự nhiên lớn hơn 10}
{Bằng tử số}
{Nhỏ hơn tử số}
{\True Là số tự nhiên khác 0}
\loigiai{
Mẫu số phải là số tự nhiên khác 0.\\
\textbf{Đáp án của câu hỏi là D}
}
\end{ex}

\begin{ex} %[4D10H53]
Một cái bánh pizza được chia đều thành 8 miếng bằng nhau. Bạn An đã ăn 3 miếng. Phân số chỉ phần bánh còn lại là
\choice
{$\dfrac{3}{8}$}
{\True $\dfrac{5}{8}$}
{$\dfrac{8}{3}$}
{$\dfrac{5}{3}$}
\loigiai{
Phần còn lại: $8 - 3 = 5$ miếng $\Rightarrow \dfrac{5}{8}$ cái bánh.\\
\textbf{Đáp án của câu hỏi là B}
}
\end{ex}

\begin{ex} %[4D10V53]
Một hộp có 12 quả bóng gồm 5 bóng đỏ và 7 bóng xanh. Phân số chỉ số bóng đỏ so với tổng số bóng trong hộp là
\choice
{\True $\dfrac{5}{12}$}
{$\dfrac{7}{12}$}
{$\dfrac{5}{7}$}
{$\dfrac{7}{5}$}
\loigiai{
Phân số là $\dfrac{5}{12}$.\\
\textbf{Đáp án của câu hỏi là A}
}
\end{ex}

\begin{ex} %[4D10V53]
Một lớp học có 35 học sinh, trong đó có 18 học sinh nữ. Phân số chỉ số học sinh nam so với cả lớp là
\choice
{$\dfrac{18}{35}$}
{$\dfrac{17}{18}$}
{$\dfrac{18}{17}$}
{\True $\dfrac{17}{35}$}
\loigiai{
Số nam: $35 - 18 = 17$ học sinh $\Rightarrow \dfrac{17}{35}$.\\
\textbf{Đáp án của câu hỏi là D}
}
\end{ex}

\begin{ex} %[4D10V53]
Trên một đoạn thẳng $AB$ dài $10\text{ cm}$, điểm $M$ chia đoạn thẳng $AB$ thành 5 phần bằng nhau. Điểm $C$ nằm ở vạch chia thứ 2 tính từ $A$. Phân số chỉ độ dài đoạn $AC$ so với $AB$ là
\choice
{\True $\dfrac{2}{5}$}
{$\dfrac{3}{5}$}
{$\dfrac{2}{10}$}
{$\dfrac{1}{5}$}
\loigiai{
$AC$ chiếm 2 phần trong tổng 5 phần $\Rightarrow \dfrac{2}{5}$.\\
\textbf{Đáp án của câu hỏi là A}
}
\end{ex}

\begin{ex} %[4D10V53]
Viết phân số có tử số là số chẵn bé nhất có 2 chữ số và mẫu số là số lớn nhất có 2 chữ số.
\choice
{$\dfrac{12}{99}$}
{\True $\dfrac{10}{99}$}
{$\dfrac{10}{98}$}
{$\dfrac{2}{99}$}
\loigiai{
Tử số là 10, mẫu số là 99 $\Rightarrow \dfrac{10}{99}$.\\
\textbf{Đáp án của câu hỏi là B}
}
\end{ex}

\begin{ex} %[4D10V53]
Một can đựng đầy $10\text{ l}$ nước. Người ta đã dùng $7\text{ l}$ nước. Phân số chỉ số lít nước đã dùng là
\choice
{\True $\dfrac{7}{10}$}
{$\dfrac{3}{10}$}
{$\dfrac{7}{3}$}
{$\dfrac{10}{7}$}
\loigiai{
Đã dùng $\dfrac{7}{10}$ can nước.\\
\textbf{Đáp án của câu hỏi là A}
}
\end{ex}

\begin{ex} %[4D10C53]
Có bao nhiêu phân số có mẫu số bằng 7 mà tử số là các số tự nhiên khác 0 và nhỏ hơn mẫu số?
\choice
{$5$ phân số}
{$7$ phân số}
{\True $6$ phân số}
{$8$ phân số}
\loigiai{
Các tử số là $1, 2, 3, 4, 5, 6$ (có 6 phân số: $\dfrac{1}{7}, \dots, \dfrac{6}{7}$).\\
\textbf{Đáp án của câu hỏi là C}
}
\end{ex}

\textbf{Phần 2. Trắc nghiệm đúng sai}
\setcounterpageref{ex}{0}

\begin{ex} %[4D10H53]
Xét tính đúng sai về cấu tạo của phân số $\dfrac{4}{9}$:
\choiceTF
{\True Tử số là $4$, mẫu số là $9$.}
{\True Phân số này đọc là Bốn phần chín.}
{Mẫu số chỉ số phần bằng nhau được lấy đi.}
{\True Mẫu số chỉ tổng số phần bằng nhau được chia ra.}
\loigiai{
\begin{itemize}
    \item Tử số là 4, mẫu số là 9: Đúng.
    \item Đọc là Bốn phần chín: Đúng.
    \item Tử số mới là số phần lấy đi (không phải mẫu số). Sai.
    \item Mẫu số chỉ tổng số phần chia đều: Đúng.
\end{itemize}
\begin{itemize}
    \item a) \textbf{Đúng}.
    \item b) \textbf{Đúng}.
    \item c) \textbf{Sai}.
    \item d) \textbf{Đúng}.
\end{itemize}
}
\end{ex}

\begin{ex} %[4D10V53]
Một trang trại có 50 con gà, 30 con vịt và 20 con ngan. Xét tính đúng sai:
\choiceTF
{\True Tổng số gia cầm của trang trại là 100 con.}
{\True Phân số chỉ số gà là $\dfrac{50}{100}$.}
{Phân số chỉ số vịt là $\dfrac{30}{50}$.}
{\True Phân số chỉ số ngan là $\dfrac{20}{100}$.}
\loigiai{
\begin{itemize}
    \item Tổng: $50 + 30 + 20 = 100$ con. Đúng.
    \item Gà: $\dfrac{50}{100}$. Đúng.
    \item Vịt so với tổng: $\dfrac{30}{100} \neq \dfrac{30}{50}$. Sai.
    \item Ngan: $\dfrac{20}{100}$. Đúng.
\end{itemize}
\begin{itemize}
    \item a) \textbf{Đúng}.
    \item b) \textbf{Đúng}.
    \item c) \textbf{Sai}.
    \item d) \textbf{Đúng}.
\end{itemize}
}
\end{ex}

\textbf{Phần 3. Trả lời ngắn}
\setcounterpageref{ex}{0}

\begin{ex} %[4D10H53]
Tìm tử số của phân số $\dfrac{8}{15}$.
\shortans{8}
\loigiai{
\textbf{Đáp số:} $8$\\
Tử số là 8.
}
\end{ex}

\begin{ex} %[4D10H53]
Tìm mẫu số của phân số $\dfrac{11}{24}$.
\shortans{24}
\loigiai{
\textbf{Đáp số:} $24$\\
Mẫu số là 24.
}
\end{ex}

\begin{ex} %[4D10V53]
Một tuần lễ có 7 ngày. Phân số chỉ số ngày đi học trong tuần (từ thứ Hai đến thứ Sáu - 5 ngày) là bao nhiêu (viết dạng a/b)?
\shortans{5/7}
\loigiai{
\textbf{Đáp số:} $5/7$\\
$5/7$.
}
\end{ex}

\begin{ex} %[4D10V53]
Một hộp có 20 viên bi, trong đó có 9 viên bi đỏ, còn lại là bi xanh. Viết phân số chỉ số bi xanh so với tổng số bi (dạng a/b).
\shortans{11/20}
\loigiai{
\textbf{Đáp số:} $11/20$\\
Bi xanh: $20 - 9 = 11 \Rightarrow 11/20$.
}
\end{ex}

\begin{ex} %[4D10V53]
Viết phân số có tử số là số lẻ lớn nhất có 1 chữ số và mẫu số là 10 (dạng a/b).
\shortans{9/10}
\loigiai{
\textbf{Đáp số:} $9/10$\\
$9/10$.
}
\end{ex}

\begin{ex} %[4D10C53]
Có bao nhiêu phân số có tổng của tử số và mẫu số bằng 5 (với tử số và mẫu số là các số tự nhiên khác 0)?
\shortans{4}
\loigiai{
\textbf{Đáp số:} $4$\\
Các cặp (tử, mẫu): (1,4), (2,3), (3,2), (4,1) $\Rightarrow 4$ phân số.
}
\end{ex}

\newpage

\TaoPhieuKiemTra{Chủ đề 10 -- Phân số}{54}{Phân số và phép chia số tự nhiên}{https://loptoanthaylinh.vn/thi-online/toan-4/bai-54}
\textbf{Phần 1. Trắc nghiệm lựa chọn}
\setcounterpageref{ex}{0}

\begin{ex} %[4D10N54]
Thương của phép chia $3 : 4$ được viết dưới dạng phân số là
\choice
{$\dfrac{4}{3}$}
{$3{,}4$}
{$4{,}3$}
{\True $\dfrac{3}{4}$}
\loigiai{
$3 : 4 = \dfrac{3}{4}$.\\
\textbf{Đáp án của câu hỏi là D}
}
\end{ex}

\begin{ex} %[4D10N54]
Phân số $\dfrac{5}{1}$ có giá trị bằng
\choice
{$1$}
{\True $5$}
{$0$}
{$10$}
\loigiai{
$\dfrac{5}{1} = 5 : 1 = 5$.\\
\textbf{Đáp án của câu hỏi là B}
}
\end{ex}

\begin{ex} %[4D10H54]
Phân số $\dfrac{0}{8}$ có giá trị bằng
\choice
{$8$}
{$1$}
{\True $0$}
{Không xác định}
\loigiai{
$\dfrac{0}{8} = 0 : 8 = 0$.\\
\textbf{Đáp án của câu hỏi là C}
}
\end{ex}

\begin{ex} %[4D10H54]
Mỗi số tự nhiên $a$ đều có thể viết dưới dạng phân số có mẫu số bằng 1 là
\choice
{\True $\dfrac{a}{1}$}
{$\dfrac{1}{a}$}
{$\dfrac{a}{a}$}
{$\dfrac{0}{a}$}
\loigiai{
$a = \dfrac{a}{1}$.\\
\textbf{Đáp án của câu hỏi là A}
}
\end{ex}

\begin{ex} %[4D10H54]
Chia đều 5 cái bánh cho 4 bạn. Mỗi bạn nhận được số phần cái bánh là
\choice
{$\dfrac{4}{5}$ cái bánh}
{$1$ cái bánh}
{$2$ cái bánh}
{\True $\dfrac{5}{4}$ cái bánh}
\loigiai{
$5 : 4 = \dfrac{5}{4}$ cái bánh.\\
\textbf{Đáp án của câu hỏi là D}
}
\end{ex}

\begin{ex} %[4D10H54]
Phân số nào sau đây lớn hơn 1?
\choice
{$\dfrac{3}{5}$}
{$\dfrac{4}{4}$}
{\True $\dfrac{7}{4}$}
{$\dfrac{1}{2}$}
\loigiai{
Phân số có tử số lớn hơn mẫu số thì lớn hơn 1: $\dfrac{7}{4} > 1$.\\
\textbf{Đáp án của câu hỏi là C}
}
\end{ex}

\begin{ex} %[4D10V54]
Phân số nào sau đây bằng 1?
\choice
{\True $\dfrac{8}{8}$}
{$\dfrac{5}{6}$}
{$\dfrac{9}{7}$}
{$\dfrac{0}{8}$}
\loigiai{
Tử số bằng mẫu số thì phân số bằng 1: $\dfrac{8}{8} = 1$.\\
\textbf{Đáp án của câu hỏi là A}
}
\end{ex}

\begin{ex} %[4D10V54]
Có 3 lít nước rót đều vào 5 chai. Mỗi chai có số lít nước là
\choice
{$\dfrac{5}{3}\text{ l}$}
{$1\text{ l}$}
{\True $\dfrac{3}{5}\text{ l}$}
{$0{,}5\text{ l}$}
\loigiai{
$3 : 5 = \dfrac{3}{5}\text{ l}$.\\
\textbf{Đáp án của câu hỏi là C}
}
\end{ex}

\begin{ex} %[4D10V54]
Trong các phân số sau, phân số nào bé hơn 1?
\choice
{$\dfrac{4}{3}$}
{\True $\dfrac{5}{7}$}
{$\dfrac{9}{9}$}
{$\dfrac{11}{8}$}
\loigiai{
Tử số nhỏ hơn mẫu số: $\dfrac{5}{7} < 1$.\\
\textbf{Đáp án của câu hỏi là B}
}
\end{ex}

\begin{ex} %[4D10V54]
Viết thương $7 : 10$ dưới dạng phân số:
\choice
{$\dfrac{10}{7}$}
{$0{,}7$}
{$70$}
{\True $\dfrac{7}{10}$}
\loigiai{
$7 : 10 = \dfrac{7}{10}$.\\
\textbf{Đáp án của câu hỏi là D}
}
\end{ex}

\begin{ex} %[4D10V54]
Chia đều $12\text{ kg}$ gạo vào 5 túi. Mỗi túi có bao nhiêu ki-lô-gam gạo?
\choice
{$\dfrac{5}{12}\text{ kg}$}
{\True $\dfrac{12}{5}\text{ kg}$ ($2{,}4\text{ kg}$)}
{$2\text{ kg}$}
{$3\text{ kg}$}
\loigiai{
$12 : 5 = \dfrac{12}{5}\text{ kg}$.\\
\textbf{Đáp án của câu hỏi là B}
}
\end{ex}

\begin{ex} %[4D10C54]
Có bao nhiêu phân số có tử số và mẫu số là các số tự nhiên có một chữ số và có giá trị bằng 1?
\choice
{\True $9$ phân số}
{$8$ phân số}
{$10$ phân số}
{$7$ phân số}
\loigiai{
Các phân số là: $\dfrac{1}{1}, \dfrac{2}{2}, \dfrac{3}{3}, \dfrac{4}{4}, \dfrac{5}{5}, \dfrac{6}{6}, \dfrac{7}{7}, \dfrac{8}{8}, \dfrac{9}{9}$ (9 phân số, mẫu số khác 0).\\
\textbf{Đáp án của câu hỏi là A}
}
\end{ex}

\textbf{Phần 2. Trắc nghiệm đúng sai}
\setcounterpageref{ex}{0}

\begin{ex} %[4D10H54]
Xét tính đúng sai về phân số và phép chia số tự nhiên:
\choiceTF
{\True Thương của phép chia số tự nhiên cho số tự nhiên khác 0 có thể viết dưới dạng phân số.}
{\True Phân số có tử số lớn hơn mẫu số thì phân số đó lớn hơn 1.}
{Phân số $\dfrac{6}{6}$ bé hơn 1.}
{\True Số 0 có thể viết dưới dạng phân số là $\dfrac{0}{5}$ (hoặc mẫu số bất kì khác 0).}
\loigiai{
\begin{itemize}
    \item Thương viết thành phân số: Đúng.
    \item Tử $>$ mẫu $\Rightarrow > 1$: Đúng.
    \item $\dfrac{6}{6} = 1$ (không bé hơn 1). Sai.
    \item $0 = \dfrac{0}{5}$: Đúng.
\end{itemize}
\begin{itemize}
    \item a) \textbf{Đúng}.
    \item b) \textbf{Đúng}.
    \item c) \textbf{Sai}.
    \item d) \textbf{Đúng}.
\end{itemize}
}
\end{ex}

\begin{ex} %[4D10V54]
Cho các phân số: $\dfrac{3}{4}, \dfrac{5}{5}, \dfrac{7}{2}, \dfrac{9}{10}, \dfrac{12}{11}$. Xét tính đúng sai:
\choiceTF
{\True Có 2 phân số bé hơn 1 (là $\dfrac{3}{4}$ và $\dfrac{9}{10}$).}
{\True Có 1 phân số bằng 1 (là $\dfrac{5}{5}$).}
{Có 3 phân số lớn hơn 1.}
{\True Có 2 phân số lớn hơn 1 (là $\dfrac{7}{2}$ và $\dfrac{12}{11}$).}
\loigiai{
\begin{itemize}
    \item Bé hơn 1: $\dfrac{3}{4}, \dfrac{9}{10}$ (2 phân số). Đúng.
    \item Bằng 1: $\dfrac{5}{5}$ (1 phân số). Đúng.
    \item Lớn hơn 1: $\dfrac{7}{2}, \dfrac{12}{11}$ (chỉ có 2 phân số, không phải 3). Ý 3 Sai, Ý 4 Đúng.
\end{itemize}
\begin{itemize}
    \item a) \textbf{Đúng}.
    \item b) \textbf{Đúng}.
    \item c) \textbf{Sai}.
    \item d) \textbf{Đúng}.
\end{itemize}
}
\end{ex}

\textbf{Phần 3. Trả lời ngắn}
\setcounterpageref{ex}{0}

\begin{ex} %[4D10H54]
Viết thương $4 : 7$ dưới dạng phân số (dạng a/b).
\shortans{4/7}
\loigiai{
\textbf{Đáp số:} $4/7$\\
$4/7$.
}
\end{ex}

\begin{ex} %[4D10H54]
Viết số tự nhiên 9 dưới dạng phân số có mẫu số bằng 1 (dạng a/b).
\shortans{9/1}
\loigiai{
\textbf{Đáp số:} $9/1$\\
$9/1$.
}
\end{ex}

\begin{ex} %[4D10V54]
Chia 7 quả cam cho 4 người. Viết số cam mỗi người nhận được dưới dạng phân số (dạng a/b).
\shortans{7/4}
\loigiai{
\textbf{Đáp số:} $7/4$\\
$7/4$ quả cam.
}
\end{ex}

\begin{ex} %[4D10V54]
Có bao nhiêu phân số nhỏ hơn 1 có mẫu số bằng 5 (tử số là số tự nhiên khác 0)?
\shortans{4}
\loigiai{
\textbf{Đáp số:} $4$\\
$1/5, 2/5, 3/5, 4/5$ (4 phân số).
}
\end{ex}

\begin{ex} %[4D10V54]
Tìm số tự nhiên $x$ biết $\dfrac{x}{8} = 1$.
\shortans{8}
\loigiai{
\textbf{Đáp số:} $8$\\
$x = 8$.
}
\end{ex}

\begin{ex} %[4D10C54]
Có 5 đoạn dây, mỗi đoạn dài $2\text{ m}$. Người ta nối lại rồi chia đều thành 8 phần bằng nhau. Viết độ dài mỗi phần theo mét dưới dạng phân số rút gọn (dạng a/b).
\shortans{5/4}
\loigiai{
\textbf{Đáp số:} $5/4$\\
Tổng chiều dài: $5 \times 2 = 10\text{ m}$. Mỗi phần: $10 : 8 = 10/8 = 5/4\text{ m}$.
}
\end{ex}

\newpage

\TaoPhieuKiemTra{Chủ đề 10 -- Phân số}{55}{Tính chất cơ bản của phân số}{https://loptoanthaylinh.vn/thi-online/toan-4/bai-55}
\textbf{Phần 1. Trắc nghiệm lựa chọn}
\setcounterpageref{ex}{0}

\begin{ex} %[4D10N55]
Khi nhân cả tử số và mẫu số của một phân số với cùng một số tự nhiên khác 0 thì ta được
\choice
{\True Phân số bằng phân số ban đầu}
{Phân số lớn hơn phân số ban đầu}
{Phân số nhỏ hơn phân số ban đầu}
{Số tự nhiên}
\loigiai{
Được một phân số bằng phân số đã cho.\\
\textbf{Đáp án của câu hỏi là A}
}
\end{ex}

\begin{ex} %[4D10N55]
Khi chia cả tử số và mẫu số của một phân số cho cùng một số tự nhiên khác 0 (và cùng chia hết) ta được
\choice
{Phân số gấp đôi}
{Phân số nhỏ hơn}
{\True Phân số bằng phân số ban đầu}
{Phân số nghịch đảo}
\loigiai{
Được một phân số bằng phân số đã cho.\\
\textbf{Đáp án của câu hỏi là C}
}
\end{ex}

\begin{ex} %[4D10H55]
Điền số thích hợp vào chỗ chấm: $\dfrac{2}{3} = \dfrac{\dots}{9}$.
\choice
{$4$}
{$8$}
{$5$}
{\True $6$}
\loigiai{
Nhân cả tử và mẫu với 3: $2 \times 3 = 6$.\\
\textbf{Đáp án của câu hỏi là D}
}
\end{ex}

\begin{ex} %[4D10H55]
Điền số thích hợp vào chỗ chấm: $\dfrac{15}{20} = \dfrac{3}{\dots}$.
\choice
{\True $4$}
{$3$}
{$5$}
{$6$}
\loigiai{
Chia cả tử và mẫu cho 5: $20 : 5 = 4$.\\
\textbf{Đáp án của câu hỏi là A}
}
\end{ex}

\begin{ex} %[4D10H55]
Phân số nào sau đây bằng phân số $\dfrac{3}{5}$?
\choice
{$\dfrac{6}{15}$}
{\True $\dfrac{9}{15}$}
{$\dfrac{5}{3}$}
{$\dfrac{6}{8}$}
\loigiai{
$\dfrac{3 \times 3}{5 \times 3} = \dfrac{9}{15}$.\\
\textbf{Đáp án của câu hỏi là B}
}
\end{ex}

\begin{ex} %[4D10H55]
Phân số nào sau đây bằng phân số $\dfrac{12}{18}$?
\choice
{$\dfrac{3}{4}$}
{$\dfrac{4}{5}$}
{$\dfrac{1}{2}$}
{\True $\dfrac{2}{3}$}
\loigiai{
$\dfrac{12 : 6}{18 : 6} = \dfrac{2}{3}$.\\
\textbf{Đáp án của câu hỏi là D}
}
\end{ex}

\begin{ex} %[4D10V55]
Tìm $x$ biết: $\dfrac{x}{8} = \dfrac{15}{24}$.
\choice
{$3$}
{$6$}
{\True $5$}
{$4$}
\loigiai{
$24 : 3 = 8 \Rightarrow x = 15 : 3 = 5$.\\
\textbf{Đáp án của câu hỏi là C}
}
\end{ex}

\begin{ex} %[4D10V55]
Tìm $y$ biết: $\dfrac{4}{y} = \dfrac{28}{35}$.
\choice
{$4$}
{$6$}
{$7$}
{\True $5$}
\loigiai{
$28 : 7 = 4 \Rightarrow y = 35 : 7 = 5$.\\
\textbf{Đáp án của câu hỏi là D}
}
\end{ex}

\begin{ex} %[4D10V55]
Trong các phân số sau, phân số nào KHÔNG bằng $\dfrac{4}{6}$?
\choice
{$\dfrac{2}{3}$}
{\True $\dfrac{8}{10}$}
{$\dfrac{8}{12}$}
{$\dfrac{12}{18}$}
\loigiai{
$\dfrac{8}{10} = \dfrac{4}{5} \neq \dfrac{4}{6}$.\\
\textbf{Đáp án của câu hỏi là B}
}
\end{ex}

\begin{ex} %[4D10V55]
Viết phân số $\dfrac{3}{4}$ thành phân số có mẫu số bằng 100:
\choice
{\True $\dfrac{75}{100}$}
{$\dfrac{30}{100}$}
{$\dfrac{80}{100}$}
{$\dfrac{70}{100}$}
\loigiai{
$\dfrac{3 \times 25}{4 \times 25} = \dfrac{75}{100}$.\\
\textbf{Đáp án của câu hỏi là A}
}
\end{ex}

\begin{ex} %[4D10V55]
Viết phân số $\dfrac{1}{5}$ thành phân số có mẫu số là 20:
\choice
{$\dfrac{2}{20}$}
{$\dfrac{5}{20}$}
{\True $\dfrac{4}{20}$}
{$\dfrac{3}{20}$}
\loigiai{
$\dfrac{1 \times 4}{5 \times 4} = \dfrac{4}{20}$.\\
\textbf{Đáp án của câu hỏi là C}
}
\end{ex}

\begin{ex} %[4D10C55]
Cho phân số $\dfrac{7}{13}$. Hỏi phải cộng thêm vào cả tử số và mẫu số cùng một số tự nhiên nào để được phân số bằng $\dfrac{2}{3}$?
\choice
{$3$}
{\True $5$}
{$4$}
{$6$}
\loigiai{
Hiệu mẫu và tử: $13 - 7 = 6$.\\
Phân số mới có hiệu mẫu và tử là $3 - 2 = 1$ phần $= 6$.\\
Tử số mới: $2 \times 6 = 12 \Rightarrow$ Số cần cộng: $12 - 7 = 5$.\\
\textbf{Đáp án của câu hỏi là B}
}
\end{ex}

\textbf{Phần 2. Trắc nghiệm đúng sai}
\setcounterpageref{ex}{0}

\begin{ex} %[4D10H55]
Xét tính đúng sai của các khẳng định về tính chất cơ bản phân số:
\choiceTF
{\True $\dfrac{a}{b} = \dfrac{a \times m}{b \times m}$ ($m$ là số tự nhiên khác 0).}
{\True $\dfrac{a}{b} = \dfrac{a : n}{b : n}$ ($n$ là ước chung khác 0 của $a$ và $b$).}
{$\dfrac{a}{b} = \dfrac{a + m}{b + m}$.}
{\True $\dfrac{3}{4} = \dfrac{6}{8} = \dfrac{9}{12}$.}
\loigiai{
\begin{itemize}
    \item Nhân cùng 1 số: Đúng.
    \item Chia cùng 1 số: Đúng.
    \item Cộng cùng 1 số thì phân số thay đổi giá trị (không bằng nhau). Sai.
    \item Dãy phân số bằng nhau đúng.
\end{itemize}
\begin{itemize}
    \item a) \textbf{Đúng}.
    \item b) \textbf{Đúng}.
    \item c) \textbf{Sai}.
    \item d) \textbf{Đúng}.
\end{itemize}
}
\end{ex}

\begin{ex} %[4D10V55]
Cho phân số $\dfrac{18}{24}$. Xét tính đúng sai:
\choiceTF
{\True Rút gọn chia cho 6 được phân số $\dfrac{3}{4}$.}
{\True Phân số $\dfrac{18}{24}$ bằng phân số $\dfrac{9}{12}$.}
{Phân số $\dfrac{18}{24}$ bằng phân số $\dfrac{6}{7}$.}
{\True Nhân cả tử và mẫu với 2 được phân số $\dfrac{36}{48}$.}
\loigiai{
\begin{itemize}
    \item $18:6 / 24:6 = 3/4$. Đúng.
    \item $18:2 / 24:2 = 9/12$. Đúng.
    \item $3/4 \neq 6/7$. Sai.
    \item $18\times 2 / 24\times 2 = 36/48$. Đúng.
\end{itemize}
\begin{itemize}
    \item a) \textbf{Đúng}.
    \item b) \textbf{Đúng}.
    \item c) \textbf{Sai}.
    \item d) \textbf{Đúng}.
\end{itemize}
}
\end{ex}

\textbf{Phần 3. Trả lời ngắn}
\setcounterpageref{ex}{0}

\begin{ex} %[4D10H55]
Điền số thích hợp vào chỗ chấm: $\dfrac{3}{5} = \dfrac{15}{\dots}$.
\shortans{25}
\loigiai{
\textbf{Đáp số:} $25$\\
$5 \times 5 = 25$.
}
\end{ex}

\begin{ex} %[4D10H55]
Điền số thích hợp vào chỗ chấm: $\dfrac{24}{36} = \dfrac{\dots}{3}$.
\shortans{2}
\loigiai{
\textbf{Đáp số:} $2$\\
$36 : 12 = 3 \Rightarrow 24 : 12 = 2$.
}
\end{ex}

\begin{ex} %[4D10V55]
Tìm $x$ biết: $\dfrac{x}{12} = \dfrac{3}{4}$.
\shortans{9}
\loigiai{
\textbf{Đáp số:} $9$\\
$x = 3 \times 3 = 9$.
}
\end{ex}

\begin{ex} %[4D10V55]
Tìm $y$ biết: $\dfrac{35}{y} = \dfrac{5}{7}$.
\shortans{49}
\loigiai{
\textbf{Đáp số:} $49$\\
$y = 7 \times 7 = 49$.
}
\end{ex}

\begin{ex} %[4D10V55]
Viết phân số $\dfrac{2}{5}$ thành phân số có mẫu số bằng 50. Tìm tử số của phân số mới.
\shortans{20}
\loigiai{
\textbf{Đáp số:} $20$\\
$2 \times 10 = 20$.
}
\end{ex}

\begin{ex} %[4D10C55]
Tìm một phân số bằng $\dfrac{3}{4}$ biết tổng của tử số và mẫu số bằng 35. Tìm tử số của phân số đó.
\shortans{15}
\loigiai{
\textbf{Đáp số:} $15$\\
Tổng số phần: $3 + 4 = 7$ phần. Giá trị 1 phần: $35 : 7 = 5$.\\
Tử số là: $5 \times 3 = 15$.
}
\end{ex}

\newpage

\TaoPhieuKiemTra{Chủ đề 10 -- Phân số}{56}{Rút gọn phân số}{https://loptoanthaylinh.vn/thi-online/toan-4/bai-56}
\textbf{Phần 1. Trắc nghiệm lựa chọn}
\setcounterpageref{ex}{0}

\begin{ex} %[4D10N56]
Rút gọn phân số là
\choice
{Nhân cả tử và mẫu với một số}
{\True Chia cả tử số và mẫu số cho cùng một số tự nhiên lớn hơn 1 để được phân số đơn giản hơn và bằng phân số đã cho}
{Cộng thêm vào tử và mẫu}
{Đổi chỗ tử số và mẫu số}
\loigiai{
Rút gọn phân số là chia cả tử và mẫu cho ước chung lớn hơn 1.\\
\textbf{Đáp án của câu hỏi là B}
}
\end{ex}

\begin{ex} %[4D10N56]
Phân số tối giản là phân số
\choice
{\True Tử số và mẫu số không cùng chia hết cho một số tự nhiên nào lớn hơn 1}
{Có tử số bằng 1}
{Có mẫu số bằng 10}
{Có tử số nhỏ hơn 10}
\loigiai{
Phân số tối giản không thể rút gọn được nữa.\\
\textbf{Đáp án của câu hỏi là A}
}
\end{ex}

\begin{ex} %[4D10H56]
Rút gọn phân số $\dfrac{6}{8}$ về dạng tối giản ta được
\choice
{$\dfrac{2}{4}$}
{\True $\dfrac{3}{4}$}
{$\dfrac{3}{2}$}
{$\dfrac{1}{2}$}
\loigiai{
$\dfrac{6 : 2}{8 : 2} = \dfrac{3}{4}$.\\
\textbf{Đáp án của câu hỏi là B}
}
\end{ex}

\begin{ex} %[4D10H56]
Trong các phân số sau, phân số nào là phân số tối giản?
\choice
{$\dfrac{4}{6}$}
{$\dfrac{9}{12}$}
{$\dfrac{15}{20}$}
{\True $\dfrac{5}{7}$}
\loigiai{
$\dfrac{5}{7}$ là phân số tối giản.\\
\textbf{Đáp án của câu hỏi là D}
}
\end{ex}

\begin{ex} %[4D10H56]
Rút gọn phân số $\dfrac{18}{27}$ ta được phân số tối giản là
\choice
{$\dfrac{6}{9}$}
{$\dfrac{3}{4}$}
{\True $\dfrac{2}{3}$}
{$\dfrac{1}{3}$}
\loigiai{
$\dfrac{18 : 9}{27 : 9} = \dfrac{2}{3}$.\\
\textbf{Đáp án của câu hỏi là C}
}
\end{ex}

\begin{ex} %[4D10H56]
Phân số nào sau đây CHƯA tối giản?
\choice
{\True $\dfrac{14}{21}$}
{$\dfrac{3}{5}$}
{$\dfrac{7}{9}$}
{$\dfrac{11}{13}$}
\loigiai{
$\dfrac{14}{21}$ cùng chia hết cho 7 (rút gọn được thành $\dfrac{2}{3}$).\\
\textbf{Đáp án của câu hỏi là A}
}
\end{ex}

\begin{ex} %[4D10V56]
Rút gọn phân số $\dfrac{45}{60}$ về tối giản ta được
\choice
{$\dfrac{9}{12}$}
{$\dfrac{15}{20}$}
{$\dfrac{4}{5}$}
{\True $\dfrac{3}{4}$}
\loigiai{
$\dfrac{45 : 15}{60 : 15} = \dfrac{3}{4}$.\\
\textbf{Đáp án của câu hỏi là D}
}
\end{ex}

\begin{ex} %[4D10V56]
Rút gọn phân số $\dfrac{75}{100}$ ta được
\choice
{$\dfrac{15}{20}$}
{$\dfrac{7}{10}$}
{\True $\dfrac{3}{4}$}
{$\dfrac{4}{5}$}
\loigiai{
$\dfrac{75 : 25}{100 : 25} = \dfrac{3}{4}$.\\
\textbf{Đáp án của câu hỏi là C}
}
\end{ex}

\begin{ex} %[4D10V56]
Cho phân số $\dfrac{36}{48}$. Số tự nhiên lớn nhất có thể dùng để rút gọn ngay phân số này về tối giản là
\choice
{\True $12$}
{$6$}
{$4$}
{$8$}
\loigiai{
Ước chung lớn nhất của 36 và 48 là 12.\\
\textbf{Đáp án của câu hỏi là A}
}
\end{ex}

\begin{ex} %[4D10V56]
Trong các phân số $\dfrac{12}{15}, \dfrac{7}{9}, \dfrac{20}{25}, \dfrac{11}{14}$, có bao nhiêu phân số tối giản?
\choice
{$1$ phân số}
{\True $2$ phân số}
{$3$ phân số}
{$4$ phân số}
\loigiai{
Các phân số tối giản là: $\dfrac{7}{9}$ và $\dfrac{11}{14}$ (có 2 phân số).\\
\textbf{Đáp án của câu hỏi là B}
}
\end{ex}

\begin{ex} %[4D10V56]
Rút gọn phân số $\dfrac{120}{150}$ về dạng tối giản:
\choice
{$\dfrac{12}{15}$}
{$\dfrac{6}{7}$}
{\True $\dfrac{4}{5}$}
{$\dfrac{2}{3}$}
\loigiai{
$\dfrac{120 : 30}{150 : 30} = \dfrac{4}{5}$.\\
\textbf{Đáp án của câu hỏi là C}
}
\end{ex}

\begin{ex} %[4D10C56]
Tính giá trị biểu thức sau bằng cách rút gọn: $\dfrac{2 \times 3 \times 5 \times 7}{3 \times 5 \times 7 \times 8}$.
\choice
{$\dfrac{1}{2}$}
{$\dfrac{2}{8}$}
{$1$}
{\True $\dfrac{1}{4}$}
\loigiai{
Cùng chia tử và mẫu cho $3 \times 5 \times 7$: được $\dfrac{2}{8} = \dfrac{1}{4}$.\\
\textbf{Đáp án của câu hỏi là D}
}
\end{ex}

\textbf{Phần 2. Trắc nghiệm đúng sai}
\setcounterpageref{ex}{0}

\begin{ex} %[4D10H56]
Xét tính đúng sai về rút gọn phân số:
\choiceTF
{\True $\dfrac{8}{12}$ rút gọn được thành $\dfrac{2}{3}$.}
{\True $\dfrac{13}{17}$ là phân số tối giản.}
{Phân số $\dfrac{15}{35}$ là phân số tối giản.}
{\True $\dfrac{15}{35}$ rút gọn được thành $\dfrac{3}{7}$.}
\loigiai{
\begin{itemize}
    \item $8/12 = 2/3$. Đúng.
    \item $13/17$ tối giản: Đúng.
    \item $15/35$ chia hết cho 5 nên chưa tối giản. Sai.
    \item $15/35 = 3/7$. Đúng.
\end{itemize}
\begin{itemize}
    \item a) \textbf{Đúng}.
    \item b) \textbf{Đúng}.
    \item c) \textbf{Sai}.
    \item d) \textbf{Đúng}.
\end{itemize}
}
\end{ex}

\begin{ex} %[4D10V56]
Cho phân số $\dfrac{40}{60}$. Xét tính đúng sai:
\choiceTF
{\True Chia cả tử và mẫu cho 10 được $\dfrac{4}{6}$.}
{\True Rút gọn về phân số tối giản được $\dfrac{2}{3}$.}
{Rút gọn về phân số tối giản được $\dfrac{4}{6}$.}
{\True Phân số $\dfrac{40}{60}$ bằng phân số $\dfrac{10}{15}$.}
\loigiai{
\begin{itemize}
    \item Chia cho 10 được 4/6: Đúng.
    \item Tối giản là 2/3: Đúng.
    \item 4/6 chưa phải tối giản. Sai.
    \item $40/60 = 10/15 = 2/3$. Đúng.
\end{itemize}
\begin{itemize}
    \item a) \textbf{Đúng}.
    \item b) \textbf{Đúng}.
    \item c) \textbf{Sai}.
    \item d) \textbf{Đúng}.
\end{itemize}
}
\end{ex}

\textbf{Phần 3. Trả lời ngắn}
\setcounterpageref{ex}{0}

\begin{ex} %[4D10H56]
Rút gọn phân số $\dfrac{10}{15}$ về tối giản (dạng a/b).
\shortans{2/3}
\loigiai{
\textbf{Đáp số:} $2/3$\\
$2/3$.
}
\end{ex}

\begin{ex} %[4D10H56]
Rút gọn phân số $\dfrac{24}{32}$ về tối giản (dạng a/b).
\shortans{3/4}
\loigiai{
\textbf{Đáp số:} $3/4$\\
$3/4$.
}
\end{ex}

\begin{ex} %[4D10V56]
Rút gọn phân số $\dfrac{42}{56}$ về tối giản (dạng a/b).
\shortans{3/4}
\loigiai{
\textbf{Đáp số:} $3/4$\\
$42 : 14 / 56 : 14 = 3/4$.
}
\end{ex}

\begin{ex} %[4D10V56]
Rút gọn phân số $\dfrac{50}{125}$ về tối giản (dạng a/b).
\shortans{2/5}
\loigiai{
\textbf{Đáp số:} $2/5$\\
$50 : 25 / 125 : 25 = 2/5$.
}
\end{ex}

\begin{ex} %[4D10V56]
Tính giá trị rút gọn của phân số $\dfrac{16 \times 5}{25 \times 4}$ (dạng a/b).
\shortans{4/5}
\loigiai{
\textbf{Đáp số:} $4/5$\\
$(16 : 4) / (25 : 5) = 4/5$.
}
\end{ex}

\begin{ex} %[4D10C56]
Rút gọn biểu thức: $\dfrac{12 \times 15 \times 18}{24 \times 30 \times 36}$ về phân số tối giản (dạng a/b).
\shortans{1/8}
\loigiai{
\textbf{Đáp số:} $1/8$\\
$(1/2) \times (1/2) \times (1/2) = 1/8$.
}
\end{ex}

\newpage

\TaoPhieuKiemTra{Chủ đề 10 -- Phân số}{57}{Quy đồng mẫu số các phân số}{https://loptoanthaylinh.vn/thi-online/toan-4/bai-57}
\textbf{Phần 1. Trắc nghiệm lựa chọn}
\setcounterpageref{ex}{0}

\begin{ex} %[4D10N57]
Quy đồng mẫu số hai phân số là
\choice
{Cộng mẫu số của hai phân số lại}
{Đổi chỗ hai mẫu số}
{Lấy tử số cộng mẫu số}
{\True Biến đổi hai phân số thành hai phân số mới có cùng mẫu số và lần lượt bằng các phân số đã cho}
\loigiai{
Quy đồng là đưa các phân số về cùng mẫu số.\\
\textbf{Đáp án của câu hỏi là D}
}
\end{ex}

\begin{ex} %[4D10N57]
Mẫu số chung của hai phân số $\dfrac{1}{3}$ và $\dfrac{2}{5}$ là
\choice
{\True $15$}
{$8$}
{$12$}
{$10$}
\loigiai{
$3 \times 5 = 15$.\\
\textbf{Đáp án của câu hỏi là A}
}
\end{ex}

\begin{ex} %[4D10H57]
Quy đồng mẫu số hai phân số $\dfrac{2}{3}$ và $\dfrac{3}{4}$ với mẫu số chung là 12 ta được
\choice
{$\dfrac{6}{12}$ và $\dfrac{9}{12}$}
{\True $\dfrac{8}{12}$ và $\dfrac{9}{12}$}
{$\dfrac{8}{12}$ và $\dfrac{6}{12}$}
{$\dfrac{4}{12}$ và $\dfrac{9}{12}$}
\loigiai{
$\dfrac{2 \times 4}{3 \times 4} = \dfrac{8}{12}$ và $\dfrac{3 \times 3}{4 \times 3} = \dfrac{9}{12}$.\\
\textbf{Đáp án của câu hỏi là B}
}
\end{ex}

\begin{ex} %[4D10H57]
Khi quy đồng mẫu số hai phân số $\dfrac{5}{6}$ và $\dfrac{7}{12}$, mẫu số chung nhỏ nhất thuận tiện là
\choice
{$72$}
{$18$}
{\True $12$}
{$24$}
\loigiai{
Vì 12 chia hết cho 6 nên chọn 12 làm MSC.\\
\textbf{Đáp án của câu hỏi là C}
}
\end{ex}

\begin{ex} %[4D10H57]
Quy đồng mẫu số hai phân số $\dfrac{3}{5}$ và $\dfrac{4}{7}$ ta được hai phân số là
\choice
{$\dfrac{21}{35}$ và $\dfrac{15}{35}$}
{$\dfrac{15}{35}$ và $\dfrac{20}{35}$}
{$\dfrac{21}{35}$ và $\dfrac{28}{35}$}
{\True $\dfrac{21}{35}$ và $\dfrac{20}{35}$}
\loigiai{
$\dfrac{3 \times 7}{5 \times 7} = \dfrac{21}{35}$ và $\dfrac{4 \times 5}{7 \times 5} = \dfrac{20}{35}$.\\
\textbf{Đáp án của câu hỏi là D}
}
\end{ex}

\begin{ex} %[4D10H57]
Mẫu số chung của hai phân số $\dfrac{1}{4}$ và $\dfrac{5}{8}$ là
\choice
{$12$}
{\True $8$}
{$32$}
{$16$}
\loigiai{
8 chia hết cho 4 nên MSC là 8.\\
\textbf{Đáp án của câu hỏi là B}
}
\end{ex}

\begin{ex} %[4D10V57]
Quy đồng mẫu số ba phân số $\dfrac{1}{2}, \dfrac{2}{3}, \dfrac{3}{4}$ với mẫu số chung bằng 12 ta được
\choice
{$\dfrac{6}{12}, \dfrac{4}{12}, \dfrac{9}{12}$}
{$\dfrac{3}{12}, \dfrac{8}{12}, \dfrac{9}{12}$}
{\True $\dfrac{6}{12}, \dfrac{8}{12}, \dfrac{9}{12}$}
{$\dfrac{6}{12}, \dfrac{8}{12}, \dfrac{6}{12}$}
\loigiai{
$\dfrac{6}{12}, \dfrac{8}{12}, \dfrac{9}{12}$.\\
\textbf{Đáp án của câu hỏi là C}
}
\end{ex}

\begin{ex} %[4D10V57]
Quy đồng mẫu số hai phân số $\dfrac{7}{10}$ và $\dfrac{4}{15}$ với mẫu số chung nhỏ nhất là 30 ta được
\choice
{\True $\dfrac{21}{30}$ và $\dfrac{8}{30}$}
{$\dfrac{14}{30}$ và $\dfrac{8}{30}$}
{$\dfrac{21}{30}$ và $\dfrac{12}{30}$}
{$\dfrac{70}{30}$ và $\dfrac{40}{30}$}
\loigiai{
$\dfrac{7 \times 3}{10 \times 3} = \dfrac{21}{30}$ và $\dfrac{4 \times 2}{15 \times 2} = \dfrac{8}{30}$.\\
\textbf{Đáp án của câu hỏi là A}
}
\end{ex}

\begin{ex} %[4D10V57]
Cho hai phân số $\dfrac{5}{8}$ và $\dfrac{3}{4}$. Sau khi quy đồng mẫu số chung là 8 thì tổng hai tử số mới là
\choice
{$10$}
{$12$}
{\True $11$}
{$13$}
\loigiai{
$\dfrac{5}{8}$ và $\dfrac{6}{8}$. Tổng hai tử số: $5 + 6 = 11$.\\
\textbf{Đáp án của câu hỏi là C}
}
\end{ex}

\begin{ex} %[4D10V57]
Quy đồng mẫu số của số 2 và phân số $\dfrac{3}{5}$ ta được
\choice
{\True $\dfrac{10}{5}$ và $\dfrac{3}{5}$}
{$\dfrac{5}{5}$ và $\dfrac{3}{5}$}
{$\dfrac{2}{5}$ và $\dfrac{3}{5}$}
{$\dfrac{6}{5}$ và $\dfrac{3}{5}$}
\loigiai{
$2 = \dfrac{2}{1} = \dfrac{10}{5}$.\\
\textbf{Đáp án của câu hỏi là A}
}
\end{ex}

\begin{ex} %[4D10V57]
Tìm mẫu số chung nhỏ nhất của hai phân số $\dfrac{5}{12}$ và $\dfrac{7}{18}$:
\choice
{$216$}
{$72$}
{$48$}
{\True $36$}
\loigiai{
$36$ chia hết cho cả 12 và 18.\\
\textbf{Đáp án của câu hỏi là D}
}
\end{ex}

\begin{ex} %[4D10C57]
Quy đồng mẫu số ba phân số $\dfrac{1}{3}, \dfrac{1}{4}, \dfrac{1}{5}$ được ba phân số có mẫu số chung nhỏ nhất là 60. Tổng các tử số mới là
\choice
{$45$}
{\True $47$}
{$49$}
{$51$}
\loigiai{
Ba phân số: $\dfrac{20}{60}, \dfrac{15}{60}, \dfrac{12}{60}$. Tổng tử số: $20 + 15 + 12 = 47$.\\
\textbf{Đáp án của câu hỏi là B}
}
\end{ex}

\textbf{Phần 2. Trắc nghiệm đúng sai}
\setcounterpageref{ex}{0}

\begin{ex} %[4D10H57]
Xét tính đúng sai khi quy đồng mẫu số hai phân số $\dfrac{3}{4}$ và $\dfrac{5}{6}$:
\choiceTF
{\True Có thể chọn mẫu số chung là $24$ ($4 \times 6$).}
{\True Có thể chọn mẫu số chung nhỏ nhất là $12$.}
{Khi quy đồng với MSC là $12$ ta được $\dfrac{6}{12}$ và $\dfrac{10}{12}$.}
{\True Khi quy đồng với MSC là $12$ ta được $\dfrac{9}{12}$ và $\dfrac{10}{12}$.}
\loigiai{
\begin{itemize}
    \item MSC 24: Đúng.
    \item MSC nhỏ nhất 12: Đúng.
    \item $3/4 = 9/12 \neq 6/12$. Ý 3 Sai, Ý 4 Đúng.
\end{itemize}
\begin{itemize}
    \item a) \textbf{Đúng}.
    \item b) \textbf{Đúng}.
    \item c) \textbf{Sai}.
    \item d) \textbf{Đúng}.
\end{itemize}
}
\end{ex}

\begin{ex} %[4D10V57]
Cho hai phân số $\dfrac{2}{5}$ và $\dfrac{4}{15}$. Xét tính đúng sai:
\choiceTF
{\True Mẫu số $15$ chia hết cho mẫu số $5$.}
{\True Chọn $15$ làm mẫu số chung.}
{Quy đồng được $\dfrac{8}{15}$ và $\dfrac{4}{15}$.}
{\True Quy đồng được $\dfrac{6}{15}$ và $\dfrac{4}{15}$.}
\loigiai{
\begin{itemize}
    \item 15 chia hết cho 5: Đúng.
    \item Chọn 15 làm MSC: Đúng.
    \item $2/5 = 6/15 \neq 8/15$. Ý 3 Sai, Ý 4 Đúng.
\end{itemize}
\begin{itemize}
    \item a) \textbf{Đúng}.
    \item b) \textbf{Đúng}.
    \item c) \textbf{Sai}.
    \item d) \textbf{Đúng}.
\end{itemize}
}
\end{ex}

\textbf{Phần 3. Trả lời ngắn}
\setcounterpageref{ex}{0}

\begin{ex} %[4D10H57]
Tìm mẫu số chung của hai phân số $\dfrac{2}{3}$ và $\dfrac{1}{7}$ (lấy tích hai mẫu số).
\shortans{21}
\loigiai{
\textbf{Đáp số:} $21$\\
$3 \times 7 = 21$.
}
\end{ex}

\begin{ex} %[4D10H57]
Quy đồng mẫu số phân số $\dfrac{3}{4}$ với mẫu số chung 20. Tìm tử số mới.
\shortans{15}
\loigiai{
\textbf{Đáp số:} $15$\\
$3 \times 5 = 15$.
}
\end{ex}

\begin{ex} %[4D10V57]
Quy đồng mẫu số hai phân số $\dfrac{5}{6}$ và $\dfrac{3}{8}$ với mẫu số chung nhỏ nhất là 24. Tìm tổng hai tử số mới.
\shortans{29}
\loigiai{
\textbf{Đáp số:} $29$\\
$5 \times 4 + 3 \times 3 = 20 + 9 = 29$.
}
\end{ex}

\begin{ex} %[4D10V57]
Tìm mẫu số chung nhỏ nhất của hai phân số $\dfrac{7}{12}$ và $\dfrac{5}{16}$.
\shortans{48}
\loigiai{
\textbf{Đáp số:} $48$\\
$48$.
}
\end{ex}

\begin{ex} %[4D10V57]
Quy đồng mẫu số của 3 và $\dfrac{2}{7}$ với mẫu số là 7. Tìm tử số của số 3 sau khi quy đồng.
\shortans{21}
\loigiai{
\textbf{Đáp số:} $21$\\
$3 \times 7 = 21$.
}
\end{ex}

\begin{ex} %[4D10C57]
Tìm mẫu số chung nhỏ nhất của 3 phân số: $\dfrac{1}{6}, \dfrac{1}{8}, \dfrac{1}{12}$.
\shortans{24}
\loigiai{
\textbf{Đáp số:} $24$\\
$24$.
}
\end{ex}

\newpage

\TaoPhieuKiemTra{Chủ đề 10 -- Phân số}{58}{So sánh hai phân số}{https://loptoanthaylinh.vn/thi-online/toan-4/bai-58}
\textbf{Phần 1. Trắc nghiệm lựa chọn}
\setcounterpageref{ex}{0}

\begin{ex} %[4D10N58]
Trong hai phân số có cùng mẫu số dương, phân số nào lớn hơn?
\choice
{Phân số có tử số bé hơn}
{\True Phân số có tử số lớn hơn}
{Hai phân số luôn bằng nhau}
{Không so sánh được}
\loigiai{
Cùng mẫu số: tử lớn hơn thì phân số lớn hơn.\\
\textbf{Đáp án của câu hỏi là B}
}
\end{ex}

\begin{ex} %[4D10N58]
Điền dấu thích hợp vào chỗ chấm: $\dfrac{3}{7} \dots \dfrac{5}{7}$.
\choice
{\True $<$}
{$>$}
{$=$}
{Không so sánh được}
\loigiai{
$3 < 5 \Rightarrow \dfrac{3}{7} < \dfrac{5}{7}$.\\
\textbf{Đáp án của câu hỏi là A}
}
\end{ex}

\begin{ex} %[4D10H58]
Trong hai phân số có cùng tử số dương, phân số nào lớn hơn?
\choice
{Phân số có mẫu số lớn hơn}
{Hai phân số luôn bằng nhau}
{\True Phân số có mẫu số bé hơn}
{Không so sánh được}
\loigiai{
Cùng tử số: mẫu bé hơn thì phân số lớn hơn.\\
\textbf{Đáp án của câu hỏi là C}
}
\end{ex}

\begin{ex} %[4D10H58]
Điền dấu thích hợp vào chỗ chấm: $\dfrac{4}{5} \dots \dfrac{4}{9}$.
\choice
{$<$}
{$=$}
{Không so sánh được}
{\True $>$}
\loigiai{
$5 < 9 \Rightarrow \dfrac{4}{5} > \dfrac{4}{9}$.\\
\textbf{Đáp án của câu hỏi là D}
}
\end{ex}

\begin{ex} %[4D10H58]
So sánh hai phân số $\dfrac{2}{3}$ và $\dfrac{3}{4}$:
\choice
{\True $\dfrac{2}{3} < \dfrac{3}{4}$}
{$\dfrac{2}{3} > \dfrac{3}{4}$}
{$\dfrac{2}{3} = \dfrac{3}{4}$}
{Không so sánh được}
\loigiai{
$\dfrac{2}{3} = \dfrac{8}{12} < \dfrac{9}{12} = \dfrac{3}{4}$.\\
\textbf{Đáp án của câu hỏi là A}
}
\end{ex}

\begin{ex} %[4D10H58]
Phân số nào sau đây lớn hơn 1?
\choice
{$\dfrac{5}{6}$}
{$\dfrac{7}{7}$}
{\True $\dfrac{8}{5}$}
{$\dfrac{3}{4}$}
\loigiai{
$\dfrac{8}{5} > 1$ (tử số lớn hơn mẫu số).\\
\textbf{Đáp án của câu hỏi là C}
}
\end{ex}

\begin{ex} %[4D10V58]
Sắp xếp các phân số $\dfrac{1}{2}, \dfrac{1}{3}, \dfrac{1}{5}, \dfrac{1}{4}$ theo thứ tự từ bé đến lớn:
\choice
{$\dfrac{1}{2} < \dfrac{1}{3} < \dfrac{1}{4} < \dfrac{1}{5}$}
{$\dfrac{1}{5} < \dfrac{1}{3} < \dfrac{1}{4} < \dfrac{1}{2}$}
{$\dfrac{1}{2} < \dfrac{1}{4} < \dfrac{1}{3} < \dfrac{1}{5}$}
{\True $\dfrac{1}{5} < \dfrac{1}{4} < \dfrac{1}{3} < \dfrac{1}{2}$}
\loigiai{
Cùng tử số: mẫu càng lớn phân số càng bé $\Rightarrow \dfrac{1}{5} < \dfrac{1}{4} < \dfrac{1}{3} < \dfrac{1}{2}$.\\
\textbf{Đáp án của câu hỏi là D}
}
\end{ex}

\begin{ex} %[4D10V58]
So sánh hai phân số $\dfrac{7}{8}$ và $\dfrac{8}{9}$ bằng cách so sánh phần bù tới 1:
\choice
{\True $\dfrac{7}{8} < \dfrac{8}{9}$}
{$\dfrac{7}{8} > \dfrac{8}{9}$}
{$\dfrac{7}{8} = \dfrac{8}{9}$}
{Không so sánh được}
\loigiai{
Phần bù: $1 - \dfrac{7}{8} = \dfrac{1}{8} > \dfrac{1}{9} = 1 - \dfrac{8}{9} \Rightarrow \dfrac{7}{8} < \dfrac{8}{9}$.\\
\textbf{Đáp án của câu hỏi là A}
}
\end{ex}

\begin{ex} %[4D10V58]
Bạn An ăn $\dfrac{1}{3}$ cái bánh, bạn Bình ăn $\dfrac{2}{5}$ cái bánh. Bạn nào ăn nhiều bánh hơn?
\choice
{Bạn An}
{\True Bạn Bình}
{Hai bạn ăn bằng nhau}
{Không so sánh được}
\loigiai{
$\dfrac{1}{3} = \dfrac{5}{15} < \dfrac{6}{15} = \dfrac{2}{5}$ nên bạn Bình ăn nhiều hơn.\\
\textbf{Đáp án của câu hỏi là B}
}
\end{ex}

\begin{ex} %[4D10V58]
Phân số nào sau đây nằm giữa $\dfrac{1}{3}$ và $\dfrac{2}{3}$?
\choice
{$\dfrac{1}{6}$}
{$\dfrac{5}{6}$}
{\True $\dfrac{1}{2}$}
{$\dfrac{3}{4}$}
\loigiai{
$\dfrac{1}{3} = \dfrac{2}{6} < \dfrac{1}{2} = \dfrac{3}{6} < \dfrac{2}{3} = \dfrac{4}{6}$.\\
\textbf{Đáp án của câu hỏi là C}
}
\end{ex}

\begin{ex} %[4D10V58]
Tìm phân số lớn nhất trong các phân số: $\dfrac{3}{4}, \dfrac{5}{6}, \dfrac{2}{3}, \dfrac{7}{12}$.
\choice
{$\dfrac{3}{4}$}
{\True $\dfrac{5}{6}$}
{$\dfrac{2}{3}$}
{$\dfrac{7}{12}$}
\loigiai{
MSC là 12: $\dfrac{9}{12}, \dfrac{10}{12}, \dfrac{8}{12}, \dfrac{7}{12}$. Lớn nhất là $\dfrac{10}{12} = \dfrac{5}{6}$.\\
\textbf{Đáp án của câu hỏi là B}
}
\end{ex}

\begin{ex} %[4D10C58]
So sánh hai phân số: $A = \dfrac{2025}{2026}$ và $B = \dfrac{2026}{2027}$.
\choice
{$A > B$}
{$A = B$}
{Không so sánh được}
{\True $A < B$}
\loigiai{
Phần bù: $1 - A = \dfrac{1}{2026} > \dfrac{1}{2027} = 1 - B \Rightarrow A < B$.\\
\textbf{Đáp án của câu hỏi là D}
}
\end{ex}

\textbf{Phần 2. Trắc nghiệm đúng sai}
\setcounterpageref{ex}{0}

\begin{ex} %[4D10H58]
Xét tính đúng sai của các phép so sánh phân số:
\choiceTF
{\True $\dfrac{5}{8} < \dfrac{7}{8}$.}
{\True $\dfrac{3}{4} > \dfrac{3}{7}$.}
{$\dfrac{5}{4} < 1$.}
{\True $\dfrac{8}{8} = 1$.}
\loigiai{
\begin{itemize}
    \item $5/8 < 7/8$: Đúng.
    \item $3/4 > 3/7$: Đúng.
    \item $5/4 > 1$ (không bé hơn 1). Sai.
    \item $8/8 = 1$: Đúng.
\end{itemize}
\begin{itemize}
    \item a) \textbf{Đúng}.
    \item b) \textbf{Đúng}.
    \item c) \textbf{Sai}.
    \item d) \textbf{Đúng}.
\end{itemize}
}
\end{ex}

\begin{ex} %[4D10V58]
Cho 3 phân số: $\dfrac{3}{5}, \dfrac{2}{3}, \dfrac{7}{10}$. Xét tính đúng sai:
\choiceTF
{\True Quy đồng với mẫu số chung 30 được $\dfrac{18}{30}, \dfrac{20}{30}, \dfrac{21}{30}$.}
{\True Phân số lớn nhất là $\dfrac{7}{10}$.}
{Phân số bé nhất là $\dfrac{2}{3}$.}
{\True Sắp xếp tăng dần: $\dfrac{3}{5} < \dfrac{2}{3} < \dfrac{7}{10}$.}
\loigiai{
\begin{itemize}
    \item $3/5 = 18/30; 2/3 = 20/30; 7/10 = 21/30$. Đúng.
    \item $7/10$ lớn nhất: Đúng.
    \item Bé nhất là $3/5$ (không phải $2/3$). Sai.
    \item Thứ tự tăng dần: $3/5 < 2/3 < 7/10$. Đúng.
\end{itemize}
\begin{itemize}
    \item a) \textbf{Đúng}.
    \item b) \textbf{Đúng}.
    \item c) \textbf{Sai}.
    \item d) \textbf{Đúng}.
\end{itemize}
}
\end{ex}

\textbf{Phần 3. Trả lời ngắn}
\setcounterpageref{ex}{0}

\begin{ex} %[4D10H58]
Trong hai phân số $\dfrac{4}{9}$ và $\dfrac{7}{9}$, phân số nào lớn hơn? (Viết dạng a/b).
\shortans{7/9}
\loigiai{
\textbf{Đáp số:} $7/9$\\
$7/9$.
}
\end{ex}

\begin{ex} %[4D10H58]
Trong hai phân số $\dfrac{5}{6}$ và $\dfrac{5}{8}$, phân số nào lớn hơn? (Viết dạng a/b).
\shortans{5/6}
\loigiai{
\textbf{Đáp số:} $5/6$\\
$5/6$.
}
\end{ex}

\begin{ex} %[4D10V58]
Tìm phân số lớn nhất trong 3 phân số: $\dfrac{1}{2}, \dfrac{2}{3}, \dfrac{3}{4}$ (dạng a/b).
\shortans{3/4}
\loigiai{
\textbf{Đáp số:} $3/4$\\
$3/4$.
}
\end{ex}

\begin{ex} %[4D10V58]
Tìm phân số bé nhất trong 3 phân số: $\dfrac{2}{5}, \dfrac{1}{3}, \dfrac{4}{15}$ (dạng a/b).
\shortans{4/15}
\loigiai{
\textbf{Đáp số:} $4/15$\\
$6/15, 5/15, 4/15 \Rightarrow 4/15$ bé nhất.
}
\end{ex}

\begin{ex} %[4D10V58]
Có bao nhiêu phân số có mẫu số bằng 8 lớn hơn $\dfrac{2}{8}$ và bé hơn $\dfrac{7}{8}$?
\shortans{4}
\loigiai{
\textbf{Đáp số:} $4$\\
$3/8, 4/8, 5/8, 6/8$ (4 phân số).
}
\end{ex}

\begin{ex} %[4D10C58]
Tìm chữ số tự nhiên $x$ sao cho $\dfrac{3}{5} < \dfrac{x}{10} < \dfrac{4}{5}$.
\shortans{7}
\loigiai{
\textbf{Đáp số:} $7$\\
$6/10 < x/10 < 8/10 \Rightarrow x = 7$.
}
\end{ex}

\newpage

\TaoPhieuKiemTra{Chủ đề 10 -- Phân số}{59}{Luyện tập chung phân số}{https://loptoanthaylinh.vn/thi-online/toan-4/bai-59}
\textbf{Phần 1. Trắc nghiệm lựa chọn}
\setcounterpageref{ex}{0}

\begin{ex} %[4D10N59]
Phân số chỉ số phần không tô màu trong một hình chia làm 5 phần bằng nhau, đã tô màu 2 phần là
\choice
{$\dfrac{2}{5}$}
{$\dfrac{5}{2}$}
{\True $\dfrac{3}{5}$}
{$\dfrac{3}{2}$}
\loigiai{
Phần chưa tô màu: $5 - 2 = 3$ phần $\Rightarrow \dfrac{3}{5}$.\\
\textbf{Đáp án của câu hỏi là C}
}
\end{ex}

\begin{ex} %[4D10N59]
Thương của phép chia $5 : 8$ viết dưới dạng phân số là
\choice
{$\dfrac{8}{5}$}
{\True $\dfrac{5}{8}$}
{$5{,}8$}
{$8{,}5$}
\loigiai{
$5 : 8 = \dfrac{5}{8}$.\\
\textbf{Đáp án của câu hỏi là B}
}
\end{ex}

\begin{ex} %[4D10H59]
Phân số $\dfrac{16}{20}$ rút gọn về phân số tối giản là
\choice
{$\dfrac{8}{10}$}
{$\dfrac{2}{3}$}
{$\dfrac{3}{4}$}
{\True $\dfrac{4}{5}$}
\loigiai{
$\dfrac{16 : 4}{20 : 4} = \dfrac{4}{5}$.\\
\textbf{Đáp án của câu hỏi là D}
}
\end{ex}

\begin{ex} %[4D10H59]
Quy đồng mẫu số hai phân số $\dfrac{2}{5}$ và $\dfrac{3}{10}$ với mẫu số chung là 10 ta được
\choice
{$\dfrac{2}{10}$ và $\dfrac{3}{10}$}
{$\dfrac{6}{10}$ và $\dfrac{3}{10}$}
{\True $\dfrac{4}{10}$ và $\dfrac{3}{10}$}
{$\dfrac{5}{10}$ và $\dfrac{3}{10}$}
\loigiai{
$\dfrac{2 \times 2}{5 \times 2} = \dfrac{4}{10}$ và $\dfrac{3}{10}$.\\
\textbf{Đáp án của câu hỏi là C}
}
\end{ex}

\begin{ex} %[4D10H59]
Phân số nào sau đây lớn nhất?
\choice
{\True $\dfrac{6}{7}$}
{$\dfrac{3}{7}$}
{$\dfrac{5}{7}$}
{$\dfrac{2}{7}$}
\loigiai{
Cùng mẫu số: $\dfrac{6}{7}$ lớn nhất.\\
\textbf{Đáp án của câu hỏi là A}
}
\end{ex}

\begin{ex} %[4D10H59]
Phân số nào sau đây bé hơn 1?
\choice
{$\dfrac{7}{5}$}
{$\dfrac{9}{9}$}
{$\dfrac{11}{10}$}
{\True $\dfrac{5}{8}$}
\loigiai{
$\dfrac{5}{8} < 1$ (tử số nhỏ hơn mẫu số).\\
\textbf{Đáp án của câu hỏi là D}
}
\end{ex}

\begin{ex} %[4D10V59]
Tìm $x$ biết: $\dfrac{x}{15} = \dfrac{4}{5}$.
\choice
{\True $12$}
{$10$}
{$14$}
{$16$}
\loigiai{
$x = 4 \times 3 = 12$.\\
\textbf{Đáp án của câu hỏi là A}
}
\end{ex}

\begin{ex} %[4D10V59]
Sắp xếp các phân số $\dfrac{2}{3}, \dfrac{3}{4}, \dfrac{5}{6}$ theo thứ tự từ bé đến lớn:
\choice
{$\dfrac{5}{6} < \dfrac{3}{4} < \dfrac{2}{3}$}
{$\dfrac{3}{4} < \dfrac{2}{3} < \dfrac{5}{6}$}
{\True $\dfrac{2}{3} < \dfrac{3}{4} < \dfrac{5}{6}$}
{$\dfrac{2}{3} < \dfrac{5}{6} < \dfrac{3}{4}$}
\loigiai{
MSC 12: $\dfrac{8}{12} < \dfrac{9}{12} < \dfrac{10}{12} \Rightarrow \dfrac{2}{3} < \dfrac{3}{4} < \dfrac{5}{6}$.\\
\textbf{Đáp án của câu hỏi là C}
}
\end{ex}

\begin{ex} %[4D10V59]
Một lớp học có 32 học sinh, trong đó có 24 học sinh thích môn Toán. Phân số chỉ phần học sinh thích môn Toán được rút gọn tối giản là
\choice
{$\dfrac{6}{8}$}
{\True $\dfrac{3}{4}$}
{$\dfrac{12}{16}$}
{$\dfrac{2}{3}$}
\loigiai{
$\dfrac{24 : 8}{32 : 8} = \dfrac{3}{4}$.\\
\textbf{Đáp án của câu hỏi là B}
}
\end{ex}

\begin{ex} %[4D10V59]
Điền dấu thích hợp vào chỗ chấm: $\dfrac{5}{9} \dots \dfrac{5}{11}$.
\choice
{\True $>$}
{$<$}
{$=$}
{Không so sánh được}
\loigiai{
Cùng tử số: $9 < 11 \Rightarrow \dfrac{5}{9} > \dfrac{5}{11}$.\\
\textbf{Đáp án của câu hỏi là A}
}
\end{ex}

\begin{ex} %[4D10V59]
Tìm phân số bằng phân số $\dfrac{2}{3}$ có mẫu số là 24.
\choice
{$\dfrac{14}{24}$}
{\True $\dfrac{16}{24}$}
{$\dfrac{18}{24}$}
{$\dfrac{12}{24}$}
\loigiai{
$\dfrac{2 \times 8}{3 \times 8} = \dfrac{16}{24}$.\\
\textbf{Đáp án của câu hỏi là B}
}
\end{ex}

\begin{ex} %[4D10C59]
Tìm phân số tối giản $\dfrac{a}{b}$ biết nếu cộng thêm 6 vào tử số thì được phân số mới có giá trị bằng 1, và phân số $\dfrac{a}{b}$ có giá trị bằng $\dfrac{3}{5}$.
\choice
{$\dfrac{6}{10}$}
{$\dfrac{12}{20}$}
{$\dfrac{3}{5}$}
{\True $\dfrac{9}{15}$}
\loigiai{
Cộng 6 vào tử được phân số bằng 1 tức là $b - a = 6$.\\
Hiệu số phần: $5 - 3 = 2$ phần $= 6 \Rightarrow 1$ phần $= 3$.\\
Tử số $a = 3 \times 3 = 9$, Mẫu số $b = 5 \times 3 = 15$. Phân số ban đầu rút gọn tối giản là $\dfrac{3}{5}$ (hoặc phân số cần tìm là $\dfrac{9}{15}$).\\
\textbf{Đáp án của câu hỏi là D}
}
\end{ex}

\textbf{Phần 2. Trắc nghiệm đúng sai}
\setcounterpageref{ex}{0}

\begin{ex} %[4D10H59]
Xét tính đúng sai về các tính chất phân số:
\choiceTF
{\True Phân số $\dfrac{12}{18}$ rút gọn tối giản là $\dfrac{2}{3}$.}
{\True $\dfrac{4}{7} = \dfrac{8}{14} = \dfrac{12}{21}$.}
{Phân số $\dfrac{8}{9}$ lớn hơn phân số $\dfrac{9}{8}$.}
{\True $\dfrac{8}{9} < 1 < \dfrac{9}{8}$.}
\loigiai{
\begin{itemize}
    \item $12/18 = 2/3$. Đúng.
    \item Phân số bằng nhau đúng.
    \item $8/9 < 1 < 9/8 \Rightarrow 8/9 < 9/8$. Sai.
    \item So sánh qua 1 đúng.
\end{itemize}
\begin{itemize}
    \item a) \textbf{Đúng}.
    \item b) \textbf{Đúng}.
    \item c) \textbf{Sai}.
    \item d) \textbf{Đúng}.
\end{itemize}
}
\end{ex}

\begin{ex} %[4D10V59]
Cho các phân số: $\dfrac{3}{5}, \dfrac{5}{4}, \dfrac{7}{7}, \dfrac{2}{9}$. Xét tính đúng sai:
\choiceTF
{\True Phân số lớn hơn 1 là $\dfrac{5}{4}$.}
{\True Phân số bằng 1 là $\dfrac{7}{7}$.}
{Có 3 phân số bé hơn 1.}
{\True Có 2 phân số bé hơn 1 (là $\dfrac{3}{5}$ và $\dfrac{2}{9}$).}
\loigiai{
\begin{itemize}
    \item Lớn hơn 1: 5/4. Đúng.
    \item Bằng 1: 7/7. Đúng.
    \item Bé hơn 1: chỉ có 2 phân số ($3/5, 2/9$). Ý 3 Sai, Ý 4 Đúng.
\end{itemize}
\begin{itemize}
    \item a) \textbf{Đúng}.
    \item b) \textbf{Đúng}.
    \item c) \textbf{Sai}.
    \item d) \textbf{Đúng}.
\end{itemize}
}
\end{ex}

\textbf{Phần 3. Trả lời ngắn}
\setcounterpageref{ex}{0}

\begin{ex} %[4D10H59]
Rút gọn phân số $\dfrac{18}{24}$ về tối giản (dạng a/b).
\shortans{3/4}
\loigiai{
\textbf{Đáp số:} $3/4$\\
$3/4$.
}
\end{ex}

\begin{ex} %[4D10H59]
Tìm $x$ biết $\dfrac{x}{20} = \dfrac{3}{5}$.
\shortans{12}
\loigiai{
\textbf{Đáp số:} $12$\\
$x = 3 \times 4 = 12$.
}
\end{ex}

\begin{ex} %[4D10V59]
Quy đồng mẫu số hai phân số $\dfrac{3}{4}$ và $\dfrac{5}{6}$ với mẫu số chung nhỏ nhất 12. Tìm tổng hai tử số mới.
\shortans{19}
\loigiai{
\textbf{Đáp số:} $19$\\
$3 \times 3 + 5 \times 2 = 9 + 10 = 19$.
}
\end{ex}

\begin{ex} %[4D10V59]
Tìm phân số lớn nhất trong các phân số: $\dfrac{2}{3}, \dfrac{3}{5}, \dfrac{7}{15}$ (dạng a/b).
\shortans{2/3}
\loigiai{
\textbf{Đáp số:} $2/3$\\
$10/15, 9/15, 7/15 \Rightarrow 2/3$.
}
\end{ex}

\begin{ex} %[4D10V59]
Tìm phân số bé nhất trong các phân số: $\dfrac{4}{5}, \dfrac{4}{7}, \dfrac{4}{9}$ (dạng a/b).
\shortans{4/9}
\loigiai{
\textbf{Đáp số:} $4/9$\\
$4/9$.
}
\end{ex}

\begin{ex} %[4D10C59]
Tìm một phân số bằng $\dfrac{2}{5}$ có mẫu số hơn tử số 18 đơn vị. Tìm tử số của phân số đó.
\shortans{12}
\loigiai{
\textbf{Đáp số:} $12$\\
Hiệu số phần: $5 - 2 = 3$ phần $= 18 \Rightarrow 1$ phần $= 6$.\\
Tử số: $6 \times 2 = 12$.
}
\end{ex}

\newpage

% =========================================================================
% CHỦ ĐỀ 11 - PHÉP CỘNG, PHÉP TRỪ PHÂN SỐ
% =========================================================================
\section{Phép cộng, phép trừ phân số}

\TaoPhieuKiemTra{Chủ đề 11 -- Phép cộng, phép trừ phân số}{60}{Phép cộng phân số}{https://loptoanthaylinh.vn/thi-online/toan-4/bai-60}
\textbf{Phần 1. Trắc nghiệm lựa chọn}
\setcounterpageref{ex}{0}

\begin{ex} %[4D11N60]
Muốn cộng hai phân số cùng mẫu số, ta
\choice
{Cộng tử số với tử số, cộng mẫu số với mẫu số}
{Nhân hai tử số với nhau}
{\True Cộng hai tử số với nhau và giữ nguyên mẫu số}
{Quy đồng tử số}
\loigiai{
Cộng hai tử số với nhau và giữ nguyên mẫu số.\\
\textbf{Đáp án của câu hỏi là C}
}
\end{ex}

\begin{ex} %[4D11N60]
Kết quả của phép tính $\dfrac{2}{7} + \dfrac{3}{7}$ là
\choice
{$\dfrac{5}{14}$}
{$\dfrac{6}{7}$}
{$\dfrac{1}{7}$}
{\True $\dfrac{5}{7}$}
\loigiai{
$\dfrac{2 + 3}{7} = \dfrac{5}{7}$.\\
\textbf{Đáp án của câu hỏi là D}
}
\end{ex}

\begin{ex} %[4D11H60]
Kết quả của phép tính $\dfrac{1}{3} + \dfrac{1}{6}$ là
\choice
{$\dfrac{2}{9}$}
{\True $\dfrac{1}{2}$}
{$\dfrac{3}{6}$}
{$\dfrac{2}{6}$}
\loigiai{
$\dfrac{2}{6} + \dfrac{1}{6} = \dfrac{3}{6} = \dfrac{1}{2}$.\\
\textbf{Đáp án của câu hỏi là B}
}
\end{ex}

\begin{ex} %[4D11H60]
Kết quả của phép tính $\dfrac{3}{4} + \dfrac{2}{5}$ là
\choice
{\True $\dfrac{23}{20}$}
{$\dfrac{5}{9}$}
{$\dfrac{11}{20}$}
{$\dfrac{6}{20}$}
\loigiai{
$\dfrac{15}{20} + \dfrac{8}{20} = \dfrac{23}{20}$.\\
\textbf{Đáp án của câu hỏi là A}
}
\end{ex}

\begin{ex} %[4D11H60]
Kết quả của phép tính $1 + \dfrac{2}{5}$ là
\choice
{$\dfrac{3}{5}$}
{$\dfrac{2}{5}$}
{\True $\dfrac{7}{5}$}
{$\dfrac{8}{5}$}
\loigiai{
$\dfrac{5}{5} + \dfrac{2}{5} = \dfrac{7}{5}$.\\
\textbf{Đáp án của câu hỏi là C}
}
\end{ex}

\begin{ex} %[4D11H60]
Tính thuận tiện: $\dfrac{1}{4} + \dfrac{2}{7} + \dfrac{3}{4}$.
\choice
{\True $\dfrac{9}{7}$}
{$1$}
{$\dfrac{8}{7}$}
{$2$}
\loigiai{
$\left(\dfrac{1}{4} + \dfrac{3}{4}\right) + \dfrac{2}{7} = 1 + \dfrac{2}{7} = \dfrac{9}{7}$.\\
\textbf{Đáp án của câu hỏi là A}
}
\end{ex}

\begin{ex} %[4D11V60]
Một người thợ ngày thứ nhất dệt được $\dfrac{2}{7}$ tấm vải, ngày thứ hai dệt được $\dfrac{3}{7}$ tấm vải. Cả hai ngày người đó dệt được số phần tấm vải là
\choice
{$\dfrac{5}{14}$ tấm vải}
{\True $\dfrac{5}{7}$ tấm vải}
{$\dfrac{6}{7}$ tấm vải}
{$1$ tấm vải}
\loigiai{
$\dfrac{2}{7} + \dfrac{3}{7} = \dfrac{5}{7}$ tấm vải.\\
\textbf{Đáp án của câu hỏi là B}
}
\end{ex}

\begin{ex} %[4D11V60]
Một vòi nước giờ thứ nhất chảy được $\dfrac{1}{3}$ bể nước, giờ thứ hai chảy được $\dfrac{2}{5}$ bể nước. Cả hai giờ vòi chảy được số phần bể nước là
\choice
{$\dfrac{3}{8}$ bể}
{$\dfrac{7}{15}$ bể}
{$\dfrac{2}{15}$ bể}
{\True $\dfrac{11}{15}$ bể}
\loigiai{
$\dfrac{1}{3} + \dfrac{2}{5} = \dfrac{5 + 6}{15} = \dfrac{11}{15}$ bể.\\
\textbf{Đáp án của câu hỏi là D}
}
\end{ex}

\begin{ex} %[4D11V60]
Tìm $x$ biết: $x - \dfrac{1}{4} = \dfrac{5}{8}$.
\choice
{\True $\dfrac{7}{8}$}
{$\dfrac{6}{8}$}
{$\dfrac{3}{8}$}
{$1$}
\loigiai{
$x = \dfrac{5}{8} + \dfrac{2}{8} = \dfrac{7}{8}$.\\
\textbf{Đáp án của câu hỏi là A}
}
\end{ex}

\begin{ex} %[4D11V60]
Tính giá trị biểu thức: $\dfrac{1}{2} + \dfrac{1}{3} + \dfrac{1}{6}$.
\choice
{$\dfrac{5}{6}$}
{$\dfrac{7}{6}$}
{$\dfrac{3}{11}$}
{\True $1$}
\loigiai{
$\dfrac{3}{6} + \dfrac{2}{6} + \dfrac{1}{6} = \dfrac{6}{6} = 1$.\\
\textbf{Đáp án của câu hỏi là D}
}
\end{ex}

\begin{ex} %[4D11V60]
Một hình tam giác có độ dài 3 cạnh lần lượt là $\dfrac{3}{5}\text{ m}, \dfrac{4}{5}\text{ m}, \dfrac{2}{5}\text{ m}$. Chu vi tam giác đó là
\choice
{$\dfrac{9}{15}\text{ m}$}
{\True $\dfrac{9}{5}\text{ m}$}
{$2\text{ m}$}
{$\dfrac{7}{5}\text{ m}$}
\loigiai{
$\dfrac{3 + 4 + 2}{5} = \dfrac{9}{5}\text{ m}$.\\
\textbf{Đáp án của câu hỏi là B}
}
\end{ex}

\begin{ex} %[4D11C60]
Tính tổng: $S = \dfrac{1}{1 \times 2} + \dfrac{1}{2 \times 3} + \dfrac{1}{3 \times 4} + \dfrac{1}{4 \times 5}$.
\choice
{$\dfrac{3}{5}$}
{$\dfrac{5}{6}$}
{\True $\dfrac{4}{5}$}
{$1$}
\loigiai{
$S = 1 - \dfrac{1}{2} + \dfrac{1}{2} - \dfrac{1}{3} + \dots - \dfrac{1}{5} = 1 - \dfrac{1}{5} = \dfrac{4}{5}$.\\
\textbf{Đáp án của câu hỏi là C}
}
\end{ex}

\textbf{Phần 2. Trắc nghiệm đúng sai}
\setcounterpageref{ex}{0}

\begin{ex} %[4D11H60]
Xét tính đúng sai của các phép tính cộng phân số:
\choiceTF
{\True $\dfrac{1}{5} + \dfrac{2}{5} = \dfrac{3}{5}$.}
{$\dfrac{1}{2} + \dfrac{1}{3} = \dfrac{2}{5}$.}
{\True $\dfrac{1}{2} + \dfrac{1}{3} = \dfrac{5}{6}$.}
{\True $\dfrac{3}{8} + \dfrac{5}{8} = 1$.}
\loigiai{
\begin{itemize}
    \item Cùng mẫu số: $1/5 + 2/5 = 3/5$. Đúng.
    \item $1/2 + 1/3 = 5/6 \neq 2/5$. Ý 2 Sai, Ý 3 Đúng.
    \item $3/8 + 5/8 = 8/8 = 1$. Đúng.
\end{itemize}
\begin{itemize}
    \item a) \textbf{Đúng}.
    \item b) \textbf{Sai}.
    \item c) \textbf{Đúng}.
    \item d) \textbf{Đúng}.
\end{itemize}
}
\end{ex}

\begin{ex} %[4D11V60]
Hai vòi nước cùng chảy vào một bể cạn: Vòi 1 chảy được $\dfrac{2}{7}$ bể, vòi 2 chảy được $\dfrac{3}{7}$ bể. Xét tính đúng sai:
\choiceTF
{\True Cả hai vòi cùng chảy được $\dfrac{5}{7}$ bể nước.}
{\True Vòi 2 chảy được nhiều hơn vòi 1 là $\dfrac{1}{7}$ bể nước.}
{Sau khi hai vòi chảy thì bể đã đầy nước.}
{\True Còn lại $\dfrac{2}{7}$ bể chưa có nước.}
\loigiai{
\begin{itemize}
    \item Tổng: $2/7 + 3/7 = 5/7$ bể. Đúng.
    \item Hiệu: $3/7 - 2/7 = 1/7$ bể. Đúng.
    \item $5/7 < 1$ nên chưa đầy. Sai.
    \item Còn lại: $1 - 5/7 = 2/7$ bể. Đúng.
\end{itemize}
\begin{itemize}
    \item a) \textbf{Đúng}.
    \item b) \textbf{Đúng}.
    \item c) \textbf{Sai}.
    \item d) \textbf{Đúng}.
\end{itemize}
}
\end{ex}

\textbf{Phần 3. Trả lời ngắn}
\setcounterpageref{ex}{0}

\begin{ex} %[4D11H60]
Tính tổng: $\dfrac{4}{9} + \dfrac{2}{9}$ (rút gọn tối giản dạng a/b).
\shortans{2/3}
\loigiai{
\textbf{Đáp số:} $2/3$\\
$6/9 = 2/3$.
}
\end{ex}

\begin{ex} %[4D11H60]
Tính tổng: $\dfrac{1}{4} + \dfrac{3}{8}$ (dạng a/b).
\shortans{5/8}
\loigiai{
\textbf{Đáp số:} $5/8$\\
$2/8 + 3/8 = 5/8$.
}
\end{ex}

\begin{ex} %[4D11V60]
Tính giá trị: $1 + \dfrac{3}{4}$ (dạng a/b).
\shortans{7/4}
\loigiai{
\textbf{Đáp số:} $7/4$\\
$7/4$.
}
\end{ex}

\begin{ex} %[4D11V60]
Tìm $x$ biết: $x - \dfrac{2}{5} = \dfrac{1}{2}$ (dạng a/b).
\shortans{9/10}
\loigiai{
\textbf{Đáp số:} $9/10$\\
$x = 1/2 + 2/5 = 5/10 + 4/10 = 9/10$.
}
\end{ex}

\begin{ex} %[4D11V60]
Tính thuận tiện: $\dfrac{2}{9} + \dfrac{5}{7} + \dfrac{7}{9}$ (dạng a/b).
\shortans{12/7}
\loigiai{
\textbf{Đáp số:} $12/7$\\
$1 + 5/7 = 12/7$.
}
\end{ex}

\begin{ex} %[4D11C60]
Tính tổng: $\dfrac{1}{2} + \dfrac{1}{4} + \dfrac{1}{8} + \dfrac{1}{16}$ (dạng a/b).
\shortans{15/16}
\loigiai{
\textbf{Đáp số:} $15/16$\\
$8/16 + 4/16 + 2/16 + 1/16 = 15/16$.
}
\end{ex}

\newpage

\TaoPhieuKiemTra{Chủ đề 11 -- Phép cộng, phép trừ phân số}{61}{Phép trừ phân số}{https://loptoanthaylinh.vn/thi-online/toan-4/bai-61}
\textbf{Phần 1. Trắc nghiệm lựa chọn}
\setcounterpageref{ex}{0}

\begin{ex} %[4D11N61]
Muốn trừ hai phân số cùng mẫu số, ta
\choice
{Trừ tử số với tử số, trừ mẫu số với mẫu số}
{\True Trừ tử số của phân số thứ nhất cho tử số của phân số thứ hai và giữ nguyên mẫu số}
{Nhân hai mẫu số với nhau}
{Đổi chỗ hai phân số}
\loigiai{
Trừ hai tử số và giữ nguyên mẫu số.\\
\textbf{Đáp án của câu hỏi là B}
}
\end{ex}

\begin{ex} %[4D11N61]
Kết quả của phép tính $\dfrac{5}{8} - \dfrac{3}{8}$ là
\choice
{$\dfrac{2}{0}$}
{$\dfrac{8}{8}$}
{\True $\dfrac{1}{4}$ (tối giản từ $\dfrac{2}{8}$)}
{$\dfrac{1}{8}$}
\loigiai{
$\dfrac{5 - 3}{8} = \dfrac{2}{8} = \dfrac{1}{4}$.\\
\textbf{Đáp án của câu hỏi là C}
}
\end{ex}

\begin{ex} %[4D11H61]
Kết quả của phép trừ $\dfrac{3}{4} - \dfrac{1}{2}$ là
\choice
{\True $\dfrac{1}{4}$}
{$\dfrac{2}{2}$}
{$\dfrac{1}{8}$}
{$\dfrac{2}{4}$}
\loigiai{
$\dfrac{3}{4} - \dfrac{2}{4} = \dfrac{1}{4}$.\\
\textbf{Đáp án của câu hỏi là A}
}
\end{ex}

\begin{ex} %[4D11H61]
Kết quả của phép tính $1 - \dfrac{3}{5}$ là
\choice
{$\dfrac{3}{5}$}
{\True $\dfrac{2}{5}$}
{$\dfrac{4}{5}$}
{$\dfrac{1}{5}$}
\loigiai{
$\dfrac{5}{5} - \dfrac{3}{5} = \dfrac{2}{5}$.\\
\textbf{Đáp án của câu hỏi là B}
}
\end{ex}

\begin{ex} %[4D11H61]
Kết quả của phép tính $\dfrac{5}{6} - \dfrac{2}{5}$ là
\choice
{$\dfrac{3}{1}$}
{$\dfrac{11}{30}$}
{$\dfrac{7}{30}$}
{\True $\dfrac{13}{30}$}
\loigiai{
$\dfrac{25}{30} - \dfrac{12}{30} = \dfrac{13}{30}$.\\
\textbf{Đáp án của câu hỏi là D}
}
\end{ex}

\begin{ex} %[4D11H61]
Tìm $x$ biết: $\dfrac{7}{9} - x = \dfrac{2}{9}$.
\choice
{$\dfrac{9}{9}$}
{$\dfrac{4}{9}$}
{\True $\dfrac{5}{9}$}
{$\dfrac{1}{9}$}
\loigiai{
$x = \dfrac{7}{9} - \dfrac{2}{9} = \dfrac{5}{9}$.\\
\textbf{Đáp án của câu hỏi là C}
}
\end{ex}

\begin{ex} %[4D11V61]
Một sợi dây dài $\dfrac{9}{10}\text{ m}$, người ta cắt đi $\dfrac{3}{5}\text{ m}$. Độ dài đoạn dây còn lại là
\choice
{$\dfrac{6}{5}\text{ m}$}
{\True $\dfrac{3}{10}\text{ m}$}
{$\dfrac{6}{10}\text{ m}$}
{$\dfrac{1}{10}\text{ m}$}
\loigiai{
$\dfrac{9}{10} - \dfrac{6}{10} = \dfrac{3}{10}\text{ m}$.\\
\textbf{Đáp án của câu hỏi là B}
}
\end{ex}

\begin{ex} %[4D11V61]
Một bể nước có $\dfrac{7}{8}$ bể nước. Người ta đã dùng $\dfrac{1}{2}$ bể nước. Lượng nước còn lại trong bể là
\choice
{$\dfrac{6}{6}$ bể}
{$\dfrac{5}{8}$ bể}
{\True $\dfrac{3}{8}$ bể}
{$\dfrac{1}{4}$ bể}
\loigiai{
$\dfrac{7}{8} - \dfrac{4}{8} = \dfrac{3}{8}$ bể.\\
\textbf{Đáp án của câu hỏi là C}
}
\end{ex}

\begin{ex} %[4D11V61]
Tính giá trị biểu thức: $2 - \dfrac{3}{4} - \dfrac{1}{2}$.
\choice
{\True $\dfrac{3}{4}$}
{$1$}
{$\dfrac{1}{4}$}
{$\dfrac{1}{2}$}
\loigiai{
$\dfrac{8}{4} - \dfrac{3}{4} - \dfrac{2}{4} = \dfrac{3}{4}$.\\
\textbf{Đáp án của câu hỏi là A}
}
\end{ex}

\begin{ex} %[4D11V61]
Một mảnh đất hình chữ nhật có nửa chu vi là $\dfrac{4}{5}\text{ km}$, chiều rộng là $\dfrac{1}{4}\text{ km}$. Chiều dài của mảnh đất là
\choice
{$\dfrac{3}{20}\text{ km}$}
{$\dfrac{7}{20}\text{ km}$}
{$\dfrac{9}{20}\text{ km}$}
{\True $\dfrac{11}{20}\text{ km}$}
\loigiai{
$\dfrac{4}{5} - \dfrac{1}{4} = \dfrac{16 - 5}{20} = \dfrac{11}{20}\text{ km}$.\\
\textbf{Đáp án của câu hỏi là D}
}
\end{ex}

\begin{ex} %[4D11V61]
Tìm $x$ biết: $x + \dfrac{3}{7} = \dfrac{11}{14}$.
\choice
{\True $\dfrac{5}{14}$}
{$\dfrac{4}{14}$}
{$\dfrac{6}{14}$}
{$\dfrac{8}{14}$}
\loigiai{
$x = \dfrac{11}{14} - \dfrac{6}{14} = \dfrac{5}{14}$.\\
\textbf{Đáp án của câu hỏi là A}
}
\end{ex}

\begin{ex} %[4D11C61]
Tính giá trị biểu thức: $A = \left(1 - \dfrac{1}{2}\right) + \left(1 - \dfrac{1}{3}\right) + \left(1 - \dfrac{1}{6}\right)$.
\choice
{$1$}
{$\dfrac{3}{2}$}
{$\dfrac{5}{2}$}
{\True $2$}
\loigiai{
$A = \dfrac{1}{2} + \dfrac{2}{3} + \dfrac{5}{6} = \dfrac{3 + 4 + 5}{6} = \dfrac{12}{6} = 2$.\\
\textbf{Đáp án của câu hỏi là D}
}
\end{ex}

\textbf{Phần 2. Trắc nghiệm đúng sai}
\setcounterpageref{ex}{0}

\begin{ex} %[4D11H61]
Xét tính đúng sai của các phép trừ phân số:
\choiceTF
{\True $\dfrac{7}{9} - \dfrac{4}{9} = \dfrac{3}{9} = \dfrac{1}{3}$.}
{$1 - \dfrac{1}{4} = \dfrac{1}{4}$.}
{\True $1 - \dfrac{1}{4} = \dfrac{3}{4}$.}
{\True $\dfrac{5}{6} - \dfrac{1}{3} = \dfrac{3}{6} = \dfrac{1}{2}$.}
\loigiai{
\begin{itemize}
    \item $7/9 - 4/9 = 1/3$. Đúng.
    \item $1 - 1/4 = 3/4 \neq 1/4$. Ý 2 Sai, Ý 3 Đúng.
    \item $5/6 - 2/6 = 3/6 = 1/2$. Đúng.
\end{itemize}
\begin{itemize}
    \item a) \textbf{Đúng}.
    \item b) \textbf{Sai}.
    \item c) \textbf{Đúng}.
    \item d) \textbf{Đúng}.
\end{itemize}
}
\end{ex}

\begin{ex} %[4D11V61]
Một tấm vải dài $1\text{ m}$. Lần đầu cắt $\dfrac{1}{5}\text{ m}$, lần sau cắt $\dfrac{2}{5}\text{ m}$. Xét tính đúng sai:
\choiceTF
{\True Cả hai lần cắt được $\dfrac{3}{5}\text{ m}$ vải.}
{\True Tấm vải còn lại $\dfrac{2}{5}\text{ m}$.}
{Lần sau cắt nhiều hơn lần đầu $\dfrac{2}{5}\text{ m}$.}
{\True Lần sau cắt nhiều hơn lần đầu $\dfrac{1}{5}\text{ m}$.}
\loigiai{
\begin{itemize}
    \item Đã cắt: $1/5 + 2/5 = 3/5\text{ m}$. Đúng.
    \item Còn lại: $1 - 3/5 = 2/5\text{ m}$. Đúng.
    \item Hiệu: $2/5 - 1/5 = 1/5\text{ m} \neq 2/5\text{ m}$. Ý 3 Sai, Ý 4 Đúng.
\end{itemize}
\begin{itemize}
    \item a) \textbf{Đúng}.
    \item b) \textbf{Đúng}.
    \item c) \textbf{Sai}.
    \item d) \textbf{Đúng}.
\end{itemize}
}
\end{ex}

\textbf{Phần 3. Trả lời ngắn}
\setcounterpageref{ex}{0}

\begin{ex} %[4D11H61]
Tính hiệu: $\dfrac{8}{11} - \dfrac{3}{11}$ (dạng a/b).
\shortans{5/11}
\loigiai{
\textbf{Đáp số:} $5/11$\\
$5/11$.
}
\end{ex}

\begin{ex} %[4D11H61]
Tính hiệu: $\dfrac{5}{6} - \dfrac{1}{2}$ (rút gọn tối giản dạng a/b).
\shortans{1/3}
\loigiai{
\textbf{Đáp số:} $1/3$\\
$5/6 - 3/6 = 2/6 = 1/3$.
}
\end{ex}

\begin{ex} %[4D11V61]
Tính giá trị: $1 - \dfrac{5}{8}$ (dạng a/b).
\shortans{3/8}
\loigiai{
\textbf{Đáp số:} $3/8$\\
$3/8$.
}
\end{ex}

\begin{ex} %[4D11V61]
Tìm $x$ biết: $\dfrac{9}{10} - x = \dfrac{2}{5}$ (dạng a/b).
\shortans{1/2}
\loigiai{
\textbf{Đáp số:} $1/2$\\
$x = 9/10 - 4/10 = 5/10 = 1/2$.
}
\end{ex}

\begin{ex} %[4D11V61]
Tính giá trị biểu thức: $2 - \dfrac{2}{3}$ (dạng a/b).
\shortans{4/3}
\loigiai{
\textbf{Đáp số:} $4/3$\\
$6/3 - 2/3 = 4/3$.
}
\end{ex}

\begin{ex} %[4D11C61]
Tính hiệu: $1 - \dfrac{1}{2} - \dfrac{1}{4} - \dfrac{1}{8}$ (dạng a/b).
\shortans{1/8}
\loigiai{
\textbf{Đáp số:} $1/8$\\
$8/8 - 4/8 - 2/8 - 1/8 = 1/8$.
}
\end{ex}

\newpage

\TaoPhieuKiemTra{Chủ đề 11 -- Phép cộng, phép trừ phân số}{62}{Luyện tập chung cộng trừ phân số}{https://loptoanthaylinh.vn/thi-online/toan-4/bai-62}
\textbf{Phần 1. Trắc nghiệm lựa chọn}
\setcounterpageref{ex}{0}

\begin{ex} %[4D11N62]
Kết quả của phép tính $\dfrac{3}{5} + \dfrac{1}{5}$ là
\choice
{$\dfrac{4}{10}$}
{$\dfrac{2}{5}$}
{\True $\dfrac{4}{5}$}
{$\dfrac{3}{10}$}
\loigiai{
$\dfrac{3 + 1}{5} = \dfrac{4}{5}$.\\
\textbf{Đáp án của câu hỏi là C}
}
\end{ex}

\begin{ex} %[4D11N62]
Kết quả của phép trừ $\dfrac{7}{9} - \dfrac{2}{9}$ là
\choice
{$\dfrac{5}{0}$}
{$\dfrac{9}{9}$}
{$\dfrac{4}{9}$}
{\True $\dfrac{5}{9}$}
\loigiai{
$\dfrac{7 - 2}{9} = \dfrac{5}{9}$.\\
\textbf{Đáp án của câu hỏi là D}
}
\end{ex}

\begin{ex} %[4D11H62]
Kết quả của phép tính $\dfrac{1}{2} + \dfrac{3}{4}$ là
\choice
{\True $\dfrac{5}{4}$}
{$\dfrac{4}{6}$}
{$1$}
{$\dfrac{3}{8}$}
\loigiai{
$\dfrac{2}{4} + \dfrac{3}{4} = \dfrac{5}{4}$.\\
\textbf{Đáp án của câu hỏi là A}
}
\end{ex}

\begin{ex} %[4D11H62]
Kết quả của phép tính $1 - \dfrac{4}{7}$ là
\choice
{$\dfrac{4}{7}$}
{$\dfrac{5}{7}$}
{$\dfrac{2}{7}$}
{\True $\dfrac{3}{7}$}
\loigiai{
$\dfrac{7}{7} - \dfrac{4}{7} = \dfrac{3}{7}$.\\
\textbf{Đáp án của câu hỏi là D}
}
\end{ex}

\begin{ex} %[4D11H62]
Tính thuận tiện: $\dfrac{3}{8} + \dfrac{5}{12} + \dfrac{5}{8}$.
\choice
{$1$}
{$\dfrac{13}{12}$}
{\True $\dfrac{17}{12}$}
{$2$}
\loigiai{
$\left(\dfrac{3}{8} + \dfrac{5}{8}\right) + \dfrac{5}{12} = 1 + \dfrac{5}{12} = \dfrac{17}{12}$.\\
\textbf{Đáp án của câu hỏi là C}
}
\end{ex}

\begin{ex} %[4D11H62]
Tìm $x$ biết: $x + \dfrac{1}{3} = \dfrac{5}{6}$.
\choice
{\True $\dfrac{1}{2}$}
{$\dfrac{4}{6}$}
{$\dfrac{2}{3}$}
{$\dfrac{1}{6}$}
\loigiai{
$x = \dfrac{5}{6} - \dfrac{2}{6} = \dfrac{3}{6} = \dfrac{1}{2}$.\\
\textbf{Đáp án của câu hỏi là A}
}
\end{ex}

\begin{ex} %[4D11V62]
Một mảnh vườn hình chữ nhật có chiều dài $\dfrac{4}{5}\text{ m}$, chiều rộng ngắn hơn chiều dài $\dfrac{1}{3}\text{ m}$. Chiều rộng mảnh vườn là
\choice
{$\dfrac{3}{2}\text{ m}$}
{\True $\dfrac{7}{15}\text{ m}$}
{$\dfrac{1}{2}\text{ m}$}
{$\dfrac{2}{15}\text{ m}$}
\loigiai{
$\dfrac{4}{5} - \dfrac{1}{3} = \dfrac{12 - 5}{15} = \dfrac{7}{15}\text{ m}$.\\
\textbf{Đáp án của câu hỏi là B}
}
\end{ex}

\begin{ex} %[4D11V62]
Tính giá trị biểu thức: $\dfrac{7}{10} + \dfrac{1}{2} - \dfrac{3}{5}$.
\choice
{\True $\dfrac{3}{5}$}
{$\dfrac{5}{10}$}
{$\dfrac{4}{5}$}
{$1$}
\loigiai{
$\dfrac{7}{10} + \dfrac{5}{10} - \dfrac{6}{10} = \dfrac{6}{10} = \dfrac{3}{5}$.\\
\textbf{Đáp án của câu hỏi là A}
}
\end{ex}

\begin{ex} %[4D11V62]
Hai tổ công nhân cùng sửa một đoạn đường: Tổ 1 sửa được $\dfrac{3}{8}$ đoạn đường, tổ 2 sửa được $\dfrac{2}{5}$ đoạn đường. Cả hai tổ sửa được số phần đoạn đường là
\choice
{$\dfrac{5}{13}$}
{\True $\dfrac{31}{40}$}
{$\dfrac{29}{40}$}
{$\dfrac{1}{8}$}
\loigiai{
$\dfrac{3}{8} + \dfrac{2}{5} = \dfrac{15 + 16}{40} = \dfrac{31}{40}$.\\
\textbf{Đáp án của câu hỏi là B}
}
\end{ex}

\begin{ex} %[4D11V62]
Một hộp kẹo có 24 cái. Bạn An lấy $\dfrac{1}{3}$ số kẹo, bạn Bình lấy $\dfrac{1}{4}$ số kẹo. Phân số chỉ số kẹo cả hai bạn đã lấy là
\choice
{$\dfrac{2}{7}$}
{$\dfrac{5}{12}$}
{\True $\dfrac{7}{12}$}
{$\dfrac{1}{12}$}
\loigiai{
$\dfrac{1}{3} + \dfrac{1}{4} = \dfrac{7}{12}$.\\
\textbf{Đáp án của câu hỏi là C}
}
\end{ex}

\begin{ex} %[4D11V62]
Tìm $x$ biết: $\dfrac{7}{4} - x = \dfrac{3}{4} + \dfrac{1}{2}$.
\choice
{$\dfrac{3}{4}$}
{$1$}
{$\dfrac{1}{4}$}
{\True $\dfrac{1}{2}$}
\loigiai{
$\dfrac{7}{4} - x = \dfrac{5}{4} \Rightarrow x = \dfrac{7}{4} - \dfrac{5}{4} = \dfrac{2}{4} = \dfrac{1}{2}$.\\
\textbf{Đáp án của câu hỏi là D}
}
\end{ex}

\begin{ex} %[4D11C62]
Tính giá trị biểu thức: $S = \dfrac{1}{2} + \dfrac{1}{6} + \dfrac{1}{12} + \dfrac{1}{20} + \dfrac{1}{30}$.
\choice
{$\dfrac{4}{5}$}
{\True $\dfrac{5}{6}$}
{$\dfrac{6}{7}$}
{$1$}
\loigiai{
$S = \left(1 - \dfrac{1}{2}\right) + \left(\dfrac{1}{2} - \dfrac{1}{3}\right) + \dots + \left(\dfrac{1}{5} - \dfrac{1}{6}\right) = 1 - \dfrac{1}{6} = \dfrac{5}{6}$.\\
\textbf{Đáp án của câu hỏi là B}
}
\end{ex}

\textbf{Phần 2. Trắc nghiệm đúng sai}
\setcounterpageref{ex}{0}

\begin{ex} %[4D11H62]
Xét tính đúng sai của các phép tính:
\choiceTF
{\True $\dfrac{2}{3} + \dfrac{1}{6} = \dfrac{4}{6} + \dfrac{1}{6} = \dfrac{5}{6}$.}
{\True $1 - \dfrac{2}{5} = \dfrac{3}{5}$.}
{$\dfrac{3}{4} - \dfrac{1}{3} = \dfrac{2}{1} = 2$.}
{\True $\dfrac{3}{4} - \dfrac{1}{3} = \dfrac{9 - 4}{12} = \dfrac{5}{12}$.}
\loigiai{
\begin{itemize}
    \item $2/3 + 1/6 = 5/6$. Đúng.
    \item $1 - 2/5 = 3/5$. Đúng.
    \item Phép trừ phân số phải quy đồng, không thể lấy tử trừ tử mẫu trừ mẫu. Ý 3 Sai, Ý 4 Đúng.
\end{itemize}
\begin{itemize}
    \item a) \textbf{Đúng}.
    \item b) \textbf{Đúng}.
    \item c) \textbf{Sai}.
    \item d) \textbf{Đúng}.
\end{itemize}
}
\end{ex}

\begin{ex} %[4D11V62]
Một chai đựng $1\text{ l}$ nước ngọt. Bạn Lan uống $\dfrac{1}{4}\text{ l}$, bạn Mai uống $\dfrac{3}{8}\text{ l}$. Xét tính đúng sai:
\choiceTF
{\True Cả hai bạn uống hết $\dfrac{5}{8}\text{ l}$ nước ngọt.}
{\True Bạn Mai uống nhiều hơn bạn Lan là $\dfrac{1}{8}\text{ l}$.}
{Trong chai còn lại $\dfrac{1}{2}\text{ l}$ nước ngọt.}
{\True Trong chai còn lại $\dfrac{3}{8}\text{ l}$ nước ngọt.}
\loigiai{
\begin{itemize}
    \item Tổng uống: $2/8 + 3/8 = 5/8$ l. Đúng.
    \item Hiệu: $3/8 - 2/8 = 1/8$ l. Đúng.
    \item Còn lại: $1 - 5/8 = 3/8$ l $\neq 1/2$ l. Ý 3 Sai, Ý 4 Đúng.
\end{itemize}
\begin{itemize}
    \item a) \textbf{Đúng}.
    \item b) \textbf{Đúng}.
    \item c) \textbf{Sai}.
    \item d) \textbf{Đúng}.
\end{itemize}
}
\end{ex}

\textbf{Phần 3. Trả lời ngắn}
\setcounterpageref{ex}{0}

\begin{ex} %[4D11H62]
Tính tổng: $\dfrac{1}{3} + \dfrac{2}{5}$ (dạng a/b).
\shortans{11/15}
\loigiai{
\textbf{Đáp số:} $11/15$\\
$5/15 + 6/15 = 11/15$.
}
\end{ex}

\begin{ex} %[4D11H62]
Tính hiệu: $\dfrac{5}{6} - \dfrac{1}{4}$ (dạng a/b).
\shortans{7/12}
\loigiai{
\textbf{Đáp số:} $7/12$\\
$10/12 - 3/12 = 7/12$.
}
\end{ex}

\begin{ex} %[4D11V62]
Tính giá trị: $1 + \dfrac{2}{7} - \dfrac{3}{7}$ (dạng a/b).
\shortans{6/7}
\loigiai{
\textbf{Đáp số:} $6/7$\\
$7/7 + 2/7 - 3/7 = 6/7$.
}
\end{ex}

\begin{ex} %[4D11V62]
Tìm $x$ biết: $x - \dfrac{3}{10} = \dfrac{1}{2}$ (dạng a/b).
\shortans{4/5}
\loigiai{
\textbf{Đáp số:} $4/5$\\
$x = 5/10 + 3/10 = 8/10 = 4/5$.
}
\end{ex}

\begin{ex} %[4D11V62]
Một hình tam giác có chu vi $\dfrac{9}{10}\text{ m}$, hai cạnh có độ dài là $\dfrac{1}{5}\text{ m}$ và $\dfrac{3}{10}\text{ m}$. Tính độ dài cạnh thứ ba (rút gọn tối giản dạng a/b).
\shortans{2/5}
\loigiai{
\textbf{Đáp số:} $2/5$\\
$9/10 - 2/10 - 3/10 = 4/10 = 2/5\text{ m}$.
}
\end{ex}

\begin{ex} %[4D11C62]
Tính nhanh: $\dfrac{1}{3} + \dfrac{1}{4} + \dfrac{2}{3} + \dfrac{3}{4}$.
\shortans{2}
\loigiai{
\textbf{Đáp số:} $2$\\
$(1/3 + 2/3) + (1/4 + 3/4) = 1 + 1 = 2$.
}
\end{ex}

\newpage

% =========================================================================
% CHỦ ĐỀ 12 - PHÉP NHÂN, PHÉP CHIA PHÂN SỐ
% =========================================================================
\section{Phép nhân, phép chia phân số}

\TaoPhieuKiemTra{Chủ đề 12 -- Phép nhân, phép chia phân số}{63}{Phép nhân phân số}{https://loptoanthaylinh.vn/thi-online/toan-4/bai-63}
\textbf{Phần 1. Trắc nghiệm lựa chọn}
\setcounterpageref{ex}{0}

\begin{ex} %[4D12N63]
Muốn nhân hai phân số, ta
\choice
{Quy đồng mẫu số rồi nhân tử số}
{\True Lấy tử số nhân với tử số, mẫu số nhân với mẫu số}
{Lấy tử số nhân mẫu số chéo}
{Giữ nguyên mẫu số, nhân các tử số}
\loigiai{
Lấy tử số nhân với tử số, mẫu số nhân với mẫu số.\\
\textbf{Đáp án của câu hỏi là B}
}
\end{ex}

\begin{ex} %[4D12N63]
Kết quả của phép tính $\dfrac{2}{3} \times \dfrac{4}{5}$ là
\choice
{$\dfrac{6}{8}$}
{$\dfrac{8}{8}$}
{\True $\dfrac{8}{15}$}
{$\dfrac{10}{12}$}
\loigiai{
$\dfrac{2 \times 4}{3 \times 5} = \dfrac{8}{15}$.\\
\textbf{Đáp án của câu hỏi là C}
}
\end{ex}

\begin{ex} %[4D12H63]
Kết quả của phép tính $\dfrac{3}{8} \times 4$ là
\choice
{$\dfrac{12}{32}$}
{$\dfrac{7}{8}$}
{$3$}
{\True $\dfrac{3}{2}$}
\loigiai{
$\dfrac{3 \times 4}{8} = \dfrac{12}{8} = \dfrac{3}{2}$.\\
\textbf{Đáp án của câu hỏi là D}
}
\end{ex}

\begin{ex} %[4D12H63]
Kết quả của phép tính $\dfrac{5}{7} \times \dfrac{7}{5}$ là
\choice
{\True $1$}
{$\dfrac{35}{10}$}
{$0$}
{$\dfrac{25}{49}$}
\loigiai{
$\dfrac{5 \times 7}{7 \times 5} = 1$.\\
\textbf{Đáp án của câu hỏi là A}
}
\end{ex}

\begin{ex} %[4D12H63]
Một hình vuông có cạnh dài $\dfrac{3}{4}\text{ m}$. Diện tích hình vuông đó là
\choice
{$3\text{ m}^2$}
{$\dfrac{6}{8}\text{ m}^2$}
{\True $\dfrac{9}{16}\text{ m}^2$}
{$\dfrac{3}{16}\text{ m}^2$}
\loigiai{
$\dfrac{3}{4} \times \dfrac{3}{4} = \dfrac{9}{16}\text{ m}^2$.\\
\textbf{Đáp án của câu hỏi là C}
}
\end{ex}

\begin{ex} %[4D12H63]
Tính thuận tiện: $\dfrac{2}{5} \times \dfrac{3}{7} \times \dfrac{5}{2}$.
\choice
{$1$}
{\True $\dfrac{3}{7}$}
{$\dfrac{6}{35}$}
{$\dfrac{15}{14}$}
\loigiai{
$\left(\dfrac{2}{5} \times \dfrac{5}{2}\right) \times \dfrac{3}{7} = 1 \times \dfrac{3}{7} = \dfrac{3}{7}$.\\
\textbf{Đáp án của câu hỏi là B}
}
\end{ex}

\begin{ex} %[4D12V63]
Một hình chữ nhật có chiều dài $\dfrac{4}{5}\text{ m}$, chiều rộng $\dfrac{3}{8}\text{ m}$. Diện tích hình chữ nhật đó là
\choice
{$\dfrac{7}{13}\text{ m}^2$}
{$\dfrac{12}{40}\text{ m}^2$}
{$\dfrac{3}{5}\text{ m}^2$}
{\True $\dfrac{3}{10}\text{ m}^2$}
\loigiai{
$\dfrac{4}{5} \times \dfrac{3}{8} = \dfrac{12}{40} = \dfrac{3}{10}\text{ m}^2$.\\
\textbf{Đáp án của câu hỏi là D}
}
\end{ex}

\begin{ex} %[4D12V63]
Một tấm kính hình thoi có độ dài hai đường chéo là $\dfrac{3}{4}\text{ m}$ và $\dfrac{2}{5}\text{ m}$. Diện tích tấm kính đó là
\choice
{$\dfrac{3}{10}\text{ m}^2$}
{\True $\dfrac{3}{20}\text{ m}^2$}
{$\dfrac{6}{20}\text{ m}^2$}
{$\dfrac{3}{40}\text{ m}^2$}
\loigiai{
$S = \dfrac{1}{2} \times \dfrac{3}{4} \times \dfrac{2}{5} = \dfrac{3}{20}\text{ m}^2$.\\
\textbf{Đáp án của câu hỏi là B}
}
\end{ex}

\begin{ex} %[4D12V63]
Tính giá trị biểu thức: $\dfrac{3}{4} \times \dfrac{5}{6} + \dfrac{3}{4} \times \dfrac{1}{6}$.
\choice
{\True $\dfrac{3}{4}$}
{$\dfrac{1}{2}$}
{$1$}
{$\dfrac{5}{8}$}
\loigiai{
$\dfrac{3}{4} \times \left(\dfrac{5}{6} + \dfrac{1}{6}\right) = \dfrac{3}{4} \times 1 = \dfrac{3}{4}$.\\
\textbf{Đáp án của câu hỏi là A}
}
\end{ex}

\begin{ex} %[4D12V63]
Một ô tô mỗi phút đi được $\dfrac{4}{5}\text{ km}$. Hỏi trong $15$ phút ô tô đó đi được bao nhiêu ki-lô-mét?
\choice
{$10\text{ km}$}
{$14\text{ km}$}
{$15\text{ km}$}
{\True $12\text{ km}$}
\loigiai{
$\dfrac{4}{5} \times 15 = 12\text{ km}$.\\
\textbf{Đáp án của câu hỏi là D}
}
\end{ex}

\begin{ex} %[4D12V63]
Tính: $\dfrac{2 \times 3 \times 4}{3 \times 4 \times 5}$.
\choice
{$\dfrac{1}{5}$}
{$\dfrac{3}{5}$}
{\True $\dfrac{2}{5}$}
{$\dfrac{4}{5}$}
\loigiai{
Rút gọn cho $3 \times 4$ ta được $\dfrac{2}{5}$.\\
\textbf{Đáp án của câu hỏi là C}
}
\end{ex}

\begin{ex} %[4D12C63]
Tính tích: $P = \left(1 - \dfrac{1}{2}\right) \times \left(1 - \dfrac{1}{3}\right) \times \left(1 - \dfrac{1}{4}\right) \times \dots \times \left(1 - \dfrac{1}{10}\right)$.
\choice
{\True $\dfrac{1}{10}$}
{$\dfrac{1}{5}$}
{$\dfrac{1}{9}$}
{$\dfrac{2}{10}$}
\loigiai{
$P = \dfrac{1}{2} \times \dfrac{2}{3} \times \dfrac{3}{4} \times \dots \times \dfrac{9}{10} = \dfrac{1}{10}$.\\
\textbf{Đáp án của câu hỏi là A}
}
\end{ex}

\textbf{Phần 2. Trắc nghiệm đúng sai}
\setcounterpageref{ex}{0}

\begin{ex} %[4D12H63]
Xét tính đúng sai của các phép nhân phân số:
\choiceTF
{\True $\dfrac{1}{3} \times \dfrac{2}{5} = \dfrac{2}{15}$.}
{\True $\dfrac{3}{4} \times 0 = 0$.}
{$\dfrac{5}{7} \times 1 = \dfrac{7}{5}$.}
{\True $\dfrac{5}{7} \times 1 = \dfrac{5}{7}$.}
\loigiai{
\begin{itemize}
    \item $1/3 \times 2/5 = 2/15$. Đúng.
    \item Nhân với 0 bằng 0: Đúng.
    \item Nhân với 1 bằng chính nó $\neq 7/5$. Ý 3 Sai, Ý 4 Đúng.
\end{itemize}
\begin{itemize}
    \item a) \textbf{Đúng}.
    \item b) \textbf{Đúng}.
    \item c) \textbf{Sai}.
    \item d) \textbf{Đúng}.
\end{itemize}
}
\end{ex}

\begin{ex} %[4D12V63]
Một mảnh bìa hình chữ nhật có chiều dài $\dfrac{3}{4}\text{ m}$, chiều rộng $\dfrac{2}{3}\text{ m}$. Xét tính đúng sai:
\choiceTF
{\True Chu vi mảnh bìa là $\dfrac{17}{6}\text{ m}$.}
{\True Diện tích mảnh bìa là $\dfrac{1}{2}\text{ m}^2$.}
{Diện tích mảnh bìa là $\dfrac{6}{7}\text{ m}^2$.}
{\True Nếu tăng chiều dài gấp đôi thì diện tích tăng gấp đôi ($1\text{ m}^2$).}
\loigiai{
\begin{itemize}
    \item Chu vi: $(3/4 + 2/3) \times 2 = 17/12 \times 2 = 17/6\text{ m}$. Đúng.
    \item Diện tích: $3/4 \times 2/3 = 6/12 = 1/2\text{ m}^2$. Đúng.
    \item $1/2 \neq 6/7$. Sai.
    \item Gấp đôi dài thì diện tích gấp đôi $= 1\text{ m}^2$: Đúng.
\end{itemize}
\begin{itemize}
    \item a) \textbf{Đúng}.
    \item b) \textbf{Đúng}.
    \item c) \textbf{Sai}.
    \item d) \textbf{Đúng}.
\end{itemize}
}
\end{ex}

\textbf{Phần 3. Trả lời ngắn}
\setcounterpageref{ex}{0}

\begin{ex} %[4D12H63]
Tính tích: $\dfrac{3}{5} \times \dfrac{5}{9}$ (rút gọn tối giản dạng a/b).
\shortans{1/3}
\loigiai{
\textbf{Đáp số:} $1/3$\\
$15/45 = 1/3$.
}
\end{ex}

\begin{ex} %[4D12H63]
Tính tích: $\dfrac{7}{8} \times 16$.
\shortans{14}
\loigiai{
\textbf{Đáp số:} $14$\\
$7 \times 2 = 14$.
}
\end{ex}

\begin{ex} %[4D12V63]
Một hình vuông có cạnh dài $\dfrac{2}{5}\text{ m}$. Tính chu vi hình vuông đó (dạng a/b).
\shortans{8/5}
\loigiai{
\textbf{Đáp số:} $8/5$\\
$2/5 \times 4 = 8/5\text{ m}$.
}
\end{ex}

\begin{ex} %[4D12V63]
Tính diện tích hình chữ nhật có kích thước $\dfrac{5}{6}\text{ m}$ và $\dfrac{3}{10}\text{ m}$ (rút gọn tối giản dạng a/b).
\shortans{1/4}
\loigiai{
\textbf{Đáp số:} $1/4$\\
$15/60 = 1/4\text{ m}^2$.
}
\end{ex}

\begin{ex} %[4D12V63]
Tính thuận tiện: $\dfrac{4}{9} \times \dfrac{7}{11} + \dfrac{4}{9} \times \dfrac{4}{11}$ (dạng a/b).
\shortans{4/9}
\loigiai{
\textbf{Đáp số:} $4/9$\\
$4/9 \times 1 = 4/9$.
}
\end{ex}

\begin{ex} %[4D12C63]
Tính tích: $\dfrac{1}{2} \times \dfrac{2}{3} \times \dfrac{3}{4} \times \dfrac{4}{5} \times \dfrac{5}{6}$ (dạng a/b).
\shortans{1/6}
\loigiai{
\textbf{Đáp số:} $1/6$\\
$1/6$.
}
\end{ex}

\newpage

\TaoPhieuKiemTra{Chủ đề 12 -- Phép nhân, phép chia phân số}{64}{Phép chia phân số}{https://loptoanthaylinh.vn/thi-online/toan-4/bai-64}
\textbf{Phần 1. Trắc nghiệm lựa chọn}
\setcounterpageref{ex}{0}

\begin{ex} %[4D12N64]
Phân số đảo ngược của phân số $\dfrac{3}{5}$ là
\choice
{$\dfrac{3}{5}$}
{$-\dfrac{3}{5}$}
{\True $\dfrac{5}{3}$}
{$\dfrac{1}{5}$}
\loigiai{
Đảo ngược tử số và mẫu số: $\dfrac{5}{3}$.\\
\textbf{Đáp án của câu hỏi là C}
}
\end{ex}

\begin{ex} %[4D12N64]
Muốn chia một phân số cho một phân số, ta
\choice
{Lấy tử số chia tử số, mẫu số chia mẫu số}
{Lấy mẫu số nhân tử số}
{Chia hai tử số với nhau}
{\True Lấy phân số thứ nhất nhân với phân số thứ hai đảo ngược}
\loigiai{
Nhân phân số thứ nhất với phân số thứ hai đảo ngược.\\
\textbf{Đáp án của câu hỏi là D}
}
\end{ex}

\begin{ex} %[4D12H64]
Kết quả của phép tính $\dfrac{2}{3} : \dfrac{3}{4}$ là
\choice
{$\dfrac{6}{12}$}
{\True $\dfrac{8}{9}$}
{$\dfrac{1}{2}$}
{$\dfrac{9}{8}$}
\loigiai{
$\dfrac{2}{3} \times \dfrac{4}{3} = \dfrac{8}{9}$.\\
\textbf{Đáp án của câu hỏi là B}
}
\end{ex}

\begin{ex} %[4D12H64]
Kết quả của phép tính $\dfrac{4}{5} : 2$ là
\choice
{$\dfrac{8}{5}$}
{$\dfrac{4}{10}$}
{\True $\dfrac{2}{5}$}
{$\dfrac{1}{5}$}
\loigiai{
$\dfrac{4}{5} \times \dfrac{1}{2} = \dfrac{4}{10} = \dfrac{2}{5}$.\\
\textbf{Đáp án của câu hỏi là C}
}
\end{ex}

\begin{ex} %[4D12H64]
Kết quả của phép tính $3 : \dfrac{1}{2}$ là
\choice
{$\dfrac{3}{2}$}
{\True $6$}
{$\dfrac{1}{6}$}
{$1$}
\loigiai{
$3 \times \dfrac{2}{1} = 6$.\\
\textbf{Đáp án của câu hỏi là B}
}
\end{ex}

\begin{ex} %[4D12H64]
Tìm $x$ biết: $x \times \dfrac{2}{5} = \dfrac{4}{7}$.
\choice
{\True $\dfrac{10}{7}$}
{$\dfrac{8}{35}$}
{$\dfrac{7}{10}$}
{$\dfrac{6}{12}$}
\loigiai{
$x = \dfrac{4}{7} : \dfrac{2}{5} = \dfrac{4}{7} \times \dfrac{5}{2} = \dfrac{10}{7}$.\\
\textbf{Đáp án của câu hỏi là A}
}
\end{ex}

\begin{ex} %[4D12V64]
Một hình chữ nhật có diện tích $\dfrac{2}{3}\text{ m}^2$, chiều dài $\dfrac{4}{5}\text{ m}$. Chiều rộng của hình chữ nhật là
\choice
{$\dfrac{8}{15}\text{ m}$}
{$\dfrac{6}{5}\text{ m}$}
{$\dfrac{1}{3}\text{ m}$}
{\True $\dfrac{5}{6}\text{ m}$}
\loigiai{
$\dfrac{2}{3} : \dfrac{4}{5} = \dfrac{2}{3} \times \dfrac{5}{4} = \dfrac{5}{6}\text{ m}$.\\
\textbf{Đáp án của câu hỏi là D}
}
\end{ex}

\begin{ex} %[4D12V64]
Có $\dfrac{7}{2}\text{ l}$ dầu chia đều vào các chai, mỗi chai chứa $\dfrac{1}{2}\text{ l}$ dầu. Số chai dầu rót được là
\choice
{\True $7$ chai}
{$6$ chai}
{$8$ chai}
{$9$ chai}
\loigiai{
$\dfrac{7}{2} : \dfrac{1}{2} = \dfrac{7}{2} \times 2 = 7$ chai.\\
\textbf{Đáp án của câu hỏi là A}
}
\end{ex}

\begin{ex} %[4D12V64]
Một tấm vải dài $\dfrac{8}{3}\text{ m}$ được cắt thành các đoạn dài $\dfrac{2}{3}\text{ m}$. Số đoạn vải cắt được là
\choice
{$3$ đoạn}
{\True $4$ đoạn}
{$5$ đoạn}
{$6$ đoạn}
\loigiai{
$\dfrac{8}{3} : \dfrac{2}{3} = 4$ đoạn.\\
\textbf{Đáp án của câu hỏi là B}
}
\end{ex}

\begin{ex} %[4D12V64]
Tìm $x$ biết: $x : \dfrac{3}{8} = \dfrac{4}{9}$.
\choice
{$\dfrac{32}{27}$}
{$\dfrac{7}{17}$}
{$\dfrac{12}{72}$}
{\True $\dfrac{1}{6}$}
\loigiai{
$x = \dfrac{4}{9} \times \dfrac{3}{8} = \dfrac{12}{72} = \dfrac{1}{6}$.\\
\textbf{Đáp án của câu hỏi là D}
}
\end{ex}

\begin{ex} %[4D12V64]
Tính giá trị biểu thức: $\dfrac{3}{5} : \dfrac{2}{3} \times \dfrac{4}{9}$.
\choice
{$\dfrac{1}{5}$}
{$\dfrac{3}{5}$}
{\True $\dfrac{2}{5}$}
{$\dfrac{4}{5}$}
\loigiai{
$\dfrac{3}{5} \times \dfrac{3}{2} \times \dfrac{4}{9} = \dfrac{9}{10} \times \dfrac{4}{9} = \dfrac{4}{10} = \dfrac{2}{5}$.\\
\textbf{Đáp án của câu hỏi là C}
}
\end{ex}

\begin{ex} %[4D12C64]
Một hình bình hành có diện tích $\dfrac{5}{8}\text{ m}^2$, chiều cao $\dfrac{3}{4}\text{ m}$. Độ dài đáy của hình bình hành là
\choice
{\True $\dfrac{5}{6}\text{ m}$}
{$\dfrac{15}{32}\text{ m}$}
{$\dfrac{6}{5}\text{ m}$}
{$\dfrac{1}{2}\text{ m}$}
\loigiai{
$a = S : h = \dfrac{5}{8} : \dfrac{3}{4} = \dfrac{5}{8} \times \dfrac{4}{3} = \dfrac{5}{6}\text{ m}$.\\
\textbf{Đáp án của câu hỏi là A}
}
\end{ex}

\textbf{Phần 2. Trắc nghiệm đúng sai}
\setcounterpageref{ex}{0}

\begin{ex} %[4D12H64]
Xét tính đúng sai của các phép chia phân số:
\choiceTF
{\True Phân số đảo ngược của $\dfrac{4}{7}$ là $\dfrac{7}{4}$.}
{\True $\dfrac{3}{5} : \dfrac{3}{5} = 1$.}
{$\dfrac{2}{3} : 2 = \dfrac{4}{3}$.}
{\True $\dfrac{2}{3} : 2 = \dfrac{1}{3}$.}
\loigiai{
\begin{itemize}
    \item Đảo ngược đúng.
    \item Chia cho chính nó bằng 1: Đúng.
    \item $2/3 : 2 = 1/3 \neq 4/3$. Ý 3 Sai, Ý 4 Đúng.
\end{itemize}
\begin{itemize}
    \item a) \textbf{Đúng}.
    \item b) \textbf{Đúng}.
    \item c) \textbf{Sai}.
    \item d) \textbf{Đúng}.
\end{itemize}
}
\end{ex}

\begin{ex} %[4D12V64]
Người ta chia một sợi dây dài $\dfrac{15}{4}\text{ m}$ thành các đoạn nhỏ bằng nhau:
\choiceTF
{\True Nếu chia thành 3 đoạn bằng nhau thì mỗi đoạn dài $\dfrac{5}{4}\text{ m}$.}
{\True Nếu chia thành các đoạn dài $\dfrac{3}{4}\text{ m}$ thì được 5 đoạn.}
{Nếu chia đôi thì mỗi đoạn dài $\dfrac{15}{2}\text{ m}$.}
{\True Nếu chia đôi thì mỗi đoạn dài $\dfrac{15}{8}\text{ m}$.}
\loigiai{
\begin{itemize}
    \item $15/4 : 3 = 5/4\text{ m}$. Đúng.
    \item $15/4 : 3/4 = 5$ đoạn. Đúng.
    \item Chia đôi: $15/4 : 2 = 15/8\text{ m} \neq 15/2\text{ m}$. Ý 3 Sai, Ý 4 Đúng.
\end{itemize}
\begin{itemize}
    \item a) \textbf{Đúng}.
    \item b) \textbf{Đúng}.
    \item c) \textbf{Sai}.
    \item d) \textbf{Đúng}.
\end{itemize}
}
\end{ex}

\textbf{Phần 3. Trả lời ngắn}
\setcounterpageref{ex}{0}

\begin{ex} %[4D12H64]
Tính thương: $\dfrac{3}{7} : \dfrac{5}{7}$ (dạng a/b).
\shortans{3/5}
\loigiai{
\textbf{Đáp số:} $3/5$\\
$3/5$.
}
\end{ex}

\begin{ex} %[4D12H64]
Tính thương: $\dfrac{4}{9} : \dfrac{2}{3}$ (rút gọn tối giản dạng a/b).
\shortans{2/3}
\loigiai{
\textbf{Đáp số:} $2/3$\\
$4/9 \times 3/2 = 2/3$.
}
\end{ex}

\begin{ex} %[4D12V64]
Tìm $x$ biết: $x \times \dfrac{3}{4} = \dfrac{9}{16}$ (dạng a/b).
\shortans{3/4}
\loigiai{
\textbf{Đáp số:} $3/4$\\
$x = 9/16 : 3/4 = 3/4$.
}
\end{ex}

\begin{ex} %[4D12V64]
Một hình chữ nhật có diện tích $\dfrac{3}{5}\text{ m}^2$, chiều rộng $\dfrac{1}{2}\text{ m}$. Tính chiều dài (dạng a/b).
\shortans{6/5}
\loigiai{
\textbf{Đáp số:} $6/5$\\
$3/5 : 1/2 = 6/5\text{ m}$.
}
\end{ex}

\begin{ex} %[4D12V64]
Có $6\text{ kg}$ đường chia đều vào các túi $\dfrac{3}{4}\text{ kg}$. Tính số túi đường đóng được.
\shortans{8}
\loigiai{
\textbf{Đáp số:} $8$\\
$6 : (3/4) = 8$ túi.
}
\end{ex}

\begin{ex} %[4D12C64]
Tính giá trị biểu thức: $\dfrac{1}{2} : \dfrac{2}{3} : \dfrac{3}{4} : \dfrac{4}{5}$ (dạng a/b).
\shortans{5/4}
\loigiai{
\textbf{Đáp số:} $5/4$\\
$1/2 \times 3/2 \times 4/3 \times 5/4 = 5/4$.
}
\end{ex}

\newpage

\TaoPhieuKiemTra{Chủ đề 12 -- Phép nhân, phép chia phân số}{65}{Tìm phân số của một số}{https://loptoanthaylinh.vn/thi-online/toan-4/bai-65}
\textbf{Phần 1. Trắc nghiệm lựa chọn}
\setcounterpageref{ex}{0}

\begin{ex} %[4D12N65]
Muốn tìm phân số $\dfrac{a}{b}$ của số $M$, ta tính
\choice
{\True $M \times \dfrac{a}{b}$ (hoặc lấy $M : b \times a$)}
{$M : a \times b$}
{$M + \dfrac{a}{b}$}
{$M - \dfrac{a}{b}$}
\loigiai{
Ta lấy số $M$ nhân với phân số $\dfrac{a}{b}$.\\
\textbf{Đáp án của câu hỏi là A}
}
\end{ex}

\begin{ex} %[4D12N65]
$\dfrac{2}{3}$ của $18$ là
\choice
{$6$}
{\True $12$}
{$9$}
{$15$}
\loigiai{
$18 \times \dfrac{2}{3} = 12$.\\
\textbf{Đáp án của câu hỏi là B}
}
\end{ex}

\begin{ex} %[4D12H65]
$\dfrac{3}{5}$ của $100\text{ kg}$ là
\choice
{\True $60\text{ kg}$}
{$50\text{ kg}$}
{$70\text{ kg}$}
{$40\text{ kg}$}
\loigiai{
$100 \times \dfrac{3}{5} = 60\text{ kg}$.\\
\textbf{Đáp án của câu hỏi là A}
}
\end{ex}

\begin{ex} %[4D12H65]
$\dfrac{3}{4}$ của $1$ giờ là bao nhiêu phút?
\choice
{$40$ phút}
{$50$ phút}
{\True $45$ phút}
{$30$ phút}
\loigiai{
$60 \times \dfrac{3}{4} = 45$ phút.\\
\textbf{Đáp án của câu hỏi là C}
}
\end{ex}

\begin{ex} %[4D12H65]
Một lớp học có 36 học sinh, trong đó $\dfrac{4}{9}$ số học sinh là nữ. Số học sinh nữ của lớp đó là
\choice
{$14$ học sinh}
{$18$ học sinh}
{$20$ học sinh}
{\True $16$ học sinh}
\loigiai{
$36 \times \dfrac{4}{9} = 16$ học sinh.\\
\textbf{Đáp án của câu hỏi là D}
}
\end{ex}

\begin{ex} %[4D12H65]
$\dfrac{5}{6}$ của $240\text{ m}$ là
\choice
{$180\text{ m}$}
{\True $200\text{ m}$}
{$220\text{ m}$}
{$160\text{ m}$}
\loigiai{
$240 \times \dfrac{5}{6} = 200\text{ m}$.\\
\textbf{Đáp án của câu hỏi là B}
}
\end{ex}

\begin{ex} %[4D12V65]
Một cửa hàng có $150\text{ kg}$ gạo, đã bán được $\dfrac{2}{5}$ số gạo đó. Số gạo cửa hàng còn lại là
\choice
{\True $90\text{ kg}$}
{$60\text{ kg}$}
{$100\text{ kg}$}
{$80\text{ kg}$}
\loigiai{
Đã bán: $150 \times \dfrac{2}{5} = 60\text{ kg}$.\\
Còn lại: $150 - 60 = 90\text{ kg}$.\\
\textbf{Đáp án của câu hỏi là A}
}
\end{ex}

\begin{ex} %[4D12V65]
Một đàn gà có 48 con. Số gà trống bằng $\dfrac{3}{8}$ tổng số gà. Số gà mái là
\choice
{$18$ con}
{\True $30$ con}
{$24$ con}
{$28$ con}
\loigiai{
Gà trống: $48 \times \dfrac{3}{8} = 18$ con.\\
Gà mái: $48 - 18 = 30$ con.\\
\textbf{Đáp án của câu hỏi là B}
}
\end{ex}

\begin{ex} %[4D12V65]
Một mảnh đất hình chữ nhật có chiều dài $40\text{ m}$, chiều rộng bằng $\dfrac{3}{4}$ chiều dài. Diện tích mảnh đất đó là
\choice
{$1000\text{ m}^2$}
{$1400\text{ m}^2$}
{\True $1200\text{ m}^2$}
{$1600\text{ m}^2$}
\loigiai{
Chiều rộng: $40 \times \dfrac{3}{4} = 30\text{ m}$.\\
Diện tích: $40 \times 30 = 1200\text{ m}^2$.\\
\textbf{Đáp án của câu hỏi là C}
}
\end{ex}

\begin{ex} %[4D12V65]
Bạn Nam có $120\,000$ đồng. Nam mua sách hết $\dfrac{2}{3}$ số tiền đó. Số tiền còn lại của Nam là
\choice
{$30\,000$ đồng}
{$50\,000$ đồng}
{$80\,000$ đồng}
{\True $40\,000$ đồng}
\loigiai{
Tiền mua sách: $120\,000 \times \dfrac{2}{3} = 80\,000$ đồng.\\
Còn lại: $120\,000 - 80\,000 = 40\,000$ đồng.\\
\textbf{Đáp án của câu hỏi là D}
}
\end{ex}

\begin{ex} %[4D12V65]
Một quãng đường dài $120\text{ km}$. Ô tô đã đi được $\dfrac{5}{8}$ quãng đường. Quãng đường ô tô còn phải đi tiếp là
\choice
{$40\text{ km}$}
{$50\text{ km}$}
{\True $45\text{ km}$}
{$75\text{ km}$}
\loigiai{
Đã đi: $120 \times \dfrac{5}{8} = 75\text{ km}$. Còn lại: $120 - 75 = 45\text{ km}$.\\
\textbf{Đáp án của câu hỏi là C}
}
\end{ex}

\begin{ex} %[4D12C65]
Một rổ cam có 60 quả. Lần thứ nhất bán $\dfrac{1}{3}$ số cam, lần thứ hai bán $\dfrac{1}{2}$ số cam CÒN LẠI. Số cam còn lại trong rổ sau 2 lần bán là
\choice
{$15$ quả}
{$25$ quả}
{$10$ quả}
{\True $20$ quả}
\loigiai{
Lần 1 bán: $60 \times \dfrac{1}{3} = 20$ quả $\Rightarrow$ còn 40 quả.\\
Lần 2 bán: $40 \times \dfrac{1}{2} = 20$ quả.\\
Sau 2 lần còn: $40 - 20 = 20$ quả.\\
\textbf{Đáp án của câu hỏi là D}
}
\end{ex}

\textbf{Phần 2. Trắc nghiệm đúng sai}
\setcounterpageref{ex}{0}

\begin{ex} %[4D12H65]
Xét tính đúng sai về tìm phân số của một số:
\choiceTF
{\True $\dfrac{3}{4}$ của 24 là 18.}
{\True $\dfrac{2}{5}$ của $1\text{ tấn}$ là $400\text{ kg}$.}
{$\dfrac{1}{4}$ thế kỉ là $20$ năm.}
{\True $\dfrac{1}{4}$ thế kỉ là $25$ năm.}
\loigiai{
\begin{itemize}
    \item $24 \times 3/4 = 18$. Đúng.
    \item $1000 \times 2/5 = 400\text{ kg}$. Đúng.
    \item $100 : 4 = 25$ năm $\neq 20$ năm. Ý 3 Sai, Ý 4 Đúng.
\end{itemize}
\begin{itemize}
    \item a) \textbf{Đúng}.
    \item b) \textbf{Đúng}.
    \item c) \textbf{Sai}.
    \item d) \textbf{Đúng}.
\end{itemize}
}
\end{ex}

\begin{ex} %[4D12V65]
Một khối lớp 4 có 160 học sinh, trong đó có $\dfrac{3}{5}$ số học sinh đạt học lực Xuất sắc, $\dfrac{1}{4}$ đạt học lực Tốt, còn lại là học lực Hoàn thành. Xét tính đúng sai:
\choiceTF
{\True Số học sinh Xuất sắc là 96 em.}
{\True Số học sinh Tốt là 40 em.}
{Số học sinh Hoàn thành là 30 em.}
{\True Số học sinh Hoàn thành là 24 em.}
\loigiai{
\begin{itemize}
    \item Xuất sắc: $160 \times 3/5 = 96$ em. Đúng.
    \item Tốt: $160 \times 1/4 = 40$ em. Đúng.
    \item Hoàn thành: $160 - (96 + 40) = 24$ em $\neq 30$ em. Ý 3 Sai, Ý 4 Đúng.
\end{itemize}
\begin{itemize}
    \item a) \textbf{Đúng}.
    \item b) \textbf{Đúng}.
    \item c) \textbf{Sai}.
    \item d) \textbf{Đúng}.
\end{itemize}
}
\end{ex}

\textbf{Phần 3. Trả lời ngắn}
\setcounterpageref{ex}{0}

\begin{ex} %[4D12H65]
Tìm $\dfrac{3}{5}$ của $45$.
\shortans{27}
\loigiai{
\textbf{Đáp số:} $27$\\
$45 \times 3/5 = 27$.
}
\end{ex}

\begin{ex} %[4D12H65]
Tìm $\dfrac{5}{8}$ của $64\text{ kg}$.
\shortans{40}
\loigiai{
\textbf{Đáp số:} $40$\\
$64 \times 5/8 = 40\text{ kg}$.
}
\end{ex}

\begin{ex} %[4D12V65]
Tìm $\dfrac{2}{3}$ của $1$ giờ theo phút.
\shortans{40}
\loigiai{
\textbf{Đáp số:} $40$\\
$60 \times 2/3 = 40$ phút.
}
\end{ex}

\begin{ex} %[4D12V65]
Một mảnh đất hình chữ nhật dài $50\text{ m}$, chiều rộng bằng $\dfrac{3}{5}$ chiều dài. Tính diện tích mảnh đất theo mét vuông.
\shortans{1500}
\loigiai{
\textbf{Đáp số:} $1500$\\
Rộng: $30\text{ m}$. Diện tích: $50 \times 30 = 1500\text{ m}^2$.
}
\end{ex}

\begin{ex} %[4D12V65]
Một rổ có 40 quả táo, mẹ đã biếu bà $\dfrac{3}{8}$ số táo đó. Hỏi trong rổ còn lại bao nhiêu quả táo?
\shortans{25}
\loigiai{
\textbf{Đáp số:} $25$\\
Biếu bà: $40 \times 3/8 = 15$ quả. Còn lại: $40 - 15 = 25$ quả.
}
\end{ex}

\begin{ex} %[4D12C65]
Một bình nước có 24 lít. Người ta lấy ra lần 1 là $\dfrac{1}{4}$ lượng nước trong bình, lần 2 lấy tiếp $\dfrac{1}{3}$ lượng nước ban đầu. Tính lượng nước còn lại trong bình theo lít.
\shortans{10}
\loigiai{
\textbf{Đáp số:} $10$\\
Lấy ra: $24 \times (1/4 + 1/3) = 24 \times 7/12 = 14$ lít. Còn lại: $24 - 14 = 10$ lít.
}
\end{ex}

\newpage

\TaoPhieuKiemTra{Chủ đề 12 -- Phép nhân, phép chia phân số}{66}{Luyện tập chung nhân chia phân số}{https://loptoanthaylinh.vn/thi-online/toan-4/bai-66}
\textbf{Phần 1. Trắc nghiệm lựa chọn}
\setcounterpageref{ex}{0}

\begin{ex} %[4D12N66]
Kết quả của phép nhân $\dfrac{3}{4} \times \dfrac{2}{5}$ là
\choice
{$\dfrac{5}{9}$}
{\True $\dfrac{3}{10}$}
{$\dfrac{6}{9}$}
{$\dfrac{5}{20}$}
\loigiai{
$\dfrac{6}{20} = \dfrac{3}{10}$.\\
\textbf{Đáp án của câu hỏi là B}
}
\end{ex}

\begin{ex} %[4D12N66]
Kết quả của phép chia $\dfrac{4}{7} : \dfrac{2}{7}$ là
\choice
{$\dfrac{8}{49}$}
{$\dfrac{1}{2}$}
{\True $2$}
{$\dfrac{6}{7}$}
\loigiai{
$\dfrac{4}{7} \times \dfrac{7}{2} = 2$.\\
\textbf{Đáp án của câu hỏi là C}
}
\end{ex}

\begin{ex} %[4D12H66]
$\dfrac{3}{5}$ của $40\text{ kg}$ là
\choice
{\True $24\text{ kg}$}
{$20\text{ kg}$}
{$28\text{ kg}$}
{$30\text{ kg}$}
\loigiai{
$40 \times \dfrac{3}{5} = 24\text{ kg}$.\\
\textbf{Đáp án của câu hỏi là A}
}
\end{ex}

\begin{ex} %[4D12H66]
Kết quả của phép tính $3 \times \dfrac{2}{9}$ là
\choice
{$\dfrac{6}{27}$}
{$\dfrac{5}{9}$}
{\True $\dfrac{2}{3}$}
{$2$}
\loigiai{
$\dfrac{6}{9} = \dfrac{2}{3}$.\\
\textbf{Đáp án của câu hỏi là C}
}
\end{ex}

\begin{ex} %[4D12H66]
Tìm $x$ biết: $x : \dfrac{2}{3} = \dfrac{3}{4}$.
\choice
{$\dfrac{9}{8}$}
{$\dfrac{5}{7}$}
{$\dfrac{6}{12}$}
{\True $\dfrac{1}{2}$}
\loigiai{
$x = \dfrac{3}{4} \times \dfrac{2}{3} = \dfrac{6}{12} = \dfrac{1}{2}$.\\
\textbf{Đáp án của câu hỏi là D}
}
\end{ex}

\begin{ex} %[4D12H66]
Tính thuận tiện: $\dfrac{4}{7} \times \dfrac{5}{9} + \dfrac{4}{7} \times \dfrac{4}{9}$.
\choice
{$1$}
{$\dfrac{20}{63}$}
{\True $\dfrac{4}{7}$}
{$\dfrac{8}{7}$}
\loigiai{
$\dfrac{4}{7} \times \left(\dfrac{5}{9} + \dfrac{4}{9}\right) = \dfrac{4}{7} \times 1 = \dfrac{4}{7}$.\\
\textbf{Đáp án của câu hỏi là C}
}
\end{ex}

\begin{ex} %[4D12V66]
Một hình bình hành có đáy dài $\dfrac{3}{2}\text{ m}$, chiều cao $\dfrac{4}{5}\text{ m}$. Diện tích hình bình hành là
\choice
{$\dfrac{7}{7}\text{ m}^2$}
{\True $\dfrac{6}{5}\text{ m}^2$}
{$\dfrac{12}{5}\text{ m}^2$}
{$\dfrac{3}{5}\text{ m}^2$}
\loigiai{
$S = \dfrac{3}{2} \times \dfrac{4}{5} = \dfrac{12}{10} = \dfrac{6}{5}\text{ m}^2$.\\
\textbf{Đáp án của câu hỏi là B}
}
\end{ex}

\begin{ex} %[4D12V66]
Một cửa hàng có $360\text{ kg}$ gạo, buổi sáng bán được $\dfrac{1}{3}$ số gạo, buổi chiều bán được $\dfrac{2}{5}$ số gạo. Cả ngày cửa hàng bán được số ki-lô-gam gạo là
\choice
{\True $264\text{ kg}$}
{$240\text{ kg}$}
{$270\text{ kg}$}
{$280\text{ kg}$}
\loigiai{
$360 \times \left(\dfrac{1}{3} + \dfrac{2}{5}\right) = 360 \times \dfrac{11}{15} = 264\text{ kg}$.\\
\textbf{Đáp án của câu hỏi là A}
}
\end{ex}

\begin{ex} %[4D12V66]
Một tấm bìa hình chữ nhật có diện tích $\dfrac{4}{5}\text{ m}^2$, chiều dài $\dfrac{8}{5}\text{ m}$. Chiều rộng tấm bìa là
\choice
{$\dfrac{32}{25}\text{ m}$}
{$2\text{ m}$}
{$\dfrac{2}{5}\text{ m}$}
{\True $\dfrac{1}{2}\text{ m}$}
\loigiai{
$\dfrac{4}{5} : \dfrac{8}{5} = \dfrac{4}{5} \times \dfrac{5}{8} = \dfrac{4}{8} = \dfrac{1}{2}\text{ m}$.\\
\textbf{Đáp án của câu hỏi là D}
}
\end{ex}

\begin{ex} %[4D12V66]
Tìm $x$ biết: $x \times \dfrac{3}{5} = \dfrac{9}{10} - \dfrac{3}{10}$.
\choice
{$\dfrac{3}{5}$}
{\True $1$}
{$\dfrac{2}{5}$}
{$\dfrac{6}{10}$}
\loigiai{
$x \times \dfrac{3}{5} = \dfrac{6}{10} = \dfrac{3}{5} \Rightarrow x = 1$.\\
\textbf{Đáp án của câu hỏi là B}
}
\end{ex}

\begin{ex} %[4D12V66]
Một bình đựng 2 lít nước, người ta đã uống $\dfrac{3}{4}$ lượng nước đó. Lượng nước còn lại trong bình là
\choice
{\True $\dfrac{1}{2}\text{ l}$}
{$1\text{ l}$}
{$\dfrac{3}{2}\text{ l}$}
{$\dfrac{1}{4}\text{ l}$}
\loigiai{
Uống: $2 \times \dfrac{3}{4} = \dfrac{3}{2}$ lít. Còn lại: $2 - \dfrac{3}{2} = \dfrac{1}{2}$ lít.\\
\textbf{Đáp án của câu hỏi là A}
}
\end{ex}

\begin{ex} %[4D12C66]
Tính tích: $M = \left(1 + \dfrac{1}{2}\right) \times \left(1 + \dfrac{1}{3}\right) \times \left(1 + \dfrac{1}{4}\right) \times \left(1 + \dfrac{1}{5}\right)$.
\choice
{$2$}
{$4$}
{$\dfrac{5}{2}$}
{\True $3$}
\loigiai{
$M = \dfrac{3}{2} \times \dfrac{4}{3} \times \dfrac{5}{4} \times \dfrac{6}{5} = \dfrac{6}{2} = 3$.\\
\textbf{Đáp án của câu hỏi là D}
}
\end{ex}

\textbf{Phần 2. Trắc nghiệm đúng sai}
\setcounterpageref{ex}{0}

\begin{ex} %[4D12H66]
Xét tính đúng sai của các phép toán phân số:
\choiceTF
{\True $\dfrac{2}{5} \times \dfrac{5}{2} = 1$.}
{\True $\dfrac{3}{4} : \dfrac{1}{2} = \dfrac{3}{2}$.}
{$\dfrac{3}{4} \times 0 = \dfrac{3}{4}$.}
{\True $\dfrac{2}{3}$ của $30$ là $20$.}
\loigiai{
\begin{itemize}
    \item Hai phân số đảo ngược nhân nhau bằng 1: Đúng.
    \item $3/4 \times 2/1 = 3/2$. Đúng.
    \item Nhân với 0 luôn bằng 0. Sai.
    \item $30 \times 2/3 = 20$. Đúng.
\end{itemize}
\begin{itemize}
    \item a) \textbf{Đúng}.
    \item b) \textbf{Đúng}.
    \item c) \textbf{Sai}.
    \item d) \textbf{Đúng}.
\end{itemize}
}
\end{ex}

\begin{ex} %[4D12V66]
Một hình chữ nhật có chiều dài $\dfrac{5}{4}\text{ m}$, chiều rộng $\dfrac{4}{5}\text{ m}$. Xét tính đúng sai:
\choiceTF
{\True Diện tích hình chữ nhật bằng $1\text{ m}^2$.}
{\True Chu vi hình chữ nhật bằng $\dfrac{41}{10}\text{ m}$.}
{Chiều dài ngắn hơn chiều rộng.}
{\True Chiều dài dài hơn chiều rộng $\dfrac{9}{20}\text{ m}$.}
\loigiai{
\begin{itemize}
    \item $S = 5/4 \times 4/5 = 1\text{ m}^2$. Đúng.
    \item Chu vi: $(5/4 + 4/5) \times 2 = 41/20 \times 2 = 41/10\text{ m}$. Đúng.
    \item Dài $5/4 > 1 > 4/5$ rộng. Sai.
    \item Hiệu: $5/4 - 4/5 = 9/20\text{ m}$. Đúng.
\end{itemize}
\begin{itemize}
    \item a) \textbf{Đúng}.
    \item b) \textbf{Đúng}.
    \item c) \textbf{Sai}.
    \item d) \textbf{Đúng}.
\end{itemize}
}
\end{ex}

\textbf{Phần 3. Trả lời ngắn}
\setcounterpageref{ex}{0}

\begin{ex} %[4D12H66]
Tính tích: $\dfrac{4}{7} \times \dfrac{21}{16}$ (dạng a/b).
\shortans{3/4}
\loigiai{
\textbf{Đáp số:} $3/4$\\
$3/4$.
}
\end{ex}

\begin{ex} %[4D12H66]
Tính thương: $\dfrac{5}{8} : \dfrac{15}{16}$ (dạng a/b).
\shortans{2/3}
\loigiai{
\textbf{Đáp số:} $2/3$\\
$5/8 \times 16/15 = 2/3$.
}
\end{ex}

\begin{ex} %[4D12V66]
Tìm $\dfrac{4}{5}$ của $100$.
\shortans{80}
\loigiai{
\textbf{Đáp số:} $80$\\
$100 \times 4/5 = 80$.
}
\end{ex}

\begin{ex} %[4D12V66]
Tìm $x$ biết: $x : \dfrac{4}{5} = 15$.
\shortans{12}
\loigiai{
\textbf{Đáp số:} $12$\\
$x = 15 \times 4/5 = 12$.
}
\end{ex}

\begin{ex} %[4D12V66]
Một hình vuông có chu vi $\dfrac{8}{3}\text{ m}$. Tính diện tích hình vuông đó theo mét vuông (dạng a/b).
\shortans{4/9}
\loigiai{
\textbf{Đáp số:} $4/9$\\
Cạnh: $8/3 : 4 = 2/3\text{ m}$. Diện tích: $2/3 \times 2/3 = 4/9\text{ m}^2$.
}
\end{ex}

\begin{ex} %[4D12C66]
Tính giá trị biểu thức: $A = \dfrac{2}{3} \times \dfrac{4}{5} : \dfrac{8}{15}$.
\shortans{1}
\loigiai{
\textbf{Đáp số:} $1$\\
$8/15 : 8/15 = 1$.
}
\end{ex}

\newpage

% =========================================================================
% CHỦ ĐỀ 13 - ÔN TẬP CUỐI NĂM
% =========================================================================
\section{Ôn tập cuối năm}

\TaoPhieuKiemTra{Chủ đề 13 -- Ôn tập cuối năm}{67}{Ôn tập số tự nhiên}{https://loptoanthaylinh.vn/thi-online/toan-4/bai-67}
\textbf{Phần 1. Trắc nghiệm lựa chọn}
\setcounterpageref{ex}{0}

\begin{ex} %[4D13N67]
Số \textit{Một trăm hai mươi lăm triệu} viết là
\choice
{$12\,500\,000$}
{\True $125\,000\,000$}
{$1\,250\,000$}
{$125\,000$}
\loigiai{
Số đó là $125\,000\,000$.\\
\textbf{Đáp án của câu hỏi là B}
}
\end{ex}

\begin{ex} %[4D13N67]
Trong số $45\,678\,900$, chữ số 5 ở hàng nào?
\choice
{Hàng chục triệu}
{Hàng trăm nghìn}
{\True Hàng triệu}
{Hàng chục nghìn}
\loigiai{
Chữ số 5 ở hàng triệu.\\
\textbf{Đáp án của câu hỏi là C}
}
\end{ex}

\begin{ex} %[4D13H67]
Làm tròn số $3\,480\,000$ đến hàng trăm nghìn ta được
\choice
{\True $3\,500\,000$}
{$3\,400\,000$}
{$3\,000\,000$}
{$4\,000\,000$}
\loigiai{
Hàng chục nghìn là 8 nên làm tròn lên $3\,500\,000$.\\
\textbf{Đáp án của câu hỏi là A}
}
\end{ex}

\begin{ex} %[4D13H67]
Số liền sau của số $99\,999\,999$ là
\choice
{$99\,999\,998$}
{$100\,000\,001$}
{$10\,000\,000$}
{\True $100\,000\,000$}
\loigiai{
$99\,999\,999 + 1 = 100\,000\,000$.\\
\textbf{Đáp án của câu hỏi là D}
}
\end{ex}

\begin{ex} %[4D13H67]
Số lớn nhất có 8 chữ số khác nhau là
\choice
{$99\,999\,999$}
{$98\,765\,430$}
{\True $98\,765\,432$}
{$87\,654\,321$}
\loigiai{
Số lớn nhất có 8 chữ số khác nhau là $98\,765\,432$.\\
\textbf{Đáp án của câu hỏi là C}
}
\end{ex}

\begin{ex} %[4D13H67]
Số nào sau đây chia hết cho cả 2 và 5?
\choice
{$1235$}
{$2468$}
{$3579$}
{\True $4560$}
\loigiai{
Số tận cùng bằng 0 chia hết cho cả 2 và 5.\\
\textbf{Đáp án của câu hỏi là D}
}
\end{ex}

\begin{ex} %[4D13V67]
Sắp xếp các số sau theo thứ tự từ lớn đến bé: $45\,020\,000; 45\,200\,000; 45\,002\,000; 45\,220\,000$.
\choice
{\True $45\,220\,000 > 45\,200\,000 > 45\,020\,000 > 45\,002\,000$}
{$45\,002\,000 > 45\,020\,000 > 45\,200\,000 > 45\,220\,000$}
{$45\,200\,000 > 45\,220\,000 > 45\,020\,000 > 45\,002\,000$}
{$45\,220\,000 > 45\,020\,000 > 45\,200\,000 > 45\,002\,000$}
\loigiai{
Thứ tự: $45\,220\,000 > 45\,200\,000 > 45\,020\,000 > 45\,002\,000$.\\
\textbf{Đáp án của câu hỏi là A}
}
\end{ex}

\begin{ex} %[4D13V67]
Tìm chữ số $x$ để số $\overline{35x,420} < 352\,420$:
\choice
{$x = 2$}
{\True $x \in \{0; 1\}$}
{$x = 3$}
{$x = 4$}
\loigiai{
$x < 2 \Rightarrow x \in \{0; 1\}$.\\
\textbf{Đáp án của câu hỏi là B}
}
\end{ex}

\begin{ex} %[4D13V67]
Dãy số tự nhiên chẵn liên tiếp bắt đầu từ 2 đến 100 có bao nhiêu số hạng?
\choice
{$49$ số}
{$51$ số}
{$100$ số}
{\True $50$ số}
\loigiai{
$(100 - 2) : 2 + 1 = 50$ số.\\
\textbf{Đáp án của câu hỏi là D}
}
\end{ex}

\begin{ex} %[4D13V67]
Tính tổng của dãy số: $1 + 2 + 3 + \dots + 30$.
\choice
{$450$}
{$480$}
{\True $465$}
{$435$}
\loigiai{
$(1 + 30) \times 30 : 2 = 465$.\\
\textbf{Đáp án của câu hỏi là C}
}
\end{ex}

\begin{ex} %[4D13V67]
Trong số $789\,123\,456$, tổng các chữ số thuộc lớp triệu là
\choice
{$18$}
{\True $24$}
{$15$}
{$21$}
\loigiai{
Lớp triệu gồm $7, 8, 9$. Tổng: $7 + 8 + 9 = 24$.\\
\textbf{Đáp án của câu hỏi là B}
}
\end{ex}

\begin{ex} %[4D13C67]
Có bao nhiêu số có 6 chữ số khác nhau được viết từ các chữ số $0, 1, 2, 3, 4, 5$?
\choice
{\True $600$ số}
{$720$ số}
{$120$ số}
{$500$ số}
\loigiai{
Chữ số đầu tiên có 5 cách chọn (khác 0). Các vị trí tiếp theo có $5 \times 4 \times 3 \times 2 \times 1 = 120$ cách. Tổng cộng: $5 \times 120 = 600$ số.\\
\textbf{Đáp án của câu hỏi là A}
}
\end{ex}

\textbf{Phần 2. Trắc nghiệm đúng sai}
\setcounterpageref{ex}{0}

\begin{ex} %[4D13H67]
Xét tính đúng sai về số tự nhiên:
\choiceTF
{\True Số 0 là số tự nhiên bé nhất.}
{\True Không có số tự nhiên lớn nhất.}
{Hai số tự nhiên liên tiếp hơn kém nhau 2 đơn vị.}
{\True Số liền trước của 1 triệu là $999\,999$.}
\loigiai{
\begin{itemize}
    \item 0 là số tự nhiên bé nhất: Đúng.
    \item Không có số lớn nhất: Đúng.
    \item Hai số tự nhiên liên tiếp hơn kém nhau 1 đơn vị (không phải 2). Sai.
    \item $1\,000\,000 - 1 = 999\,999$: Đúng.
\end{itemize}
\begin{itemize}
    \item a) \textbf{Đúng}.
    \item b) \textbf{Đúng}.
    \item c) \textbf{Sai}.
    \item d) \textbf{Đúng}.
\end{itemize}
}
\end{ex}

\begin{ex} %[4D13V67]
Cho số $A = 50\,000\,000 + 700\,000 + 4000 + 8$. Xét tính đúng sai:
\choiceTF
{\True Số $A$ được viết là $50\,704\,008$.}
{\True Lớp triệu của số $A$ là $50$.}
{Chữ số 7 có giá trị là $70\,000$.}
{\True Chữ số 7 có giá trị là $700\,000$.}
\loigiai{
\begin{itemize}
    \item Viết số đúng: $50\,704\,008$.
    \item Lớp triệu: 50. Đúng.
    \item Chữ số 7 ở hàng trăm nghìn nên có giá trị $700\,000 \neq 70\,000$. Ý 3 Sai, Ý 4 Đúng.
\end{itemize}
\begin{itemize}
    \item a) \textbf{Đúng}.
    \item b) \textbf{Đúng}.
    \item c) \textbf{Sai}.
    \item d) \textbf{Đúng}.
\end{itemize}
}
\end{ex}

\textbf{Phần 3. Trả lời ngắn}
\setcounterpageref{ex}{0}

\begin{ex} %[4D13H67]
Viết số: Tám mươi tư triệu ba trăm nghìn.
\shortans{84300000}
\loigiai{
\textbf{Đáp số:} $84300000$\\
$84\,300\,000$.
}
\end{ex}

\begin{ex} %[4D13H67]
Làm tròn số $7\,650\,000$ đến hàng trăm nghìn.
\shortans{7700000}
\loigiai{
\textbf{Đáp số:} $7700000$\\
$7\,700\,000$.
}
\end{ex}

\begin{ex} %[4D13V67]
Tìm số chẵn lớn nhất có 7 chữ số khác nhau.
\shortans{9876542}
\loigiai{
\textbf{Đáp số:} $9876542$\\
$9\,876\,542$.
}
\end{ex}

\begin{ex} %[4D13V67]
Trong số $456\,789\,123$, tính tổng các chữ số thuộc lớp nghìn.
\shortans{24}
\loigiai{
\textbf{Đáp số:} $24$\\
$7 + 8 + 9 = 24$.
}
\end{ex}

\begin{ex} %[4D13V67]
Dãy số $5, 10, 15, \dots, 100$ có bao nhiêu số hạng?
\shortans{20}
\loigiai{
\textbf{Đáp số:} $20$\\
$(100 - 5) : 5 + 1 = 20$.
}
\end{ex}

\begin{ex} %[4D13C67]
Tính tổng: $1 + 2 + 3 + \dots + 100$.
\shortans{5050}
\loigiai{
\textbf{Đáp số:} $5050$\\
$(1 + 100) \times 100 : 2 = 5050$.
}
\end{ex}

\newpage

\TaoPhieuKiemTra{Chủ đề 13 -- Ôn tập cuối năm}{68}{Ôn tập các phép tính với số tự nhiên}{https://loptoanthaylinh.vn/thi-online/toan-4/bai-68}
\textbf{Phần 1. Trắc nghiệm lựa chọn}
\setcounterpageref{ex}{0}

\begin{ex} %[4D13N68]
Kết quả của phép tính $350\,000 + 450\,000$ là
\choice
{$700\,000$}
{\True $800\,000$}
{$900\,000$}
{$850\,000$}
\loigiai{
$350\,000 + 450\,000 = 800\,000$.\\
\textbf{Đáp án của câu hỏi là B}
}
\end{ex}

\begin{ex} %[4D13N68]
Kết quả của phép nhân $125 \times 8$ là
\choice
{\True $1000$}
{$100$}
{$10\,000$}
{$800$}
\loigiai{
$125 \times 8 = 1000$.\\
\textbf{Đáp án của câu hỏi là A}
}
\end{ex}

\begin{ex} %[4D13H68]
Kết quả của phép chia $480\,000 : 60$ là
\choice
{$800$}
{$80\,000$}
{\True $8000$}
{$80$}
\loigiai{
$480\,000 : 60 = 8000$.\\
\textbf{Đáp án của câu hỏi là C}
}
\end{ex}

\begin{ex} %[4D13H68]
Tính thuận tiện: $25 \times 49 \times 4$.
\choice
{$4800$}
{$5000$}
{$4700$}
{\True $4900$}
\loigiai{
$(25 \times 4) \times 49 = 100 \times 49 = 4900$.\\
\textbf{Đáp án của câu hỏi là D}
}
\end{ex}

\begin{ex} %[4D13H68]
Tính thuận tiện: $45 \times 38 + 45 \times 62$.
\choice
{\True $4500$}
{$4400$}
{$4600$}
{$5000$}
\loigiai{
$45 \times (38 + 62) = 45 \times 100 = 4500$.\\
\textbf{Đáp án của câu hỏi là A}
}
\end{ex}

\begin{ex} %[4D13H68]
Tìm hai số biết tổng bằng $100$ và hiệu bằng $24$. Số lớn là
\choice
{$60$}
{$64$}
{$38$}
{\True $62$}
\loigiai{
$(100 + 24) : 2 = 62$.\\
\textbf{Đáp án của câu hỏi là D}
}
\end{ex}

\begin{ex} %[4D13V68]
Trung bình cộng của ba số: $45, 55, 80$ là
\choice
{\True $60$}
{$55$}
{$65$}
{$70$}
\loigiai{
$(45 + 55 + 80) : 3 = 60$.\\
\textbf{Đáp án của câu hỏi là A}
}
\end{ex}

\begin{ex} %[4D13V68]
Một mảnh vườn hình chữ nhật có chu vi $120\text{ m}$, chiều dài hơn chiều rộng $12\text{ m}$. Diện tích mảnh vườn đó là
\choice
{$850\text{ m}^2$}
{$880\text{ m}^2$}
{$900\text{ m}^2$}
{\True $864\text{ m}^2$}
\loigiai{
Nửa chu vi: $60\text{ m}$. Dài: $36\text{ m}$, Rộng: $24\text{ m}$.\\
Diện tích: $36 \times 24 = 864\text{ m}^2$.\\
\textbf{Đáp án của câu hỏi là D}
}
\end{ex}

\begin{ex} %[4D13V68]
Mua 5 quyển vở hết $45\,000$ đồng. Mua 12 quyển vở như thế hết số tiền là
\choice
{$100\,000$ đồng}
{$110\,000$ đồng}
{\True $108\,000$ đồng}
{$120\,000$ đồng}
\loigiai{
$45\,000 : 5 \times 12 = 108\,000$ đồng.\\
\textbf{Đáp án của câu hỏi là C}
}
\end{ex}

\begin{ex} %[4D13V68]
Tìm $x$ biết: $x \times 25 = 1500$.
\choice
{$50$}
{\True $60$}
{$70$}
{$80$}
\loigiai{
$x = 1500 : 25 = 60$.\\
\textbf{Đáp án của câu hỏi là B}
}
\end{ex}

\begin{ex} %[4D13V68]
Tính giá trị biểu thức: $120\,000 - 45\,000 \times 2 + 15\,000$.
\choice
{$40\,000$}
{$50\,000$}
{\True $45\,000$}
{$60\,000$}
\loigiai{
$120\,000 - 90\,000 + 15\,000 = 45\,000$.\\
\textbf{Đáp án của câu hỏi là C}
}
\end{ex}

\begin{ex} %[4D13C68]
Tổng của 3 số là $360$. Số thứ nhất bằng số thứ hai, số thứ ba gấp đôi số thứ nhất. Số thứ ba là
\choice
{$90$}
{\True $180$}
{$120$}
{$150$}
\loigiai{
Tổng số phần: $1 + 1 + 2 = 4$ phần.\\
Số thứ ba: $360 : 4 \times 2 = 180$.\\
\textbf{Đáp án của câu hỏi là B}
}
\end{ex}

\textbf{Phần 2. Trắc nghiệm đúng sai}
\setcounterpageref{ex}{0}

\begin{ex} %[4D13H68]
Xét tính đúng sai của các phép toán:
\choiceTF
{\True $a \times (b + c) = a \times b + a \times c$.}
{\True $250 \times 4 = 1000$.}
{$1\,000\,000 - 450\,000 = 650\,000$.}
{\True $1\,000\,000 - 450\,000 = 550\,000$.}
\loigiai{
\begin{itemize}
    \item Tính chất phân phối: Đúng.
    \item $250 \times 4 = 1000$. Đúng.
    \item $1\,000\,000 - 450\,000 = 550\,000 \neq 650\,000$. Ý 3 Sai, Ý 4 Đúng.
\end{itemize}
\begin{itemize}
    \item a) \textbf{Đúng}.
    \item b) \textbf{Đúng}.
    \item c) \textbf{Sai}.
    \item d) \textbf{Đúng}.
\end{itemize}
}
\end{ex}

\begin{ex} %[4D13V68]
Hai thùng chứa tất cả $180\text{ l}$ dầu. Thùng 1 chứa nhiều hơn thùng 2 là $30\text{ l}$. Xét tính đúng sai:
\choiceTF
{\True Thùng 1 chứa $105\text{ l}$ dầu.}
{\True Thùng 2 chứa $75\text{ l}$ dầu.}
{Nếu chuyển $15\text{ l}$ từ thùng 1 sang thùng 2 thì hai thùng bằng nhau ($90\text{ l}$ mỗi thùng).}
{\True Trung bình mỗi thùng chứa $90\text{ l}$ dầu.}
\loigiai{
\begin{itemize}
    \item Thùng 1: $(180 + 30) : 2 = 105\text{ l}$. Đúng.
    \item Thùng 2: $180 - 105 = 75\text{ l}$. Đúng.
    \item Chuyển 15 l sang thùng 2: cả 2 có 90 l. Đúng.
    \item TBC: $180 : 2 = 90\text{ l}$. Đúng.
\end{itemize}
\begin{itemize}
    \item a) \textbf{Đúng}.
    \item b) \textbf{Đúng}.
    \item c) \textbf{Đúng}.
    \item d) \textbf{Đúng}.
\end{itemize}
}
\end{ex}

\textbf{Phần 3. Trả lời ngắn}
\setcounterpageref{ex}{0}

\begin{ex} %[4D13H68]
Tính thuận tiện: $25 \times 84 \times 4$.
\shortans{8400}
\loigiai{
\textbf{Đáp số:} $8400$\\
$8400$.
}
\end{ex}

\begin{ex} %[4D13H68]
Tính tổng: $125\,000 + 875\,000$.
\shortans{1000000}
\loigiai{
\textbf{Đáp số:} $1000000$\\
$1\,000\,000$.
}
\end{ex}

\begin{ex} %[4D13V68]
Hai số có tổng là $140$, hiệu là $36$. Tìm số bé.
\shortans{52}
\loigiai{
\textbf{Đáp số:} $52$\\
$(140 - 36) : 2 = 52$.
}
\end{ex}

\begin{ex} %[4D13V68]
Tính nhanh: $67 \times 101 - 67$.
\shortans{6700}
\loigiai{
\textbf{Đáp số:} $6700$\\
$67 \times 100 = 6700$.
}
\end{ex}

\begin{ex} %[4D13V68]
Tìm $x$ biết: $x : 18 = 250$.
\shortans{4500}
\loigiai{
\textbf{Đáp số:} $4500$\\
$x = 250 \times 18 = 4500$.
}
\end{ex}

\begin{ex} %[4D13C68]
Tính tổng: $2 + 4 + 6 + \dots + 40$.
\shortans{420}
\loigiai{
\textbf{Đáp số:} $420$\\
Số số hạng: $(40 - 2) : 2 + 1 = 20$. Tổng: $(2 + 40) \times 20 : 2 = 420$.
}
\end{ex}

\newpage

\TaoPhieuKiemTra{Chủ đề 13 -- Ôn tập cuối năm}{69}{Ôn tập phân số}{https://loptoanthaylinh.vn/thi-online/toan-4/bai-69}
\textbf{Phần 1. Trắc nghiệm lựa chọn}
\setcounterpageref{ex}{0}

\begin{ex} %[4D13N69]
Phân số nào sau đây có tử số bằng 5 và mẫu số bằng 8?
\choice
{\True $\dfrac{5}{8}$}
{$\dfrac{8}{5}$}
{$5{,}8$}
{$8{,}5$}
\loigiai{
Phân số đó là $\dfrac{5}{8}$.\\
\textbf{Đáp án của câu hỏi là A}
}
\end{ex}

\begin{ex} %[4D13N69]
Phân số $\dfrac{12}{16}$ rút gọn về phân số tối giản là
\choice
{$\dfrac{6}{8}$}
{$\dfrac{2}{3}$}
{\True $\dfrac{3}{4}$}
{$\dfrac{4}{5}$}
\loigiai{
$\dfrac{12 : 4}{16 : 4} = \dfrac{3}{4}$.\\
\textbf{Đáp án của câu hỏi là C}
}
\end{ex}

\begin{ex} %[4D13H69]
Phân số nào sau đây bằng phân số $\dfrac{2}{3}$?
\choice
{$\dfrac{4}{9}$}
{\True $\dfrac{8}{12}$}
{$\dfrac{6}{10}$}
{$\dfrac{3}{2}$}
\loigiai{
$\dfrac{2 \times 4}{3 \times 4} = \dfrac{8}{12}$.\\
\textbf{Đáp án của câu hỏi là B}
}
\end{ex}

\begin{ex} %[4D13H69]
Trong các phân số sau, phân số nào lớn hơn 1?
\choice
{$\dfrac{3}{4}$}
{$\dfrac{5}{5}$}
{$\dfrac{7}{9}$}
{\True $\dfrac{9}{7}$}
\loigiai{
$\dfrac{9}{7} > 1$ (tử số lớn hơn mẫu số).\\
\textbf{Đáp án của câu hỏi là D}
}
\end{ex}

\begin{ex} %[4D13H69]
Điền dấu thích hợp vào chỗ chấm: $\dfrac{4}{7} \dots \dfrac{6}{7}$.
\choice
{\True $<$}
{$>$}
{$=$}
{Không so sánh được}
\loigiai{
$4 < 6 \Rightarrow \dfrac{4}{7} < \dfrac{6}{7}$.\\
\textbf{Đáp án của câu hỏi là A}
}
\end{ex}

\begin{ex} %[4D13H69]
Mẫu số chung nhỏ nhất của hai phân số $\dfrac{3}{4}$ và $\dfrac{5}{6}$ là
\choice
{$24$}
{$18$}
{\True $12$}
{$36$}
\loigiai{
12 chia hết cho cả 4 và 6.\\
\textbf{Đáp án của câu hỏi là C}
}
\end{ex}

\begin{ex} %[4D13V69]
Sắp xếp các phân số $\dfrac{1}{2}, \dfrac{2}{3}, \dfrac{3}{4}$ theo thứ tự từ bé đến lớn:
\choice
{$\dfrac{3}{4} < \dfrac{2}{3} < \dfrac{1}{2}$}
{\True $\dfrac{1}{2} < \dfrac{2}{3} < \dfrac{3}{4}$}
{$\dfrac{2}{3} < \dfrac{1}{2} < \dfrac{3}{4}$}
{$\dfrac{1}{2} < \dfrac{3}{4} < \dfrac{2}{3}$}
\loigiai{
$\dfrac{6}{12} < \dfrac{8}{12} < \dfrac{9}{12} \Rightarrow \dfrac{1}{2} < \dfrac{2}{3} < \dfrac{3}{4}$.\\
\textbf{Đáp án của câu hỏi là B}
}
\end{ex}

\begin{ex} %[4D13V69]
Tìm $x$ biết: $\dfrac{x}{18} = \dfrac{2}{3}$.
\choice
{$10$}
{$14$}
{$16$}
{\True $12$}
\loigiai{
$x = 2 \times 6 = 12$.\\
\textbf{Đáp án của câu hỏi là D}
}
\end{ex}

\begin{ex} %[4D13V69]
Một lớp có 30 học sinh, trong đó có 18 bạn nữ. Phân số chỉ số bạn nam so với cả lớp rút gọn tối giản là
\choice
{\True $\dfrac{2}{5}$}
{$\dfrac{3}{5}$}
{$\dfrac{12}{30}$}
{$\dfrac{2}{3}$}
\loigiai{
Nam: $30 - 18 = 12$ bạn $\Rightarrow \dfrac{12}{30} = \dfrac{2}{5}$.\\
\textbf{Đáp án của câu hỏi là A}
}
\end{ex}

\begin{ex} %[4D13V69]
Phân số nào sau đây tối giản?
\choice
{$\dfrac{15}{25}$}
{\True $\dfrac{9}{14}$}
{$\dfrac{14}{21}$}
{$\dfrac{8}{12}$}
\loigiai{
$\dfrac{9}{14}$ là phân số tối giản.\\
\textbf{Đáp án của câu hỏi là B}
}
\end{ex}

\begin{ex} %[4D13V69]
Viết phân số $\dfrac{4}{5}$ thành phân số có mẫu số bằng 100:
\choice
{$\dfrac{40}{100}$}
{$\dfrac{85}{100}$}
{\True $\dfrac{80}{100}$}
{$\dfrac{75}{100}$}
\loigiai{
$\dfrac{4 \times 20}{5 \times 20} = \dfrac{80}{100}$.\\
\textbf{Đáp án của câu hỏi là C}
}
\end{ex}

\begin{ex} %[4D13C69]
Tìm phân số tối giản $\dfrac{a}{b}$ bằng $\dfrac{4}{7}$ biết mẫu số hơn tử số 15 đơn vị.
\choice
{$\dfrac{12}{27}$}
{$\dfrac{16}{31}$}
{$\dfrac{24}{39}$}
{\True $\dfrac{20}{35}$}
\loigiai{
Hiệu số phần: $7 - 4 = 3$ phần $= 15 \Rightarrow 1$ phần $= 5$.\\
Tử: $4 \times 5 = 20$, Mẫu: $7 \times 5 = 35 \Rightarrow \dfrac{20}{35}$.\\
\textbf{Đáp án của câu hỏi là D}
}
\end{ex}

\textbf{Phần 2. Trắc nghiệm đúng sai}
\setcounterpageref{ex}{0}

\begin{ex} %[4D13H69]
Xét tính đúng sai về phân số:
\choiceTF
{\True $\dfrac{3}{5} = \dfrac{6}{10} = \dfrac{9}{15}$.}
{\True Phân số $\dfrac{7}{8}$ bé hơn 1.}
{Phân số $\dfrac{5}{4}$ bé hơn phân số $\dfrac{4}{5}$.}
{\True Phân số $\dfrac{5}{4}$ lớn hơn phân số $\dfrac{4}{5}$.}
\loigiai{
\begin{itemize}
    \item Dãy phân số bằng nhau đúng.
    \item $7/8 < 1$: Đúng.
    \item $5/4 > 1 > 4/5$. Ý 3 Sai, Ý 4 Đúng.
\end{itemize}
\begin{itemize}
    \item a) \textbf{Đúng}.
    \item b) \textbf{Đúng}.
    \item c) \textbf{Sai}.
    \item d) \textbf{Đúng}.
\end{itemize}
}
\end{ex}

\begin{ex} %[4D13V69]
Cho 3 phân số: $\dfrac{3}{4}, \dfrac{5}{6}, \dfrac{7}{12}$. Xét tính đúng sai:
\choiceTF
{\True Mẫu số chung nhỏ nhất của 3 phân số là 12.}
{\True Quy đồng được $\dfrac{9}{12}, \dfrac{10}{12}, \dfrac{7}{12}$.}
{Phân số lớn nhất là $\dfrac{3}{4}$.}
{\True Phân số lớn nhất là $\dfrac{5}{6}$.}
\loigiai{
\begin{itemize}
    \item MSC là 12: Đúng.
    \item $9/12, 10/12, 7/12$: Đúng.
    \item Lớn nhất là $5/6 = 10/12 \neq 3/4$. Ý 3 Sai, Ý 4 Đúng.
\end{itemize}
\begin{itemize}
    \item a) \textbf{Đúng}.
    \item b) \textbf{Đúng}.
    \item c) \textbf{Sai}.
    \item d) \textbf{Đúng}.
\end{itemize}
}
\end{ex}

\textbf{Phần 3. Trả lời ngắn}
\setcounterpageref{ex}{0}

\begin{ex} %[4D13H69]
Rút gọn phân số $\dfrac{25}{35}$ về tối giản (dạng a/b).
\shortans{5/7}
\loigiai{
\textbf{Đáp số:} $5/7$\\
$5/7$.
}
\end{ex}

\begin{ex} %[4D13H69]
Tìm $x$ biết $\dfrac{x}{15} = \dfrac{3}{5}$.
\shortans{9}
\loigiai{
\textbf{Đáp số:} $9$\\
$x = 3 \times 3 = 9$.
}
\end{ex}

\begin{ex} %[4D13V69]
Quy đồng mẫu số hai phân số $\dfrac{2}{3}$ và $\dfrac{3}{5}$ với mẫu số chung 15. Tìm tổng hai tử số mới.
\shortans{19}
\loigiai{
\textbf{Đáp số:} $19$\\
$10 + 9 = 19$.
}
\end{ex}

\begin{ex} %[4D13V69]
Tìm phân số lớn nhất trong các phân số: $\dfrac{3}{5}, \dfrac{4}{5}, \dfrac{2}{5}$ (dạng a/b).
\shortans{4/5}
\loigiai{
\textbf{Đáp số:} $4/5$\\
$4/5$.
}
\end{ex}

\begin{ex} %[4D13V69]
Tìm phân số bé nhất trong các phân số: $\dfrac{5}{6}, \dfrac{5}{7}, \dfrac{5}{8}$ (dạng a/b).
\shortans{5/8}
\loigiai{
\textbf{Đáp số:} $5/8$\\
$5/8$.
}
\end{ex}

\begin{ex} %[4D13C69]
Tìm chữ số $x$ để $\dfrac{1}{2} < \dfrac{x}{8} < \dfrac{3}{4}$.
\shortans{5}
\loigiai{
\textbf{Đáp số:} $5$\\
$4/8 < x/8 < 6/8 \Rightarrow x = 5$.
}
\end{ex}

\newpage

\TaoPhieuKiemTra{Chủ đề 13 -- Ôn tập cuối năm}{70}{Ôn tập các phép tính với phân số}{https://loptoanthaylinh.vn/thi-online/toan-4/bai-70}
\textbf{Phần 1. Trắc nghiệm lựa chọn}
\setcounterpageref{ex}{0}

\begin{ex} %[4D13N70]
Kết quả của phép tính $\dfrac{1}{4} + \dfrac{2}{4}$ là
\choice
{$\dfrac{3}{8}$}
{\True $\dfrac{3}{4}$}
{$\dfrac{2}{4}$}
{$\dfrac{1}{2}$}
\loigiai{
$\dfrac{1 + 2}{4} = \dfrac{3}{4}$.\\
\textbf{Đáp án của câu hỏi là B}
}
\end{ex}

\begin{ex} %[4D13N70]
Kết quả của phép tính $\dfrac{3}{5} \times \dfrac{2}{7}$ là
\choice
{\True $\dfrac{6}{35}$}
{$\dfrac{5}{12}$}
{$\dfrac{6}{12}$}
{$\dfrac{5}{35}$}
\loigiai{
$\dfrac{3 \times 2}{5 \times 7} = \dfrac{6}{35}$.\\
\textbf{Đáp án của câu hỏi là A}
}
\end{ex}

\begin{ex} %[4D13H70]
Kết quả của phép trừ $\dfrac{5}{6} - \dfrac{1}{2}$ là
\choice
{$\dfrac{4}{4}$}
{\True $\dfrac{1}{3}$}
{$\dfrac{2}{3}$}
{$\dfrac{1}{6}$}
\loigiai{
$\dfrac{5}{6} - \dfrac{3}{6} = \dfrac{2}{6} = \dfrac{1}{3}$.\\
\textbf{Đáp án của câu hỏi là B}
}
\end{ex}

\begin{ex} %[4D13H70]
Kết quả của phép chia $\dfrac{4}{9} : \dfrac{2}{3}$ là
\choice
{$\dfrac{8}{27}$}
{$\dfrac{3}{2}$}
{$\dfrac{1}{3}$}
{\True $\dfrac{2}{3}$}
\loigiai{
$\dfrac{4}{9} \times \dfrac{3}{2} = \dfrac{2}{3}$.\\
\textbf{Đáp án của câu hỏi là D}
}
\end{ex}

\begin{ex} %[4D13H70]
$\dfrac{2}{3}$ của $30\text{ kg}$ là
\choice
{$15\text{ kg}$}
{\True $20\text{ kg}$}
{$25\text{ kg}$}
{$10\text{ kg}$}
\loigiai{
$30 \times \dfrac{2}{3} = 20\text{ kg}$.\\
\textbf{Đáp án của câu hỏi là B}
}
\end{ex}

\begin{ex} %[4D13H70]
Tính thuận tiện: $\dfrac{3}{5} \times \dfrac{7}{9} + \dfrac{3}{5} \times \dfrac{2}{9}$.
\choice
{\True $\dfrac{3}{5}$}
{$1$}
{$\dfrac{9}{9}$}
{$\dfrac{6}{5}$}
\loigiai{
$\dfrac{3}{5} \times \left(\dfrac{7}{9} + \dfrac{2}{9}\right) = \dfrac{3}{5} \times 1 = \dfrac{3}{5}$.\\
\textbf{Đáp án của câu hỏi là A}
}
\end{ex}

\begin{ex} %[4D13V70]
Một mảnh đất hình chữ nhật có chiều dài $\dfrac{4}{5}\text{ m}$, chiều rộng $\dfrac{1}{2}\text{ m}$. Diện tích mảnh đất là
\choice
{$\dfrac{5}{7}\text{ m}^2$}
{$\dfrac{4}{10}\text{ m}^2$}
{\True $\dfrac{2}{5}\text{ m}^2$}
{$\dfrac{3}{10}\text{ m}^2$}
\loigiai{
$\dfrac{4}{5} \times \dfrac{1}{2} = \dfrac{4}{10} = \dfrac{2}{5}\text{ m}^2$.\\
\textbf{Đáp án của câu hỏi là C}
}
\end{ex}

\begin{ex} %[4D13V70]
Một người có $60\text{ kg}$ gạo, đã bán $\dfrac{3}{4}$ số gạo đó. Số gạo còn lại là
\choice
{\True $15\text{ kg}$}
{$45\text{ kg}$}
{$20\text{ kg}$}
{$30\text{ kg}$}
\loigiai{
Đã bán: $60 \times \dfrac{3}{4} = 45\text{ kg}$. Còn lại: $60 - 45 = 15\text{ kg}$.\\
\textbf{Đáp án của câu hỏi là A}
}
\end{ex}

\begin{ex} %[4D13V70]
Tìm $x$ biết: $x - \dfrac{1}{3} = \dfrac{2}{5}$.
\choice
{$\dfrac{1}{2}$}
{$\dfrac{3}{8}$}
{$\dfrac{1}{15}$}
{\True $\dfrac{11}{15}$}
\loigiai{
$x = \dfrac{2}{5} + \dfrac{1}{3} = \dfrac{6 + 5}{15} = \dfrac{11}{15}$.\\
\textbf{Đáp án của câu hỏi là D}
}
\end{ex}

\begin{ex} %[4D13V70]
Tìm $x$ biết: $x : \dfrac{3}{4} = \dfrac{8}{9}$.
\choice
{$\dfrac{32}{27}$}
{$\dfrac{5}{12}$}
{\True $\dfrac{2}{3}$}
{$\dfrac{6}{12}$}
\loigiai{
$x = \dfrac{8}{9} \times \dfrac{3}{4} = \dfrac{24}{36} = \dfrac{2}{3}$.\\
\textbf{Đáp án của câu hỏi là C}
}
\end{ex}

\begin{ex} %[4D13V70]
Tính giá trị biểu thức: $\left(\dfrac{1}{2} + \dfrac{1}{3}\right) \times \dfrac{6}{5}$.
\choice
{$\dfrac{5}{6}$}
{$\dfrac{6}{5}$}
{$2$}
{\True $1$}
\loigiai{
$\dfrac{5}{6} \times \dfrac{6}{5} = 1$.\\
\textbf{Đáp án của câu hỏi là D}
}
\end{ex}

\begin{ex} %[4D13C70]
Tính giá trị biểu thức: $A = \dfrac{1}{2} \times \dfrac{2}{3} \times \dfrac{3}{4} \times \dots \times \dfrac{19}{20}$.
\choice
{$\dfrac{1}{10}$}
{$\dfrac{1}{19}$}
{\True $\dfrac{1}{20}$}
{$\dfrac{2}{20}$}
\loigiai{
$A = \dfrac{1}{20}$.\\
\textbf{Đáp án của câu hỏi là C}
}
\end{ex}

\textbf{Phần 2. Trắc nghiệm đúng sai}
\setcounterpageref{ex}{0}

\begin{ex} %[4D13H70]
Xét tính đúng sai của các phép toán phân số:
\choiceTF
{\True $\dfrac{1}{3} + \dfrac{2}{3} = 1$.}
{\True $\dfrac{3}{4} \times \dfrac{4}{3} = 1$.}
{$\dfrac{5}{6} - \dfrac{1}{6} = \dfrac{4}{0}$.}
{\True $\dfrac{5}{6} - \dfrac{1}{6} = \dfrac{4}{6} = \dfrac{2}{3}$.}
\loigiai{
\begin{itemize}
    \item $1/3 + 2/3 = 1$. Đúng.
    \item $3/4 \times 4/3 = 1$. Đúng.
    \item Mẫu số không thể bằng 0. Sai.
    \item $5/6 - 1/6 = 2/3$. Đúng.
\end{itemize}
\begin{itemize}
    \item a) \textbf{Đúng}.
    \item b) \textbf{Đúng}.
    \item c) \textbf{Sai}.
    \item d) \textbf{Đúng}.
\end{itemize}
}
\end{ex}

\begin{ex} %[4D13V70]
Một thùng có $24\text{ l}$ dầu: Lần 1 lấy $\dfrac{1}{3}$ thùng, lần 2 lấy $\dfrac{1}{4}$ thùng. Xét tính đúng sai:
\choiceTF
{\True Lần 1 lấy được $8\text{ l}$ dầu.}
{\True Lần 2 lấy được $6\text{ l}$ dầu.}
{Cả hai lần lấy được tất cả $15\text{ l}$ dầu.}
{\True Trong thùng còn lại $10\text{ l}$ dầu.}
\loigiai{
\begin{itemize}
    \item Lần 1: $24 \times 1/3 = 8\text{ l}$. Đúng.
    \item Lần 2: $24 \times 1/4 = 6\text{ l}$. Đúng.
    \item Tổng lấy: $8 + 6 = 14\text{ l} \neq 15\text{ l}$. Ý 3 Sai, Ý 4 Đúng ($24 - 14 = 10\text{ l}$).
\end{itemize}
\begin{itemize}
    \item a) \textbf{Đúng}.
    \item b) \textbf{Đúng}.
    \item c) \textbf{Sai}.
    \item d) \textbf{Đúng}.
\end{itemize}
}
\end{ex}

\textbf{Phần 3. Trả lời ngắn}
\setcounterpageref{ex}{0}

\begin{ex} %[4D13H70]
Tính tổng: $\dfrac{2}{5} + \dfrac{1}{2}$ (dạng a/b).
\shortans{9/10}
\loigiai{
\textbf{Đáp số:} $9/10$\\
$4/10 + 5/10 = 9/10$.
}
\end{ex}

\begin{ex} %[4D13H70]
Tính hiệu: $\dfrac{7}{8} - \dfrac{3}{4}$ (dạng a/b).
\shortans{1/8}
\loigiai{
\textbf{Đáp số:} $1/8$\\
$7/8 - 6/8 = 1/8$.
}
\end{ex}

\begin{ex} %[4D13V70]
Tính tích: $\dfrac{5}{9} \times \dfrac{3}{10}$ (rút gọn tối giản dạng a/b).
\shortans{1/6}
\loigiai{
\textbf{Đáp số:} $1/6$\\
$15/90 = 1/6$.
}
\end{ex}

\begin{ex} %[4D13V70]
Tính thương: $\dfrac{3}{5} : \dfrac{9}{10}$ (rút gọn tối giản dạng a/b).
\shortans{2/3}
\loigiai{
\textbf{Đáp số:} $2/3$\\
$3/5 \times 10/9 = 2/3$.
}
\end{ex}

\begin{ex} %[4D13V70]
Tìm $\dfrac{3}{7}$ của $42$.
\shortans{18}
\loigiai{
\textbf{Đáp số:} $18$\\
$42 \times 3/7 = 18$.
}
\end{ex}

\begin{ex} %[4D13C70]
Tính nhanh: $\dfrac{3}{4} \times \dfrac{5}{7} + \dfrac{3}{4} \times \dfrac{2}{7}$ (dạng a/b).
\shortans{3/4}
\loigiai{
\textbf{Đáp số:} $3/4$\\
$3/4 \times 1 = 3/4$.
}
\end{ex}

\newpage

\TaoPhieuKiemTra{Chủ đề 13 -- Ôn tập cuối năm}{71}{Ôn tập hình học và đo lường}{https://loptoanthaylinh.vn/thi-online/toan-4/bai-71}
\textbf{Phần 1. Trắc nghiệm lựa chọn}
\setcounterpageref{ex}{0}

\begin{ex} %[4H13N71]
Góc bẹt có số đo bằng
\choice
{\True $180^\circ$}
{$90^\circ$}
{$60^\circ$}
{$120^\circ$}
\loigiai{
Góc bẹt bằng $180^\circ$.\\
\textbf{Đáp án của câu hỏi là A}
}
\end{ex}

\begin{ex} %[4H13N71]
$1$ tấn bằng bao nhiêu ki-lô-gam?
\choice
{$100\text{ kg}$}
{$10\text{ kg}$}
{\True $1000\text{ kg}$}
{$10\,000\text{ kg}$}
\loigiai{
$1$ tấn $= 1000\text{ kg}$.\\
\textbf{Đáp án của câu hỏi là C}
}
\end{ex}

\begin{ex} %[4H13H71]
$1\text{ m}^2$ bằng bao nhiêu đề-xi-mét vuông?
\choice
{$10\text{ dm}^2$}
{\True $100\text{ dm}^2$}
{$1000\text{ dm}^2$}
{$10\,000\text{ dm}^2$}
\loigiai{
$1\text{ m}^2 = 100\text{ dm}^2$.\\
\textbf{Đáp án của câu hỏi là B}
}
\end{ex}

\begin{ex} %[4H13H71]
Hình thoi có
\choice
{4 góc vuông}
{2 đường chéo bằng nhau}
{Chỉ có 1 cặp cạnh song song}
{\True 4 cạnh bằng nhau và 2 đường chéo vuông góc với nhau}
\loigiai{
Hình thoi có 4 cạnh bằng nhau và hai đường chéo vuông góc.\\
\textbf{Đáp án của câu hỏi là D}
}
\end{ex}

\begin{ex} %[4H13H71]
Diện tích hình bình hành có độ dài đáy $a$ và chiều cao tương ứng $h$ là
\choice
{$S = (a + h) \times 2$}
{\True $S = a \times h$}
{$S = a \times h : 2$}
{$S = a + h$}
\loigiai{
$S = a \times h$.\\
\textbf{Đáp án của câu hỏi là B}
}
\end{ex}

\begin{ex} %[4H13H71]
Năm $2000$ thuộc thế kỉ nào?
\choice
{Thế kỉ XIX}
{Thế kỉ XXI}
{Thế kỉ XVIII}
{\True Thế kỉ XX}
\loigiai{
Năm $2000$ là năm cuối thế kỉ XX.\\
\textbf{Đáp án của câu hỏi là D}
}
\end{ex}

\begin{ex} %[4H13V71]
Một hình bình hành có đáy dài $25\text{ cm}$, chiều cao $14\text{ cm}$. Diện tích hình bình hành là
\choice
{$175\text{ cm}^2$}
{$78\text{ cm}^2$}
{\True $350\text{ cm}^2$}
{$300\text{ cm}^2$}
\loigiai{
$S = 25 \times 14 = 350\text{ cm}^2$.\\
\textbf{Đáp án của câu hỏi là C}
}
\end{ex}

\begin{ex} %[4H13V71]
Một hình thoi có chu vi bằng $68\text{ dm}$. Độ dài cạnh hình thoi là
\choice
{\True $17\text{ dm}$}
{$16\text{ dm}$}
{$18\text{ dm}$}
{$34\text{ dm}$}
\loigiai{
$68 : 4 = 17\text{ dm}$.\\
\textbf{Đáp án của câu hỏi là A}
}
\end{ex}

\begin{ex} %[4H13V71]
Đổi $3$ tấn $50\text{ kg}$ ra ki-lô-gam:
\choice
{$350\text{ kg}$}
{$3500\text{ kg}$}
{$30\,050\text{ kg}$}
{\True $3050\text{ kg}$}
\loigiai{
$3 \times 1000 + 50 = 3050\text{ kg}$.\\
\textbf{Đáp án của câu hỏi là D}
}
\end{ex}

\begin{ex} %[4H13V71]
Một căn phòng hình chữ nhật dài $9\text{ m}$, rộng $6\text{ m}$. Diện tích sàn phòng là
\choice
{\True $54\text{ m}^2$}
{$30\text{ m}^2$}
{$540\text{ dm}^2$}
{$108\text{ m}^2$}
\loigiai{
$9 \times 6 = 54\text{ m}^2$.\\
\textbf{Đáp án của câu hỏi là A}
}
\end{ex}

\begin{ex} %[4H13V71]
Lúc 3 giờ đúng, góc tạo bởi kim giờ và kim phút là góc gì?
\choice
{Góc nhọn}
{Góc tù}
{\True Góc vuông ($90^\circ$)}
{Góc bẹt}
\loigiai{
Lúc 3 giờ tạo thành góc vuông $90^\circ$.\\
\textbf{Đáp án của câu hỏi là C}
}
\end{ex}

\begin{ex} %[4H13C71]
Một thửa ruộng hình bình hành có đáy $40\text{ m}$, chiều cao $25\text{ m}$. Người ta trồng lúa trên thửa ruộng đó, cứ $100\text{ m}^2$ thu được $60\text{ kg}$ thóc. Cả thửa ruộng thu hoạch được số tạ thóc là
\choice
{$5\text{ tạ}$}
{\True $6\text{ tạ}$}
{$7\text{ tạ}$}
{$60\text{ tạ}$}
\loigiai{
Diện tích: $40 \times 25 = 1000\text{ m}^2$.\\
Số thóc: $1000 : 100 \times 60 = 600\text{ kg} = 6$ tạ.\\
\textbf{Đáp án của câu hỏi là B}
}
\end{ex}

\textbf{Phần 2. Trắc nghiệm đúng sai}
\setcounterpageref{ex}{0}

\begin{ex} %[4H13H71]
Xét tính đúng sai về hình học và đo lường:
\choiceTF
{\True Hai đường thẳng song song không bao giờ cắt nhau.}
{\True $1$ thế kỉ $= 100$ năm.}
{$2\text{ m}^2 5\text{ dm}^2 = 250\text{ dm}^2$.}
{\True $2\text{ m}^2 5\text{ dm}^2 = 205\text{ dm}^2$.}
\loigiai{
\begin{itemize}
    \item Song song không cắt nhau: Đúng.
    \item 1 thế kỉ $= 100$ năm: Đúng.
    \item $2\text{ m}^2 5\text{ dm}^2 = 205\text{ dm}^2 \neq 250\text{ dm}^2$. Ý 3 Sai, Ý 4 Đúng.
\end{itemize}
\begin{itemize}
    \item a) \textbf{Đúng}.
    \item b) \textbf{Đúng}.
    \item c) \textbf{Sai}.
    \item d) \textbf{Đúng}.
\end{itemize}
}
\end{ex}

\begin{ex} %[4H13V71]
Một mảnh đất hình chữ nhật có chu vi $100\text{ m}$, chiều dài hơn chiều rộng $10\text{ m}$. Xét tính đúng sai:
\choiceTF
{\True Chiều dài mảnh đất là $30\text{ m}$.}
{\True Chiều rộng mảnh đất là $20\text{ m}$.}
{Diện tích mảnh đất là $500\text{ m}^2$.}
{\True Diện tích mảnh đất là $600\text{ m}^2$.}
\loigiai{
\begin{itemize}
    \item Nửa chu vi: $50\text{ m}$. Dài: $(50 + 10) : 2 = 30\text{ m}$. Đúng.
    \item Rộng: $50 - 30 = 20\text{ m}$. Đúng.
    \item Diện tích: $30 \times 20 = 600\text{ m}^2 \neq 500\text{ m}^2$. Ý 3 Sai, Ý 4 Đúng.
\end{itemize}
\begin{itemize}
    \item a) \textbf{Đúng}.
    \item b) \textbf{Đúng}.
    \item c) \textbf{Sai}.
    \item d) \textbf{Đúng}.
\end{itemize}
}
\end{ex}

\textbf{Phần 3. Trả lời ngắn}
\setcounterpageref{ex}{0}

\begin{ex} %[4H13H71]
Tính diện tích hình bình hành có đáy $16\text{ cm}$, chiều cao $10\text{ cm}$.
\shortans{160}
\loigiai{
\textbf{Đáp số:} $160$\\
$16 \times 10 = 160\text{ cm}^2$.
}
\end{ex}

\begin{ex} %[4H13H71]
Tính chu vi hình thoi có cạnh $15\text{ cm}$.
\shortans{60}
\loigiai{
\textbf{Đáp số:} $60$\\
$15 \times 4 = 60\text{ cm}$.
}
\end{ex}

\begin{ex} %[4H13V71]
Đổi $4$ tấn $6$ tạ ra tạ.
\shortans{46}
\loigiai{
\textbf{Đáp số:} $46$\\
$46$ tạ.
}
\end{ex}

\begin{ex} %[4H13V71]
Đổi $5\text{ m}^2 30\text{ dm}^2$ ra đề-xi-mét vuông.
\shortans{530}
\loigiai{
\textbf{Đáp số:} $530$\\
$530\text{ dm}^2$.
}
\end{ex}

\begin{ex} %[4H13V71]
Năm $1975$ thuộc thế kỉ thứ mấy (ghi bằng số Ả Rập)?
\shortans{20}
\loigiai{
\textbf{Đáp số:} $20$\\
Thế kỉ 20 (XX).
}
\end{ex}

\begin{ex} %[4H13C71]
Một hình thoi có độ dài hai đường chéo là $12\text{ cm}$ và $18\text{ cm}$. Tính diện tích hình thoi theo xăng-ti-mét vuông.
\shortans{108}
\loigiai{
\textbf{Đáp số:} $108$\\
$12 \times 18 : 2 = 108\text{ cm}^2$.
}
\end{ex}

\newpage

\TaoPhieuKiemTra{Chủ đề 13 -- Ôn tập cuối năm}{72}{Ôn tập thống kê và xác suất}{https://loptoanthaylinh.vn/thi-online/toan-4/bai-72}
\textbf{Phần 1. Trắc nghiệm lựa chọn}
\setcounterpageref{ex}{0}

\begin{ex} %[4X13N72]
Cho dãy số liệu về số cây xanh 4 tổ trồng: $15, 18, 20, 17$. Tổ trồng nhiều cây nhất trồng được
\choice
{$18\text{ cây}$}
{$17\text{ cây}$}
{\True $20\text{ cây}$}
{$15\text{ cây}$}
\loigiai{
$20$ cây là số lớn nhất trong dãy.\\
\textbf{Đáp án của câu hỏi là C}
}
\end{ex}

\begin{ex} %[4X13N72]
Khi tung một đồng xu 1 lần, các khả năng có thể xuất hiện ở mặt trên là
\choice
{\True Mặt ngửa hoặc mặt sấp}
{Mặt ngửa hoặc mặt nghiêng}
{Chắc chắn mặt ngửa}
{Chắc chắn mặt sấp}
\loigiai{
Đồng xu có 2 mặt: ngửa hoặc sấp.\\
\textbf{Đáp án của câu hỏi là A}
}
\end{ex}

\begin{ex} %[4X13H72]
Số trung bình cộng của dãy số liệu: $10, 20, 30, 40$ là
\choice
{$20$}
{\True $25$}
{$30$}
{$35$}
\loigiai{
$(10 + 20 + 30 + 40) : 4 = 25$.\\
\textbf{Đáp án của câu hỏi là B}
}
\end{ex}

\begin{ex} %[4X13H72]
Biểu đồ cột số học sinh yêu thích các môn học: Toán ($24$), Tiếng Việt ($20$), Tiếng Anh ($26$). Môn học được nhiều học sinh yêu thích nhất là
\choice
{\True Tiếng Anh}
{Toán}
{Tiếng Việt}
{Cả ba môn như nhau}
\loigiai{
Tiếng Anh có 26 bạn, cao nhất.\\
\textbf{Đáp án của câu hỏi là A}
}
\end{ex}

\begin{ex} %[4X13H72]
Tung đồng xu 20 lần, có 11 lần xuất hiện mặt ngửa. Số lần xuất hiện mặt sấp là
\choice
{$8$ lần}
{\True $9$ lần}
{$10$ lần}
{$11$ lần}
\loigiai{
$20 - 11 = 9$ lần.\\
\textbf{Đáp án của câu hỏi là B}
}
\end{ex}

\begin{ex} %[4X13H72]
Trong một hộp kín có 10 viên bi đỏ và 5 viên bi xanh. Lấy ngẫu nhiên 1 viên bi. Khả năng lấy được bi đỏ
\choice
{Bé hơn khả năng lấy bi xanh}
{Bằng khả năng lấy bi xanh}
{Không thể xảy ra}
{\True Lớn hơn khả năng lấy bi xanh}
\loigiai{
10 bi đỏ $> 5$ bi xanh nên khả năng lấy bi đỏ cao hơn.\\
\textbf{Đáp án của câu hỏi là D}
}
\end{ex}

\begin{ex} %[4X13V72]
Biểu đồ số sách đọc được trong học kì 2 của 4 bạn: An ($12$), Bình ($15$), Cường ($10$), Dung ($18$). Tổng số sách 4 bạn đã đọc là
\choice
{$50$ cuốn}
{\True $55$ cuốn}
{$60$ cuốn}
{$52$ cuốn}
\loigiai{
$12 + 15 + 10 + 18 = 55$ cuốn.\\
\textbf{Đáp án của câu hỏi là B}
}
\end{ex}

\begin{ex} %[4X13V72]
Gieo một con xúc xắc 30 lần, thấy có 6 lần xuất hiện mặt 1 chấm. Số lần xuất hiện các mặt còn lại là
\choice
{$20$ lần}
{$25$ lần}
{\True $24$ lần}
{$26$ lần}
\loigiai{
$30 - 6 = 24$ lần.\\
\textbf{Đáp án của câu hỏi là C}
}
\end{ex}

\begin{ex} %[4X13V72]
Biểu đồ số điểm 10 của 3 tổ: Tổ 1 ($35$), Tổ 2 ($40$), Tổ 3 ($45$). Trung bình mỗi tổ đạt được số điểm 10 là
\choice
{$38$ điểm}
{$42$ điểm}
{$45$ điểm}
{\True $40$ điểm}
\loigiai{
$(35 + 40 + 45) : 3 = 40$ điểm 10.\\
\textbf{Đáp án của câu hỏi là D}
}
\end{ex}

\begin{ex} %[4X13V72]
Cho dãy số liệu điểm thi: $6, 7, 8, 9, 8, 10, 8, 7, 8$. Số điểm 8 xuất hiện bao nhiêu lần?
\choice
{\True $4$ lần}
{$3$ lần}
{$5$ lần}
{$2$ lần}
\loigiai{
Điểm 8 xuất hiện 4 lần.\\
\textbf{Đáp án của câu hỏi là A}
}
\end{ex}

\begin{ex} %[4X13V72]
Quay vòng quay có 3 màu (Xanh, Đỏ, Vàng) 40 lần: kim chỉ vào màu Đỏ 15 lần, màu Xanh 17 lần. Số lần kim chỉ vào màu Vàng là
\choice
{$7$ lần}
{$9$ lần}
{$10$ lần}
{\True $8$ lần}
\loigiai{
$40 - (15 + 17) = 8$ lần.\\
\textbf{Đáp án của câu hỏi là D}
}
\end{ex}

\begin{ex} %[4X13C72]
Một nhóm 5 học sinh có điểm kiểm tra trung bình là 8 điểm. Nếu thêm điểm của bạn thứ 6 là 9 điểm thì điểm trung bình của cả 6 bạn là bao nhiêu (viết dạng số thập phân)?
\choice
{$8{,}1$}
{$8{,}2$}
{\True $8{,}17$ (hoặc $(40+9)/6 = 49/6$)}
{$8{,}5$}
\loigiai{
Tổng 5 bạn: $8 \times 5 = 40$. Tổng 6 bạn: $40 + 9 = 49$. Trung bình: $49 : 6 \approx 8{,}17$.\\
\textbf{Đáp án của câu hỏi là C}
}
\end{ex}

\textbf{Phần 2. Trắc nghiệm đúng sai}
\setcounterpageref{ex}{0}

\begin{ex} %[4X13H72]
Xét tính đúng sai về các yếu tố thống kê và xác suất:
\choiceTF
{\True Biểu đồ cột giúp trực quan hóa việc so sánh các số liệu.}
{\True Số trung bình cộng nằm giữa giá trị nhỏ nhất và lớn nhất của dãy số liệu.}
{Tung đồng xu 100 lần thì chắc chắn số lần ngửa luôn bằng 50 lần.}
{\True Khi gieo xúc xắc 6 mặt, có 6 khả năng xảy ra đối với số chấm ở mặt trên.}
\loigiai{
\begin{itemize}
    \item Biểu đồ cột trực quan: Đúng.
    \item TBC nằm giữa min và max: Đúng.
    \item Số lần ngửa là ngẫu nhiên, không chắc chắn bằng 50. Sai.
    \item 6 mặt từ 1 đến 6 chấm: Đúng.
\end{itemize}
\begin{itemize}
    \item a) \textbf{Đúng}.
    \item b) \textbf{Đúng}.
    \item c) \textbf{Sai}.
    \item d) \textbf{Đúng}.
\end{itemize}
}
\end{ex}

\begin{ex} %[4X13V72]
Biểu đồ số cây 3 lớp trồng được: 4A ($40\text{ cây}$), 4B ($50\text{ cây}$), 4C ($60\text{ cây}$). Xét tính đúng sai:
\choiceTF
{\True Lớp 4C trồng được nhiều cây nhất.}
{\True Lớp 4A trồng được ít cây nhất.}
{Cả 3 lớp trồng được tất cả $160\text{ cây}$.}
{\True Trung bình mỗi lớp trồng được $50\text{ cây}$.}
\loigiai{
\begin{itemize}
    \item 4C nhiều nhất (60 cây): Đúng.
    \item 4A ít nhất (40 cây): Đúng.
    \item Tổng: $40 + 50 + 60 = 150\text{ cây} \neq 160\text{ cây}$. Sai.
    \item TBC: $150 : 3 = 50\text{ cây}$. Đúng.
\end{itemize}
\begin{itemize}
    \item a) \textbf{Đúng}.
    \item b) \textbf{Đúng}.
    \item c) \textbf{Sai}.
    \item d) \textbf{Đúng}.
\end{itemize}
}
\end{ex}

\textbf{Phần 3. Trả lời ngắn}
\setcounterpageref{ex}{0}

\begin{ex} %[4X13H72]
Tính số trung bình cộng của 3 số: $30, 40, 50$.
\shortans{40}
\loigiai{
\textbf{Đáp số:} $40$\\
$40$.
}
\end{ex}

\begin{ex} %[4X13H72]
Tung đồng xu 40 lần có 22 lần ngửa. Tìm số lần sấp.
\shortans{18}
\loigiai{
\textbf{Đáp số:} $18$\\
$40 - 22 = 18$ lần.
}
\end{ex}

\begin{ex} %[4X13V72]
Cho dãy số liệu: $15, 20, 18, 25, 30$. Tìm số lớn nhất trong dãy.
\shortans{30}
\loigiai{
\textbf{Đáp số:} $30$\\
$30$.
}
\end{ex}

\begin{ex} %[4X13V72]
Biểu đồ số học sinh giỏi 3 lớp: $12, 15, 18$. Tính tổng số học sinh giỏi cả 3 lớp.
\shortans{45}
\loigiai{
\textbf{Đáp số:} $45$\\
$12 + 15 + 18 = 45$.
}
\end{ex}

\begin{ex} %[4X13V72]
Gieo xúc xắc 50 lần có 8 lần ra mặt 6 chấm. Hỏi các mặt còn lại xuất hiện bao nhiêu lần?
\shortans{42}
\loigiai{
\textbf{Đáp số:} $42$\\
$50 - 8 = 42$ lần.
}
\end{ex}

\begin{ex} %[4X13C72]
Bốn bạn có chiều cao trung bình là $135\text{ cm}$. Bạn thứ năm có chiều cao $140\text{ cm}$. Tính chiều cao trung bình của cả 5 bạn theo xăng-ti-mét.
\shortans{136}
\loigiai{
\textbf{Đáp số:} $136$\\
$(135 \times 4 + 140) : 5 = (540 + 140) : 5 = 680 : 5 = 136\text{ cm}$.
}
\end{ex}

\newpage

\TaoPhieuKiemTra{Chủ đề 13 -- Ôn tập cuối năm}{73}{Ôn tập chung cuối năm}{https://loptoanthaylinh.vn/thi-online/toan-4/bai-73}
\textbf{Phần 1. Trắc nghiệm lựa chọn}
\setcounterpageref{ex}{0}

\begin{ex} %[4D13N73]
Số lớn nhất có 7 chữ số là
\choice
{$1\,000\,000$}
{$9\,876\,543$}
{$10\,000\,000$}
{\True $9\,999\,999$}
\loigiai{
Số lớn nhất có 7 chữ số là $9\,999\,999$.\\
\textbf{Đáp án của câu hỏi là D}
}
\end{ex}

\begin{ex} %[4D13N73]
Kết quả của phép tính $\dfrac{2}{5} + \dfrac{3}{5}$ là
\choice
{\True $1$}
{$\dfrac{5}{10}$}
{$\dfrac{6}{5}$}
{$\dfrac{1}{5}$}
\loigiai{
$\dfrac{2 + 3}{5} = \dfrac{5}{5} = 1$.\\
\textbf{Đáp án của câu hỏi là A}
}
\end{ex}

\begin{ex} %[4D13H73]
Kết quả của phép tính $\dfrac{3}{4} \times \dfrac{2}{3}$ là
\choice
{$\dfrac{5}{7}$}
{\True $\dfrac{1}{2}$}
{$\dfrac{6}{7}$}
{$\dfrac{2}{3}$}
\loigiai{
$\dfrac{6}{12} = \dfrac{1}{2}$.\\
\textbf{Đáp án của câu hỏi là B}
}
\end{ex}

\begin{ex} %[4D13H73]
Một hình bình hành có đáy $30\text{ cm}$, chiều cao $15\text{ cm}$. Diện tích của nó là
\choice
{$225\text{ cm}^2$}
{$90\text{ cm}^2$}
{\True $450\text{ cm}^2$}
{$45\text{ cm}^2$}
\loigiai{
$S = 30 \times 15 = 450\text{ cm}^2$.\\
\textbf{Đáp án của câu hỏi là C}
}
\end{ex}

\begin{ex} %[4D13H73]
$2$ tấn $5$ tạ bằng bao nhiêu tạ?
\choice
{$205\text{ tạ}$}
{\True $25\text{ tạ}$}
{$250\text{ tạ}$}
{$2500\text{ tạ}$}
\loigiai{
$2 \times 10 + 5 = 25$ tạ.\\
\textbf{Đáp án của câu hỏi là B}
}
\end{ex}

\begin{ex} %[4D13H73]
Hai số có tổng bằng $120$ và hiệu bằng $20$. Số lớn là
\choice
{$60$}
{$80$}
{$50$}
{\True $70$}
\loigiai{
$(120 + 20) : 2 = 70$.\\
\textbf{Đáp án của câu hỏi là D}
}
\end{ex}

\begin{ex} %[4D13V73]
Một mảnh đất hình chữ nhật có chu vi $140\text{ m}$, chiều dài hơn chiều rộng $10\text{ m}$. Diện tích mảnh đất là
\choice
{$1100\text{ m}^2$}
{$1300\text{ m}^2$}
{\True $1200\text{ m}^2$}
{$1400\text{ m}^2$}
\loigiai{
Nửa chu vi: $70\text{ m}$. Dài: $40\text{ m}$, Rộng: $30\text{ m}$.\\
Diện tích: $40 \times 30 = 1200\text{ m}^2$.\\
\textbf{Đáp án của câu hỏi là C}
}
\end{ex}

\begin{ex} %[4D13V73]
Một người có $120\,000$ đồng, đã tiêu hết $\dfrac{3}{4}$ số tiền. Số tiền còn lại là
\choice
{$40\,000$ đồng}
{\True $30\,000$ đồng}
{$50\,000$ đồng}
{$90\,000$ đồng}
\loigiai{
Đã tiêu: $120\,000 \times \dfrac{3}{4} = 90\,000$ đồng. Còn lại: $30\,000$ đồng.\\
\textbf{Đáp án của câu hỏi là B}
}
\end{ex}

\begin{ex} %[4D13V73]
Tính thuận tiện: $25 \times 18 \times 4$.
\choice
{$1600$}
{$2000$}
{\True $1800$}
{$1700$}
\loigiai{
$(25 \times 4) \times 18 = 100 \times 18 = 1800$.\\
\textbf{Đáp án của câu hỏi là C}
}
\end{ex}

\begin{ex} %[4D13V73]
Một hình thoi có chu vi $80\text{ cm}$. Độ dài cạnh hình thoi là
\choice
{\True $20\text{ cm}$}
{$15\text{ cm}$}
{$25\text{ cm}$}
{$40\text{ cm}$}
\loigiai{
$80 : 4 = 20\text{ cm}$.\\
\textbf{Đáp án của câu hỏi là A}
}
\end{ex}

\begin{ex} %[4D13V73]
Tính giá trị biểu thức: $\dfrac{3}{5} + \dfrac{1}{2} - \dfrac{7}{10}$.
\choice
{$\dfrac{1}{5}$}
{$\dfrac{3}{5}$}
{$\dfrac{4}{5}$}
{\True $\dfrac{2}{5}$}
\loigiai{
$\dfrac{6 + 5 - 7}{10} = \dfrac{4}{10} = \dfrac{2}{5}$.\\
\textbf{Đáp án của câu hỏi là D}
}
\end{ex}

\begin{ex} %[4D13C73]
Tính tổng: $S = \dfrac{1}{1 \times 2} + \dfrac{1}{2 \times 3} + \dots + \dfrac{1}{9 \times 10}$.
\choice
{\True $\dfrac{9}{10}$}
{$\dfrac{8}{9}$}
{$1$}
{$\dfrac{10}{11}$}
\loigiai{
$S = 1 - \dfrac{1}{10} = \dfrac{9}{10}$.\\
\textbf{Đáp án của câu hỏi là A}
}
\end{ex}

\textbf{Phần 2. Trắc nghiệm đúng sai}
\setcounterpageref{ex}{0}

\begin{ex} %[4D13H73]
Xét tính đúng sai của các khẳng định toàn bộ chương trình Toán 4:
\choiceTF
{\True $1$ tấn $= 1000\text{ kg}$.}
{\True Hai đường chéo hình thoi vuông góc với nhau.}
{$\dfrac{4}{5} \times \dfrac{5}{4} = 0$.}
{\True $\dfrac{4}{5} \times \dfrac{5}{4} = 1$.}
\loigiai{
\begin{itemize}
    \item 1 tấn $= 1000$ kg: Đúng.
    \item Đường chéo hình thoi vuông góc: Đúng.
    \item Tích bằng 1 (không phải 0). Ý 3 Sai, Ý 4 Đúng.
\end{itemize}
\begin{itemize}
    \item a) \textbf{Đúng}.
    \item b) \textbf{Đúng}.
    \item c) \textbf{Sai}.
    \item d) \textbf{Đúng}.
\end{itemize}
}
\end{ex}

\begin{ex} %[4D13V73]
Một hình chữ nhật có chiều dài $20\text{ m}$, chiều rộng bằng $\dfrac{3}{4}$ chiều dài. Xét tính đúng sai:
\choiceTF
{\True Chiều rộng hình chữ nhật là $15\text{ m}$.}
{\True Chu vi hình chữ nhật là $70\text{ m}$.}
{Diện tích hình chữ nhật là $250\text{ m}^2$.}
{\True Diện tích hình chữ nhật là $300\text{ m}^2$.}
\loigiai{
\begin{itemize}
    \item Rộng: $20 \times 3/4 = 15\text{ m}$. Đúng.
    \item Chu vi: $(20 + 15) \times 2 = 70\text{ m}$. Đúng.
    \item Diện tích: $20 \times 15 = 300\text{ m}^2 \neq 250\text{ m}^2$. Ý 3 Sai, Ý 4 Đúng.
\end{itemize}
\begin{itemize}
    \item a) \textbf{Đúng}.
    \item b) \textbf{Đúng}.
    \item c) \textbf{Sai}.
    \item d) \textbf{Đúng}.
\end{itemize}
}
\end{ex}

\textbf{Phần 3. Trả lời ngắn}
\setcounterpageref{ex}{0}

\begin{ex} %[4D13H73]
Tính tổng: $450\,000 + 550\,000$.
\shortans{1000000}
\loigiai{
\textbf{Đáp số:} $1000000$\\
$1\,000\,000$.
}
\end{ex}

\begin{ex} %[4D13H73]
Tính tích: $\dfrac{3}{7} \times \dfrac{14}{9}$ (dạng a/b).
\shortans{2/3}
\loigiai{
\textbf{Đáp số:} $2/3$\\
$2/3$.
}
\end{ex}

\begin{ex} %[4D13V73]
Một hình thoi có cạnh $12\text{ cm}$. Tính chu vi hình thoi theo xăng-ti-mét.
\shortans{48}
\loigiai{
\textbf{Đáp số:} $48$\\
$12 \times 4 = 48\text{ cm}$.
}
\end{ex}

\begin{ex} %[4D13V73]
Tìm số trung bình cộng của 3 số: $40, 60, 80$.
\shortans{60}
\loigiai{
\textbf{Đáp số:} $60$\\
$60$.
}
\end{ex}

\begin{ex} %[4D13V73]
Tìm $\dfrac{3}{5}$ của $50\text{ kg}$.
\shortans{30}
\loigiai{
\textbf{Đáp số:} $30$\\
$50 \times 3/5 = 30\text{ kg}$.
}
\end{ex}

\begin{ex} %[4D13C73]
Tính tổng: $1 + 3 + 5 + 7 + \dots + 19$.
\shortans{100}
\loigiai{
\textbf{Đáp số:} $100$\\
Số số hạng: $(19 - 1) : 2 + 1 = 10$. Tổng: $(1 + 19) \times 10 : 2 = 100$.
}
\end{ex}

\newpage
\end{document}
